% Options for packages loaded elsewhere
\PassOptionsToPackage{unicode}{hyperref}
\PassOptionsToPackage{hyphens}{url}
\documentclass[
]{article}
\usepackage{xcolor}
\usepackage{amsmath,amssymb}
\setcounter{secnumdepth}{-\maxdimen} % remove section numbering
\usepackage{iftex}
\ifPDFTeX
  \usepackage[T1]{fontenc}
  \usepackage[utf8]{inputenc}
  \usepackage{textcomp} % provide euro and other symbols
\else % if luatex or xetex
  \usepackage{unicode-math} % this also loads fontspec
  \defaultfontfeatures{Scale=MatchLowercase}
  \defaultfontfeatures[\rmfamily]{Ligatures=TeX,Scale=1}
  % unicode-math maps \setminus to U+29F5 (REVERSE SOLIDUS OPERATOR), which
  % latinmodern-math.otf lacks (tectonic/lualatex emit "Missing character
  % U+29F5"). Remap to U+2216 (SET MINUS) -- the glyph the font provides and
  % that matches the pdfTeX/amssymb rendering. Output-preserving; silences the
  % warning. Applies to all \setminus occurrences globally.
  \AtBeginDocument{\renewcommand{\setminus}{\mathbin{\char"2216}}}
\fi
\usepackage{lmodern}
\ifPDFTeX\else
  % xetex/luatex font selection
\fi
% Use upquote if available, for straight quotes in verbatim environments
\IfFileExists{upquote.sty}{\usepackage{upquote}}{}
\IfFileExists{microtype.sty}{% use microtype if available
  \usepackage[]{microtype}
  \UseMicrotypeSet[protrusion]{basicmath} % disable protrusion for tt fonts
}{}
\makeatletter
\@ifundefined{KOMAClassName}{% if non-KOMA class
  \IfFileExists{parskip.sty}{%
    \usepackage{parskip}
  }{% else
    \setlength{\parindent}{0pt}
    \setlength{\parskip}{6pt plus 2pt minus 1pt}}
}{% if KOMA class
  \KOMAoptions{parskip=half}}
\makeatother
\setlength{\emergencystretch}{3em} % prevent overfull lines
\providecommand{\tightlist}{%
  \setlength{\itemsep}{0pt}\setlength{\parskip}{0pt}}
\usepackage{bookmark}
\IfFileExists{xurl.sty}{\usepackage{xurl}}{} % add URL line breaks if available
\urlstyle{same}
\hypersetup{
  hidelinks,
  pdfcreator={LaTeX via pandoc}}

\author{}
\date{}

\begin{document}

\textbf{Ratushnyi Neural Repair Coding, Part I: Framework, Constructions,
and Complexity}

\emph{A Framework for Lossless Neural Compression via Side‑Information
Repair: Stream, Positional, and Tensor Entropy‑Field Coding}

Pavlo Ratushnyi

\emph{Independent Research --- Draft v0.8, May 2026}

\textbf{Abstract}

We introduce Ratushnyi Neural Repair Coding (RNR Coding), a framework
for lossless compression in which a deterministic,
decoder-reproducible neural or hybrid model produces side information
about a source block, and a short repair description reconstructs the
exact original from that side information. The framework is organized
along the dimensionality and target of prediction. Type-I performs
one-dimensional stream repair; Type-II performs two-dimensional
positional symbol coding over a partition of positions among symbol
classes; Type-III performs multi-dimensional tensor entropy-field
coding, divided into III-A (entropy-field meta-prediction), III-B
(byte-tensor coding of joint nibble structure), and III-C
(token-tensor coding over a learned hierarchically indexed
dictionary). For Types I, II, III-A, and III-B we give concrete
constructions and achievability bounds; for Type-III-C we give a
greedy marginal-MDL admission criterion (Theorem 6.5), a
bits-back rate analysis (Theorem 6.6), and a standalone
achievability bound for the learned hierarchical
dictionary/token-cover construction (Theorem 6.6b): when the
dictionary is built by a universal grammar transform (an LZ78
incremental parser or an irreducible-grammar transform in the sense of
Kieffer--Yang 2000), the construction attains $N \cdot h(X) + o(N)$ for
stationary ergodic sources, the hierarchical index being rate-neutral
(Lemma 6.6a) and the $K = O(N/\log N)$ dictionary entries adding only
$o(N)$ overhead; whether the specific greedy marginal-MDL builder of
Theorem 6.5 is itself a universal transform is partially resolved: it is
universal once its literals are entropy-coded under a universal predictor
and a one-bit baseline guard is added (Lemma 6.6d), leaving only the bare
flat-$8\ell$ literal rule open (Remarks 6.6b, 6.6e).
Matching converses are given for Type-I (Theorem 4.2), Type-II
(Theorem 5.3b), and Type-III-B (Theorem 6.3b, the byte coder cannot beat
the joint conditional entropy $H(H,L\mid C)$, so the $I(H;L\mid C)$
advantage of Theorem 6.3 is exactly tight). For Type-II the base bound
$\mathbb{E}[L] \geq H(X \mid Y)$ holds for every code and is the operative
converse, met (up to coder redundancy) by ideal joint coding of the
partition and also by the count-plus-rank construction of
Theorems 5.2--5.3 once its enumerative index is entropy-coded; the
flat-enumerative candidate-set cost (the multinomial baseline when no
support is used) is a converse only for the uniform-index scheme,
equalling $H(X \mid Y)$ exactly when the conditional law is uniform on the
decoder-available candidate set (so Theorem 5.3 is entropy-optimal there)
and otherwise exceeding it by the recoverable within-candidate-set
uniform deficit, a KL divergence that is the price of uniform indexing
rather than of the partition representation. Type-III-C is
\emph{asymptotically rate-optimal} under a universal dictionary builder:
the achievability bound (Theorem 6.6b) is matched by the entropy-rate
converse (Theorem 6.6c), so $\tfrac{1}{N}\mathbb{E}[L_{\mathrm{IIIC}}]
\to h(X)$ whenever the dictionary is built universally (the bare
greedy-MDL builder's own universality being the one residual sub-question,
Remark 6.6b, partially resolved in augmented form by Lemma 6.6d). Type-III-A has both a \emph{relative} converse --- the
$2\varepsilon$ approximate-oracle gap is tight (Theorem 6.2) --- and an
\emph{absolute-rate} converse (Theorem 6.2b): $\mathbb{E}[L_{\mathrm{IIIA}}]
\ge H(X\mid Y)$, attained within $2\sum_t\varepsilon_t$ iff the modes are
per-tile-expressive \emph{and} causally context-sufficient (each tile
coded from its full decoder-reproducible past), the slack otherwise being
the context deficit --- maximal, and equal to the inter-tile total
correlation $\mathrm{TC}(X_1,\ldots,X_T\mid Y)$, for the block-independent
variant, but removable by autoregressive mode conditioning. Complexity
results include
NP-hardness of the
varying-parameter learned-code labeling problem (Theorem 4.3), a
polynomial-time vs.\ NP-hard dichotomy in semi-non-factorized
Type-II support selection (Lemma 5.7a polynomial-time for
permutation-like joint laws; Lemma 5.7c NP-hard for general joint
laws via CLIQUE reduction), and NP-hardness of a restricted
normal-form dictionary-only case of Type-III-C (Theorem 6.7), plus
an APX-hardness transfer for the same restricted case under the
unbounded-alphabet variant with explicit inapproximability constant
$\gamma_{\mathrm{SGP}} = 8569/8568$ inherited from Charikar et al.\
(2005, Theorem 1) via the symbol-count bijection (Theorem 6.8); the
fixed-byte-alphabet case ($|\Sigma| = 256$ as in canonical RNR)
requires a separate gap-preserving construction and is open.
Hardness of the unrestricted joint (model, dictionary, repair)
problem and the analogous inapproximability for $E_{12}$ are also
open. For the continuous (Gaussian-field) Type-III-C limit the
support-selection problem is characterised completely by complexity
regime: with the field model fixed the per-block residual is a
Schur-complement log-determinant whose minimisation is the
maximum-entropy-sampling problem, hence APX-hard with constant $5/4$
under $\mathrm{P} \neq \mathrm{NP}$ (Lemma 6.9, via Ko--Lee--Queyranne
and Ohsaka; \textsc{Partial}, the gap surviving the manifold
covariance restriction at high SNR), whereas with the predictor free
and amortised the same objective is in $\mathrm{P}$ via a
sample-covariance maximum-likelihood plug-in (Remark 6.9a), the
crossover budget being characterised in Remark 6.9c.
Worked numerical examples are
given for Type-I, Type-II, and Type-III-B; III-A and III-C are
illustrated through their composition and dictionary-construction
sections respectively. Two structural results identify the
information-theoretic quantity governing factorized-vs-joint
separation: a $\Theta(N)$ amortized separation for a specific
linear-code family in semi-non-factorized Type-II (Theorem 5.4), and
the identity that for an arbitrary source $D_N$ the gap equals
multi-information $\mathrm{TC}(D_N)$ (Theorem 5.5, where the gap may
vanish on sources without joint structure). Type-III-A admits a tight
$2\varepsilon$ approximate-oracle bound (Theorem 6.1, with matching
lower bound Theorem 6.2) and an absolute-rate converse
$\mathbb{E}[L_{\mathrm{IIIA}}]\ge H(X\mid Y)$ whose slack is the
per-tile context deficit --- at most the inter-tile total correlation
$\mathrm{TC}(X_1,\ldots,X_T\mid Y)$ (the block-independent worst case),
and zero under causal-context modes (Theorem 6.2b).
Bit-exact cross-platform reproduction is
reduced to a precision-and-determinism specification (Theorem 10.1);
published integer-arithmetic methods (integer SGD, signSGD, integer
quantized inference) provide candidate components that an implementer
must verify against the full specification on a chosen target. We do not claim universal
dominance --- no lossless code can map every input of length $n$ to a
strictly shorter one (Theorem 8.1) --- and the framework is evaluated
conditionally, in the precise sense of the conditionally-advantageous
predicate of Definition 2.1: on data classes for which side
information reduces residual uncertainty by more than the model,
repair, verification, and metadata cost, and only in the
model-amortization regime that makes the model-description cost
payable (§10.2). The qualifying data-class precondition is assembled
as one explicit inequality in §10.2.1; which real corpora satisfy it
is an open empirical question (§11). Companion parts develop the random-access and deviation theory ([RNR-II]) and the operational rate--distortion dispersion ([RNR-III]).


\emph{Numbering follows the unified RNR series; this Part contains
Sections 1--6, the composition and mode-selection lead of Section 7
(Theorems 7.1/7.1b--7.1d and the runtime-parameter text of §7.2),
Sections 8--9, the practical core §§10.1--10.6, and Sections 11--13.
Results proved in a companion Part are cited as [RNR-II]/[RNR-III]
(Parts II/III): Part II develops the random-access and deviation theory
(Theorems 7.2--7.33, 7.35, 7.36; §§10.7--10.8), and Part III the
operational rate--distortion dispersion (the §7.34 family). Section and
theorem numbers are preserved across the series, so a companion Part is
referenced by its original number (e.g.\ Theorem 7.15 of [RNR-II]); the
§7.2 runtime-parameter subsection of this Part and Theorem 7.2 of [RNR-II]
are distinct objects that share a number by this convention.}

\textbf{Keywords:} lossless compression, neural compression, side
information, Slepian--Wolf coding, syndrome coding, positional coding,
enumerative coding, entropy‑field coding, tensor coding, Minimum
Description Length, bits‑back coding, RNR Coding.

\textbf{1. Introduction}

\textbf{1.1 Motivation}

Classical lossless compressors --- LZ77/LZ78 derivatives, LZMA, the
Burrows--Wheeler family, arithmetic‑coded context mixers such as PAQ and
Cmix, and modern entropy coders based on arithmetic coding or asymmetric
numeral systems --- are highly effective on a wide range of inputs
because they exploit repeated substrings, local contexts, block sorting,
and adaptive probability models. Their operational question is
sequential: given the bytes seen so far, what is the next byte, or where
did this fragment occur previously? This question is well‑matched to
data whose redundancy is dominated by local repetition or short‑range
statistics.

Many practically important data classes are not well described by this
question. Database pages, raw or lightly preprocessed image sensor data,
telemetry grids, structured binary sections, scientific tensors, and
certain instruction streams contain structure that is positional,
geometric, multi-scale, or coupled across non-adjacent positions. For
such data, the conjecture motivating this paper is that sequential
redundancy is weaker than non-local structural regularity, and
dictionary-based methods leave entropy unexploited that more structural
predictors could in principle reach. The conjecture is empirical, not
proven here: §11 lists falsifiable bit-per-byte targets for specific
corpora, each cross-referenced to a numbered, pre-registered hypothesis
(H1--H7, H15--H16) in the experimental-design companion document, which
specifies the pre-registered measurement protocol and its named
baselines (including the seekable formats bgzip and zstd-seekable and
the deployed format squashfs+zstd). Within this paper we make no
quantitative claim about absolute compression ratios on any specific
corpus.

Neural lossless compressors such as Bellard's NNCP, PixelCNN‑style
autoregressive models, integer discrete flows, and bits‑back schemes
(BB‑ANS, Bit‑Swap, HiLLoC) demonstrate that learned predictors can drive
arithmetic or numeral‑system coders to bit‑rates below those of strong
classical baselines on certain corpora. These methods almost universally
adopt the same operational question as classical coders --- predict the
next symbol --- and replace the predictor with a neural network.

A second obstruction blocks neural compressors from one of the largest
practical deployment niches for lossless coding: archive and
time-series storage engines inside database servers. Compressed
read-only filesystems (squashfs, EROFS), columnar archive engines, and
time-partitioned metric stores rely on \emph{random access} ---
reading a byte at offset $p$ in time polylogarithmic in archive
size $|X|$. Dictionary-based seekable formats (bgzip, zstd seekable)
achieve this operationally via independent compressed frames plus a
seek index. Autoregressive neural byte models (NNCP and similar) do
not, because their predictors require sequential evaluation of all
bytes $[0, p)$ before byte $p$ is available; the latent-variable lossless schemes
(IDF, Bit-Swap, HiLLoC) decode in block units that can be larger than
the desired random-access granularity. The author's engineering
experience with a Rust MySQL-compatible server motivated this
obstruction concretely (compressed-archive and time-series engines
across 20 storage engines): when ratio matters and seek latency
matters jointly, neither existing class fits the niche. This anecdotal
motivation is what suggested treating the neural model output as side
information rather than as a sequential predictor; the question
addressed in this paper is whether such a recasting can support random
access at sync-point granularity without losing the ratio advantage
(an empirical question deferred to §11).

This work proposes a different organizing question. Treat the neural
model output as side information about the source, not as a candidate
reconstruction. A neural generator alone cannot losslessly compress
data, because any single bit of disagreement between generator output
and ground truth is fatal. But if the model output is treated as a noisy
or partial description of the data --- a known damaged version, a
positional probability field, or a tensor of predicted symbol
probabilities and reliability --- then the archive needs to store only
the repair information that, combined with the model output, recovers
the exact original.

The Ratushnyi Neural Repair Coding (RNR) framework formalizes this idea
along a single axis: the dimensionality and target of prediction. Type‑I
predicts a one‑dimensional stream; Type‑II predicts a two‑dimensional
symbol‑position partition; Type‑III predicts a multi‑dimensional tensor
that may include the predicted entropy of prediction itself. The central
claim is conditional, not universal: there is no claim of universal
dominance (Theorem 8.1), and ``conditional'' here means exactly the
conditionally-advantageous predicate of Definition 2.1
($\mathbb{E}[L(\mathsf{Enc}(X))] < \mathbb{E}[L(B(X))]$). For data
classes in which neural side information reduces residual uncertainty
by more than the cost of describing the model, the repair, the
verification, and the metadata --- and only in the model-amortization
regime in which the per-archive model cost $L(M)/m$ is payable
(§10.2) --- RNR can in principle outperform sequential dictionary or
entropy‑only baselines. The qualifying precondition is stated as a
single inequality with each cost term named and §-anchored in §10.2.1;
the structural data-class families plausibly satisfying it are listed
there as a conjecture, and which real corpora actually qualify is an
open empirical question deferred to §11.

\textbf{1.2 Position relative to existing work}

RNR draws on five established lines of work. Slepian and Wolf
established that a source can be encoded losslessly at a rate determined
by its conditional entropy given correlated side information at the
decoder. Syndrome‑based and parity‑check constructions, including LDPC
and BCH codes, give practical machinery for treating prediction error as
a virtual channel noise process. Cover's enumerative source coding shows
how subsets and sequences under combinatorial constraints can be encoded
by their rank among admissible possibilities. Rissanen's Minimum
Description Length principle provides the criterion under which a
learned representation justifies its own overhead. Arithmetic coding and
asymmetric numeral systems provide near‑optimal entropy coders for given
probability models.

Within neural lossless compression, the closest reference points are
NNCP (autoregressive byte‑level transformer with arithmetic coding),
PixelCNN++ and PixelSNAIL (autoregressive factorized image models with
categorical likelihoods), the bits‑back family (BB‑ANS, Bit‑Swap,
HiLLoC, encoding latent variables of a variational model losslessly via
ANS), integer discrete flows (IDF, IDF++), and L3C for hierarchical
image priors. Among classical context‑mixing coders, PAQ and Cmix
represent the upper end of what mixture‑of‑experts entropy coding
achieves without neural networks.

RNR is not a replacement for any of these. It is a taxonomy and a set
of constructions that, in our reading, naturally organize them ---
though the mappings below are interpretive (not theorems) and are
included to fix the vocabulary of later sections. Per-position
arithmetic coding over a factorized predicted distribution
(NNCP-style direct arithmetic coding) is \emph{rate-equivalent} to
the ideal Type-I class+payload coder under a matching predictive
law, but is not literally a special case of the Type-I construction
as defined here (§4.5 gives the precise relationship). PixelCNN++
over images is autoregressive over pixels conditioned on prior
pixels (categorical or discretized-mixture likelihoods, depending
on the head), and the relation to Type-II is analogy not
equivalence: Type-II as defined in §3.2 is partition-based and does
not subsume PixelCNN++'s pixel-wise softmax. Bits-back coding
on a continuous latent fits into Type-III-C as the limit where the
discrete dictionary is replaced by a variational posterior, with
exact cross-entropy plus KL accounting (Theorem 6.6). PAQ-style
switching context mixing can be viewed as a hand-designed
Type-III-A-like scheduler with a continuous-weight logistic mix
(rather than a discrete argmin), so the mapping is analogy rather
than instantiation.
The framework's intended contribution is to give a unified vocabulary
and a set of constructions for three regimes that the existing methods
treat ad hoc or not at all --- semi-non-factorized positional models,
predicted entropy fields, and structured token tensors. We do not
claim that existing methods are empirically weak in these regimes;
Theorem 5.4 specifically separates only factorized-marginal coders
from semi-non-factorized ones, and a true autoregressive coder can in
principle match the joint-coding rate on the same source. Whether the
RNR constructions outperform existing methods on natural data classes
is the empirical question of §11; the present paper provides only the
theoretical scaffolding for the comparison.

\textbf{1.3 Contributions}

This article makes the following contributions.

\begin{itemize}
\item
  A framework definition (Section 3) that decouples the choice of
  predictor (the model M) from the choice of representation target
  (stream, partition, tensor) and the choice of repair encoding (mask,
  residual, rank, selector).
\item
  Concrete constructions for each of the five types: Type-I with an
  explicit Code12 systematic shortened-Hamming-like byte-to-12-bit
  mapping (Section 4); semi-non-factorized Type-II positional coding
  with disjoint-partition and count constraints (Section 5); and
  Type-III-A/B/C with entropy-field, byte-tensor, and hierarchical
  $32 \times 32$ token-tensor constructions (Section 6).
\item
  Achievability and converse bounds for Type-I (Theorems 4.1, 4.2):
  the \emph{ideal} Type-I rate (true conditional class+payload
  model, Theorem 4.2 part (A)) matches the conditional source-coding
  lower bound $H(X \mid Y)$ to first order in $N$; the concrete
  Theorem 4.1 construction, which uses a fixed $H(P_E)$ class model
  and a uniform within-class payload, achieves
  $N \cdot H(P_E) + N \cdot \ell_{\text{cond}} + \varepsilon_N$
  bits and matches $H(X \mid Y)$ only when both the class model and
  payload model match the true conditional law on $(C_E, Y)$ and
  $(E \mid C_E, Y)$ respectively; otherwise the gap is the
  cross-entropy excess plus KL, which can be $\Theta(N)$.
\item
  Complexity boundary for Type-I codeword selection: the fixed
  Code12 selection at $n_0 = 256, m_0 = 12$ is a finite constant-size
  optimization (Remark 4.3a), while the varying-parameter
  learned-code labeling problem (alphabet size $n$ and codeword
  length $m$ as input) is NP-hard via reduction from Wagner and
  Corneil's (1990) tree-to-hypercube subgraph embedding (Theorem 4.3).
\item
  A corrected positional-advantage theorem (Theorem 5.3) with a
  multinomial baseline and an explicit escape-mode cost term,
  addressing two errors in earlier drafts that overstated the
  Type-II advantage.
\item
  A matching Type-II converse (Theorem 5.3b), completing the Type-II
  proof matrix in the shape of the Type-I converse (Theorem 4.2(C)):
  the base bound $\mathbb{E}[L_{\mathrm{II}}] \geq H(X \mid Y)$ holds
  for every Type-II code and is the operative converse, met by ideal
  joint coding of the partition and equally by the count-plus-rank
  construction of Theorems 5.2--5.3 once its enumerative index is
  entropy-coded. The flat-enumerative candidate-set cost
  $\log_2 |\mathcal{C}(Y,\mathbf{k})|$ (the multinomial baseline only when
  no support is used) is a converse only for the uniform-index scheme; it
  equals $H(X \mid Y)$ exactly when the conditional law is uniform on the
  decoder-available candidate set (so Theorem 5.3 is entropy-optimal
  there) and otherwise exceeds it by the recoverable within-candidate-set
  uniform deficit (a KL divergence, with source-dependent support metadata
  charged). Three worked counterexamples ($N=4$, $\mathbf{k}=(1,3)$) show
  this overhead is the price of uniform indexing on the candidate set ---
  not of the partition representation, not measured against the full count
  class once a support is used, and not evadable by source-dependent
  support metadata (which is charged); it is thus not misattributed to
  Type-II as such.
\item
  An amortized separation theorem (Theorem 5.4) for
  semi-non-factorized Type-II via a specific linear-code
  construction, with the exact information-theoretic quantity
  governing factorized-versus-joint coding identified as
  multi-information $\mathrm{TC}(D_N)$ (Theorem 5.5; the gap may be
  zero on sources without joint structure). Four explicit
  $\Theta(N)$-TC classes are established within the §5.6 framework:
  Lemma 5.6a gives $\mathrm{TC} = N \cdot c(\alpha) + O(\log N)$ for
  the uniqueness-constrained class (column-permutation invariance);
  Lemma 5.6b gives the exact identity
  $\mathrm{TC} = N \cdot I(X_1; Y) - H(Y) + H(Y \mid X^N)$ for
  hidden-class iid sources (latent-topic mixture); Lemma 5.6c
  gives $\mathrm{TC} = N \cdot E - \delta_N$ for stationary
  processes, where $E = H_{\mathrm{marg}} - h(X)$ is the per-position
  excess entropy and $\delta_N$ is bounded for order-$W$ Markov
  chains (temporal correlation); and Lemma 5.6d extends to hidden
  Markov models with the Baum--Petrie filter-recursion entropy rate
  $h_{\mathrm{HMM}}$ (combined hidden-and-temporal structure).
  Together these four lemmas cover the principal mechanisms by which
  natural data accumulates linear TC. Whether $\mathrm{TC}(D_N) =
  \Theta(N)$ extends to further classes (factor graphs over
  expanders, continuous-state HMMs, non-stationary processes)
  remains an open empirical question.
\item
  A status remark on Type-II support selection
  (Remark 5.7): the natural marginal-$Q$-prefix heuristic is not
  provably optimal in the semi-non-factorized regime in general ---
  expected cost depends on the full joint law of the partition sets,
  not just marginals --- and a concrete counterexample (at $N=7$,
  $k_v=2$, $s=2$) is included in the paper's verification scripts.
  Lemma 5.7a recovers optimality of top-$Q$ for the
  \emph{permutation-like} (conditional Bernoulli) class of joint laws
  where atom weights factorize as $\Pr(M_v = T) \propto
  \prod_{i \in T} q_i$ (covering the uniform-on-fixed-cardinality
  and log-likelihood-weighted regimes), via a stochastic dominance
  coupling. Lemma 5.7c then closes the complementary direction:
  optimal fixed-size support selection for arbitrary joint laws is
  NP-hard, via a CLIQUE-to-support-selection reduction at $N=3s+1$
  where the algebraic identity $A_s(1) = A_s(2)$ collapses the
  cost to a linear function of the densest-$k$-subgraph value
  $n_0(S)$. Together Lemma 5.7a and Lemma 5.7c form a polynomial-time
  vs.\ NP-hard dichotomy in the joint law's permutation-like structure.
\item
  An approximate-oracle mode-selection bound for Type-III-A (Theorem
  6.1) and a tight matching lower bound (Theorem 6.2) showing the
  $2\varepsilon$ factor cannot be improved without additional
  assumptions on the prediction error structure.
\item
  A mutual-information characterization of the Type-III-B byte-tensor
  advantage (Theorem 6.3) with a matching Shannon converse
  (Theorem 6.3b) proving the $I(H;L\mid C)$ saving is exactly the
  maximum achievable by any uniquely decodable byte coder, and an
  invariance result (Theorem 6.4) showing all rooted binary tree
  decompositions of the 256-byte alphabet achieve the same
  information-theoretic rate.
\item
  Six results on Type-III-C: the greedy MDL admission criterion
  (Theorem 6.5), the bits-back integration with an exact
  cross-entropy plus KL-divergence accounting (Theorem 6.6), a
  standalone entropy-rate achievability bound $N \cdot h(X) + o(N)$ for
  the learned hierarchical dictionary under a universal grammar
  transform (LZ78 or irreducible-grammar, Theorem 6.6b), with the
  hierarchical family/variant index shown rate-neutral (Lemma 6.6a) and
  the greedy-MDL builder's own universality partially resolved (universal
  under an entropy-coded literal fallback and baseline guard, Lemma 6.6d;
  the bare flat-$8\ell$ rule open, Remark 6.6b),
  NP-hardness of the restricted (constant-predictor, normal-form
  recursive dictionary, symbol-count objective) dictionary-learning
  problem (Theorem 6.7), APX-hardness of the same restricted
  case under the \emph{unbounded-alphabet} variant with explicit
  constant $\gamma_{\mathrm{SGP}} = 8569/8568$ inherited from
  Charikar et al.\ (2005, Theorem 1) via the symbol-count bijection
  (Theorem 6.8), and an extension of the same APX-hardness constant
  to the \emph{bit-cost} objective $L^{\mathrm{bit}} =
  L_{\mathrm{sym}} \cdot \lceil \log_2(|\Sigma|+K) \rceil$ within the
  Charikar normal-form grammar class on Berman-Karpinski-hard
  instances (Lemma 6.8b), via a ratio-preserving reduction exploiting
  monotonicity of $L^{\mathrm{bit}}$ in the vertex-cover size
  $|\sigma|$. The 8569/8568 constant derives from the
  Berman-Karpinski (1999) bound on Vertex-Cover-3 inapproximability
  propagated through the $15|V|+3|E|+k$ encoding overhead. The
  fixed-byte-alphabet case requires a fresh gap-preserving
  construction and we do not claim a specific constant for it;
  unrestricted bit-cost APX-hardness over all (non-Charikar-format)
  grammars and unrestricted-predictor joint hardness remain open.
\item
  A complete complexity characterization of the \emph{continuous}
  (Gaussian-field) Type-III-C support-selection problem by regime
  (§13.4.4): with the field covariance $\Sigma$ fixed, the per-block
  residual rate is the Schur-complement conditional log-determinant
  $\tfrac12\log_2\det\Sigma_{S\mid\bar S}$, whose minimisation over
  size-$s$ supports is the maximum-entropy-sampling / $D$-optimal
  design problem, hence APX-hard with constant $5/4$ under
  $\mathrm{P}\neq\mathrm{NP}$ (Lemma 6.9, via Ko--Lee--Queyranne 1995
  and Ohsaka 2022), with the gap surviving the Lemma 5.6j manifold
  covariance restriction $\Sigma=U\Lambda U^\top+\sigma^2 I$ at high
  SNR --- a genuinely non-count gap (the log-det is the actual cost,
  not a substituted functional) that, unlike the Type-II import,
  does not dilute. This is \textsc{Partial}: with the predictor free
  and amortised ($L(M)/m\to0$) the same objective is in $\mathrm{P}$,
  the optimum attained by a decoder-reproducible sample-covariance
  maximum-likelihood plug-in (Remark 6.9a), so the fixed-field
  hardness is the sharp boundary; the budget regime at which the
  multiplicative gap appears is delimited in Remark 6.9c, and the
  $\Sigma\succeq I$ noise-floor calibration of the greedy log-det
  guarantee is shown necessary in Remark 6.9d. Corollary 6.9b records
  the operational consequence: an additive-regret Gaussian-on-manifold
  RNR encoder synthesised from Lemma 5.6j and Lemma 6.9.
\item
  The discrete, deployed counterpart of the continuous Type-III-C
  support-selection characterization above --- its instantiation for a
  transformer's discrete autoregressive field, together with the
  associated exact finite-block entropy complexity arc --- is developed
  in the companion Part [RNR-II] (the Remark 7.15a--7.15i sequence).
\item
  Additive rule-cost results for the discrete Type-III-C
  smallest-grammar problem (§13.4.1): the additive rule-size SGP
  objective inherits Charikar APX-hardness for all $\alpha$
  (ceiling-free) via a ratio-preserving reduction (Lemma 6.8k),
  matched by a poly-time $O(\log)$-approximation algorithm
  (Lemma 6.8k-approx); a $K^{\mathrm{poly}}$ one-sided upper bound on
  the joint Type-III-C cost (Lemma 6.8g, making no hardness claim);
  and a K-monotone isolation observation (Observation 6.8f) that
  retains discrete bit-cost hardness only by excluding the
  $K$-escape obstacle by fiat.
\item
  A no-universal-dominance result (Theorem 8.1) establishing the
  framework's claim as conditional, and a verification-discipline
  theorem (Theorem 9.1) ruling out cross-subblock error propagation
  in autoregressive modes.
\item
  Sufficient conditions for bit-exact decoder reproduction under
  fixed predictors (Theorem 10.1) and adaptive predictors (Theorem
  10.2 with Definition 10.1, the compatible-online-optimizer class),
  reducing reproducibility from a software-engineering question to a
  conformance specification. Published low-precision-training methods
  (integer SGD, signSGD, integer Adam, integer-quantized inference)
  are candidate components but their conformance against the full
  specification on any specific hardware target requires implementation
  and testing.
\item
  Random-access (seek-and-read) capability with cost
  $O(\log L(R) + (K + k) \cdot \mathrm{cost}(M))$ for uniform
  sync-point spacing and \emph{per-position sequential} coding
  modes (non-uniform spacing replaces the direct lookup with binary
  search at cost
  $O(\log(|X|/K) \cdot (\log |X| + \log L(R)))$ bits;
  whole-block Type-II/III-B partition-rank modes get only
  sub-block-granular access with a separate
  $O(K \cdot \mathrm{unrank}(K))$ per sub-block, Remark 10.3a)
  --- polylogarithmic in archive size $|X|$ once sync-point spacing
  $K$, read window $k$, predictor cost $\mathrm{cost}(M)$, and
  $\log L(R)$ are held fixed (or polylog in $|X|$) --- via Theorem 10.3
  with Definitions 10.2, 10.3 of $W$-bounded predictors and sync points.
  The factor $\mathrm{cost}(M)$ is a single per-position neural-predictor
  evaluation: non-universal (architecture-dependent; characterized below
  via Lemma 5.1b) and the dominant term of the per-query cost, so the
  ``polylog in $|X|$'' headline bounds the dependence on archive size
  with $\mathrm{cost}(M)$ fixed, not the per-query cost in absolute terms.
  Theorem 10.4 covers non-$W$-bounded predictors via state snapshots.
  Lemma 5.1b strengthens this under exponentially decaying
  conditional mutual information of the form
  $I(X_n; X_{1:n-W-1} \mid X_{n-W:n-1}) \leq C \rho^W$ (eq.\ 5.8;
  strictly satisfied by order-$W'$ Markov chains for any finite
  $W'$, and plausible under additional positivity assumptions for
  hidden Markov models with geometrically ergodic hidden chains).
  For these, choosing $W = O(\log N)$ brings the per-position
  sequential rate $H_W(D_N)$ within $O(1)$ of $H(D_N)$. The total
  expected length is $H(D_N) + O(1) + \Delta_{\mathrm{sync,total}} +
  \varepsilon_N N$, where the $O(1)$ is the truncated-window rate
  gap and $\Delta_{\mathrm{sync,total}} = m \cdot S \cdot W \cdot
  \log_2(1/\eta) = O((|X|/K) \log |X| \cdot \log_2(1/\eta))$ for
  $m = \lceil |X|/K \rceil$ sync points. This overhead trades off
  with the per-query random-access cost via the sync-point spacing
  $K$: at $K = \log^a |X|$ for $a > 1$ the total sync overhead is
  $O(|X| \log |X|/\log^a |X|) = O(|X|/\log^{a-1} |X|)$ (sub-linear
  in $|X|$, but not polylog; at $a = 1$ the total is $\Theta(|X|)$,
  linear); at $K = |X|/\mathrm{polylog}(|X|)$ the total overhead
  drops to $O(\mathrm{polylog}(|X|))$ but per-query cost grows to
  $\Theta(|X|/\mathrm{polylog}(|X|) \cdot \mathrm{cost}(M))$.
  Lemma 5.1c additionally identifies the \emph{tight} cold-context
  excess per sync for the special case of stationary order-$W'$ Markov
  sources with predictor context $W \geq W'$ and sub-block size
  $K \geq W'$: the per-sync excess equals exactly
  $\delta_{W'} = H(X^{W'}) - W' h(X)$ (a finite source-specific
  constant), and total sync excess over $H(D_N)$ is exactly
  $(m-1) \cdot \delta_{W'}$ when $N = m K$ (full sub-blocks) —
  strictly tightening Lemma 5.1a's worst-case
  $S \cdot W \cdot \log_2(1/\eta)$ per-sync factor whenever
  $\delta_{W'}$ is smaller (typical for natural-data Markov sources).
  The simultaneous polylog regime (polylog per-query AND polylog
  total) does not exist within this trade-off. The per-query random-access cost
  is polylog in $|X|$ provided $K = \mathrm{polylog}(|X|)$ and the
  predictor architecture admits context $W = \Theta(\log N)$ at
  $\mathrm{poly}(\log N)$ per evaluation (transformers with windowed attention pay
  $\Theta(W)$ per token; lookup-table predictors over $\Sigma^W$
  have $\Theta(\sigma^W) = N^{\Theta(1)}$ \emph{model storage} so
  are excluded by model size, even though their per-evaluation
  lookup is $O(1)$).
  The sync-point compression overhead for natural data is estimated
  in §10.7; a tight worst-case bound in general is not proved. The
  comparison to existing neural compressors (NNCP-style autoregressive
  byte models, block-decoding latent-variable schemes) on either seek
  capability or ratio is empirical and is the subject of §11.
\item
  Block-device and read-only filesystem deployment as direct
  corollaries (Theorems 10.5 and 10.6 in §10.8): an OS-level driver
  can serve LBA-indexed read requests with \emph{per-query} cost
  polylogarithmic in archive size $|X|$ in the parameter regime
  where the sync-point spacing $K$, read window $k$, predictor cost
  $\mathrm{cost}(M)$, and seek-index entry width $\log L(R)$ are
  fixed or grow at most polylogarithmically in $|X|$, and
  additionally $L(R) = \mathrm{poly}(|X|)$ (Theorem 10.5). Note
  that the driver's resident memory still includes the seek index
  itself, which is $\Theta(|X|/K)$ entries for any fixed $K$ ---
  linear in $|X|$, not polylogarithmic; for $|X| = 1$ GiB and
  $K = 64$ KiB this is 16384 entries (a few hundred KiB total),
  small in absolute terms but not asymptotically negligible. The
  predictor can
  be required to
  run bit-exactly on any hardware target satisfying the conformance
  specification of Theorem 10.1 (Theorem 10.6). Specific target
  classes named in the spec (CPU scalar, CPU SIMD including
  AVX-512, GPU integer Tensor Cores, integer NPUs) are engineering
  assertions about what targets satisfy the spec, not theorems;
  conformance for each named target requires implementation and
  testing per the conformance suite (engineering-spec document).
\item
  A runtime parameter-selection framework (Section 7.2) that subsumes
  input-dependent design choices (Code-$k$ rate, alphabet size,
  support threshold, tile size, dictionary size) under the existing
  mode-selection machinery, with the same $2\varepsilon$
  approximate-oracle gap as Theorem 6.1 plus the per-tile metadata
  overhead $O(\log K_{\text{params}})$ for transmitting the chosen
  parameter values.
\item
  The random-access and deviation-theory development of Section 7
  (Theorems 7.2--7.33, 7.35, 7.36), which bridges the abstract framework
  to a practical neural-LM compression pipeline with byte-granular random
  access, source-coding extensions, and the concentration/deviation
  hierarchy for archive-size fluctuations, is the subject of the companion
  Part [RNR-II]; the operational rate--distortion dispersion (the §7.34
  family) is developed in Part [RNR-III].
\item
  A concrete attack-vector roadmap for the remaining open problems
  (§13.4), recording for each open question both the positive
  sub-results obtained and the standard reductions that demonstrably
  fail, so as to prevent re-attempts on known-dead paths. For OP2(a)
  (unrestricted bit-cost SGP) the density-shifted Berman--Karpinski,
  Dinur--Safra/FGLSS, $K$-escape-budget, $K$-uniform/Label-Cover, and
  SGP-constant-import routes are each closed negatively (§13.4.1),
  with the additive-rule-cost variant resolved positively
  (Lemmas 6.8k, 6.8k-approx). For OP2(b) (free-predictor joint
  Type-III-C) the predictor-absorption obstacle (Lemma 13.4.2a)
  absorbs the grammar-size gap into $L(M)$ so the GSGP/MSGP route
  does not \emph{directly} transfer; but that argument gives only a
  common upper bound (not a collapse of the optimum ratio), and the
  PRG construction neutralises it, reducing amortised OP2(b)-APX to
  single-string OP2(a)-APX --- so OP2(b)-APX remains open
  (Theorem 13.4.2b).
  For OP3$''$ (efficient enumerative encoders) the unambiguous-CFL
  sub-case is closed constructively with an explicit polynomial-time
  rank/unrank encoder within $1/N$ bit/symbol of the Cover optimum
  (Lemma 13.4.3a). For OP4 (general dependent-law support selection)
  the Densest-$k$-Subhypergraph and non-count entropy/log-det routes
  are closed negatively (§13.4.4), and for the $E_{12}$ constant-gap
  locus the obstacle is identified as a missing robust-NO promise
  rather than affine dilution (Remark 4.3d). The §13.4.5 refinement
  of the TC-scaling programme adds the two new natural-data TC classes
  Lemma 5.6i (combinatorial-uniform, sublinear $\mathrm{TC}$) and
  Lemma 5.6j (Gaussian-on-manifold, log-det $\mathrm{TC}$).
\item
  Worked numerical examples for Type-I (§4.7), Type-II (§5.7), and
  Type-III-B (§6.4), demonstrating concrete coding-length arithmetic
  on toy distributions. Type-III-A and Type-III-C are illustrated
  through their composition and dictionary-construction sections
  respectively.
\end{itemize}

\textbf{1.4 Organization}

Section 2 fixes notation and reviews the necessary background. Section 3
defines the framework. Sections 4, 5, and 6 develop Type-I, Type-II,
and the three Type-III subtypes. Worked numerical examples are given
in §4.7 (Type-I), §5.7 (Type-II), and §6.4 (Type-III-B); Type-III-A
and Type-III-C are illustrated through the composition discussion in
§7 and the dictionary-construction discussion in §6.5 respectively
rather than through standalone worked examples. Section 7 opens with composition and mode selection: the composition-correctness and hierarchy-limit theorems (Theorems 7.1/7.1b/7.1c and Remark 7.1d) and the runtime parameter-selection framework (§7.2), which subsumes input-dependent design choices under the same mode-selection machinery. The remainder of Section 7 --- the byte-granular random-access architecture, the source-coding extensions, the concentration and deviation hierarchy (Theorems 7.2--7.33, 7.35, 7.36), and the operational rate--distortion dispersion (the §7.34 family) --- is developed in the companion Parts [RNR-II] (random access and deviation theory, together with §§10.7--10.8) and [RNR-III] (rate--distortion dispersion). Section 8 states the
no‑universal‑dominance result. Section 9 covers verification and error
containment. Section 10 addresses practical reproducibility and proves a
sufficient condition for bit-exact decoder reproduction. Section 11
lists experimental predictions for empirical follow‑up, including the
explicit T7.15 enwik9 prediction in §11.7. Section 12 surveys related
work, and Section 13 reports the status of open problems, distinguishing
resolved results (§13.1), remaining mathematical open questions (§13.2),
items explicitly deferred to engineering or empirical work (§13.3), and
a concrete attack-vector roadmap (§13.4) that for each remaining open
problem records the positive sub-results obtained and the standard
reductions shown to fail, so as to prevent re-attempts on known-dead
paths --- including the continuous Type-III-C support-selection
characterization (Lemma 6.9 fixed-field APX-hardness, Remark 6.9a
free-field membership in $\mathrm{P}$, Remarks 6.9c/6.9d and
Corollary 6.9b; §13.4.4), the additive rule-cost SGP results
(Lemmas 6.8k, 6.8k-approx; §13.4.1), and the constructive
unambiguous-CFL enumerative encoder (Lemma 13.4.3a; §13.4.3). The questions exposed by the §7 expansion are addressed within the companion Part [RNR-II].

\textbf{2. Preliminaries and notation}

\textbf{2.1 The lossless compression problem}

A lossless compressor is an injective map
$\mathsf{Enc}$ from a source alphabet (here, finite-length binary
strings or byte sequences of a fixed block length $N$) to binary
codewords, together with a decoder $\mathsf{Dec}$ such that
$\mathsf{Dec}(\mathsf{Enc}(X)) = X$ for every $X$. Injectivity is
necessary; for the multi-stream concatenations used later (model,
repair, verification, metadata), each stream must in addition be
prefix-free or carry an explicit length so that the decoder can
partition the codeword.

Throughout the paper, we fix the block length $N$ and treat $X$ as a
random variable over the finite alphabet $\{0,1\}^{8N}$ (byte block)
under a distribution $D$ on that finite set. We write $L(\cdot)$ for
the bit length of a binary description, and adopt the convention
that $0 \cdot \log 0 = 0$. For a random variable $X$ over a finite
alphabet, $H(X)$ denotes the Shannon entropy of $X$ in bits. Given
a side-information variable $Y$, $H(X \mid Y)$ denotes the
conditional entropy.

\textbf{Definition 2.1 (Conditionally advantageous on a data class).}
For a source $X \sim D$ and a baseline compressor $B$, an RNR coder
$\mathsf{Enc}$ is \emph{conditionally advantageous} on $D$ if
\[
\mathbb{E}[L(\mathsf{Enc}(X))] < \mathbb{E}[L(B(X))], \quad X \sim D.
\tag{2.1}
\]
\emph{Notation note.} The symbol $C$ in later sections denotes
\emph{public context} (the input to the predictor $M$, so
$Y = M(C)$), not the compressor; the symbol $D$ in later sections
denotes the source distribution, not the decoder. The
encoder/decoder pair is referred to as $\mathsf{Enc}/\mathsf{Dec}$
throughout to avoid this collision.

\textbf{2.2 Side-information coding}

Slepian and Wolf's distributed-source-coding theorem (Slepian and
Wolf 1973) characterizes the rate region for encoding jointly
distributed sources $X$ and $Y$ with separate encoders and a joint
decoder, with constraints $R_X \geq H(X \mid Y)$,
$R_Y \geq H(Y \mid X)$, and $R_X + R_Y \geq H(X, Y)$ in the
finite-alphabet memoryless setting. The corner of this region we use
is the asymmetric case in which $Y$ is already available to the
decoder, so only $X$ needs to be coded at rate
$\geq H(X \mid Y)$. The motivating intuition --- that decoder-side
information drives down the rate to a conditional entropy --- is
Slepian-Wolf's; the achievability in our setting, however, reduces
to ordinary \emph{conditional source coding} (Cover and Thomas 2006,
Theorem 5.4.1) and does not require the binning constructions that
Slepian-Wolf invokes when the encoder lacks $Y$.

Our setting differs from the classical distributed Slepian-Wolf
setting in three ways. First, we consider a single sender, not two
correlated sources at different physical locations. Second, the side
information $Y$ is not a separately observed correlated random
variable but the deterministic output of a public model $M$ applied
to a public context $C$: $Y = M(C)$. Third, this side information is
available to both encoder and decoder, so the encoder can exploit it
directly; the relevant achievability is conditional source coding,
and Slepian-Wolf binning is not needed.

We use the term ``side information'' in this broader sense
throughout: any deterministic, decoder-reproducible function of
public state that constrains the conditional distribution of $X$.

\textbf{2.3 Entropy coding}

We assume the availability of an entropy coder --- arithmetic coding
(Witten, Neal, Cleary 1987; Rissanen and Langdon 1979) or asymmetric
numeral systems (Duda 2009) --- that, given an explicit probability
model $Q$ with $Q(x) > 0$ on the support of the source and finite-
precision arithmetic of $b$ bits per state, encodes a sequence $x$ of
length $n$ in $\lceil -\log_2 Q(x) \rceil + \varepsilon_n$ bits,
where the redundancy $\varepsilon_n$ accounts for end-of-stream and
finite-precision rounding (under-2-bit total overhead for the
classical Witten-Neal-Cleary integer arithmetic coder; ANS adds an
$O(\log b)$ state-initialization overhead amortized over the stream).
Taking expectation under the true source $X \sim D$ gives expected
length $\mathbb{E}_D[-\log_2 Q(X)] + \varepsilon_n$, i.e.,
cross-entropy plus redundancy. We assume $\varepsilon_n / n \to 0$
as the block length $n$ grows.

\textbf{2.4 Enumerative coding}

Cover's enumerative source coding (Cover 1973) describes an element
of a \emph{known} finite class by its rank among the elements of that
class. The information-theoretic cost is the binary logarithm of the
class size; a self-delimiting binary rank requires
$\lceil \log_2 |\mathcal{C}| \rceil$ bits (or
$\log_2 |\mathcal{C}| + O(1)$ via arithmetic coding driven by the
uniform rank distribution). Whenever we write the cost as
$\log_2 |\mathcal{C}|$ below we mean the information-theoretic rate;
the integer-representation overhead is absorbed into the
$\varepsilon_n$ entropy-coder redundancy of §2.3. Crucially, the
class itself must be decoder-reproducible: when its description
depends on data-dependent parameters (e.g., the cardinality vector
of a partition), those parameters must be transmitted separately.
Two enumerative formulas are used repeatedly in this paper.

The number of subsets of size $k$ drawn from $\{1, \ldots, N\}$ is
$\binom{N}{k}$, so a subset is described in
$\log_2 \binom{N}{k}$ bits relative to the class
$\binom{[N]}{k}$ \emph{given $k$}. When $k$ is data-dependent it is
transmitted in $O(\log N)$ bits before the rank.

The number of partitions of $\{1, \ldots, N\}$ into 16 disjoint sets
$M_0, \ldots, M_{15}$ with cardinalities $k_0, \ldots, k_{15}$ is the
multinomial coefficient
\[
M(N; k_0, \ldots, k_{15}) = \frac{N!}{k_0!\, k_1! \cdots k_{15}!},
\tag{2.2}
\]
so such a partition is described in
$\log_2 M(N; k_0, \ldots, k_{15})$ bits given the count vector
$(k_0, \ldots, k_{15})$. The count vector itself lies in a set of
size at most $(N + 1)^{16}$ and is transmitted in
$16 \log_2(N + 1) = O(\log N)$ bits by direct enumeration. The
multinomial baseline is strictly tighter than
$\sum_v \log_2 \binom{N}{k_v}$, which corresponds to encoding 16
independent subsets without enforcing disjointness, and therefore
overcounts.

\textbf{2.5 Minimum Description Length}

Rissanen's Minimum Description Length principle (Rissanen 1978)
selects the model that minimizes the total description length: model
structure, real-valued parameters, and data given the model. There
are several technical formulations --- crude two-part MDL,
prequential MDL, universal coding, normalized maximum likelihood
(NML) --- with different optimality criteria. RNR uses MDL only
heuristically, in its simplest two-part admission form: a candidate
model or token is included if its description cost plus the resulting
conditional code length beats the alternative without it. The
admission rule
\[
L(M) + L(X \mid M) < L_{\text{base}}(X)
\]
is a two-part-MDL comparison heuristic, not a theorem of optimality;
we invoke it as a design criterion at every level (Type-I model
admission, Type-III-C token admission) but do not claim that any RNR
construction is the MDL optimum under a stronger criterion such as
NML.

\textbf{2.6 Bits-back coding}

Bits-back coding is a technique for losslessly compressing data under
a latent-variable model. The encoder uses an auxiliary bit string to
sample a latent $z$ from a variational posterior $q(z \mid x)$,
encodes $(z, x)$ jointly via prior $P(z)$ and conditional
$P(x \mid z)$, and the decoder recovers $z$ and $x$ and extracts the
auxiliary bits. The expected net length under the data distribution
$X \sim D$ and the variational sampler $z \sim q(\cdot \mid X)$ is
the \emph{negative ELBO}:
\[
\mathbb{E}_{X \sim D,\, z \sim q(\cdot \mid X)} \!\left[-\log P(X, z) + \log q(z \mid X)\right]
= \mathbb{E}_{X \sim D}\!\left[-\log P_X(X)\right]
+ \mathbb{E}_{X \sim D}\!\left[\mathrm{KL}(q(\cdot \mid X) \Vert P(\cdot \mid X))\right].
\]
This equals the source entropy $H_D(X)$ only when the model marginal
matches the data ($P_X = D$) and $q$ equals the true posterior; in
general it is cross-entropy under the model marginal plus the
variational gap. Theorem 6.6 makes this accounting explicit for the
RNR Type-III-C decomposition. The lineage is: Hinton and van Camp
(1993) introduced the bits-back argument for noisy-weight
communication in a Bayesian-network MDL setting; Frey (1997, 1998)
extended it to one-to-many stochastic source codes with discrete
latents in graphical models; Townsend, Bird and Barber (2019) gave
the practical BB-ANS construction that handles continuous latents via
ANS coding (Bit-Swap and HiLLoC are subsequent extensions). We use
the term ``bits-back'' for the general principle and BB-ANS for the
specific continuous-latent realization; see Section 6.5.

\textbf{2.7 Notation summary}

We adopt the following notation throughout.
\begin{description}
\item[$N$] block length (number of bytes per source block, fixed throughout)
\item[$n$] generic symbol-sequence length (e.g., positional, when distinct from $N$)
\item[$X$] original source block (sequence of bytes, nibbles, or symbols)
\item[$D$] source distribution on $\{0,1\}^{8N}$; $X \sim D$
\item[$Y$] decoder-reproducible side information: $Y = M(C)$
\item[$M$] deterministic predictor; $M_I, M_{\text{pos}}$ for type-specific predictors
\item[$C$] public context (seed, verified previous subblock, metadata), distinct from the §2.1 compressor symbol
\item[$Q$] predicted probability field (vector, matrix, or tensor); used as the entropy-coder model in §2.3
\item[$U_m(t)$] predicted code-length / mode-cost field for tile $t$ under mode $m$ (Type-III-A)
\item[$T$] invertible representation transform; $Z = T(X)$
\item[$R$] repair / residual / rank / selector stream
\item[$V$] verification overhead (hash, checksum); $H_v(X)$ stored as $v$
\item[$\Theta$] metadata overhead (sizes, mode flags, quantization tables)
\item[$B$] baseline compressor (LZMA, zstd, NNCP, $\ldots$); $L_{\text{base}}$ is its description length
\item[$\mathsf{Enc}, \mathsf{Dec}$] encoder/decoder pair (compressor); see §2.1
\item[$\varepsilon_n$] entropy-coder redundancy over $n$ symbols (positivity + finite-precision overhead per §2.3)
\item[$\Pi(X) = (M_0, \ldots, M_{15})$] positional partition of $X$ into symbol-class position sets
\item[$k_v$] cardinality of $M_v$, i.e., count of symbol $v$ in the block
\item[$H(\cdot), H(\cdot \mid \cdot)$] Shannon entropy / conditional entropy in bits
\item[$I(\cdot\,;\cdot \mid \cdot)$] conditional mutual information
\end{description}

\textbf{3. The RNR framework}

\textbf{3.1 The repair principle}

A neural generator alone cannot losslessly compress data, because the
generator output and the source disagree on at least some bits with
non‑negligible probability, and a single bit of disagreement is fatal.
RNR resolves this by treating the model output not as a candidate
reconstruction but as side information, and storing a separate repair
description that combined with the side information reconstructs the
exact source. The framework is parameterized by three independent
choices.

First, the predictor M. This may be a small autoregressive transformer,
a positional convolutional model, a learned token predictor, a
non‑neural context‑mixing predictor (PAQ‑style), a hash‑based predictor,
or any deterministic function of public state.

Second, the representation target. The model may predict (a) a
one‑dimensional stream of bytes or codewords, (b) a two‑dimensional
positional probability field over a partition of positions among symbol
classes, or (c) a higher‑dimensional tensor that may include the model's
own predicted uncertainty as a field. These three choices define Type‑I,
Type‑II, and Type‑III respectively.

Third, the repair encoding. Given a representation target and a model
prediction, the residual that the archive stores may be (i) an error
mask between predicted and true codewords, (ii) a rank within a
constrained enumerative class, (iii) a selector from a candidate list,
(iv) a per‑position probability under which an entropy coder operates,
(v) a tile‑wise mode flag, or any composition of these. The repair
encoding is a design choice within each type.

\textbf{3.2 The five RNR types}

\textbf{Definition 3.1 (RNR Type-I).} An RNR Type-I coder predicts a
one-dimensional damaged stream $Y = M_I(C)$ and stores a repair
stream $R_I$ together with a verification value $v = H_v(X)$ as part
of the archive. The decoder forms
$\hat{X} = \mathrm{DecRepair}(Y, R_I)$ and accepts the reconstruction
only if $H_v(\hat{X}) = v$; for a valid archive, $\hat{X} = X$
exactly. The canonical instance is Code12, in which each source byte
is mapped to a 12-bit codeword and the predicted error mask is
entropy-coded by Hamming-weight class.

\textbf{Definition 3.2 (RNR Type-II).} An RNR Type-II coder predicts
a positional marginal probability field $Q_v(i)$ approximating
$\Pr(i \in M_v \mid C)$, $v \in \{0, \ldots, 15\}$,
$i \in \{1, \ldots, N\}$, with $\sum_v Q_v(i) = 1$ for each position
$i$. The archive stores a partition residual $R_\Pi$ that, combined
with $Q$, recovers the actual partition
$\Pi(X) = (M_0, \ldots, M_{15})$, where the sets $M_v$ are pairwise
disjoint subsets of $\{1, \ldots, N\}$ and
$\sum_{v=0}^{15} |M_v| = N$. The recovered partition determines the
source $X$.

\textbf{Definition 3.3 (RNR Type-III-A).} An RNR Type-III-A coder
predicts an entropy field $U_m(t)$ over tiles $t$ and coding modes
$m$, selects predicted per-tile modes
\[
\hat{m}(t) = \min \{ m : U_m(t) = \min_{m'} U_{m'}(t) \},
\]
breaking ties by lexicographic order on a fixed canonical mode
index (specified by the model description as a total ordering of
mode identifiers; e.g., the ordering in which modes are declared in
$L(M)$). For decoder reproducibility under finite-precision
arithmetic, the $U_m(t)$ values must be computed in the same
deterministic representation on encoder and decoder (e.g., as the
bit-exact integer outputs of the same conforming inference engine
per Theorem 10.6); two encoders that produce $U_m(t)$ in
floating-point and round differently would tie-break differently
and break decoder reproducibility. The selection is then
deterministic and decoder-reproducible. The archive stores the
per-tile residual streams and, when the encoder departs from the
predicted argmin, override flags as metadata.

\textbf{Definition 3.4 (RNR Type-III-B).} An RNR Type-III-B coder
predicts a normalized byte tensor $Q_{h, l}(i)$ approximating
$\Pr(h_i = h, l_i = l \mid C)$ over high-nibble $h$ and low-nibble $l$
of byte $i$, with $h, l \in \{0, \ldots, 15\}$ and
$\sum_{h, l} Q_{h, l}(i) = 1$ for each position $i$. A
decoder-reproducible mode choice, either stored as metadata or
computed by a fixed rule from $(C, Q)$, selects one of three coding
strategies (detailed in §6.3): (i) direct 256-class partition
coding; (ii) hierarchical partition coding that first encodes the
high-nibble partition and then the low-nibble partition conditional
on the high nibble; or (iii) per-byte sequential coding, which
encodes each byte's high-then-low nibble in source order via
per-position arithmetic coding under $Q^H(\cdot \mid i, C)$ and
$Q^{L \mid H}(\cdot \mid i, h_i, C)$.

\textbf{Definition 3.5 (RNR Type-III-C).} An RNR Type-III-C coder is
parameterized by a learned token dictionary
$D = \{\tau_0, \ldots, \tau_{K-1}\}$ with $K = a \cdot b$ (typically
$32 \times 32 = 1024$ or $64 \times 64 = 4096$) hierarchically indexed
by $\text{family} \in \{0, \ldots, a-1\}$ and
$\text{variant} \in \{0, \ldots, b-1\}$ as
$\text{token\_id} = b \cdot \text{family} + \text{variant}$. The
model predicts $Q(k, i)$ approximating
$\Pr(\tau_k \text{ begins at position } i \mid C)$.
The archive stores a tokenization residual $R_C$ sufficient to
reconstruct the exact token-interval cover of the block.

\textbf{3.3 The repair length decomposition}

Every RNR coder has the same general code-length decomposition:
\[
L_{\text{RNR}}(X) = L(M) + L(R \mid Y) + L(V) + L(\Theta),
\tag{3.1}
\]
where $L(M)$ is the model description length (parameters,
quantization tables, adapter weights as relevant), $L(R \mid Y)$ is
the repair stream length, $L(V)$ is verification overhead, and
$L(\Theta)$ is metadata overhead (sizes, mode flags, dictionary
headers). The conditional advantage criterion is then
\[
\mathbb{E}_D\!\left[L(M) + L(R \mid Y) + L(V) + L(\Theta)\right] < \mathbb{E}_D[B(X)].
\tag{3.2}
\]

\textbf{3.4 Amortization of model cost}

$L(M)$ is paid once and amortized across all blocks encoded by the
same model. For a single-block archive, $L(M)$ dominates and the
framework loses unless $M$ is extremely small. For an archive of $m$
blocks of average size $n$, the effective per-block model cost is
$L(M)/m$, and the criterion (3.2) becomes more permissive in $m$.
Throughout, we assume the amortized regime; Section 10 discusses the
single-block case explicitly.

\textbf{4. Type-I: Stream Repair Coding}

\textbf{4.1 Definition and operational form}

Type-I is the one-dimensional form. The source
$X = (x_1, \ldots, x_N)$ is a stream of bytes or symbols. The model
produces a predicted stream $Y = M_I(C)$ of the same length,
considered as a damaged version of $X$. The encoder, knowing both
$X$ and $Y$, computes a repair $R_I = \mathrm{EncRepair}(X, Y)$, and
the decoder recovers $\hat{X} = \mathrm{DecRepair}(Y, R_I)$. The
encoder also stores the verification value $v = H_v(X)$ in the
archive; the decoder accepts the reconstruction only if
$H_v(\hat{X}) = v$ (Definition 3.1).

\textbf{4.2 The Code12 construction}

In the Code12 variant of Type-I, each source byte
$b \in \{0, \ldots, 255\}$ is mapped through an injective function
$E_{12} : \{0, \ldots, 255\} \to \{0,1\}^{12}$ to a 12-bit codeword
$c = E_{12}(b)$. The model predicts a codeword $\hat{c}_i = M_I(C, i)$
for each position. The raw error mask is
$e_i = c_i \oplus \hat{c}_i$. The motivation is that, with a
well-chosen $E_{12}$, the error-mask distribution concentrates on
low Hamming weights when the model is approximately correct, and
Hamming weight is cheap to entropy-code.

We specify $E_{12}$ as a shortened-Hamming-like systematic mapping. Let
the source byte bits be $(b_0, \ldots, b_7)$. Use 1-based codeword
positions $c_1, \ldots, c_{12}$ and place data bits at non-parity
positions:
\[
\begin{aligned}
c_3 &= b_0, & c_5 &= b_1, & c_6 &= b_2, & c_7 &= b_3, \\
c_9 &= b_4, & c_{10} &= b_5, & c_{11} &= b_6, & c_{12} &= b_7.
\end{aligned}
\]
The parity positions $c_1, c_2, c_4, c_8$ are computed by even parity
over selected data positions:
\[
\begin{aligned}
c_1 &= c_3 \oplus c_5 \oplus c_7 \oplus c_9 \oplus c_{11}, \\
c_2 &= c_3 \oplus c_6 \oplus c_7 \oplus c_{10} \oplus c_{11}, \\
c_4 &= c_5 \oplus c_6 \oplus c_7 \oplus c_{12}, \\
c_8 &= c_9 \oplus c_{10} \oplus c_{11} \oplus c_{12}.
\end{aligned}
\]
This construction is reversible because the data bits $b_0, \ldots, b_7$
appear directly in $c$. We use it for two reasons: first, single-bit
and many two-bit codeword errors decode to small Hamming-weight masks,
which are cheap; second, the parity structure means random predictor
noise distributes errors over both data and parity positions, rather
than concentrating in the data bits. We discuss alternative $E_{12}$
mappings (Gray-ordered, learned, bucketed with candidate lists) in
Section 4.6.

\textbf{4.3 Repair stream encoding}

The simplest repair symbol at position $i$ is the error mask $e_i \in
\{0,1\}^{12}$ itself, of which there are $2^{12} = 4096$ possibilities.
To avoid paying 12 bits per position regardless of prediction quality,
we encode by Hamming-weight class:
\begin{itemize}
\item class 0: no error (1 codeword);
\item class 1: one-bit flip (12 codewords, $\log_2 12 \approx 3.585$ bits);
\item class 2: two-bit pattern ($\binom{12}{2} = 66$ codewords, $\log_2 66 \approx 6.044$ bits);
\item class candidate-$k$: correct codeword is $k$-th in predicted candidate list;
\item class escape: store raw byte (8 bits) or raw codeword (12 bits).
\end{itemize}
Let the class probabilities be $p_0, p_1, p_2, p_{\text{cand}}, p_{\text{esc}}$
summing to 1. The expected per-position repair length is

\[
\ell_{\text{rep}} = p_0 \cdot 0 + p_1 \log_2(12) + p_2 \log_2(66) + p_{\text{cand}} \cdot (\text{header} + \text{rank}) + p_{\text{esc}} \cdot 8 + H(\text{class}) + \varepsilon,
\tag{4.1}
\]

where $H(\text{class})$ is the entropy of the class label encoded with
a small adaptive entropy coder, and $\varepsilon$ is the entropy-coder
redundancy. When $p_0$ dominates, the dominant cost is $H(\text{class})$
plus a tiny contribution from the other classes; $\ell_{\text{rep}}$
can fall well below 1 bit per byte if the model is good.

\textbf{4.4 Achievability bound}

\textbf{Theorem 4.1 (Type-I achievability).} \emph{Let $X$ be a source
block of $N$ bytes, $Y = M_I(C)$ be the predicted Code12 stream, and
$E = (e_1, \ldots, e_N)$ be the error-mask sequence with empirical
class distribution $P_E$. Suppose the class entropy coder achieves
expected length at most $N \cdot H(P_E) + \varepsilon_N$. Then the
expected Type-I code length is bounded by}
\[
\mathbb{E}[L_I(X)] \leq L(M_I) + N \cdot H(P_E) + N \cdot \ell_{\text{cond}} + \varepsilon_N + L(V) + L(\Theta),
\tag{4.2}
\]
\emph{where $\ell_{\text{cond}}$ denotes the expected conditional
payload (within-class rank or escape) per position. Consequently,
Type-I is conditionally advantageous on a distribution $D$ if the
right-hand side, in expectation under $D$, is smaller than the
expected baseline length $\mathbb{E}_D[B(X)]$.}

\emph{\textbf{Proof.}} The Type-I code is the concatenation of the
model description, the class stream, the within-class payloads, the
verification block, and the metadata block. By the entropy-coder
hypothesis, the class stream has expected length at most
$N \cdot H(P_E) + \varepsilon_N$. The within-class payload is a
separate stream whose expected length, summed across positions, is
$N \cdot \ell_{\text{cond}}$ by definition. Adding the constant
overheads $L(M_I), L(V), L(\Theta)$ gives the bound. Taking
expectation under $D$ and comparing to $\mathbb{E}_D[B(X)]$ yields
the advantage condition. \hfill$\blacksquare$

\textbf{Theorem 4.2 (Type-I converse and ideal achievability).}
\emph{Let $X \sim D$ be a source over length-$N$ byte blocks and let
$Y = M_I(C)$ be decoder-reproducible side information. Fix an
\emph{injective} Code12 map
$E_{12} : \{0, \ldots, 255\} \to \{0, 1\}^{12}$ so that, given $Y$,
the error variable $E = E_{12}(X) \oplus Y$ stands in bijection with
$X$ (the decoder recovers $X$ from $(Y, E)$ via
$X = E_{12}^{-1}(E \oplus Y)$). Then:}
\begin{itemize}
\item[\textbf{(C)}] \emph{(\emph{Converse}) Any uniquely decodable lossless
code for $X$ whose decoder uses only the codeword and $Y$ has
expected length at least $H(X \mid Y)$.}
\item[\textbf{(A)}] \emph{(\emph{Ideal achievability}) An \emph{ideal}
Type-I class-and-payload coder driven by the true conditional law
of $E$ achieves expected length $H(X \mid Y) + \varepsilon_N$, plus
the constant headers $L(M_I) + L(V) + L(\Theta)$.}
\end{itemize}
\emph{A concrete coder using an approximate class model or a uniform
within-class payload pays the corresponding cross-entropy, which
matches the bound only when those models match the true
conditional law. Non-injective $E_{12}$ would
collapse the bijection: e.g., if $E_{12}(0) = E_{12}(1) = 0$ and
$Y \equiv 0$, then $E \equiv 0$ for $X \sim \mathrm{Unif}\{0, 1\}$
while $H(X \mid Y) = 1$, so no decoder can recover $X$ from $(Y, E)$.}

\emph{\textbf{Proof.}} The converse is obtained by applying Shannon's
source-coding bound to each conditional distribution
$P(X \mid Y = y)$ and averaging over $Y$. Concretely, for a
\emph{uniquely decodable} code McMillan's theorem (Cover and Thomas
2006, Theorem 5.5.1) gives the Kraft inequality, whence the
expected-length lower bound $\mathbb{E}[L \mid Y = y] \geq
H(X \mid Y = y)$ (Cover and Thomas 2006, Theorem 5.3.1); averaging
over $Y$ yields $\mathbb{E}[L] \geq H(X \mid Y)$. (Theorem 5.4.1,
the $H \leq L^\ast < H + 1$ bound, is the achievability--gap
statement, not the converse used here; see also Chapter 15 for
coding with side information.)
For achievability, fix an injective Code12 map and define
$E = E_{12}(X) \oplus Y$. Given $Y$, the map $X \leftrightarrow E$ is
bijective, so $H(E \mid Y) = H(X \mid Y)$. Encoding first the class
$C_E$ of $E$ and then the payload within that class, both under their
true conditional probabilities given $Y$, achieves
$H(C_E \mid Y) + H(E \mid C_E, Y) = H(E \mid Y) = H(X \mid Y)$ by the
chain rule for entropy (Cover and Thomas 2006, Theorem 2.6.4). An entropy coder driven by these conditional laws attains
this rate up to arithmetic-coding redundancy $\varepsilon_N$. A
specific Code12 implementation that uses an approximate $P(C_E \mid Y)$
or a uniform within-class payload (the construction of Theorem 4.1)
pays the corresponding cross-entropy in place of the entropy; it
equals the bound only when those models match the true conditional
law. The constant headers $L(M_I) + L(V) + L(\Theta)$ are $O(1)$
independent of $N$. \hfill$\blacksquare$

\textbf{4.5 Relationship to NNCP-style direct arithmetic coding}

The NNCP-style autoregressive byte coder treats the model output as
a full probability distribution $Q_i$ over the next byte and
arithmetic-codes the byte directly under $Q_i$; this is not a
special case of Type-I as defined in this paper, since Type-I
specifically requires a deterministic damaged codeword
$Y = M_I(C)$ and a repair stream encoded via the Code12 error
variable $E = E_{12}(X) \oplus Y$. The two schemes are however
\emph{rate-equivalent} in a precise sense: given the same predictor
state, the NNCP-style direct arithmetic coder achieves expected
per-position cost equal to the categorical cross-entropy plus
$\varepsilon/N$, while the ideal class+payload Type-I coder
(Theorem 4.2 part A) driven by the true conditional law of $E$
achieves $H(X \mid Y) + \varepsilon_N$. When $Y$ encodes a
deterministic predicted codeword and the true predictive
distribution induces the same conditional $P(E \mid Y)$, both
expected rates match the same source-coding bound, so neither
construction is strictly more powerful in expected rate. Type-I as
defined generalizes the deterministic-prediction architecture by
allowing (a) non-trivial $E_{12}$ mappings whose error distribution
is more concentrated than a flat softmax, and (b) class-based
entropy coding that can be cheaper than
arithmetic coding for sharply peaked predictors.

\textbf{4.6 Variants of $E_{12}$}

Four $E_{12}$ variants are worth comparing empirically.
\begin{itemize}
\item
  $E_{12}$-A: shortened Hamming(12,8)-like systematic mapping as in
  Section 4.2.
\item
  $E_{12}$-B: Gray-ordered mapping in which numerically adjacent bytes
  differ by Hamming weight 1, suitable when model errors are
  predominantly small numerical perturbations of the true byte.
\item
  $E_{12}$-C: learned mapping in which the byte-to-codeword assignment
  is optimized to minimize the expected Hamming weight of $c \oplus
  \hat{c}$ over a training corpus and predictor.
\item
  $E_{12}$-D: bucketed mapping in which bytes that the predictor
  frequently confuses share a 4-byte bucket, and the repair stores a
  2-bit within-bucket selector.
\end{itemize}

Among these, only $E_{12}$-A and $E_{12}$-B admit closed-form
construction. $E_{12}$-C and $E_{12}$-D are learned, and we now show
that finding the optimum is computationally intractable.

\textbf{Remark 4.3a (Fixed-size Code12 search).} In the deployed
Code12 setting the alphabet size and code length are fixed constants:
$n_0 = 256$ and $m_0 = 12$. For these fixed parameters, optimizing
$E_{12}$ is a finite search over $\binom{2^{12}}{256} \cdot 256!$
injective byte-to-codeword maps, a number independent of the input
description length of the empirical confusion data. Exhaustive search
is therefore astronomically large but still a constant-size
enumeration, hence polynomial in the formal asymptotic sense. Thus
the practical Code12 case is not itself an asymptotically NP-hard
problem; NP-hardness only applies to a varying-parameter family in
which $n$ or $m$ is part of the input.

\textbf{Theorem 4.3 (NP-hardness of varying-parameter learned-code
labeling).} \emph{Consider the family of instances specified by
integers $n, m$ with $n \leq 2^m$, a symmetric row-stochastic matrix
$F \in \mathbb{Q}_{\geq 0}^{n \times n}$, and a rational threshold
$B$. Given an injective labeling
$E : \{1, \ldots, n\} \to \{0, 1\}^m$, define the unnormalized
uniform-prior Hamming loss}
\[
C_F(E) = \sum_{i,j=1}^n F_{ij} \cdot d_H(E(i), E(j)).
\]
\emph{Deciding whether there exists an injective labeling $E$ with
$C_F(E) \leq B$ is NP-hard in the description size of $(F, B, n, m)$.
Consequently, optimizing the learned code labeling is NP-hard for
varying parameters.}

\emph{\textbf{Proof.}} Reduce from tree-to-hypercube subgraph
embedding. Let $Q_m$ denote the graph on $\{0,1\}^m$ in which two
vertices are adjacent exactly when their Hamming distance is one.
Wagner and Corneil (1990) showed that deciding, given a tree $T$ and
integer $m$, whether $T$ is isomorphic to a subgraph of $Q_m$ is
NP-complete. Equivalently, the question is whether there is an
injective map $\phi : V(T) \to \{0, 1\}^m$ such that
$d_H(\phi(u), \phi(v)) = 1$ for every tree edge $\{u, v\}$.

If $|V(T)| > 2^m$, then no injective embedding into $Q_m$ exists. In
that case the reduction outputs the fixed no-instance with $n' = 2$,
$m' = 1$,
\[
F' = \begin{pmatrix} 0 & 1 \\ 1 & 0 \end{pmatrix}, \qquad B' = 1,
\]
which satisfies $n' \leq 2^{m'}$ and has no injective labeling with
$C_{F'}(E) \leq B'$ (the only injective $E : \{1, 2\} \to \{0, 1\}$
gives $C_{F'}(E) = 2 > B'$). Hence we may assume for the remaining
construction that $|V(T)| \leq 2^m$.

Given such an instance, write $n = |V(T)|$ and
$\Delta = \max_{v \in V(T)} \deg_T(v)$. For nontrivial instances
$\Delta \geq 1$. Construct an $n \times n$ matrix $F$ indexed by the
vertices of $T$ as follows:
\[
F_{uv} =
\begin{cases}
1/\Delta, & u \neq v \text{ and } \{u, v\} \in E(T), \\
1 - \deg_T(u)/\Delta, & u = v, \\
0, & \text{otherwise.}
\end{cases}
\]
This construction is polynomial in the size of $T$. The matrix is
symmetric because $T$ is undirected. It is row-stochastic because, for
each vertex $u$,
\[
F_{uu} + \sum_{v : \{u, v\} \in E(T)} F_{uv}
= 1 - \deg_T(u)/\Delta + \deg_T(u)/\Delta = 1.
\]
Thus $F$ is a valid symmetric row-stochastic confusion matrix: the
off-diagonal mass follows the tree edges, and the remaining
probability is placed on the diagonal self-confusion term.

Set
\[
B = \frac{2(n - 1)}{\Delta}.
\]
For any injective labeling $E$, diagonal terms contribute zero
($d_H(E(u), E(u)) = 0$), so
\[
C_F(E) = \frac{1}{\Delta}
\sum_{\{u, v\} \in E(T)} \!\!\bigl(d_H(E(u), E(v)) + d_H(E(v), E(u))\bigr)
= \frac{2}{\Delta} \sum_{\{u, v\} \in E(T)} d_H(E(u), E(v)).
\]
Since $E$ is injective, every tree edge has distinct endpoint labels,
hence $d_H(E(u), E(v)) \geq 1$. Because a tree on $n$ vertices has
$n - 1$ edges, every feasible labeling satisfies
\[
C_F(E) \geq \frac{2(n - 1)}{\Delta} = B.
\]
Equality holds if and only if every tree edge is mapped to Hamming
distance one. Therefore $T$ embeds as a subgraph of $Q_m$ if and only
if the constructed learned-code instance has an injective labeling
with $C_F(E) \leq B$. A polynomial-time algorithm for the learned-code
decision problem, or for exact optimization followed by comparison
with $B$, would decide the NP-complete tree-to-hypercube embedding
problem of Wagner and Corneil (1990). Hence varying-parameter
learned-code labeling is NP-hard. \hfill$\blacksquare$

\textbf{Remark 4.3b (Handling asymmetric confusion matrices).} For any
nonnegative confusion matrix $F$,
\[
\sum_{i, j} F_{ij} \cdot d_H(E(i), E(j))
= \sum_{i, j} \frac{F_{ij} + F_{ji}}{2} \cdot d_H(E(i), E(j)),
\]
because $d_H$ is symmetric. Thus an asymmetric empirical confusion
matrix induces the same labeling objective as its symmetric pair-weight
matrix $(F + F^T)/2$. However, if $F$ is row-stochastic,
$(F + F^T)/2$ need not be row-stochastic, so this identity does not by
itself extend the NP-hardness statement of Theorem 4.3 to the
asymmetric row-stochastic promise.

Consequently, $E_{12}$-C and $E_{12}$-D must be obtained by
approximation algorithms. Standard relaxations include gradient
descent on a continuous relaxation of the assignment (interpreting
$E_{12}$ as a stochastic mapping with simplex constraints), simulated
annealing, and tabu search. These approximations carry no
constant-factor guarantee in the varying-parameter regime; on the
fixed Code12 instance they are evaluated empirically against the
constant-size oracle of Remark 4.3a.

\textbf{Remark 4.3c (Wagner-Corneil reduction has an inverse-linear
inapproximability gap).}
\emph{Let $n \geq 4$ and consider Theorem 4.3 instances $(T, m)$ with
$n = |V(T)|$, $\Delta = \max_v \deg_T(v)$. The reduction gives:}

\emph{(i) Lower bound on NO-instance ratio:}
\[
\frac{\mathrm{OPT}_{\mathrm{NO}}}{\mathrm{OPT}_{\mathrm{YES}}} \geq 1 + \frac{1}{n-1}.
\]

\emph{(ii) Tight NO-family witness: for every $n \geq 4$, the star
$T = K_{1, n-1}$ with $m = n - 2$ is a NO instance attaining
$\mathrm{OPT} = B + 2/(n - 1)$ exactly, so the ratio bound (i) is
tight on this family.}

\emph{(iii) Consequence: the varying-parameter learned-code labeling
problem is NP-hard to approximate (on the minimization side) within
any multiplicative factor strictly less than $1 + 1/(n - 1)$.
Equivalently: any polynomial-time minimization approximator with
guarantee $\alpha < 1 + 1/(n - 1)$ would decide tree-to-hypercube
subgraph embedding, contradicting Wagner-Corneil NP-hardness. The
gap $1 + 1/(n - 1)$ is inverse-linear in $n$ and vanishes as
$n \to \infty$; the Theorem 4.3 reduction alone does not establish
constant-gap inapproximability.}

\emph{Proof.} (i) In the Theorem 4.3 reduction, the YES case has
$C_F(E) = B = 2(n-1)/\Delta$ exactly (each tree edge maps to $d_H = 1$).
In any NO case, Hamming distances on tree edges are positive integers
$\geq 1$ for injective labelings; at least one edge has $d_H \geq 2$
(otherwise the labeling embeds $T$ as a subgraph of $Q_m$, contradicting
NO), contributing an extra $\geq 2/\Delta$ to $C_F$. Hence
$\mathrm{OPT}_{\mathrm{NO}} \geq B + 2/\Delta = (2/\Delta) \cdot
(n-1 + 1) = (2n)/\Delta$. The ratio:
\[
\frac{\mathrm{OPT}_{\mathrm{NO}}}{\mathrm{OPT}_{\mathrm{YES}}}
\geq \frac{B + 2/\Delta}{B}
= 1 + \frac{2/\Delta}{2(n-1)/\Delta}
= 1 + \frac{1}{n-1}.
\]

(ii) For the star $K_{1, n-1}$ with $\Delta = n - 1$ and $m = n - 2$,
we have $|V(T)| = n \leq 2^{n-2}$ for $n \geq 4$ (so injective
labeling exists). The center has $m = n - 2$ codeword neighbors at
$d_H = 1$ in $Q_m$, so at most $n - 2$ of the $n - 1$ leaves can be
at $d_H = 1$ from the center; the remaining leaf must be at $d_H = 2$.
This labeling is feasible (place center and $n - 2$ neighbors at
$d_H = 1$, then the remaining leaf at any $d_H = 2$ codeword). The
$C_F$ value:
\[
\mathrm{OPT} = \frac{2}{\Delta} \bigl((n - 2) \cdot 1 + 1 \cdot 2\bigr)
= \frac{2n}{n - 1} = B + \frac{2}{n - 1} = B + \frac{2}{\Delta}.
\]
Ratio $\mathrm{OPT} / B = (2n / (n-1)) / (2 \cdot (n-1) / (n-1)) =
n / (n - 1) = 1 + 1/(n-1)$ exactly. The smallest such instance is
$K_{1, 3}$ with $m = 2$ (n=4), giving ratio $4/3$. The next is
$K_{1, 4}$ with $m = 3$ (n=5), giving ratio $5/4$. Both are verified
in \texttt{scripts/verify/remark\_4\_3c\_wagner\_corneil\_gap.py}.

(iii) follows from (i) and (ii): any minimization approximator with
guarantee strictly less than $1 + 1/(n-1)$ would output cost
$\leq \alpha \cdot \mathrm{OPT}_{\mathrm{YES}} = \alpha \cdot B
< B + 2/\Delta$, which is strictly less than $\mathrm{OPT}_{\mathrm{NO}}$,
hence the approximator's output below this threshold certifies the
YES case. As $n \to \infty$ the gap shrinks to $1$. \hfill$\blacksquare$

\emph{Implications for OP4.} A constant-gap inapproximability for
$E_{12}$ requires \emph{gap amplification} beyond the inverse-linear
ratio of Theorem 4.3. Three studied amplification mechanisms in the
linear-code / lattice / hypercube literature are candidates:
(a) Nearest Codeword Problem (NCP) inapproximability via PCP-style
reductions from Label Cover, established by
Arora, Babai, Stern, Sweedyk 1997 \emph{The Hardness of Approximate
Optima in Lattices, Codes, and Systems of Linear Equations}
(JCSS 54(2): 317--331);
(b) Closest Vector Problem (CVP) inapproximability for general
lattices, established by Dinur, Kindler, Raz, Safra 2003
\emph{Approximating CVP to within almost-polynomial factors is
NP-hard} (Combinatorica 23(2): 205--243), specialized in subsequent
work to small ambient dimensions like $\mathbb{F}_2^{12}$;
(c) Bounded-Distance Decoding (BDD) hardness in the coding-theoretic
formulation introduced by Dumer, Micciancio, Sudan 2003
\emph{Hardness of Approximating the Minimum Distance of a Linear
Code} (IEEE Trans.\ Inform.\ Theory 49(1): 22--37). DMS 2003 itself
defines BDD and leaves the $\rho = 1/2$ complexity open; subsequent
literature has established BDD hardness for various lattice
formulations.
A fourth mechanism is \emph{graph-powering / tensor-product
gap amplification} in the style of Dinur 2007 \emph{The PCP theorem
by gap amplification} (JACM 54(3): 12); if a Wagner-Corneil
tree-embedding instance could be composed with a robust PCP
verifier, the inverse-linear gap could in principle be amplified to
a constant. None of these constructions is known to realize a
gap-preserving reduction from any APX-hard problem directly into
the byte-alphabet $E_{12}$ geometry; closing OP4 remains open.

The paper makes no inapproximability claim for the varying-parameter
learned-code labeling problem beyond Theorem 4.3 + Remark 4.3c
(inverse-linear gap of $1 + 1/(n-1)$); a sound constant-gap reduction
into the Hamming-restricted geometry is left as an open problem.

\textbf{Remark 4.3d (Gap-locus characterization: the $E_{12}$
obstacle is the missing dilation-$1$/robust-NO promise gap, not affine
dilution and not row-stochasticity).}
\emph{We isolate, exactly, what a constant-gap reduction into $E_{12}$
must supply, and show that the failure of the known candidate sources
(Remark 4.3c) is not the affine-dilution / row-stochastic obstacle that
blocks the continuous-cost open problems (cf.\ the Type-II support and
Type-III-C dictionary objectives), but a separate ``no robust NO
promise'' barrier inherited from the metric geometry of partial cubes.}

\emph{(i) Exact gap identity.} For any symmetric nonnegative $F$ write
the off-diagonal mass $w = \sum_{i<j} F_{ij}$, and recall every
injective labeling sends each support edge to Hamming distance
$\geq 1$. Hence the ``all support edges at distance one'' value $2w$ is
a universal lower bound on $C_F$, attained iff the weighted support
graph $G_F$ is a \emph{dilation-$1$} (subgraph) embedding into $Q_m$.
Writing $\mathrm{OPT} = \min_E C_F(E)$, the multiplicative ratio
\[
\rho(F) \;=\; \frac{\mathrm{OPT}}{2w}
\;=\; 1 \;+\; \frac{1}{2w}\!\!\sum_{i<j} 2F_{ij}\,
\bigl(d_H(E^\star(i),E^\star(j)) - 1\bigr)
\;=\; 1 + \frac{\text{forced weighted excess}}{2w},
\]
so the inapproximability factor is governed \emph{solely} by the
fraction of weighted edge mass that every injective labeling is forced
to push beyond distance one. A constant-gap inapproximability for
$E_{12}$ is therefore equivalent to an NP-hard promise gap between
\textbf{(YES)} $G_F$ embeds at dilation $1$, cost $=2w$, and
\textbf{(NO)} every injective labeling stretches a \emph{constant
fraction} $\varepsilon_0 > 0$ of the weighted mass, cost
$\geq (1+\varepsilon_0)\,2w$.

\emph{(ii) Row-stochasticity is not the obstacle.} The ratio
$\rho(F)$ is invariant under global rescaling of $F$: replacing the
row-stochastic edge weight $1/\Delta$ by a uniform weight $1$ multiplies
both $\mathrm{OPT}$ and $2w$ by $\Delta$ and leaves $\rho$ unchanged.
Row-stochasticity only enforces the mass bound $w \leq n/2$; it neither
caps the forced excess fraction nor forces it to vanish. This is in
pointed contrast to the continuous-relaxation obstacles elsewhere in
the paper, where normalization of an \emph{affine} cost dilutes any
multiplicative source gap; here the cost is a quadratic
(assignment-type) functional of $E$, the gap survives normalization,
and the per-instance excess fraction can be a genuine constant
(e.g.\ $K_{2,4}$ at $m=3$ gives $\rho = 3/2$, and $K_{3,3}$ at $m=3$
gives $\rho = 13/9$; both are normalization-invariant, verified in
\texttt{scripts/verify/remark\_4\_3d\_e12\_constant\_gap\_dilution.py}).

\emph{(iii) Why Wagner-Corneil dilutes specifically, and why the
candidate amplifications fail.} The Wagner-Corneil family is a
\emph{tree}: it is always dilation-$1$ embeddable when $2^m \geq n$,
and a NO instance arises only from the \emph{packing} constraint that
$m$ is too small to seat all neighbors of a high-degree vertex at
distance one. The forced excess is then a single bumped edge of weight
$1/\Delta$ over a baseline of $\Theta(n/\Delta)$, giving the additive
$+2/\Delta$ and hence $\rho = 1 + 1/(n-1) \to 1$. Amplifying this gap
requires a family that is simultaneously (a) dilation-$1$ embeddable in
the YES case (so a YES instance exists at all) and (b) robustly
non-embeddable --- constant fraction stretched --- in the NO case. The
two requirements collide with two literature facts. First, the class of
dilation-$1$ embeddable graphs is exactly the \emph{partial cubes}
(isometric--subgraph case) / hypercube subgraphs, recognizable in
polynomial time (Djokovi\'c 1973; Winkler 1984; quadratic-time
algorithm of Eppstein 2008); so the embeddability predicate alone is
not NP-hard, and any hardness must come (as in Wagner-Corneil) from the
\emph{dimension/packing} side, where the slack is inherently
$O(1/n)$. Second, dense graphs that \emph{would} force a constant
stretch fraction (complete bipartite, expanders) are \emph{never}
dilation-$1$ embeddable for any $m$ --- e.g.\ $K_{2,3}$ is not a
hypercube subgraph in any dimension --- so they possess no YES instance
and yield a ratio between two stretched configurations rather than a
promise gap from the $2w$ baseline. This is the same degeneracy that
prevents the no-constant-approximation status of Minimum Linear
Arrangement (Ambühl-Mastrolilli-Svensson 2011, no PTAS under
subexponential-ETH) from being a clean constant multiplicative-gap APX
statement, and it explains why the lattice/coding hardness of NCP and
CVP (Arora-Babai-Stern-Sweedyk 1997; Dinur-Kindler-Raz-Safra 2003),
though constant- and superconstant-hard, does not transfer: those
problems minimize distance from a fixed received word to a linear
\emph{code} (a coset), not a pairwise quadratic sum over an
\emph{injective labeling} into $Q_m$, and no gap-preserving map between
the two geometries is known.

\emph{(iv) Verdict (NEGATIVE, with a sharpened locus).} The
constant-gap inapproximability of $E_{12}$ (OP4) remains open, and the
obstruction is now localized: it is \emph{neither} affine dilution
\emph{nor} the row-stochastic promise, but the absence of a
gap-amplified hypercube-embedding source exhibiting a YES dilation-$1$
instance together with a constant-fraction robust-NO instance. By
(i)--(ii) any such source immediately yields the constant gap; by (iii)
the natural candidates fail for structurally distinct reasons
(partial-cube tractability on the embeddable side, absence of a YES
baseline on the dense side). Constructing a robust gap version of the
hypercube subgraph embedding problem --- or a gap-preserving reduction
from Label Cover that lands in the injective-into-$Q_m$ geometry while
preserving a dilation-$1$ YES case --- is the precise remaining step.

\textbf{4.7 Worked numerical example}

Consider a block of $N = 4096$ bytes drawn from a source where the
Code12 predictor achieves the following empirical class distribution:
$p_0 = 0.60$ (no error), $p_1 = 0.25$ (one-bit flip), $p_2 = 0.10$
(two-bit pattern), $p_{\text{cand}} = 0.04$ (within candidate list of
size 4), $p_{\text{esc}} = 0.01$ (escape to raw byte).

The class entropy is
\[
H(P_E) = -0.60 \log_2 0.60 - 0.25 \log_2 0.25 - 0.10 \log_2 0.10
- 0.04 \log_2 0.04 - 0.01 \log_2 0.01 \approx 1.527 \text{ bits.}
\]
The expected within-class payload per position is
\[
\ell_{\text{cond}} = 0.25 \log_2 12 + 0.10 \log_2 66 + 0.04 \log_2 4
+ 0.01 \cdot 8 \approx 0.896 + 0.604 + 0.080 + 0.080 = 1.660 \text{ bits.}
\]
Per-position total is therefore $1.527 + 1.660 \approx 3.187$ bits,
or $3.187 \cdot 4096 \approx 13055$ bits $\approx 1632$ bytes for
the repair stream. Adding $L(V) = 32$ bytes (a SHA-256 per 4 KiB
subblock) and $L(\Theta) = 16$ bytes of metadata gives approximately
$1680$ bytes plus the amortized model cost. For comparison, a 4096-byte block at 8 bits per byte raw
is 4096 bytes; LZMA on structured text typically achieves 1500--2000
bytes. The Type-I coder is competitive whenever the amortized $L(M_I)$
per block is below approximately 100--300 bytes, which is achievable
for a small distilled byte-level predictor.

\textbf{5. Type‑II: Positional Symbol Coding}

\textbf{5.1 From sequential to positional prediction}

Type‑I and the classical neural lossless coders share an operational
question: predict the next symbol. Type‑II changes the unit of modeling.
Instead of asking ``what is the next symbol?'', Type‑II asks: ``where
does each symbol appear in the block?'' The block is represented as a
partition of position indices among symbol classes, and the model
predicts the geometry of that partition.

This change is useful when a block has positional, geometric,
matrix‑like, periodic, or otherwise non‑local structure that is not
visible as local dictionary repetition. Examples include raw image data,
Bayer‑mosaic sensor data, database pages with structured record offsets,
instruction streams with positional opcode constraints, and numeric
matrices.

\textbf{5.2 Nibble partition representation}

Without loss of generality, we present Type-II at the nibble alphabet
$A_{16} = \{0, \ldots, 15\}$. Larger alphabets (full bytes, byte pairs,
learned tokens) generalize directly and are revisited in Section 6.
Let $X = (x_1, \ldots, x_N)$ with $x_i \in A_{16}$. For each nibble
value $v$, define the position set
\[
M_v = \{\, i \in \{1, \ldots, N\} : x_i = v \,\}.
\]
The collection $\Pi(X) = (M_0, \ldots, M_{15})$ is a partition of
$\{1, \ldots, N\}$: the $M_v$ are pairwise disjoint and their union is
the full position set. The original sequence $X$ is recovered from
$\Pi(X)$ by reading off the position memberships in order. Encoding
$X$ is therefore equivalent to encoding $\Pi(X)$.

\textbf{5.3 Factorized and semi-non-factorized positional models}

A Type-II model predicts a probability field $Q \in [0,1]^{16 \times N}$
with column normalization $\sum_v Q_v(i) = 1$. We distinguish two
regimes that affect the framework's relationship to Type-I.

In the factorized regime, $Q_v(i)$ is treated as if the position labels
were independent across $i$: the joint probability assigned to a
candidate sequence $(x_1, \ldots, x_N)$ is the product over $i$ of
$Q_{x_i}(i)$. Under this regime, optimal entropy coding achieves
length equal to $-\sum_i \log_2 Q_{x_i}(i)$, which is bit-equivalent
to per-position arithmetic coding driven by the same $Q$. Type-II
under a factorized model is therefore not operationally distinct from
Type-I with a per-position predictor. We make this explicit in
Lemma 5.1.

In the semi-non-factorized regime, the model additionally enforces
structural constraints that the factorized form violates. The minimal
such constraints are: (a) the disjoint-partition constraint that
exactly one $v$ has $x_i = v$ at each $i$; and (b) the count
constraint that the partition cardinalities $k_v = |M_v|$ are
themselves modeled (either predicted, transmitted, or constrained by a
learned distribution over count vectors). Under the
semi-non-factorized regime, the joint distribution over partitions is
not a product of per-position categoricals, and the coding gain over
the factorized regime can be substantial when the counts or the joint
geometry are highly informative. Theorem 5.4 quantifies this
separation.

\textbf{5.4 Coding strategies for the partition residual}

Given $Q$ and possible knowledge of counts $k_v$, the archive must
store a residual $R_\Pi$ sufficient to recover the exact partition.
Several coding strategies are available.
\begin{itemize}
\item
  Per-position arithmetic coding under $Q$: encode $x_i$ with
  probability $Q_{x_i}(i)$. Cost per position: $-\log_2 Q_{x_i}(i)$.
\item
  Per-symbol enumerative coding given the counts: for each $v$ in
  some order, encode $M_v$ as a subset of the remaining positions by
  its rank among $\binom{N'}{k_v}$ possibilities, where $N'$ is the
  number of positions not yet assigned. The final mask is implicit.
\item
  Per-symbol enumerative coding given a predicted support $S_v$
  containing $M_v$ of size $s_v$: encode $M_v$ by its rank among
  $\binom{s_v}{k_v}$ possibilities, with an escape mode for positions
  in $M_v \setminus S_v$.
\item
  Joint enumerative coding under the multinomial: encode the partition
  by its rank among the $M(N; k_0, \ldots, k_{15})$ valid partitions
  with the given counts.
\end{itemize}

\textbf{5.5 Baseline, lemma, and the corrected gain theorem}

A natural but incorrect baseline for the partition residual is the
per-symbol enumerative sum $\sum_v \log_2 \binom{N}{k_v}$. This
overcounts: the 16 subset choices it sums over are not constrained to
be disjoint, so the same position may be assigned to multiple $v$ in
the enumeration, and the resulting count exceeds the true number of
valid partitions. The correct baseline is the multinomial, which we
state in Theorem 5.2 below; the gain theorem in Theorem 5.3 is
expressed against it.

\textbf{Lemma 5.1.} \emph{Under the factorized regime, per-position
arithmetic coding under $Q$ and per-position Type-I arithmetic coding
under the equivalent sequential conditional $Q$ are bit-equivalent:
both produce code lengths within $\varepsilon_N$ of
$-\sum_i \log_2 Q_{x_i}(i)$.}

\emph{\textbf{Proof.}} Both coders are arithmetic coders driven by the
same per-position probability vector $Q_*(i)$. Arithmetic coding
redundancy is a constant overhead independent of which axis of the
data the coder is sweeping. The two code lengths therefore agree up to
$\varepsilon_N$. \hfill$\blacksquare$

\textbf{Lemma 5.1a (Per-position sequential recoding achieves the
$W$-Markov rate plus sync overhead, with byte-granular random
access).} \emph{Let $X \sim D_N$ be a block of $N$ symbols over a
$\sigma$-symbol alphabet, and let $M$ be a $W$-bounded predictor
(Definition 10.2) that at each position $i$ outputs the order-$W$
Markov chain-rule conditional distribution
$P(x_i \mid x_{i-W}, \dots, x_{i-1}, C)$ (necessarily depending on
at most the last $W$ bytes, by Definition 10.2). Without sync
points, per-position sequential arithmetic coding of $X$ under
these conditionals achieves expected total length}
\[
H_W(D_N) \;:=\; \sum_{i=1}^N H(X_i \mid X_{\max(1, i-W)}, \dots, X_{i-1}, C)
\;\;\leq\;\; N \log_2 \sigma,
\]
\emph{plus $\varepsilon_N \cdot N$ entropy-coder redundancy. With
sync points at spacing $K$ (Definition 10.3) inserted to enable
byte-granular random access, the coding additionally pays a
sync-overhead term $\Delta_{\mathrm{sync, total}} \leq m \cdot S
\cdot W \cdot \log_2(1/\eta)$ bits (the per-sync-point cold-context
excess analyzed in §10.7, with $S$ symbols per source byte for the
mode, $\eta$ the probability floor, and $m$ the sync-point count
as in Definition 10.3). The total expected length is therefore
$H_W(D_N) + \Delta_{\mathrm{sync, total}} + \varepsilon_N \cdot N$.
In general $H_W(D_N) \geq H(D_N)$ with equality iff $D_N$ is
order-$W$ Markov given $C$. Combined with sync points, this coding
admits byte-granular random access via Theorem 10.3 at
$O(\log L(R) + (K+k) \cdot \mathrm{cost}(M))$ per query.}

\emph{\textbf{Proof.}} Without sync points, the order-$W$ Markov
chain rule gives
$H_W(D_N) = \sum_i H(X_i \mid X_{i-W:i-1}, C)$, where the
conditional is over the at-most-$W$-bytes window. Arithmetic
coding under these conditionals achieves
$\sum_i [-\log_2 P(x_i \mid x_{i-W:i-1}, C)]$ on any realization,
which in expectation under $D_N$ equals
$H_W(D_N) + \varepsilon_{\mathrm{ac}}^{\mathrm{total}}$ where
$\varepsilon_{\mathrm{ac}}^{\mathrm{total}}$ is the
arithmetic-coder total redundancy
(at most $2$ bits per stream for canonical implementations, by
Cover \& Thomas 2006 Theorem 5.4.1; not to be confused with $\varepsilon_n$ in
§2.3 which is a separate per-symbol redundancy notation). The inequality $H_W(D_N) \geq H(D_N)$ follows from
conditioning on a strict subset of the past not increasing entropy
beyond conditioning on the full past. With sync points, each
sync-point boundary resets the predictor context (Definition 10.3
condition (a)), forcing the first $W$ positions after each sync
point to use a shortened/cold context; this contributes the
sync-overhead bound $\Delta_{\mathrm{sync, total}} \leq m \cdot S
\cdot W \cdot \log_2(1/\eta)$ derived in §10.7 (per-sync excess
bounded by $W \cdot S \cdot \log_2(1/\eta)$ times the sync count
$m$). By the $W$-bounded property, the coding is per-position
sequential in the sense of Theorem 10.3's scope; applying Theorem
10.3 directly gives byte-granular random access at the stated
per-query cost. \hfill$\blacksquare$

\emph{Consequence for Theorem 5.4.} Theorem 5.4's $\Theta(N)$
amortized separation between factorized (marginal-only) and joint
coding of the linear-code family was stated using whole-block
partition-rank coding, which gives only sub-block-granular random
access (Remark 10.3a). Lemma 5.1a shows that per-position sequential
coding under an order-$W$ Markov predictor achieves rate
$H_W(D_N)$, which is intermediate between the factorized rate
$\sum_i H_{\mathrm{marg}}(X_i)$ (Markov order 0) and the full joint
entropy $H(D_N)$ (Markov order $N-1$). For data classes that are
$W$-Markov given $C$, the full joint rate $H(D_N)$ is matched by the
W-bounded rate $H_W(D_N)$ (since $H_W = H$ for $W$-Markov sources)
and is achieved by per-position sequential coding with byte-granular
random access \emph{plus} the sync-point overhead
$\Delta_{\mathrm{sync,total}} \leq m \cdot S \cdot W \cdot \log_2(1/\eta)
= \lceil N/K \rceil \cdot S \cdot W \cdot \log_2(1/\eta)$ (worst-case
upper bound; the realized cold-context excess may be smaller), which is
$o(N)$ only when sync-point spacing $K = \omega(W)$ strictly (for
fixed $K, W, S, \eta$ the total is $\Theta(N)$; see the asymptotic
trade-off discussion accompanying Lemma 5.1b). For
data classes whose joint structure is not order-$W$ Markov, the
W-bounded rate $H_W(D_N)$ exceeds $H(D_N)$, and the
random-access tax is the multi-information $\mathrm{TC}_W(D_N) :=
H_W(D_N) - H(D_N)$ beyond order $W$. Choosing a larger $W$ reduces
this gap (at the cost of higher per-position predictor evaluation
cost). The Theorem 5.4 linear-code family is generally not order-$W$
Markov for small $W$; recovering its full $\Theta(N)$ separation with
byte-granular access requires either $W \geq \Theta(N)$ (eliminating
the polylog advantage) or sub-block restructuring (Remark 10.3a).

\textbf{Lemma 5.1b (Logarithmic $W$ suffices under exponential
conditional-MI decay).} \emph{Let $X = (X_n)_{n \geq 1}$ be a
strictly stationary process on a finite alphabet $\mathcal{X}$.
Suppose the conditional mutual information between $X_n$ and its
far past, given the recent window $X_{n-W:n-1}$, decays
exponentially in the window size: there exist constants $C > 0$ and
$\rho \in (0, 1)$ such that for every $n \geq W + 2$ and every
$W \geq 1$,}
\[
I\bigl(X_n ; X_{1:n-W-1} \,\big|\, X_{n-W:n-1}\bigr) \leq C \cdot \rho^W.
\tag{5.8}
\]
\emph{(For $n \leq W + 1$, the conditioning window already contains
the full past, so the conditional MI is identically zero by
convention.) Then for any $N \geq 1$ and any $W \geq 1$,}
\[
0 \leq H_W(D_N) - H(D_N) \leq \max(N - W - 1, 0) \cdot C \cdot \rho^W \leq N \cdot C \cdot \rho^W.
\]
\emph{In particular, choosing
$W_N := \max\bigl(1, \lceil \log_{1/\rho}(N \cdot C) \rceil\bigr)$
(clamped at $W \geq 1$) gives $H_{W_N}(D_N) - H(D_N) \leq 1$, and
the per-position sequential arithmetic-coding scheme of Lemma 5.1a
achieves total expected length}
\[
H(D_N) + O(1) + \Delta_{\mathrm{sync,total}}(W_N) + \varepsilon_N \cdot N,
\]
\emph{where $\Delta_{\mathrm{sync,total}}(W_N) = m \cdot S \cdot W_N
\cdot \log_2(1/\eta) = O((N/K) \cdot S \cdot \log N \cdot \log_2(1/\eta))$
scales with the sync-point count $m = \lceil N/K \rceil$. The total
overhead is sub-linear in $N$ only when sync-point spacing
$K = \omega(\log N)$ strictly: for $K = \log^a N$ with $a > 1$ the
overhead is $O(N/\log^{a-1} N)$ (sub-linear but \emph{not} polylog);
for $K = \Theta(\log N)$ exactly the overhead is $\Theta(N
\log_2(1/\eta))$, linear in $N$; polylog total overhead would
require $K = N/\mathrm{polylog}(N)$, which is incompatible with
polylog per-query random-access cost. For fixed $K = O(1)$ the
overhead is $\Theta(N \log N \log_2(1/\eta))$, linear in $N$ with
a log factor.
The per-position rate gap $H_{W_N}(D_N) - H(D_N) \leq 1$ from the
truncated window is itself $O(1)$ in $N$, independent of $K$.
The byte-granular random-access cost from Theorem 10.3 becomes
$O(\log L(R) + (K + k) \cdot \mathrm{cost}(M_{W_N}))$ per query. The per-position
predictor-evaluation cost $\mathrm{cost}(M_{W_N})$ depends on the
architecture: a transformer with context $W_N$ requires
$\Omega(W_N)$ attention operations per token (so $\Omega(\log N)$ at
$W_N = \Theta(\log N)$); an explicit conditional-probability lookup
table over $\Sigma^W$ contexts has $\Theta(\sigma^{W_N}) =
\Theta(N^{\log_2 \sigma / \log_2(1/\rho)})$ \emph{model storage} but
$O(1)$ per-evaluation lookup time (hashing by context);
a fixed-architecture predictor with structural context $W$
(Definition 10.2) requires changing the architecture to accommodate
$W = W_N$, not a runtime parameter scan.}

\emph{\textbf{Proof.}} For any $n \geq W + 2$ and $W \geq 1$,
chain-rule expansion gives
\[
H(X_n \mid X_{n-W:n-1}) - H(X_n \mid X_{1:n-1})
= I(X_n; X_{1:n-W-1} \mid X_{n-W:n-1}) \leq C \cdot \rho^W
\]
by the hypothesis (5.8). For $n \leq W + 1$, the window
$X_{n-W:n-1}$ already covers the full past $X_{1:n-1}$ (since
$\max(1, n - W) = 1$ when $n \leq W + 1$), so the two conditional
entropies coincide and the excess is exactly zero. Summing:
\[
H_W(D_N) - H(D_N) = \sum_{n = W + 2}^N \bigl[H(X_n \mid X_{n-W:n-1}) - H(X_n \mid X_{1:n-1})\bigr]
\leq (N - W - 1)_+ \cdot C \cdot \rho^W
\leq N \cdot C \cdot \rho^W.
\]
Setting $W = W_N = \max(1, \lceil \log_{1/\rho}(N \cdot C) \rceil)$,
we have $\rho^{W_N} \leq 1/(N C)$ when $NC \geq 1$, so
$H_{W_N}(D_N) - H(D_N) \leq 1$. (When $NC < 1$, the bound
$N \cdot C \cdot \rho \leq NC < 1$ holds already for $W = 1$.) The
total-length bound follows from Lemma 5.1a with this $W_N$, and the
random-access cost from Theorem 10.3.
\hfill$\blacksquare$

\emph{Implications.} For natural-data sources satisfying the
exponential conditional-MI decay (5.8), Lemma 5.1b shows that
byte-granular random access is achievable with rate $H(D_N) + O(1)$
excess (relative to the per-position-sequential lower bound) plus
sync-point overhead $\Delta_{\mathrm{sync,total}} = O((N/K) \log N
\log_2(1/\eta))$, which scales with the sync-point count
$m = N/K$. The hypothesis (5.8) is
\emph{strictly satisfied} by order-$W'$ Markov chains for any
finite $W'$ (where the conditional MI is identically zero for
$W \geq W'$, so any $\rho \in (0, 1)$ works) and by the convolution
of order-$W'$ Markov chains. The hypothesis is \emph{plausible} for
hidden Markov models with geometrically ergodic hidden chains and
positive-density emission distributions, by adapting filter-stability
results of Cappe, Moulines, and Ryden (\emph{Inference in Hidden
Markov Models}, Springer 2005, \S 4.3 ``Forgetting of the Initial
Condition'') from total-variation forgetting to entropy / conditional
MI under additional positivity assumptions; the formal derivation
is non-trivial and we do not include it here.

\emph{Terminology note.} The decay (5.8) is a conditional
mutual-information condition, distinct from the classical
$\beta$-mixing (or absolute regularity) coefficient that compares
the joint past-future sigma-algebra against the product law in total
variation. The two conditions are related but not equivalent: a
deterministic period-$2$ alternating chain $X_n = 1 - X_{n-1}$ has
$I(X_n ; X_{1:n-W-1} \mid X_{n-W:n-1}) = 0$ for any $W \geq 1$
(every position is a deterministic function of the previous one) but
is not $\beta$-mixing (Bradley \emph{Basic Properties of Strong
Mixing Conditions}, Probability Surveys 2 (2005), 107--144, DOI
10.1214/154957805100000104, Corollary 3.6: a Markov chain is
$\beta$-mixing iff Harris-recurrent and aperiodic).
Lemma 5.1b only requires (5.8).

\textbf{Lemma 5.1c (Tight cold-context excess for order-$W'$ Markov
with $W' \leq W$).} \emph{Let $X = (X_n)_{n \geq 1}$ be a strictly
stationary order-$W'$ Markov process on a finite alphabet, and let
$M$ be a $W$-bounded predictor (Definition 10.2) with $W \geq W'$
that outputs the true conditional $P(X_i \mid X_{i-W:i-1})$ at each
position. In the per-position sequential coding scheme of Lemma 5.1a
on sub-blocks of size $K$ starting at sync points (Definition 10.3),
the expected cold-context excess per sync point equals exactly}
\[
\Delta_{\mathrm{sync, per-pt}}
\;=\; \delta_{W'}
\;:=\; \sum_{i=1}^{W'} \bigl[ H(X_i \mid X_1, \ldots, X_{i-1}) - h(X) \bigr]
\;=\; H(X^{W'}) - W' \cdot h(X)
\;\geq\; 0,
\tag{5.9}
\]
\emph{a finite source-specific constant independent of $W$ (provided
$W \geq W'$) and independent of sub-block size $K$ (provided
$K \geq W'$). For order-1 Markov, $\delta_{W'} = \delta_1 =
H_{\mathrm{marg}} - h(X) = E$. The
total sync \emph{excess over $H(D_N)$} in the byte-granular
random-access scheme with $m$ sync points and full sub-blocks of
size $K$ (so $N = m K$) is exactly}
\[
L_{\mathrm{synced}} - H(D_N) \;=\; (m - 1) \cdot \delta_{W'},
\]
\emph{the additional cost paid by the $m-1$ sync resets after the
initial natural-start sub-block, which strictly tightens Lemma 5.1a's
worst-case bound $m \cdot S \cdot W \cdot \log_2(1/\eta)$ whenever
the source-specific constant $\delta_{W'}$ is less than
$S \cdot W \cdot \log_2(1/\eta)$. The latter is typical for
natural-data Markov sources: e.g.\ for an order-$1$ Markov source on
the 27-letter English alphabet with $E \approx 0.71$ bits/letter
(Shannon 1951 bigram estimate) and $W = O(\log N) = 20$, $S = 1$,
$\eta = 2^{-12}$, the worst-case Lemma 5.1a per-sync bound is
$S \cdot W \cdot \log_2(1/\eta) = 240$ bits/sync versus Lemma 5.1c's
tight value $\delta_1 \approx 0.71$ bits/sync, a $\sim 340\times$
improvement.}

\emph{\textbf{Proof.}} For a stationary order-$W'$ Markov source with
$W' \leq W$ and the $W$-bounded predictor outputting the true
conditional, the predictor's output at position $i > W'$ within a
sub-block (where the $W$-byte window already includes the relevant
$W'$ context bytes) is $P(X_i \mid X_{i-W':i-1})$ by the
order-$W'$ Markov property. By stationarity, the conditional entropy
$H(X_i \mid X_{i-W':i-1}) = h(X)$ is the entropy rate. Thus the
warm region ($i > W'$ within a sub-block of size $K \geq W'$)
contributes exactly $h(X)$ per position. The cold region
($i \leq W'$, where the in-block context is shorter than $W'$
bytes) contributes $\sum_{i=1}^{W'} H(X_i \mid X_{<i}) = H(X^{W'})$
by the chain rule and stationarity. The expected sub-block length is
\[
\sum_{i=1}^{K} \mathbb{E}_X[-\log_2 P(X_i \mid X_{i-W:i-1})]
\;=\; H(X^{W'}) + (K - W') \cdot h(X)
\;=\; K \cdot h(X) + \delta_{W'},
\]
matching $H(D_K)$ on the full sub-block by chain-rule expansion.
Summing over $m$ sub-blocks (with $N = m K$):
$L_{\mathrm{synced}} = m \cdot (K \cdot h(X) + \delta_{W'}) =
N \cdot h(X) + m \cdot \delta_{W'}$. Compare to the unsynced full-block
joint entropy: $H(D_N) = N \cdot h(X) + \delta_{W'}$ (one $\delta_{W'}$
correction at the natural start of the stream). The excess paid for
sync points is therefore
$L_{\mathrm{synced}} - H(D_N) = (m - 1) \cdot \delta_{W'}$, since the
first sub-block's cold-context cost is shared with the natural
stream start. \hfill$\blacksquare$

\emph{Scope.} Lemma 5.1c applies to order-$W'$ Markov sources with
predictor context bound $W \geq W'$. For general stationary mixing
sources that are not finite-order Markov, the predictor's warm rate is
$h_W := H(X_{W+1} \mid X_1^W) \geq h(X)$ with equality only in the
order-$\leq W$ Markov case. Outside the Markov hypothesis, the
sub-block length decomposes as
$\sum_{i=1}^{W} H(X_i \mid X_{<i}) + (K-W) \cdot h_W$ instead of
$K \cdot h(X) + \delta_W$, and the cold-context excess relative to
$h_W$-based steady-state is
$\sum_{i=1}^{W} [H(X_i \mid X_{<i}) - h_W]$, a different (and harder
to characterize) quantity. Lemma 5.1b's choice $W = O(\log N)$ makes
$(h_W - h) \to 0$ via the conditional-MI decay (5.8), so the difference
to the entropy rate $h(X)$ contributes only $O(1)$ in total over the
block.

\emph{History.} An earlier draft contained a Lemma 5.1c claiming
$\Delta_{\mathrm{sync, per-pt}} = \delta_W$ for \emph{arbitrary}
stationary processes, with $\delta_W$ bounded by the
Crutchfield--Feldman ``excess entropy'' under exponential mixing. That
broader claim was retracted after adversarial review pointed out the
$h$/$h_W$ conflation for non-finite-order-Markov sources. The current
Lemma 5.1c restricts hypothesis to order-$W'$ Markov ($W' \leq W$),
where the claim holds exactly. Lemma 5.1d below extends the
characterization to the broader class of exponentially mixing
stationary sources, with the Crutchfield--Feldman excess entropy
$\delta_\infty$ replacing the Markov-stabilized $\delta_{W'}$.

\textbf{Lemma 5.1d (Cold-context excess for exponentially mixing
stationary sources).} \emph{Let $X = (X_n)_{n \geq 1}$ be a strictly
stationary process on a finite alphabet satisfying the
exponential conditional-MI decay (5.8) of Lemma 5.1b with constants
$C, \rho$, and let $h_W := H(X_{W+1} \mid X_1, \ldots, X_W)$ be the
$W$-bounded conditional rate. In the per-position sequential coding
scheme of Lemma 5.1a on sub-blocks of size $K \geq W$ starting at
sync points, the expected sub-block length satisfies}
\[
\mathbb{E}[L_{\mathrm{sub}}(K)]
\;=\; K \cdot h_W + \delta_W^{(h_W)}
\;\leq\; K \cdot h_W + \delta_\infty,
\]
\emph{where $\delta_W^{(h_W)} := \sum_{i=1}^W [H(X_i \mid X_{<i}) -
h_W] \geq 0$ is the cold-context excess relative to $h_W$, and
$\delta_\infty := \sum_{n=1}^\infty [H(X_n \mid X_{<n}) - h(X)]$ is
the Crutchfield--Feldman excess entropy (finite under
summable-correlation mixing, c.f.\ Lemma 5.6c (b)). With $m =
\lceil N/K \rceil$ sub-blocks and $W = W_N := \max(1, \lceil
\log_{1/\rho}(N C) \rceil)$ as in Lemma 5.1b (clamped at $W \geq 1$),
the total expected length satisfies}
\[
L_{\mathrm{synced}} \;\leq\; N \cdot h_W + m \cdot \delta_\infty,
\quad\text{with}\quad
L_{\mathrm{synced}} - H(D_N) \leq N \cdot (h_W - h) + (m-1) \cdot \delta_\infty + 1.
\]
\emph{Since $h_W - h \leq C \rho^W = O(1/N)$ at the chosen $W_N$, the
first term is $O(1)$, and the total excess over $H(D_N)$ is
$(m-1) \cdot \delta_\infty + O(1)$.}

\emph{\textbf{Proof.}} The sub-block consists of a cold region
($i \in [1, W]$, where the predictor's $W$-byte window is partially
populated) and a warm region ($i \in [W+1, K]$, fully populated).
The cold region contributes $\sum_{i=1}^W H(X_i \mid X_{<i})$ via
arithmetic coding on the true in-block conditional. The warm region
contributes $(K - W) \cdot h_W$, since for $i > W$, the conditioning
window is $X_{i-W:i-1}$, a $W$-byte block, and by stationarity
$\mathbb{E}[-\log_2 P(X_i \mid X_{i-W:i-1})] = h_W$. Adding:
\[
\mathbb{E}[L_{\mathrm{sub}}(K)]
= \sum_{i=1}^W H(X_i \mid X_{<i}) + (K-W) \cdot h_W
= K \cdot h_W + \underbrace{\sum_{i=1}^W [H(X_i \mid X_{<i}) - h_W]}_{=: \delta_W^{(h_W)}}.
\]
For the $\delta_W^{(h_W)} \leq \delta_\infty$ bound: by stationarity,
$H(X_i \mid X_{<i})$ for $i \leq W$ equals $H(X_W \mid X_1^{W-1+i-W})$,
which conditions on $i-1 \leq W-1$ symbols of recent past, a strict
sub-context of the $W$-byte window used to define $h_W = H(X_{W+1}
\mid X_1^W)$ (by stationarity, $H(X_{W+1} \mid X_1^W) =
H(X_W \mid X_0^{W-1})$, a $W$-symbol conditioning). Conditioning on
fewer past symbols gives weakly higher entropy, so $H(X_i \mid X_{<i})
\geq h_W$ for $i \leq W$; each summand is non-negative.
Moreover, $H(X_i \mid X_{<i}) - h_W \leq H(X_i \mid X_{<i}) - h(X)$
since $h_W \geq h(X)$. Summing: $\delta_W^{(h_W)} \leq \sum_{i=1}^W
[H(X_i \mid X_{<i}) - h(X)] = \delta_W \leq \delta_\infty$ (by the
summability hypothesis, Lemma 5.6c (b)). Multiplying by $m$ and comparing to the finite identity
$H(D_N) = N \cdot h(X) + \delta_N$ from Lemma 5.6c, with
$\delta_N \to \delta_\infty$ as $N \to \infty$ under summable
correlation: the excess is $N \cdot (h_W - h) +
m \cdot \delta_W^{(h_W)} - \delta_\infty + o(1)$. With $W = W_N$:
$h_W - h \leq C \rho^{W_N} \leq C / (NC) = 1/N$, so $N \cdot (h_W -
h) \leq 1$. The total excess simplifies to $(m-1) \cdot \delta_\infty
+ O(1)$.
\hfill$\blacksquare$

\emph{Scope.} Lemma 5.1d extends Lemma 5.1c from order-$W'$ Markov
sources (where the excess is exactly $\delta_{W'}$, a finite constant
in the source) to exponentially mixing stationary sources (where the
excess is at most $\delta_\infty$, a generally finite but possibly
larger constant). Lemma 5.1c is the special case $\delta_\infty =
\delta_{W'}$ when $X$ is order-$W'$ Markov.

\textbf{Lemma 5.1e (Matching information-theoretic lower bound for
sub-block-based byte-granular random access).} \emph{Let $X$ be a
stationary source on a finite alphabet, and consider any byte-granular
random-access coding scheme that (i) encodes the archive as $m$
independently-decodable sub-blocks of size $K$ (so $N = mK$); (ii)
each sub-block's payload is a uniquely decodable codeword for the
realized sub-block data alone, without reference to other sub-blocks'
data; (iii) any shared metadata is source-INDEPENDENT (model constants,
fixed structural parameters, but \emph{not} a source-dependent global
header that could amortize sub-block content). Then}
\[
L_{\mathrm{synced}} \;\geq\; m \cdot H(D_K) \;=\; m \cdot H(X_1, \ldots, X_K),
\]
\emph{matched by Lemma 5.1c's construction at $L_{\mathrm{synced}} =
m \cdot (K \cdot h(X) + \delta_{W'})$ for order-$W'$ Markov $X$ with
$K \geq W'$. In this case the upper and lower bounds coincide:
$L_{\mathrm{synced}} = m \cdot H(D_K)$ exactly, and the total excess
$L_{\mathrm{synced}} - H(D_N) = m \cdot \delta_{W'} - \delta_{W'} =
(m-1) \cdot \delta_{W'}$ is information-theoretically tight.}

\emph{\textbf{Proof.}} Conditions (i)-(iii) force each sub-block to
be encoded as a uniquely decodable codeword for its realized data
under a fixed (source-independent) per-sub-block coding scheme. By
Shannon's converse to source coding (Cover \& Thomas 2006 Theorem
5.4.1, applied to the per-sub-block uniquely-decodable codes), the
expected length of any such encoding of a single sub-block of size
$K$ drawn from $D_K$ is at least the joint entropy $H(D_K) =
H(X_1, \ldots, X_K)$. With $m$ sub-blocks (which are correlated under
the global stationary law but encoded independently per (ii)):
$L_{\mathrm{synced}} \geq m \cdot H(D_K)$. Substituting the chain-rule
expansion $H(D_K) = K \cdot h(X) + \delta_W$ for order-$W'$ Markov
$X$ with $K \geq W' = W$ (Lemma 5.6c case (a) on the sub-block law),
$L_{\mathrm{synced}} \geq m \cdot K \cdot h(X) + m \cdot \delta_{W'}
= N \cdot h(X) + m \cdot \delta_{W'}$. The full-block joint entropy
$H(D_N) = N \cdot h(X) + \delta_{W'}$ also from Lemma 5.6c. Excess
$L_{\mathrm{synced}} - H(D_N) \geq (m - 1) \cdot \delta_{W'}$.
\hfill$\blacksquare$

\emph{Source-dependent shared header excluded.} Condition (iii) is
essential. Consider a degenerate stationary source $X_n = Y$ for all
$n$, where $Y$ is a single uniformly distributed byte ($h(X) = 0$,
$H(D_K) = 8$ bits for any $K$). A source-dependent global header
storing $Y$ in $8$ bits, plus empty sub-block payloads, supports
byte-granular access in $L_{\mathrm{synced}} = 8$ bits, violating the
naive $m \cdot H(D_K) = 8m$ lower bound. The lemma's source-independent
metadata restriction excludes such cross-block amortization: each
sub-block must independently encode its realized data, no global
covariate-extraction trick allowed.

\emph{Implication.} Lemma 5.1c's construction achieves
$(m-1) \cdot \delta_{W'}$ excess by per-position sequential coding
within each sub-block. Lemma 5.1e shows this is
\emph{information-theoretically optimal}: no other byte-granular
random-access scheme on the same sub-block partition can do strictly
better. The $(m-1) \cdot \delta_{W'}$ excess is the unavoidable
``cost of byte-granular random access'' for order-$W'$ Markov sources
under the sub-block partition. Schemes that relax sub-block
independence (e.g., share state across sync points via Theorem 10.4
state snapshots) operate outside this framework and are not bounded
by this lemma.

\textbf{Corollary 5.1f (Open Problem 1 resolution within sub-block
paradigm).} \emph{Combining Lemmas 5.1a, 5.1b, 5.1c, 5.1d, 5.1e: for
strictly stationary finite-alphabet sources satisfying exponential
conditional-MI decay (5.8), under the sub-block-based byte-granular
random-access framework (source-independent metadata, independently
decodable sub-blocks per Lemma 5.1e (ii)-(iii)):}

\emph{(a) Achievability.} \emph{The per-position sequential
arithmetic-coding scheme of Lemma 5.1a, with window choice
$W = W_N = O(\log N)$ from Lemma 5.1b and sub-block size $K \geq W$,
achieves total expected length}
\[
L_{\mathrm{synced}}^{\mathrm{achieve}} \;\leq\; H(D_N) + (m-1) \cdot \delta_{W} + N \cdot (h_W - h(X)) + \varepsilon_{\mathrm{ac}}^{\mathrm{total}},
\]
\emph{with the $(h_W - h)$ term bounded by $O(N^{-c})$ at $W = W_N$
(Lemma 5.1d), so the leading-order excess over $H(D_N)$ is}
$(m-1) \cdot \delta_W$.

\emph{(b) Information-theoretic optimality.} \emph{Lemma 5.1e
establishes that any sub-block-based byte-granular random-access
scheme (under conditions (i)-(iii) of Lemma 5.1e) satisfies}
$L_{\mathrm{synced}} \geq m \cdot H(D_K) = N \cdot h(X) + m \cdot \delta_K$
\emph{(by Shannon's source coding theorem per sub-block plus the
finite identity $H(D_K) = K h(X) + \delta_K$ of Lemma 5.6c). For
order-$W'$ Markov sources with $W' \leq W$ and $K \geq W'$, the
lower bound becomes $H(D_N) + (m-1) \cdot \delta_{W'}$ exactly
(matching (a) ignoring the $\varepsilon_{\mathrm{ac}}^{\mathrm{total}}$
coder redundancy, since $h_W = h$ for order-$W'$ Markov). For general
exp-mixing sources satisfying (5.8), the achievability-vs-lower-bound
gap is}
\[
K (h_W - h) + \delta_W^{(h_W)} - \delta_K
= \sum_{i=W+1}^K I(X_i; X_{1:i-W-1} \mid X_{i-W:i-1}) \;\geq\; 0,
\]
\emph{(by the chain-rule expansion of conditional entropies), bounded
by $K \cdot C \rho^W$ via (5.8); over $m$ sub-blocks the total gap is
$\leq N \cdot C \rho^W = O(1)$ for $W = W_N = O(\log N)$. So the
achievability of (a) and the Lemma 5.1e lower bound match up to
$O(1)$ additive slack uniformly across all sources satisfying (5.8) ---
a clean weak-matching characterization.}

\emph{(c) Open extensions.} \emph{(i) Non-sub-block schemes:
state-snapshot frameworks (Theorem 10.4) store predictor state at
sync points. Such state is source-dependent (a function of prior
source bytes), so the scheme violates Lemma 5.1e condition (ii) ---
the sub-block payload is no longer decoded from its own data alone
but from $(\mathrm{state}, \mathrm{payload})$ together. Lemma 5.1e
itself does not literally apply to state-snapshot schemes (its
hypotheses exclude source-dependent state); a more general
charged-object converse — for each sub-block $j$,
$\mathbb{E}[\ell(\mathrm{state}_j, \mathrm{payload}_j)] \geq H(D_K)$
since the sub-block data is a deterministic function of
$(\mathrm{state}_j, \mathrm{payload}_j)$ — gives the analogous bound
that the total charged archive length (state plus payload, summed
over sync points) is at least $m \cdot H(D_K)$. The payload alone
can be smaller than $H(D_K)$ because the state carries cross-block
context. The engineering
trade-off is between payload excess $(m-1) \cdot \delta_W$ (with
state size zero) and combined state-plus-payload allocation
(reducing payload but adding $m \cdot L(\mathrm{state})$ bits of
state storage). The optimum depends on the architecture's
$L(\mathrm{state})$ (e.g., transformer KV cache, RNN hidden state)
and is engineering-determined per use case.
(ii) Non-Markov stationary sources with non-summable correlations:
Lemma 5.1d's bound requires summable correlation; sources with
$\delta_\infty = \infty$ require a different framework.}

\textbf{Theorem 5.2 (Multinomial baseline for unconstrained partition
coding).} \emph{Let $N$ be a block length, and let
$(k_0, \ldots, k_{15})$ be the symbol counts with $\sum_v k_v = N$.
Without side information, the number of valid 16-way disjoint
partitions of $\{1, \ldots, N\}$ with these counts is
$M(N; k_0, \ldots, k_{15})$. The enumerative code length for the
partition is therefore $\log_2 M(N; k_0, \ldots, k_{15})$ bits, plus
the cost of describing the count vector itself.}

\emph{\textbf{Proof.}} A 16-way partition of $\{1, \ldots, N\}$ with
cardinalities $(k_0, \ldots, k_{15})$ is in bijection with a length-$N$
labeling of the positions in which label $v$ occurs $k_v$ times. The
number of such labelings is the multinomial coefficient $N!$ divided
by the product over $v$ of $k_v!$ by direct combinatorial counting.
Enumerative coding describes the partition by its rank in this class.
\hfill$\blacksquare$

\textbf{Theorem 5.3 (Type-II support-set advantage with escape).}
\emph{Let $\Pi(X) = (M_0, \ldots, M_{15})$ be the true partition with
counts $k_v$ and predicted support sets $S_v$ of size $s_v$.
Encode symbols in fixed order $v = 0, 1, \ldots, 15$. Let
$U_v = \bigcup_{u < v} M_u$ be the positions already claimed,
$A_v = \{1, \ldots, N\} \setminus U_v$ the remaining universe
($|A_v| = N - \sum_{u<v} k_u$), $S'_v = S_v \cap A_v$ the support
restricted to $A_v$ with $s'_v = |S'_v|$, and
$\alpha_v = |M_v \setminus S'_v|$ the escape count.
After transmitting the count vector $(k_v)$ and the escape-count
vector $(\alpha_v)$, encode $M_v \cap S'_v$ by its rank among
$\binom{s'_v}{k_v - \alpha_v}$ subsets of $S'_v$ of that size, and
$M_v \setminus S'_v$ by its rank among
$\binom{|A_v| - s'_v}{\alpha_v}$ subsets of $A_v \setminus S'_v$ of
that size. The total partition code length is at most}

\[
L_{R\Pi} \leq \sum_v \left[\log_2 \binom{s'_v}{k_v - \alpha_v} + \log_2 \binom{|A_v| - s'_v}{\alpha_v}\right] + L(\text{counts}) + L(\alpha) + \varepsilon_N,
\tag{5.1}
\]
\emph{where one may take
$L(\alpha) \leq \sum_v \lceil \log_2 (k_v + 1) \rceil + O(1)$ bits to
describe the escape-count vector explicitly.}

\emph{Comparing against the multinomial baseline
$\log_2 M(N; k_0, \ldots, k_{15})$ of Theorem 5.2, the positional
gain is}

\[
G_{\text{pos}} = \log_2 M(N; k_0, \ldots, k_{15}) - \sum_v \left[\log_2 \binom{s'_v}{k_v - \alpha_v} + \log_2 \binom{|A_v| - s'_v}{\alpha_v}\right].
\tag{5.2}
\]

\emph{Type-II is conditionally advantageous when $G_{\text{pos}}$
exceeds $L(M_{\text{pos}}) + L(V) + L(\Theta) + L(\alpha) +
\varepsilon_N$ plus the cost of describing the count vector and any
per-symbol support headers.}

\emph{\textbf{Proof.}} The decoder reconstructs $M_0, M_1, \ldots$ in
order. At step $v$ the available universe is $A_v$ (positions not
yet claimed by $M_0, \ldots, M_{v-1}$), so the newly decoded set is
automatically disjoint from all previous $M_u$. The transmitted
$\alpha_v$ specifies the two subset sizes; the two enumerative ranks
then identify $M_v \cap S'_v$ within $S'_v$ and $M_v \setminus S'_v$
within $A_v \setminus S'_v$ uniquely. Summing the two log-binomials
across $v$ and adding the count, escape-count, and arithmetic-coder
overheads gives (5.1). The gain identity (5.2) follows by subtracting
from the multinomial baseline. \hfill$\blacksquare$

Theorem 5.3 differs from a natural but flawed formulation in two
respects. First, it uses the multinomial coefficient as the baseline,
which is the correct entropy of an unconstrained partition with
known counts; the looser $\sum_v \log_2 \binom{N}{k_v}$ baseline
overstates the gain by allowing inconsistent subset choices. Second,
it admits the realistic case where $M_v$ is not entirely contained
in $S_v$ by introducing the per-symbol escape count $\alpha_v$,
which carries its own enumerative cost. The bound becomes
particularly clean when $\alpha_v = 0$ for all $v$: the second
log-binomial vanishes and the formula reduces to
$\sum_v \log_2 \binom{s'_v}{k_v}$ (with the chained $s'_v = |S_v \cap A_v|$,
not the raw $s_v$), recovering the idealized oracle bound on the
remaining universe at each step.

\textbf{5.5a The Type-II converse}

Theorems 5.2 and 5.3 are achievability statements: they upper-bound the
length of explicit partition codes. We now supply the matching
converse, completing the Type-II proof matrix in the same shape as the
Type-I converse (Theorem 4.2(C)). It is important to separate three
nested notions, because the multinomial / support cost is a converse for
only the narrowest of them.
\begin{itemize}
\item[(i)] \emph{Any Type-II code} (the partition $\Pi$, in bijection
with $X$, coded by any uniquely decodable scheme using $Y$).
\item[(ii)] \emph{Count-plus-rank form}: first describe the count vector
$\mathbf{k}$, then identify $\Pi$ within its count class by an
enumerative index (``rank''), the index being coded by \emph{any}
uniquely decodable scheme.
\item[(iii)] \emph{Flat-enumerative form} (the literal Theorem 5.2 /
Theorem 5.3 construction): the count-plus-rank form in which the rank is
coded \emph{uniformly}, at the fixed length $\log_2$ of the number of
candidate sets --- i.e.\ enumerative coding with no entropy coding of the
index.
\end{itemize}
The result has two layers. The first is the clean Shannon bound
$\mathbb{E}[L_{\mathrm{II}}] \geq H(X \mid Y)$ for \emph{every} Type-II
code (i); it is tight, because an ideal joint coder of the partition
reaches it with no overhead, and the count-plus-rank form (ii) also
meets it with no excess (by entropy-coding the index). The structural
content is the second layer, scoped to the \emph{flat-enumerative} form
(iii): there the flat cost is the log of the decoder-available
candidate-set size $\log_2 |\mathcal{C}(Y,\mathbf{k})|$ (the multinomial
$\log_2 M(N;\mathbf{k})$ when no support is used, and
$\sum_v \log_2 \binom{s'_v}{k_v}$ with a correct support), and it equals
the within-class entropy $H(\Pi \mid Y, \mathbf{k})$ --- hence the
information-theoretic optimum $H(X \mid Y)$ once the count cost is added
--- exactly when the conditional partition law is uniform on its
candidate set $\mathcal{C}(Y,\mathbf{k})$. In that regime Theorem 5.3 is
entropy-optimal; otherwise its flatness leaves a strictly positive gap,
the within-candidate-set uniform deficit, recovered by
\emph{entropy-coding the same enumerative index} (staying in
count-plus-rank form (ii)) or, equivalently in rate, by ideal joint
coding. We are explicit throughout about which of (i)--(iii) a given
claim addresses, so the flat-enumerative overhead is never misattributed
to Type-II, nor to the count-plus-rank form, as such.

\textbf{Theorem 5.3b (Type-II converse).} \emph{Let $X \sim D$ be a
source over length-$N$ nibble blocks with partition representation
$\Pi(X) = (M_0, \ldots, M_{15})$ (Definition 3.2; $\Pi$ is in bijection
with $X$), and let $Y = M(C)$ be decoder-reproducible side information.
Write $\mathbf{k} = (k_0, \ldots, k_{15})$, $k_v = |M_v|$, for the
count vector.}
\begin{itemize}
\item[\textbf{(C)}] (\emph{Base converse; all Type-II codes.}) Any
uniquely decodable Type-II code whose decoder uses only the codeword and
$Y$ has expected length
\[
\mathbb{E}[L_{\mathrm{II}}(X)] \;\geq\; H(X \mid Y).
\]
This bound is tight: an ideal Type-II coder that arithmetic-codes
$\Pi$ under its true conditional law $P(\Pi \mid Y)$ attains
$H(X \mid Y) + \varepsilon_N$, with no count, support, or escape
overhead. Hence the multinomial / support cost of Theorems 5.2--5.3 is
\emph{not} a converse against arbitrary Type-II codes.
\item[\textbf{(S)}] (\emph{Structural converse and the flat-enumerative
overhead.}) For the \emph{count-plus-rank form} (ii) --- describe
$\mathbf{k}$, then identify $\Pi$ within its count class by any uniquely
decodable index code (the Theorem 5.2 / Theorem 5.3 chained scheme,
support-restricted and escape variants included) --- the converse is
again the base bound:
\[
\mathbb{E}[L_{\mathrm{II}}^{\mathrm{cr}}(X)] \;\geq\; H(\mathbf{k} \mid Y) + H(\Pi \mid Y, \mathbf{k}) \;=\; H(X \mid Y),
\]
so the count-plus-rank form never beats $H(X \mid Y)$ and, by
entropy-coding the index, meets it (achievability of part T). No
multinomial / support \emph{overhead} is forced by the count-plus-rank
structure itself. The overhead appears only in the \emph{flat-enumerative
form} (iii). One must account honestly for the support metadata. Let $Z$
collect all \emph{source-dependent} side information the decoder uses to
form the candidate set --- the declared supports $S_v$ and escape counts
$\alpha_v$ when these depend on the block --- transmitted in the header;
purely source-\emph{independent} structure (fixed by $Y$ and the model)
contributes nothing to $Z$. Let $\mathcal{C}(Y, \mathbf{k}, Z)$ be the
resulting candidate set (the count-$\mathbf{k}$ partitions consistent with
$Y$ and $Z$), so flat enumeration codes the index uniformly at the fixed
length $\log_2 |\mathcal{C}(Y, \mathbf{k}, Z)|$, and the header charges
$Z$ at expected cost $\geq H(Z \mid Y, \mathbf{k})$. Provided the support
is correct ($\Pi \in \mathcal{C}(Y, \mathbf{k}, Z)$ always, which the
escape mechanism of Theorem 5.3 guarantees for the realized block), the
flat payload decomposes \emph{conditioned on $Z$}:
\[
\underbrace{\mathbb{E}\!\left[\log_2 |\mathcal{C}(Y, \mathbf{k}, Z)|\right]}_{\text{flat rank payload}} = \underbrace{H(\Pi \mid Y, \mathbf{k}, Z)}_{\text{conditional index entropy}} + G_{\mathrm{within}}^{Z},
\qquad
G_{\mathrm{within}}^{Z} = \mathbb{E}\!\left[D\!\left(P_{\Pi \mid Y, \mathbf{k}, Z} \,\big\|\, U_{\mathcal{C}(Y, \mathbf{k}, Z)}\right)\right] \;\geq\; 0,
\]
where $U_{\mathcal{C}}$ is the uniform law on the candidate set and
$G_{\mathrm{within}}^{Z}$ is the \emph{within-candidate-set uniform
deficit}. Crucially, the rank payload alone can be made small (even $0$)
by pushing the partition information into $Z$, but then the header pays
for it: since $Z$ is a deterministic function of $\Pi$ (it is computed
from the block), the chain rule gives total
$\geq H(\mathbf{k} \mid Y) + H(Z \mid Y, \mathbf{k}) + H(\Pi \mid Y, \mathbf{k}, Z) = H(\mathbf{k}, Z, \Pi \mid Y) = H(\Pi \mid Y) = H(X \mid Y)$,
recovering the base converse no matter how the support is chosen. The
\emph{support-free} special case takes $Z$ empty,
$\mathcal{C}(Y, \mathbf{k}) = \{$all partitions with counts $\mathbf{k}\}$,
$|\mathcal{C}| = M(N; \mathbf{k})$ (the multinomial of Theorem 5.2), and
$G_{\mathrm{within}}^{\varnothing} = \mathbb{E}[D(P_{\Pi \mid Y, \mathbf{k}} \| U_{\mathbf{k}})]$.
A correct \emph{source-independent} support (one fixed by $Y$, so still
$Z$ empty) shrinks $\mathcal{C}$ and lowers both the rank payload and the
deficit for free; a source-\emph{dependent} support shrinks the rank
payload but adds $H(Z \mid Y, \mathbf{k})$ to the header, with the two
trading off so that the total never drops below $H(X \mid Y)$. The flat
scheme is entropy-optimal ($G_{\mathrm{within}}^{Z} = 0$) exactly when the
conditional law is uniform on its candidate set
$\mathcal{C}(Y, \mathbf{k}, Z)$. Thus the flat scheme's total cost is
$H(X \mid Y) + G_{\mathrm{within}}^{Z}$ (with the count and metadata
headers charged at their entropies and $\varepsilon_N$); its only excess
over the optimum is $G_{\mathrm{within}}^{Z}$, the price of \emph{uniform}
indexing on the candidate set, not of the count-plus-rank decomposition
--- it is removed by entropy-coding the index, without leaving form (ii).
\item[\textbf{(T)}] (\emph{Tightness / order-optimality.}) The
flat-enumerative scheme (iii) is entropy-optimal ---
$G_{\mathrm{within}}^{Z} = 0$, flat rank payload
$= H(\Pi \mid Y, \mathbf{k}, Z)$ --- exactly when the conditional law is
uniform on its candidate set $\mathcal{C}(Y, \mathbf{k}, Z)$. Two clean
cases realize this with zero total excess. (a) \emph{No support} ($Z$
empty), with the conditional partition law uniform on the whole count
class, where $|\mathcal{C}| = M(N; \mathbf{k})$ and
\[
\mathbb{E}_{Y, \mathbf{k}}\!\left[\log_2 M(N; \mathbf{k})\right] = H(\Pi \mid Y, \mathbf{k}),
\qquad H(\mathbf{k} \mid Y) + \mathbb{E}_{Y, \mathbf{k}}\!\left[\log_2 M(N; \mathbf{k})\right] = H(X \mid Y).
\]
(b) A correct \emph{source-independent} support (one fixed by $Y$, so $Z$
is still empty) that restricts $\mathcal{C}$ to a subclass on which the
conditional law is uniform, where $\log_2 |\mathcal{C}| < \log_2 M$ but
still $= H(\Pi \mid Y, \mathbf{k})$, again with no header cost for the
support. A source-\emph{dependent} support reduces the rank payload but
charges $H(Z \mid Y, \mathbf{k})$ in the header; the total is still
$H(X \mid Y) + G_{\mathrm{within}}^{Z}$, entropy-optimal iff the law is
uniform on $\mathcal{C}(Y, \mathbf{k}, Z)$. In all entropy-optimal cases
the flat Theorem 5.3 code meets the converse up to the arithmetic-coder
redundancy $\varepsilon_N$ and the header cost. Otherwise
($G_{\mathrm{within}}^{Z} > 0$), the flat scheme overpays by exactly
$G_{\mathrm{within}}^{Z}$, but the count-plus-rank form recovers the gap
by entropy-coding the index under $P(\Pi \mid Y, \mathbf{k}, Z)$, costing
$H(\Pi \mid Y, \mathbf{k}, Z)$ in expectation (Shannon achievability), so
with the headers $\mathbb{E}[L^{\mathrm{cr}}] = H(\mathbf{k} \mid Y) + H(Z \mid Y, \mathbf{k}) + H(\Pi \mid Y, \mathbf{k}, Z) + \varepsilon_N = H(X \mid Y) + \varepsilon_N$.
Thus Type-II is order-optimal in general (part C), entropy-optimal via
the flat scheme exactly on the uniform-candidate-set regime, and
entropy-optimal always within the count-plus-rank form by entropy-coding
the index; the flat scheme's $G_{\mathrm{within}}^{Z}$ excess is the price
of \emph{uniform} indexing on the candidate set rather than an
unavoidable cost of the partition representation.
\end{itemize}

\emph{Sharpness of the scope (a counterexample to over-claiming).} The
restriction of the $\log_2 M(N;\mathbf{k})$ lower bound to the
\emph{flat-enumerative} form (iii), rather than to the broader
count-plus-rank form (ii), is necessary. Take $N = 4$, no side
information, alphabet $\{0,1\}$, and a deterministic count vector
$\mathbf{k} = (1,3)$, so $H(\mathbf{k}) = 0$ and the count class has
$M(4;\mathbf{k}) = \binom{4}{1} = 4$ partitions (indexed by the position
of the single $0$). Let the conditional law over the four enumerative
ranks be $(\tfrac12, \tfrac14, \tfrac18, \tfrac18)$, so
$H(\Pi \mid \mathbf{k}) = 1.75 < 2 = \log_2 M$ and (support-free)
$G_{\mathrm{within}}^{\varnothing} = 0.25$. The count-plus-rank code that
sends an empty count prefix and then the Huffman index code
$\{0 \mapsto 0,\, 1 \mapsto 10,\, 2 \mapsto 110,\, 3 \mapsto 111\}$ is
prefix-free, decodes by unranking, and has expected length
$\tfrac12\cdot1 + \tfrac14\cdot2 + \tfrac18\cdot3 + \tfrac18\cdot3 = 1.75$
bits $= H(X)$. It stays strictly within the count-plus-rank form yet
beats the flat multinomial payload $\log_2 M = 2$ by exactly
$G_{\mathrm{within}}^{\varnothing}$. Hence $\log_2 M(N;\mathbf{k})$ is
\emph{not} a converse for form (ii); the structural lower bound (part S)
is the entropy $H(\Pi \mid Y, \mathbf{k})$, and the $\log_2 M$ value is a
converse only for the support-free flat scheme.

A second variant shows that even the \emph{flat} cost is the candidate-set
size $\log_2 |\mathcal{C}|$, not $\log_2 M(N;\mathbf{k})$, once a
\emph{source-independent} support is used. Keep $N = 4$,
$\mathbf{k} = (1,3)$, but let the single $0$ be uniform on positions
$\{1,2\}$ only: $P(0111) = P(1011) = \tfrac12$. Relative to the full count
class $H(\Pi \mid \mathbf{k}) = 1 < 2 = \log_2 M$, a deficit of $1$. Yet a
flat code with the fixed support $S_0 = \{1,2\}$ --- the same set for every
block, so $Z$ is empty and no support header is paid --- has $\mathcal{C}$
of size $\binom{2}{1} = 2$ and pays $\log_2 2 = 1 = H(\Pi \mid \mathbf{k})$,
entropy-optimal: it does \emph{not} overpay by the full-class deficit.
This is why part (S) measures the deficit against the uniform law on the
decoder-available candidate set $\mathcal{C}$, with $\log_2 M$ recovered
only in the support-free special case.

A third variant shows why the support metadata $Z$ must be charged. Same
source, but now let the code declare a \emph{source-dependent} support
$S_0 = \{$the actual zero position$\}$. For every realized block the
support is correct and $|\mathcal{C}(Y, \mathbf{k}, Z)| = 1$, so the flat
\emph{rank} payload is $\log_2 1 = 0$ --- seemingly beating even the
entropy. But $Z$ now identifies the zero's position, so
$H(Z \mid \mathbf{k}) = 1$ bit must be transmitted in the header, and the
total is $0 + 1 = 1 = H(X)$ again. The apparent free lunch is illusory:
the partition information has merely moved from the rank into $Z$. This is
why part (S) conditions the candidate set and the deficit
$G_{\mathrm{within}}^{Z}$ on $Z$ and charges $H(Z \mid Y, \mathbf{k})$
separately, exactly the source-dependent-header accounting that Lemma
5.1e enforces in the sub-block setting; with $Z$ a deterministic function
of $\Pi$, the total can never drop below $H(X \mid Y)$.

\emph{\textbf{Proof.}} \textbf{(C)} Because $\Pi(X)$ is in bijection
with $X$ (§5.2: ``encoding $X$ is equivalent to encoding $\Pi(X)$''),
any Type-II code is a uniquely decodable code for $X$ whose decoder
sees only the codeword and $Y$. Shannon's source-coding converse
applied to each conditional law $P(X \mid Y = y)$ and averaged over $Y$
(Cover and Thomas 2006, Theorem 5.4.1; Chapter 15 for the
side-information form) gives $\mathbb{E}[L_{\mathrm{II}}] \geq H(X \mid Y)$,
exactly as in Theorem 4.2(C). For tightness, code $\Pi$ directly under
$P(\Pi \mid Y)$ with an arithmetic coder: since $\Pi \leftrightarrow X$,
$H(\Pi \mid Y) = H(X \mid Y)$, and the coder attains
$H(\Pi \mid Y) + \varepsilon_N$ (Cover and Thomas 2006, Theorem 5.4.1).
No count vector, support, or escape stream appears, so the overhead of
Theorems 5.2--5.3 is avoidable; it cannot be a converse against all
Type-II codes.

\textbf{(S)} A count-plus-rank code splits as a self-delimiting
description of $\mathbf{k}$ followed by an identification of $\Pi$ within
its count class (the chained per-symbol enumerative ranks of Theorem 5.3
specify $\Pi$ given $\mathbf{k}$ and the supports, all decoder-available;
$\mathbf{k}$ is logically prior since the candidate counts $\binom{\cdot}{k_v}$
require it). By linearity of expectation the expected length is
$\mathbb{E}[|c(\mathbf{k})|] + \mathbb{E}[|r(\Pi;\mathbf{k},Y)|]$. The
prefix $c$ is a uniquely decodable code for $\mathbf{k}$ given $Y$, so
$\mathbb{E}[|c|] \geq H(\mathbf{k} \mid Y)$ (Shannon converse on
$\mathbf{k}$); conditioned on $(Y, \mathbf{k})$ the remainder $r$ is a
uniquely decodable code for $\Pi$ over the count class, so
$\mathbb{E}[|r|] \geq H(\Pi \mid Y, \mathbf{k})$ (Shannon converse on
$\Pi$ within the class). Summing and using the chain rule
$H(\Pi \mid Y) = H(\mathbf{k} \mid Y) + H(\Pi \mid Y, \mathbf{k})$
(valid because $\mathbf{k}$ is a deterministic function of $\Pi$), and
$H(\Pi \mid Y) = H(X \mid Y)$ from the bijection,
\[
\mathbb{E}[L_{\mathrm{II}}^{\mathrm{cr}}] \;\geq\; H(\mathbf{k} \mid Y) + H(\Pi \mid Y, \mathbf{k}) \;=\; H(X \mid Y).
\]
This is the converse for form (ii); the only structural lower bound on
the index payload is the within-class \emph{entropy}
$H(\Pi \mid Y, \mathbf{k})$, not $\log_2 M(N;\mathbf{k})$.

The flat payload identity is separate, and must charge the support
metadata. The flat-enumerative form (iii) transmits, beyond
$\mathbf{k}$, any source-dependent support metadata $Z$ (the declared
$S_v, \alpha_v$ when block-dependent), then codes the index uniformly at
the fixed length $\log_2 |\mathcal{C}(y, \mathbf{k}, z)|$, where
$\mathcal{C}(y, \mathbf{k}, z)$ is the count-$\mathbf{k}$ partition set
consistent with $(y, z)$ (for the chained Theorem 5.3 scheme,
$|\mathcal{C}| = \prod_v \binom{s'_v}{k_v - \alpha_v}\binom{|A_v| - s'_v}{\alpha_v}$,
and $= M(N; \mathbf{k})$ when $Z$ is empty). The header pays $Z$ at
expected cost $\geq H(Z \mid Y, \mathbf{k})$ (Shannon converse on $Z$).
If the support is correct,
$\mathrm{supp}\,P_{\Pi \mid Y=y, \mathbf{k}, z} \subseteq \mathcal{C}(y, \mathbf{k}, z)$,
so by Cover and Thomas 2006 Theorem 2.6.4 the conditional entropy is at
most $\log_2 |\mathcal{C}(y, \mathbf{k}, z)|$, with deficit
\[
G_{\mathrm{within}}^{Z} = \mathbb{E}\!\left[\log_2 |\mathcal{C}(Y,\mathbf{k},Z)| - H(\Pi\mid Y,\mathbf{k},Z)\right] = \mathbb{E}\!\left[D\!\left(P_{\Pi\mid Y,\mathbf{k},Z} \,\big\|\, U_{\mathcal{C}(Y,\mathbf{k},Z)}\right)\right] \geq 0,
\]
by the identity $D(P\|U) = \log_2|\mathrm{supp}\,U| - H(P)$ for $U$
uniform on a finite set containing $\mathrm{supp}\,P$. The flat scheme's
\emph{total} expected length is then at least
\[
H(\mathbf{k} \mid Y) + H(Z \mid Y, \mathbf{k}) + \mathbb{E}\!\left[\log_2 |\mathcal{C}(Y,\mathbf{k},Z)|\right] = H(\mathbf{k} \mid Y) + H(Z \mid Y, \mathbf{k}) + H(\Pi \mid Y, \mathbf{k}, Z) + G_{\mathrm{within}}^{Z}.
\]
Since $\mathbf{k}$ is a deterministic function of $\Pi$, the chain rule
gives
$H(\mathbf{k} \mid Y) + H(Z \mid Y, \mathbf{k}) + H(\Pi \mid Y, \mathbf{k}, Z) = H(\mathbf{k}, Z, \Pi \mid Y) = H(\Pi \mid Y) + H(Z \mid \Pi, Y) \geq H(\Pi \mid Y) = H(X \mid Y)$,
with equality when $Z$ is a deterministic function of $\Pi$ (the typical
case: the support is computed from the block, so $H(Z \mid \Pi, Y) = 0$).
So the flat total is $\geq H(X \mid Y) + G_{\mathrm{within}}^{Z}$, no
matter how the support metadata is chosen --- even a \emph{randomized} $Z$
only adds the nonnegative $H(Z \mid \Pi, Y)$, so pushing partition
information into $Z$ shrinks the rank payload but inflates the header by at
least as much. (The deficit $G_{\mathrm{within}}^{Z}$ is the entropy
below the \emph{candidate-set} maximum, not the full count-class maximum
unless $Z$ is empty and the support uninformative; and it is \emph{not}
the cross-position total correlation of Theorem 5.5, since
$U_{\mathcal{C}}$ has correlated coordinates --- the counts and support
are fixed --- so $\sum_i H_{\mathrm{marg}}(X_i \mid Y,\mathbf{k},Z) \neq \log_2 |\mathcal{C}(Y,\mathbf{k},Z)|$
in general.)

\textbf{(T)} If $\Pi$ is uniform on its candidate set
$\mathcal{C}(Y, \mathbf{k}, Z)$ given $(Y, \mathbf{k}, Z)$, the maximizer
is attained: $H(\Pi \mid Y, \mathbf{k}, Z) = \mathbb{E}[\log_2 |\mathcal{C}(Y,\mathbf{k},Z)|]$,
$G_{\mathrm{within}}^{Z} = 0$, and by the collapse above the flat total
equals $H(X \mid Y)$. With $Z$ empty this is the support-free case: take
the Theorem 5.3 code with $\alpha_v = 0$ and $S_v = A_v$, under which the
chained sum telescopes,
\[
\sum_{v=0}^{15} \log_2 \binom{|A_v|}{k_v} = \log_2 \binom{N}{k_0}\binom{N - k_0}{k_1}\cdots = \log_2 \frac{N!}{\prod_v k_v!} = \log_2 M(N; \mathbf{k}),
\]
so the flat payload equals $\log_2 M(N; \mathbf{k})$; with a correct
source-independent support it is the smaller
$\sum_v \log_2 \binom{s'_v}{k_v} = \log_2 |\mathcal{C}|$ (no $Z$ header),
still $\geq H(\Pi \mid Y, \mathbf{k})$ by part (S). When $\Pi$ is uniform
on $\mathcal{C}$ this equals $H(\Pi \mid Y, \mathbf{k})$, so adding the
count description at its entropy $H(\mathbf{k} \mid Y) + O(\log N)$ and
the coder redundancy $\varepsilon_N$ matches the converse up to lower-order
terms; the flat scheme is entropy-optimal. For a general (non-uniform)
law on $\mathcal{C}$, replace the uniform index code by an arithmetic code
driven by the conditional rank law $P(\Pi \mid Y, \mathbf{k}, Z)$: this
still has count-plus-rank structure (the rank is the same enumerative
index), costs $H(\Pi \mid Y, \mathbf{k}, Z) + \varepsilon_N$ in
expectation for the index (Cover and Thomas 2006, Theorem 5.4.1), and with
the count and metadata headers gives total
$H(X \mid Y) + \varepsilon_N + O(\log N)$ by the same collapse --- matching
the converse without leaving form (ii). The flat scheme's excess
$G_{\mathrm{within}}^{Z}$ is therefore the cost of uniform indexing on the
candidate set alone. \hfill$\blacksquare$

\emph{Relation to the sub-block converse (Lemma 5.1e).} Lemma 5.1e is a
converse for a \emph{different} operational setting: byte-granular
random access via $m$ independently decodable sub-blocks under a
sequential (Type-I-style) per-position coder, giving
$L_{\mathrm{synced}} \geq m \cdot H(D_K)$ with source-independent
metadata. It bounds the \emph{sync overhead} of sequential sub-block
coding, not the partition-rank cost of Type-II. Theorem 5.3b is the
converse for the partition (Type-II) representation and is logically
independent of Lemma 5.1e: the former concerns the count-plus-rank
decomposition of a whole-block partition code; the latter concerns the
sub-block decomposition of a sequential stream code. Where they touch is
that both are instances of the same master principle --- a uniquely
decodable code for a deterministic function of the data costs at least
that function's conditional entropy --- applied to two different
factorizations.

\emph{Scope and honesty.} The converse for both the general Type-II
class (i) and the count-plus-rank form (ii) is the same: $H(X \mid Y)$,
with no structural overhead. Three coders attain it (up to
$\varepsilon_N$): ideal joint arithmetic coding of $\Pi$ (i); the
count-plus-rank coder with an entropy-coded index (ii); and, when the
conditional law is uniform on its candidate set, the flat-enumerative
Theorem 5.3 coder (iii). The flat-enumerative \emph{total}
$H(Z \mid Y, \mathbf{k}) + \log_2 |\mathcal{C}(Y,\mathbf{k},Z)|$ for
source-dependent support metadata $Z$ --- reducing to the multinomial
$\log_2 M(N;\mathbf{k})$ when $Z$ is empty and no support is used --- is a
converse \emph{only} for the flat form (iii), where it exceeds the optimum
by the within-candidate-set uniform deficit $G_{\mathrm{within}}^{Z}$,
zero exactly when the conditional law is uniform on
$\mathcal{C}(Y,\mathbf{k},Z)$. The three counterexamples above show this
overhead is genuinely specific to uniform indexing on the candidate set:
it is not forced by the count-plus-rank decomposition (entropy-coding the
index removes it), is not measured against the full count class once a
source-independent support is used, and cannot be evaded by smuggling
partition information into a source-dependent support (which is charged in
$Z$, leaving the total $\geq H(X \mid Y)$) --- and a fortiori it is not a
property of Type-II. This is the honest reading of ``matching converse'':
Theorem 5.3 is entropy-optimal exactly on the uniform-candidate-set
regime (where its flatness costs nothing), and is matched to the optimum
elsewhere after its enumerative index is entropy-coded, the residual
flat-indexing gap $G_{\mathrm{within}}^{Z}$ being fully characterized as a
KL divergence to the candidate-set uniform law. We do \emph{not} claim
that $\log_2 M(N;\mathbf{k})$ lower-bounds an arbitrary Type-II or
count-plus-rank code, nor even an arbitrary \emph{support-restricted}
flat code; for those, part (C)'s $H(X \mid Y)$ is the operative and tight
bound, and the flat total is
$H(Z \mid Y, \mathbf{k}) + \log_2 |\mathcal{C}(Y,\mathbf{k},Z)|$.

\textbf{5.6 Type-II vs Type-I: a separation theorem}

Lemma 5.1 shows that Type-II under a factorized $Q$ is operationally
identical to per-position Type-I arithmetic coding. The natural
question is whether semi-non-factorized Type-II provides a genuine
separation. The next theorem answers this in the affirmative for a
concrete family of joint distributions.

\textbf{Theorem 5.4 (Type-II semi-non-factorized separation,
amortized).} \emph{There exists a family of distributions $D_N$ over
length-$N$ nibble strings, and a positive constant $c > 0$, such that
for every sufficiently large $N$ the optimal factorized per-position
arithmetic coder driven by the true marginals
$Q^{\text{marg}}_v(i) = \Pr_{X \sim D_N}(x_i = v)$ has expected code
length at least $4N$ bits, while a semi-non-factorized Type-II coder
that has access to a single bit per position of joint side information
(membership in a structured set) achieves amortized expected length at
most $(4 - c)N + O(\log N)$ bits per block, when amortized over
$B_N = \Omega(N^2)$ blocks sharing the same generator-matrix model
header. The amortized per-block separation is therefore $\Theta(N)$.}

\emph{\textbf{Proof.}} Fix a rational parameter $c \in (0, 4)$ and
set $d := \lfloor (4 - c) N \rfloor$ (the integer message
dimension; the floor causes at most a $1/N$ additive loss in $c$
that is absorbed into the $O(\log N)$ slack of the theorem
statement and does not affect the $\Theta(N)$ amortized separation).
Let $G$ be a $d \times 4N$ generator matrix over $\mathrm{GF}(2)$
of a linear code $C_N$ satisfying the following two properties:
\begin{itemize}
\item[\textbf{(P1)}]
for every codeword position $j \in \{1, \ldots, 4N\}$, the $j$-th
coordinate of $mG$ (where $m \in \mathrm{GF}(2)^d$ is the message)
is a non-trivial $\mathrm{GF}(2)$-linear function of $m$ ---
equivalently, the $j$-th column of $G$ is non-zero;
\item[\textbf{(P2)}]
for every nibble position $i \in \{1, \ldots, N\}$, the four
columns of $G$ at codeword coordinates $4(i-1)+1, \ldots, 4i$ are
linearly independent in $\mathrm{GF}(2)^d$;
\item[\textbf{(P3)}]
$G$ has full row rank $d$.
\end{itemize}
Such matrices exist with high probability over the uniform
distribution on $d \times 4N$ matrices over $\mathrm{GF}(2)$. For
\textbf{(P1)}, each individual column is the zero vector in
$\mathrm{GF}(2)^d$ with probability $2^{-d}$; a union bound over the
$4N$ columns gives failure probability $\leq 4N \cdot 2^{-d}$. For
\textbf{(P2)}, four fixed columns are linearly dependent iff one of
the $2^4 - 1 = 15$ non-zero $\mathrm{GF}(2)$-linear combinations
yields the zero vector; each such combination of four uniform
independent random vectors in $\mathrm{GF}(2)^d$ is itself uniform
on $\mathrm{GF}(2)^d$ and vanishes with probability $2^{-d}$. Union
over $N$ nibble positions: $\leq 15 N \cdot 2^{-d}$. For
\textbf{(P3)}, a fixed non-zero row-combination
$\lambda \in \mathrm{GF}(2)^d \setminus \{0\}$ produces
$\lambda G \in \mathrm{GF}(2)^{4N}$ which vanishes identically (on
all $4N$ columns) with probability $2^{-4N}$ over the random choice
of $G$; a union bound over the $2^d - 1 < 2^d$ non-zero combinations
gives failure probability $\leq 2^{d - 4N} = 2^{-cN}$.
Combined failure probability:
$\leq 19 N \cdot 2^{-(4-c)N} + 2^{-cN}$, which is $o(1)$ as
$N \to \infty$ for any fixed $c \in (0, 4)$.

Let $S_N \subset \{0, \ldots, 15\}^N$ be the image of all $2^d$
messages under $G$, with the $4N$ output bits parsed as $N$ nibbles.
Property \textbf{(P3)} implies that $G$ has full row rank $d$, so the
encoding map $m \mapsto mG$ is injective and
$|S_N| = 2^d = 2^{\lfloor (4-c)N \rfloor}$. Define $D_N$ as the uniform distribution
on $S_N$.

\textbf{(i) Marginal uniformity.} Fix a nibble position
$i \in \{1, \ldots, N\}$. The four bits encoding the nibble at
position $i$ are values of four $\mathrm{GF}(2)$-linear forms on the
message $m \in \mathrm{GF}(2)^d$, which by \textbf{(P2)} are linearly
independent. The standard fact that $k$ linearly independent
$\mathrm{GF}(2)$-linear forms on a uniformly random vector in
$\mathrm{GF}(2)^d$ induce the uniform distribution on
$\mathrm{GF}(2)^k$ (the joint map is a surjective
$\mathrm{GF}(2)$-linear homomorphism with kernel of size $2^{d - k}$,
so each image has the same preimage size) then gives: each of the
16 bit patterns at position $i$ occurs with probability exactly
$1/16$. Therefore $Q^{\text{marg}}_v(i) = 1/16$ for all $i, v$.

\textbf{(ii) Factorized lower bound.} The cross-entropy of any
$X \sim D_N$ under the factorized model with $Q_v(i) = 1/16$ is
exactly $N \cdot \log_2 16 = 4N$ bits; this is also the optimal
expected code length of an arithmetic coder driven by
$Q^{\text{marg}}$, up to the redundancy $\varepsilon_N$.

\textbf{(iii) Joint upper bound.} A semi-non-factorized coder that
encodes membership in $S_N$ as the unique $d$-bit message index per
block, followed by deterministic expansion via $G$, achieves per-block
content length
$\log_2 |S_N| + \varepsilon_N = d + \varepsilon_N \leq (4 - c)N +
\varepsilon_N$ bits (with the $\leq$ from the floor in
$d = \lfloor (4-c)N \rfloor$, an additive loss of at most $1$ bit
absorbed into the $O(\log N)$ slack), i.e., savings of at least
$4N - (4-c)N = cN$ bits per block \emph{before} model-header
amortization. Describing the generator matrix $G$ itself costs
$L(G) = d \cdot 4N \leq 4(4-c)N^2$ bits, paid once and absorbed
into the $L(M_{\text{pos}})$ term of (3.1). When the same $G$ is
reused across $B_N$ blocks, the per-block model contribution is
$L(G)/B_N$; taking $B_N = \Omega(N^2)$ (e.g., $B_N = N^2$) makes
this $O(1)$ bits per block, absorbed into the $O(\log N)$ overhead.
For fewer blocks the per-block cost is correspondingly higher.

\textbf{(iv) Separation.} The linear-code joint coder of step (iii)
transmits only the $d$-bit message index per block; the decoder
recovers $X$ by computing $m G$ and parsing as nibbles, with no
auxiliary count vector required (the construction of $D_N$ already
fixes $|S_N| = 2^d$ and the marginals are uniform by step (i)).
Subtracting the per-block upper bound (iii) from the per-block lower
bound (ii) thus gives
\[
\Delta(N) \geq 4N - \left[(4 - c)N + \varepsilon_N + 1\right] = cN - \varepsilon_N - 1 = \Theta(N)
\]
for any fixed $c \in (0, 4)$, once the one-time generator-matrix
description $L(G) \leq 4(4-c)N^2$ is amortized across $B_N = \Omega(N^2)$
blocks (the per-block model contribution then absorbs into the
$O(1)$ overhead). The $-1$ from the floor in $d$ is absorbed.
\hfill$\blacksquare$

Theorem 5.4 establishes that semi-non-factorized Type-II strictly
separates from \emph{factorized marginal} per-position coders: coders
that produce one categorical distribution per position using only
each position's marginal $Q^{\text{marg}}_v(i)$, with no joint
conditioning. A genuine autoregressive coder (PixelRNN, PixelCNN,
NNCP) that uses chain-rule conditionals
$Q(x_i \mid x_1, \ldots, x_{i-1})$ can in principle achieve the
joint entropy $H(D_N) = \lfloor (4 - c)N \rfloor$ on the same $D_N$
and is \emph{not} separated by Theorem 5.4; the relevant boundary is
between models with and without joint conditioning, not between RNN
and CNN architectures. The separation is $\Theta(N)$, not
$\Theta(\log N)$: joint structure beyond per-position marginals can
carry a constant fraction of a bit per position even for the
marginal-conditioning case. The linear-code
family used in the proof is a worst case for factorized models
(marginals are exactly uniform) and a best case for joint coders (the
joint distribution lives on a $2^{\lfloor (4-c)N \rfloor}$-element set). Realistic
data classes interpolate between these extremes; the practical regime
in which Type-II is advantageous is one where the data lives on a
small, structurally describable subset of $\{0, \ldots, 15\}^N$ whose
per-position marginals appear approximately uniform. Count constraints
alone yield only a logarithmic separation (the Stirling correction),
so richer joint structure --- geometric constraints, algebraic
parities, sparsity patterns, sortedness --- is required to exhibit the
$\Theta(N)$ regime.

The linear-code construction of Theorem 5.4 is one extreme of a more
general phenomenon. The next result identifies the exact
information-theoretic quantity governing the separation between
factorized and joint coding for any source distribution, and the
following remark connects this quantity to the manifold hypothesis from
statistical learning theory and topological data analysis.

\textbf{Theorem 5.5 (Multi-information separation).} \emph{Let $D_N$
be any distribution over $A_{16}^N$. Define the total correlation
(multi-information, Watanabe 1960; see also Studen\'y and Vejnarov\'a
1998)
$\mathrm{TC}(D_N) = \sum_{i=1}^N H_{\text{marg}}(X_i) - H(D_N)$, where
$H_{\text{marg}}(X_i)$ is the marginal entropy at position $i$ under
$D_N$. Then:}
\begin{enumerate}
\item[\textnormal{(F)}] \emph{(Factorized rate.) The optimal factorized
per-position arithmetic coder driven by the true marginals
$Q^{\mathrm{marg}}_v(i) = \Pr_{X \sim D_N}(x_i = v)$ has expected
length \emph{at least} $\sum_i H_{\mathrm{marg}}(X_i)$ bits per block
(Shannon's converse for product-distribution coders), and a block
arithmetic coder under those marginals attains
$\sum_i H_{\mathrm{marg}}(X_i) + \varepsilon_N$ in expectation
(achievability), with $\varepsilon_N$ the per-block entropy-coder
redundancy of §2.3.}
\item[\textnormal{(J)}] \emph{(Joint rate.) Any uniquely decodable lossless
code for $X$ has expected length \emph{at least} $H(D_N)$ bits per
block (Shannon's converse), and a block arithmetic coder driven by
the joint distribution $D_N$ attains
$H(D_N) + \varepsilon_N$ in expectation (achievability).}
\item[\textnormal{(S)}] \emph{(Separation.) The gap between the optimal
factorized and the optimal joint expected per-block rates equals
$\mathrm{TC}(D_N) + O(\varepsilon_N)$.}
\end{enumerate}

\emph{\textbf{Proof.}} (F) The converse is Shannon's source coding
theorem applied to the product distribution
$\prod_i Q^{\mathrm{marg}}(\cdot \mid i)$; the achievability uses
the same product model with an arithmetic coder (Cover and Thomas
2006, Theorem 5.4.1), incurring at most $\varepsilon_N$ redundancy
per block (§2.3). The expected length under $D_N$ of any code matched
to this product model is the cross-entropy
$\mathbb{E}_{D_N}[-\sum_i \log_2 Q^{\mathrm{marg}}_{X_i}(i)] =
\sum_i H_{\mathrm{marg}}(X_i)$ (since the marginals match).
(J) The converse is Shannon's source coding theorem applied directly
to $D_N$; the achievability is block arithmetic coding under $D_N$,
again with $\varepsilon_N$ redundancy. Neither argument uses AEP /
typical-set machinery (which would require an iid asymptotic setup
that we do not assume).
(S) Subtracting the two attained rates gives
$\sum_i H_{\mathrm{marg}}(X_i) - H(D_N) = \mathrm{TC}(D_N)$ by
definition. Total correlation is non-negative, with equality iff
positions are mutually independent under $D_N$.
\hfill$\blacksquare$

Theorem 5.5 generalizes Theorem 5.4 from the uniform-marginal
linear-code construction to arbitrary source distributions. The
quantity $\mathrm{TC}(D_N)$ is the classical \emph{multi-information}
of Watanabe (1960), and the gap between factorized and joint coding
is a standard consequence of the chain rule for entropy; the novelty
is the placement of this quantity as the operational separator
between Type-I and Type-II under the RNR framework. The $\Theta(N)$
separation observed in Theorem 5.4 corresponds to
$\mathrm{TC} \geq cN - 1$ for the linear-code family (sum of
marginals $\sum_i H_{\text{marg}}(X_i) = 4N$ minus joint entropy
$H(D_N) = \lfloor (4-c)N \rfloor$). For general
$D_N$, the magnitude of $\mathrm{TC}(D_N)$ determines whether Type-II
provides a constant-factor or vanishing improvement over factorized
coding.

\textbf{Remark 5.6 (Manifold-hypothesis intuition, not a theorem).}
The manifold hypothesis from statistical learning theory and
topological data analysis (Donoho 2000; Tenenbaum, de Silva, Langford
2000; Carlsson 2009) is that high-dimensional natural data ---
images, text, audio, structured records --- concentrates near a
low-dimensional structure in the ambient space. It is tempting to
conclude that if $D_N$'s typical realizations live near a structure
of intrinsic dimension $d \ll N$, then $H(D_N) = O(d \log N)$ and
hence $\mathrm{TC}(D_N) = \Theta(N)$ under near-uniform marginals.
\emph{This implication is not a theorem.} A discrete distribution
supported near a $d$-dimensional manifold can have arbitrarily large
Shannon entropy unless additional structure is assumed: a
quantization scale $\Delta$, a bound on the manifold's reach
(local-curvature radius) and on the density (e.g., Lipschitz density
with respect to the manifold's intrinsic measure), and a bound on
the parameter complexity of the manifold itself. Under explicit such
assumptions one can derive entropy bounds of the form
$H(D_N) \leq d \log_2(1/\Delta) + C(d, \mathrm{reach}, \text{density})$;
without them, low intrinsic dimension is not enough.

The cited empirical estimates also do not support a direct
quantitative claim about whole natural images. Carlsson, Ishkhanov,
de Silva and Zomorodian (2008) found low-dimensional local topology
(specifically a Klein-bottle-like structure) in $3 \times 3$
high-contrast image patches; they do not establish
$d \in [10^2, 10^3]$ for full images in ambient dimension $10^6$.
Levina and Bickel (2004) gave a maximum-likelihood intrinsic-dimension
estimator with known biases on noisy real data. We therefore treat
the manifold connection as informal motivation only. Whether any
specific natural-data class $D_N$ satisfies $\mathrm{TC}(D_N) =
\Theta(N)$ is an empirical question (deferred to §11), and
constructing an efficient enumerative encoder for the corresponding
typical set is open.

\emph{Formal refinement: Lemma 5.6i (below in this section) makes
the limits of the manifold intuition precise by exhibiting a class
that does fit a ``low-dimensional combinatorial support'' framing
but has \emph{sub-linear} $\mathrm{TC}$: uniform $K$-sparse binary
vectors with $|\mathcal{C}| = \binom{N}{K} = 2^{O(K \log(N/K))}$
have $\mathrm{TC}(D_N) = O(\log K) + O(\log N) = o(N)$. The
contrast with uniform permutations (also a ``low-dimensional''
combinatorial support but with $\mathrm{TC} = \Theta(N)$) shows
that $\Theta(N)$ TC requires explicit pairwise dependence structure
(distinct-value constraints, Markov-style local-clique factors, or
hidden latent classes per Lemma 5.6b), not just low support
cardinality.}

For a specific structured data class --- one with explicit
combinatorial constraints rather than learned approximate
structure --- we can give an explicit linear-in-$N$ scaling.

\textbf{Lemma 5.6a (Linear $\mathrm{TC}$ for uniqueness-constrained
database pages).} \emph{Let $D_N$ be the uniform distribution over
injections $\{1, \dots, N\} \to \{1, \dots, K\}$. For any $K \geq
N+1$ (so the injection condition is feasible and non-degenerate),}
\[
\mathrm{TC}(D_N) \;=\; N \cdot \bigl[\log_2 e - (K/N - 1) \log_2\bigl((K/N)/((K/N)-1)\bigr)\bigr] + O(\log N).
\tag{5.5}
\]
\emph{When $K = \lceil \alpha N \rceil$ for a constant $\alpha > 1$,
this becomes $\mathrm{TC}(D_N) = N \cdot c(\alpha) + O(\log N)$ with
$c(\alpha) = \log_2 e - (\alpha-1) \log_2(\alpha/(\alpha-1))$
(integer-$K$ corrections absorbed into the $O(\log N)$ remainder).
In particular, $\mathrm{TC}(D_N) = \Theta(N)$ for any fixed
$\alpha > 1$. The coefficient $c(\alpha)$ is strictly positive,
satisfies $\lim_{\alpha \to 1^+} c(\alpha) = \log_2 e \approx
1.4427$ (tightest constraint at $K = N+1$, where the injections
saturate the UNIQUE constraint), and decreases monotonically to
$0$ as $\alpha \to \infty$ (the constraint becomes vacuous when
$K \gg N$). Selected values: $c(1.5) \approx 0.650$,
$c(2) = \log_2(e/2) \approx 0.4427$, $c(3) \approx 0.273$,
$c(10) \approx 0.075$.}

\emph{\textbf{Proof.}} By symmetry over rows, the marginal of any
single row is uniform on $\{1, \dots, K\}$, giving sum of marginal
entropies
$\sum_i H(\text{row } i) = N \log_2 K = N \log_2(\alpha N)$.

The joint entropy is the log of the number of valid configurations:
$H(D_N) = \log_2 \frac{K!}{(K-N)!}$.
Using Stirling
$\log_2(M!) = M \log_2 M - M \log_2 e + (1/2)\log_2(2\pi M) + O(1/M)$:
\[
\log_2 \frac{K!}{(K-N)!}
= K \log_2 K - (K-N) \log_2(K-N) - N \log_2 e + O(\log N).
\]
Substituting $K = \alpha N$ and $K-N = (\alpha-1) N$:
\[
\begin{aligned}
H(D_N) &= \alpha N \log_2(\alpha N) - (\alpha-1) N \log_2((\alpha-1) N) - N \log_2 e + O(\log N) \\
&= N \log_2 N \cdot [\alpha - (\alpha-1)] + N \cdot [\alpha \log_2 \alpha - (\alpha-1) \log_2(\alpha-1)] - N \log_2 e + O(\log N) \\
&= N \log_2 N + N \cdot [\alpha \log_2 \alpha - (\alpha-1) \log_2(\alpha-1) - \log_2 e] + O(\log N).
\end{aligned}
\]
Subtracting from the sum of marginals
$N \log_2(\alpha N) = N \log_2 \alpha + N \log_2 N$:
\[
\begin{aligned}
\mathrm{TC}(D_N) &= N \log_2 \alpha - N \cdot [\alpha \log_2 \alpha - (\alpha-1) \log_2(\alpha-1) - \log_2 e] + O(\log N) \\
&= N \cdot [\log_2 \alpha + \log_2 e - \alpha \log_2 \alpha + (\alpha-1) \log_2(\alpha-1)] + O(\log N) \\
&= N \cdot [\log_2 e - (\alpha-1)(\log_2 \alpha - \log_2(\alpha-1))] + O(\log N) \\
&= N \cdot [\log_2 e - (\alpha-1) \log_2(\alpha/(\alpha-1))] + O(\log N).
\end{aligned}
\]
The coefficient
$c(\alpha) = \log_2 e - (\alpha-1) \log_2(\alpha/(\alpha-1))$
satisfies $c(\alpha) > 0$ for all $\alpha > 1$. Differentiating with
respect to $\alpha$:
$dc/d\alpha = -\log_2(\alpha/(\alpha-1)) + (\alpha-1) \cdot \frac{1}{\alpha(\alpha-1)\ln 2} \cdot 1
= -\log_2(\alpha/(\alpha-1)) + \frac{1}{\alpha \ln 2}$.
For $\alpha > 1$, $\log_2(\alpha/(\alpha-1)) > \frac{1}{\alpha \ln 2}$
(by concavity of $\log$), so $dc/d\alpha < 0$ and $c$ is monotone
decreasing on $(1, \infty)$. As $\alpha \to 1^+$:
$(\alpha-1) \log_2(\alpha/(\alpha-1)) \to 0 \cdot \infty$ which by
L'H\^opital evaluates to $0$, so $c(1^+) = \log_2 e$. As
$\alpha \to \infty$:
$(\alpha-1) \log_2(\alpha/(\alpha-1)) \to \log_2 e$ (using
$\log(1 + 1/(\alpha-1)) \sim 1/((\alpha-1) \ln 2)$), so $c(\infty) = 0$.
\hfill$\blacksquare$

Lemma 5.6a gives the first explicit non-asymptotic $\Theta(N)$
scaling for a specific structured data class within the framework
of Theorem 5.5. The numerical verification at $N = 10000$,
$\alpha = 2$ matches the formula to within $\approx -0.5$ bits
absolute residual (the higher-order Stirling correction);
see \texttt{scripts/verify/lemma\_5\_6a\_tc\_uniqueness.py}.
The proof extends with no structural change to the multi-column
case where each row carries $c$ independent uniqueness-constrained
columns of varying $\alpha_j$, with TC additive across columns.

\textbf{Lemma 5.6b (Linear $\mathrm{TC}$ for hidden-class iid
sources).} \emph{Let $Y$ be a discrete random variable on a finite
alphabet $\mathcal{Y}$ with $|\mathcal{Y}| < \infty$, and let
$X_1, \ldots, X_N$ be conditionally independent and identically
distributed given $Y$ with each $X_i$ taking values in a finite
observation alphabet $\mathcal{X}$ under conditional law
$\Pr(X_i = x \mid Y = y) = p_y(x)$ for each $i$. Then}
\[
\mathrm{TC}(D_N) \;=\; N \cdot I(X_1; Y) - H(Y) + H(Y \mid X_1, \ldots, X_N).
\tag{5.6}
\]
\emph{In particular, if the conditional laws are pairwise
distinguishable in the sense that
$\min_{y \neq y'} D_{\mathrm{TV}}(p_y, p_{y'}) \geq \delta > 0$, then
$H(Y \mid X_1, \ldots, X_N) \to 0$ exponentially in $N$, and}
\[
\mathrm{TC}(D_N) = N \cdot I(X_1; Y) - H(Y) + o(1) = \Theta(N)
\]
\emph{whenever $I(X_1; Y) > 0$ (equivalently, whenever the joint
$(Y, X_i)$ is non-product, which by conditional iid is equivalent
for every single $i$). The pairwise hypothesis is essential:
collapsing two classes (say $p_y = p_{y'}$ for some $y \neq y'$)
makes them indistinguishable to any function of $X^N$, leaving
$H(Y \mid X^N) \to H(Y \mid \tilde Y)$ where $\tilde Y$ is the
projection of $Y$ onto its equivalence classes under $p_\cdot$;
the asymptotic identity then degrades to
$\mathrm{TC}(D_N) = N \cdot I(X_1; Y) - H(Y) + H(Y \mid \tilde Y) + o(1)$.}

\emph{\textbf{Proof.}} The marginal of any single position is the
mixture
$\Pr(X_i = x) = \sum_y \Pr(Y = y) p_y(x) =: \bar p(x)$,
independent of $i$ by the iid-given-$Y$ assumption. Hence
$H(X_i) = H(\bar p)$ for all $i$ and
\[
\sum_{i=1}^N H_{\mathrm{marg}}(X_i) = N \cdot H(\bar p).
\]

For the joint entropy, expand via the chain rule on the joint
distribution of $(Y, X_1, \ldots, X_N)$:
\[
\begin{aligned}
H(Y, X_1, \ldots, X_N)
&= H(Y) + H(X_1, \ldots, X_N \mid Y) \\
&= H(Y) + N \cdot H(X_1 \mid Y)
\end{aligned}
\]
(the second equality uses conditional iid given $Y$). Symmetrically,
\[
H(Y, X_1, \ldots, X_N) = H(X_1, \ldots, X_N) + H(Y \mid X_1, \ldots, X_N).
\]
Equating and solving for $H(X_1, \ldots, X_N)$:
\[
H(X_1, \ldots, X_N) = H(Y) + N \cdot H(X_1 \mid Y) - H(Y \mid X_1, \ldots, X_N).
\]

Combining with the sum of marginals:
\[
\begin{aligned}
\mathrm{TC}(D_N) &= \sum_i H_{\mathrm{marg}}(X_i) - H(X_1, \ldots, X_N) \\
&= N \cdot H(\bar p) - H(Y) - N \cdot H(X_1 \mid Y) + H(Y \mid X_1, \ldots, X_N) \\
&= N \cdot [H(X_1) - H(X_1 \mid Y)] - H(Y) + H(Y \mid X_1, \ldots, X_N) \\
&= N \cdot I(X_1; Y) - H(Y) + H(Y \mid X_1, \ldots, X_N).
\end{aligned}
\]

The asymptotic behavior of the posterior $H(Y \mid X_1, \ldots, X_N)$
follows from a standard Chernoff-Hellinger bound for multi-hypothesis
testing. Under the pairwise separation hypothesis
$\min_{y \neq y'} D_{\mathrm{TV}}(p_y, p_{y'}) \geq \delta > 0$, the
Bayes-optimal classifier for $Y$ given $X^N$ errs with probability
$P_e \leq O(\exp(-c \delta^2 N))$ for some constant
$c > 0$ depending on the prior $\Pr(Y)$ and the alphabet size
$|\mathcal{Y}|$; this is the multi-hypothesis Chernoff bound (the
two-hypothesis case is Cover \& Thomas \emph{Elements of Information
Theory} 2nd ed., 2006, Theorem 11.9.1; the multi-hypothesis
extension follows by a union bound over the $\binom{|\mathcal{Y}|}{2}$
pairwise tests, each at error $O(e^{-c \delta^2 N})$ by Theorem
11.9.1. The two-hypothesis testing inequalities of Bretagnolle--Huber
1979 and LeCam 1973 provide the underlying TV-to-error-probability
estimate). Fano's inequality then gives
$H(Y \mid X^N) \leq h_2(P_e) + P_e \cdot \log_2(|\mathcal{Y}| - 1)
= O(N \cdot e^{-c\delta^2 N}) \to 0$ as $N \to \infty$.
Hence the leading-order TC is $N \cdot I(X_1; Y) - H(Y)$ up to
$o(1)$ (in fact exponentially small) corrections.
\hfill$\blacksquare$

\emph{Examples.}
(i) \emph{Binary mixture}: $Y \sim \mathrm{Bern}(1/2)$,
$p_0 = \mathrm{Bern}(0.3)$, $p_1 = \mathrm{Bern}(0.7)$. The marginal
$\bar p = \mathrm{Bern}(0.5)$ has $H(\bar p) = 1$; the conditional
$H(X|Y) = h_2(0.3) \approx 0.881$. So
$I(X_1; Y) = 1 - 0.881 = 0.119$ bit, $H(Y) = 1$ bit, and for
large $N$,
$\mathrm{TC}(D_N) \approx 0.119 N - 1$ bits.
(ii) \emph{Topic-model text}: a Latent Dirichlet Allocation source
with $K$ topics, fixed per-topic vocabulary distributions, and a
single topic per document. If the topic distributions $\{p_k\}$ are
mutually well-separated, $\mathrm{TC}$ of $N$ tokens scales linearly
with rate $I(X_1; Y) = H(\bar p) - H(X|Y) \in (0, \log K]$.
(iii) \emph{Markov chain (degenerate case)}: when $Y = (X_0)$ is
the previous symbol, conditional independence given $Y$ fails (the
chain has correlations beyond the immediate past); the lemma does
not apply directly but is recoverable by augmenting $Y$ to encode
sufficient context. The simple stationary-Markov-chain analog
$\mathrm{TC}(D_N) = (N-1) \cdot I(X_1; X_2)$ requires a separate
argument (chain-rule decomposition) and is included in
\texttt{scripts/verify/lemma\_5\_6b\_hidden\_class\_tc.py} as a
companion check.

\emph{Discussion.} Lemma 5.6b provides a second explicit
$\Theta(N)$-TC class beyond the uniqueness-constrained injections
of Lemma 5.6a, capturing the latent-variable structure characteristic
of natural-language and document sources. Both lemmas use the
Theorem 5.5 identity as their formal foundation; the technical
difficulty in each case is computing the joint entropy exactly
(Stirling for Lemma 5.6a, the chain rule with posterior correction
for Lemma 5.6b).

\textbf{Lemma 5.6c (Linear $\mathrm{TC}$ for stationary processes
with positive entropy-rate gap).} \emph{Let $(X_n)_{n \geq 1}$ be a
strictly stationary discrete-time random process on a finite
alphabet $\mathcal{X}$ (so the joint law of any block
$(X_{j+1}, \ldots, X_{j+N})$ is independent of the shift $j \geq 0$).
Define the marginal entropy
$H_{\mathrm{marg}} := H(X_1)$, the entropy rate
$h(X) := \lim_{n \to \infty} H(X_n \mid X_1, \ldots, X_{n-1})$
(which exists by stationarity, Cover \& Thomas 2006 Thm 4.2.1),
and the per-position excess entropy
$E := H_{\mathrm{marg}} - h(X) \geq 0$. Then for any $N \geq 1$,}
\[
\mathrm{TC}(D_N) \;=\; N \cdot E - \delta_N,
\quad\text{where}\quad
\delta_N := \sum_{n=1}^N \bigl[H(X_n \mid X_1, \ldots, X_{n-1}) - h(X)\bigr] \geq 0.
\tag{5.7}
\]
\emph{In particular:}

\emph{(a) Order-$W$ stationary Markov chains.} If $(X_n)$ is order-$W$
Markov (so $H(X_n \mid X_1, \ldots, X_{n-1}) = h(X)$ for $n > W$),
then $\delta_N$ stabilizes at a constant
$\delta_W = \sum_{n=1}^W [H(X_n \mid X_1, \ldots, X_{n-1}) - h(X)]$
for $N \geq W$, and
\[
\mathrm{TC}(D_N) = N \cdot E - \delta_W \quad (\text{exact, for } N \geq W).
\]
This is $\Theta(N)$ if and only if $E > 0$ (the chain is not
memoryless). For order-1 Markov, $\delta_W = E$ and the formula
recovers $\mathrm{TC}(D_N) = (N-1) \cdot I(X_1; X_2)$, the companion
of Lemma 5.6b.

\emph{(b) Summable-correlation stationary processes.} If the
conditional entropy convergence is summable
($\sum_{n=1}^\infty [H(X_n \mid \mathrm{past}) - h(X)] < \infty$,
a standard mixing assumption; see Crutchfield \& Feldman 2003
\emph{Chaos} 13(1):25--54 for ``excess entropy'' in this sense),
then $\delta_N \to \delta_\infty < \infty$ as $N \to \infty$ and
\[
\mathrm{TC}(D_N) = N \cdot E - \delta_\infty + o(1).
\]
Hence $\mathrm{TC}(D_N) = \Theta(N)$ whenever $E > 0$.

\emph{(c) General stationary processes (Ces\`aro limit).} For any
strictly stationary process, $H(X_n \mid X_1, \ldots, X_{n-1}) \to h(X)$
as $n \to \infty$ (Cover \& Thomas 2006 Thm 4.2.2), so by the Ces\`aro
averaging
$\delta_N / N \to 0$
as $N \to \infty$. Consequently
\[
\mathrm{TC}(D_N) / N \to E,
\]
i.e.\ $\mathrm{TC}(D_N) = N \cdot E + o(N)$ whenever $E > 0$, even
without summability of the excess terms.

\emph{\textbf{Proof.}} By the chain rule for entropy,
$H(X_1, \ldots, X_N) = \sum_{n=1}^N H(X_n \mid X_1, \ldots, X_{n-1})$
with the convention $H(X_1 \mid \emptyset) = H(X_1)$. By
stationarity each marginal equals $H_{\mathrm{marg}}$, so the sum
of marginals is $N \cdot H_{\mathrm{marg}}$. Subtracting:
\[
\mathrm{TC}(D_N) = N \cdot H_{\mathrm{marg}} - \sum_{n=1}^N H(X_n \mid X_1, \ldots, X_{n-1}).
\]
Adding and subtracting $N \cdot h(X)$ on the right:
\[
\mathrm{TC}(D_N) = N \cdot [H_{\mathrm{marg}} - h(X)] - \sum_{n=1}^N [H(X_n \mid X_1, \ldots, X_{n-1}) - h(X)]
                 = N \cdot E - \delta_N.
\]
Non-negativity of $\delta_N$ follows from $H(X_n \mid \mathrm{past})
\geq h(X)$ for each finite $n$ (the conditional entropy is
non-increasing in the size of the conditioning set, with limit
$h(X)$; this is Cover \& Thomas 2006 Theorem 4.2.2). Equality of
$\delta_W$ to a stabilized constant under order-$W$ Markov is
immediate from the Markov property. Summable correlations imply
convergence of $\delta_N$ by Cauchy's test on positive-term series.
\hfill$\blacksquare$

\emph{Discussion.} Lemma 5.6c extends Lemmas 5.6a, 5.6b to the third
canonical natural-data pattern: temporal/sequential correlations.
Together the three lemmas cover the three foundational structural
mechanisms by which natural sources accumulate linear
$\mathrm{TC}$: combinatorial column constraints (Lemma 5.6a),
shared latent variable across positions (Lemma 5.6b), and temporal
dependence (Lemma 5.6c). Numerical illustration on English text (Shannon 1951
\emph{Bell System Technical Journal} 30(1):50, 27-letter alphabet
including space): Shannon's Table 1 gives the letter-frequency rate
$F_1 \approx 4.03$ bits/letter and the bigram-conditional rate
$F_2 \approx 3.32$ bits/letter, plus experimental bounds on the
asymptotic entropy rate from Shannon's 100-letter prediction
experiment ($0.6$ bits/letter lower bound, $1.3$ bits/letter upper
bound, with explicit sampling-error caveats). Taking
$H_{\mathrm{marg}} = F_1$, the per-position excess entropy at the
bigram level is
$E_{\mathrm{bigram}} \approx 4.03 - 3.32 = 0.71$ bits/letter, and
asymptotically $E_{\mathrm{en}} \approx 4.03 - 0.95 \approx 3.08$
bits/letter (using the midpoint of Shannon's bounds). For
$N = 2^{20}$ letters of English text, a per-position factorized
coder leaves $\approx 0.74$ Mbit (order-$1$ bigram approximation,
using $E_{\mathrm{bigram}} = 0.71$ bits/letter) to $\approx 3.23$
Mbit (asymptotic, using midpoint $E_{\mathrm{en}} = 3.08$
bits/letter; ranging from $2.86$ Mbit at Shannon's upper bound
$1.3$ to $3.59$ Mbit at lower bound $0.6$) of joint structure
unexploited. The
verification script
\texttt{scripts/verify/lemma\_5\_6c\_stationary\_markov\_tc.py}
exercises identity (5.7) for stationary order-$1$ and order-$2$
chains, including non-trivial transition matrices.

\emph{Terminology note.} Our per-position excess entropy
$E := H_{\mathrm{marg}} - h(X)$ is the \emph{slope} of TC as a
function of $N$ and should not be conflated with Crutchfield \&
Feldman's (2003) ``excess entropy'' $\mathbf{E}_{\mathrm{CF}}$,
which is defined via the block-entropy asymptote
$H(X^L) \sim h(X) \cdot L + \mathbf{E}_{\mathrm{CF}}$ as $L \to \infty$.
In our notation $\mathbf{E}_{\mathrm{CF}} = \lim_{N \to \infty}
[N \cdot H_{\mathrm{marg}} - N \cdot E - H(X^N)] = \delta_\infty$
(i.e., $\mathbf{E}_{\mathrm{CF}}$ corresponds to $\delta_\infty$).
The two are distinct non-negative quantities: $E$ is the per-position
slope of TC and determines whether $\mathrm{TC}(D_N) = \Theta(N)$,
while $\delta_\infty = \mathbf{E}_{\mathrm{CF}}$ is an asymptotic
correction constant and does not change the leading $N$-scaling.

\textbf{Lemma 5.6d (Linear $\mathrm{TC}$ for hidden Markov models).}
\emph{Let $(Y_n)_{n \geq 1}$ be a stationary Markov chain on a
finite hidden state space $\mathcal{Y}$ with transition matrix $P$
and stationary distribution $\pi$, and let $(X_n)_{n \geq 1}$ be
observations on a finite alphabet $\mathcal{X}$ conditionally
independent given $(Y_n)$ with emission laws
$p_y(x) = \Pr(X_n = x \mid Y_n = y)$. The observation process
$(X_n)$ is strictly stationary; Lemma 5.6c applies with
$E_{\mathrm{HMM}} := H(X_1) - h_{\mathrm{HMM}}$, where}
\[
H(X_1) = H\!\left(\sum_{y \in \mathcal{Y}} \pi(y) p_y\right)
\quad\text{and}\quad
h_{\mathrm{HMM}} = \lim_{n \to \infty} H(X_n \mid X_1, \ldots, X_{n-1})
\]
\emph{is the entropy rate of the observation process. The
filter recursion of Baum--Petrie (1966) /
Birch (1962) gives}
\[
h_{\mathrm{HMM}} = \lim_{n \to \infty} \mathbb{E}_{X^{n-1}}
\bigl[H(X_n \mid \mu_n)\bigr],
\]
\emph{where $\mu_n(y) := \Pr(Y_n = y \mid X_1, \ldots, X_{n-1})$ is
the prediction filter and the conditional entropy is computed
under $X_n \sim \sum_y \mu_n(y) p_y$. The HMM has
$E_{\mathrm{HMM}} > 0$ if and only if the observation process
$(X_n)$ is not iid (under its own marginal law, with $Y$ marginalized
out). Equivalently, letting $K_{y,x} := p_y(x)$ denote the emission
matrix, $E_{\mathrm{HMM}} > 0$ iff the observable prediction
$q_n := \mu_n K \in \Delta(\mathcal{X})$ does not stay almost surely
equal to the stationary marginal output law $m := \pi K$. Sufficient
conditions for $E_{\mathrm{HMM}} > 0$ are subtle: neither
``hidden chain non-trivial'' nor ``emissions distinguish states''
alone suffices, nor their conjunction (counterexample: hidden state
$(A_n, B_n)$ with $A_n$ a persistent Markov chain, $B_n$ iid, and
emission distinct for all four state pairs but
$\mathbb{E}_{B_n}[p_{(A_n, B_n)}] = \text{const in } A_n$, yielding
$X_n$ iid despite both conditions). The cleanest necessary-and-sufficient
condition is that the observation process is not iid; verifying this
from the HMM parameters $(P, \pi, \{p_y\})$ reduces to checking
that the observable prediction $\mu_n K$ moves under the
filter recursion.}

\emph{\textbf{Proof.}} Stationarity of $(X_n)$: from stationarity
of $(Y_n)$, the joint distribution of any block
$(Y_{j+1:j+N}, X_{j+1:j+N})$ is shift-invariant in $j$;
marginalizing over the hidden states preserves stationarity in
$(X_n)$. Apply Lemma 5.6c with this $E$ and $h(X) = h_{\mathrm{HMM}}$.

The entropy rate $h_{\mathrm{HMM}}$ formula in terms of the
prediction filter is standard (Baum \& Petrie 1966
\emph{Statistical Inference for Probabilistic Functions of Finite
State Markov Chains}, Annals of Mathematical Statistics 37(6):
1554--1563; Birch 1962 \emph{Approximations for the Entropy for
Functions of Markov Chains}, Annals of Mathematical Statistics
33(3): 930--938). The filter
recursion is $\mu_{n+1}(y) = \frac{1}{Z}
\sum_{y' \in \mathcal{Y}} \mu_n(y') P_{y' y} p_{y'}(X_n)$,
where $Z$ normalizes. The entropy rate
is the asymptotic per-position entropy of the prediction-filter
output under the stationary measure of $(\mu_n)$.

Informativeness: $E_{\mathrm{HMM}} = H(X_1) - h_{\mathrm{HMM}} \geq 0$ in
all cases (Cover \& Thomas 2006 Theorem 4.2.2). Equality
$E_{\mathrm{HMM}} = 0$ holds iff $H(X_n \mid X_{<n}) = H(X_n)$ for all $n$,
i.e., the observed process $(X_n)$ is iid (in a strictly stationary
process, this also forces independence at all block lengths). This can
occur in several non-equivalent ways: (a) all emission laws $\{p_y\}$
are identical, making $X_n$ iid given $Y_n$ and hence marginally;
(b) the hidden chain $(Y_n)$ is iid (i.e., all rows of $P$ in the
support of $\pi$ are equal to $\pi$); or (c) the hidden chain has non-trivial temporal structure
in components that average out in the observed prediction (e.g.,
$Y = (A, B)$ with $A$ a persistent Markov chain and $B$ an iid
``filler'' that washes out $A$'s information through emission averaging).
The cleanest necessary-and-sufficient condition for
$E_{\mathrm{HMM}} = 0$ is that the observable prediction
$\mu_n K \equiv \pi K = m$ almost surely (the filter may still move
within the affine fiber $\{\mu : \mu K = m\}$, but the observable output
prediction is stuck at the stationary marginal). \hfill$\blacksquare$

\emph{Example: binary-state binary-observation HMM.} Take
$\mathcal{Y} = \mathcal{X} = \{0, 1\}$, symmetric hidden chain
with transition probability $p \in (0, 1/2)$ (so the hidden state
flips with probability $p$ at each step), and emission $p_y =
\mathrm{Bern}(q_y)$ with $q_0 = \alpha$ and $q_1 = 1 - \alpha$
for $\alpha \in (0, 1/2)$. The stationary distribution of $Y$ is
uniform; the marginal of $X_n$ is $\mathrm{Bern}(1/2)$, so
$H(X_1) = 1$. The entropy rate $h_{\mathrm{HMM}}$ does not admit a
closed form for general $(p, \alpha)$ but is computable by
filter-recursion Monte Carlo. Selected values verified in
\texttt{scripts/verify/lemma\_5\_6d\_hmm\_tc.py} using 50 chains of
length $10^4$:
$h_{\mathrm{HMM}}(p = 0.1, \alpha = 0.2) \approx 0.910$ bits,
so $E_{\mathrm{HMM}} \approx 0.090$ bits/position.
$h_{\mathrm{HMM}}(p = 0.3, \alpha = 0.3) \approx 0.997$ bits,
so $E_{\mathrm{HMM}} \approx 0.003$ bits/position. These are
strictly between the iid-extreme ($E = 0$ when emission carries
no information about state) and the deterministic-emission case
($E = H_{\mathrm{marg}} - h(Y) = 1 - h(Y)$ for binary symmetric
emissions onto a mixing hidden chain, with $E$ approaching $1$
only in the non-mixing limit where $h(Y) \to 0$).

\emph{Scope.} Lemma 5.6d completes Open Problem 5's resolution for
the four structurally cleanest natural-data patterns:
combinatorial column constraints (5.6a), shared latent variable
(5.6b), temporal correlation (5.6c), and combined hidden-and-temporal
(5.6d). Lemma 5.6e (below) extends to the continuous-hidden-state
finite-observation case via filter-stability theory. Further classes
(factor graphs over expanders, general non-stationary processes)
admit similar treatments but with strictly more involved bookkeeping
or non-trivial functional-analytic content; they remain open as
explicit lemmas.

\textbf{Lemma 5.6e (Linear $\mathrm{TC}$ for continuous-hidden-state
finite-observation HMMs under geometric ergodicity).}
\emph{Let $(Y_n)_{n \geq 1}$ be a strictly stationary Markov chain on
a Polish state space $(\mathcal{Y}, \mathcal{B})$ with transition
kernel $P(y, dy')$ and stationary measure $\pi$, and let
$(X_n)_{n \geq 1}$ be observations on a \emph{finite} alphabet
$\mathcal{X}$, conditionally independent given $(Y_n)$ with
emission probabilities $p_y(x) = \Pr(X_n = x \mid Y_n = y)$.
Assume:}
\begin{itemize}
\item \emph{(GE) [Geometric ergodicity of the hidden chain]
$\|P^n(y, \cdot) - \pi\|_{\mathrm{TV}} \leq C_{\mathrm{erg}} \cdot
\rho_Y^n$ for some constants $C_{\mathrm{erg}} > 0$ and $\rho_Y \in
(0, 1)$, uniformly in $y \in \mathcal{Y}$.}
\item \emph{(EP) [Positive emission density] $p_y(x) \geq \rho_{\min} > 0$
for every $y \in \mathcal{Y}$ and every $x \in \mathcal{X}$.}
\end{itemize}
\emph{Then the observed process $(X_n)$ is strictly stationary on the
finite alphabet $\mathcal{X}$, satisfies the exponential
conditional-MI decay (5.8) of Lemma 5.1b (with constants
$C', \rho'_X \in (0, 1)$ inherited via filter stability from
$C_{\mathrm{erg}}, \rho_Y, \rho_{\min}$), and Lemma 5.6c (b) applies
to give}
\[
\mathrm{TC}(D_N) \;=\; N \cdot E_{\mathrm{HMM}}^{\mathrm{cont}}
                       - \delta_\infty^{\mathrm{cont}} + o(1),
\]
\emph{where $E_{\mathrm{HMM}}^{\mathrm{cont}} := H(X_1) -
h_{\mathrm{HMM}}^{\mathrm{cont}} \geq 0$ is the per-position excess
entropy, $h_{\mathrm{HMM}}^{\mathrm{cont}} := \lim_n H(X_n \mid X_{<n})$
is the entropy rate computed via the measure-valued prediction
filter $\mu_n := \Pr(Y_n \in \cdot \mid X_1, \ldots, X_{n-1})$ (the
filter recursion of Stratonovich 1960 / Baum--Petrie 1966 generalised
to Polish state space; see Capp{\'e}--Moulines--Ryd{\'e}n 2005,
\emph{Inference in Hidden Markov Models}, Ch.\ 4), and
$\delta_\infty^{\mathrm{cont}} = \sum_{n=1}^\infty
[H(X_n \mid X_{<n}) - h_{\mathrm{HMM}}^{\mathrm{cont}}] < \infty$
under (GE) and (EP) by filter stability.}

\emph{\textbf{Proof.}} \emph{Stationarity.} From stationarity of
$(Y_n)$ and conditional independence of $(X_n)$ given $(Y_n)$, the
joint law of any block of $(Y, X)$ is shift-invariant; marginalising
over $Y$ preserves stationarity in $X$.

\emph{Exponential conditional-MI decay (5.8) for $(X_n)$.}
Geometric ergodicity (GE) of $(Y_n)$ combined with positive emission
(EP) implies \emph{filter stability}: for any two prior measures
$\nu, \nu'$ on $\mathcal{Y}$ and any observation sequence
$x_1, \ldots, x_n$, the resulting prediction filters
$\mu_n^{\nu}, \mu_n^{\nu'}$ satisfy
$\|\mu_n^{\nu} - \mu_n^{\nu'}\|_{\mathrm{TV}} \leq C' \cdot \rho_X^n$
for some constants $C', \rho_X \in (0, 1)$ depending on
$C_{\mathrm{erg}}, \rho_Y, \rho_{\min}$ (van~Handel 2008,
\emph{Observability and Nonlinear Filtering},
\emph{Probab.\ Theory Relat.\ Fields} 145(1--2):35--74; see also
Atar--Zeitouni 1997 \emph{Annales de l'IHP} 33(6):697--725 for the
original exponential stability bounds). Since the observed
conditional distribution $\Pr(X_n = \cdot \mid X_{<n}) = \sum_y \mu_n(y)
p_y(\cdot)$ is a Lipschitz functional of $\mu_n$ (uniformly in $n$
because $\mathcal{X}$ is finite and $p_y$ is bounded), the conditional
mutual information $I(X_n; X_{1:n-W-1} \mid X_{n-W:n-1})$ decays
exponentially in $W$: writing the conditional MI as the relative
entropy between the true conditional and the $W$-bounded
approximation, both quantities depend on the filter $\mu_n$, and the
TV-difference of filters with truncated vs.\ full history controls
the KL-divergence by Pinsker's inequality, giving
$I(\cdot) \leq C'' \cdot (\rho_X)^W$ for some $C''$ inheriting from
the filter-stability constants. Hence (5.8) holds.

\emph{Applying Lemma 5.6c (b).} Under (5.8) the conditional-entropy
excess $\sum_n [H(X_n \mid X_{<n}) - h_{\mathrm{HMM}}^{\mathrm{cont}}]$
is summable (geometric series), so Lemma 5.6c (b) applies: $\mathrm{TC}(D_N)
= N \cdot E_{\mathrm{HMM}}^{\mathrm{cont}} - \delta_\infty^{\mathrm{cont}}
+ o(1)$.

\emph{Filter-based entropy-rate formula.} The Baum--Petrie /
Birch entropy-rate identity $h = \lim_n \mathbb{E}[H(X_n \mid \mu_n)]$
extends to continuous state space by replacing the finite-state
filter recursion with the measure-valued Bayesian update
$\mu_{n+1}(dy') = \int_{\mathcal{Y}} \mu_n(dy) P(y, dy') p_{y'}(X_n) / Z$
where $Z$ normalises (Capp{\'e}--Moulines--Ryd{\'e}n 2005 §4.2). The
ergodicity-stability machinery guarantees that the limit exists and
equals $h_{\mathrm{HMM}}^{\mathrm{cont}}$.
\hfill$\blacksquare$

\emph{Scope and significance.} Lemma 5.6e closes the
``continuous-state HMM'' open subcase of OP3 listed in Lemma 5.6d's
Scope. The technical contribution is the transfer of filter-stability
theory (functional-analytic, measure-theoretic) to the Crutchfield--Feldman
summability requirement of Lemma 5.6c (b), avoiding the need for a
direct discretization argument (which would face the non-Markovianity
of bin-discretised continuous chains). The two hypotheses
(GE)+(EP) are mild and standard in the HMM literature; (GE) is
satisfied by any uniformly ergodic chain (e.g., compact state space
with positive transition density), and (EP) holds whenever emissions
do not assign zero probability to any observable symbol.

\emph{Remaining open.} The case of continuous \emph{observation}
alphabets requires differential entropy and Radon--Nikodym derivatives;
the case of polynomially (rather than exponentially) ergodic hidden
chains substitutes a weaker $\rho_X^W \to W^{-c}$ decay and still
yields linear TC scaling provided $c > 1$ (summability), but with
$\delta_\infty < \infty$ requiring tighter analysis. Both are open
as separate lemmas.

\textbf{Lemma 5.6f (Linear $\mathrm{TC}$ for asymptotically mean
stationary (AMS) processes via Cesaro time-averaging).}
\emph{Let $(X_n)_{n \geq 1}$ be a (possibly non-stationary) stochastic
process on a finite alphabet $\mathcal{X}$. Suppose it is
\emph{asymptotically mean stationary}: there exists a stationary
measure $\bar{\mu}$ on $\mathcal{X}^{\mathbb{Z}_+}$ such that for
every measurable cylinder event $A \in \sigma(X_1, \ldots, X_k)$
(any $k$),}
\[
\bar{\mu}(A) \;=\; \lim_{N \to \infty} \tfrac{1}{N} \sum_{n=1}^N
  \Pr(\sigma^n \omega \in A),
\]
\emph{where $\sigma$ is the shift map (Gray 1990 \emph{Entropy and
Information Theory}, Springer, Ch.\ 7 ``AMS Random Processes''). Then
the joint entropy $H(D_N) = H(X_1, \ldots, X_N)$ satisfies the
$\mathrm{AMS}$ Shannon--McMillan--Breiman theorem
(Barron 1985 \emph{Ann.\ Probab.}\ 13(4):1292--1303,
\emph{The Strong Ergodic Theorem for Densities}; Gray 1990 Ch.\ 8)
$\lim_N H(D_N)/N = h(\bar{\mu})$ and the average per-position
marginal entropy converges to $H_{\mathrm{marg}}(\bar{\mu})$. Hence}
\[
\mathrm{TC}(D_N) \;=\; N \cdot E_{\bar{\mu}} + o(N),
\]
\emph{where $E_{\bar{\mu}} = H_{\mathrm{marg}}(\bar{\mu}) -
h(\bar{\mu}) \geq 0$ is the AMS-limit per-position excess entropy.
In particular, $\mathrm{TC}(D_N) = \Theta(N)$ whenever
$E_{\bar{\mu}} > 0$ (equivalently, the AMS-limit process is not iid).}

\emph{\textbf{Proof.}} The total correlation is
$\mathrm{TC}(D_N) = \sum_n H(X_n) - H(D_N)$. The first sum
divided by $N$ is the average per-position marginal entropy
$\bar{H}_N := \tfrac{1}{N} \sum_n H(X_n)$. By the AMS hypothesis
applied to single-position events,
$\bar{H}_N \to H_{\mathrm{marg}}(\bar{\mu})$ as $N \to \infty$ (Gray
1990 Theorem 7.3.1; entropy is a continuous functional of cylinder
probabilities on the finite alphabet, so AMS cylinder convergence
implies entropy convergence). The joint entropy
$H(D_N)/N \to h(\bar{\mu})$ is the AMS Shannon--McMillan--Breiman
identity (Barron 1985; Gray 1990 Theorem 8.7.1). Hence
$\mathrm{TC}(D_N)/N = \bar{H}_N - H(D_N)/N \to H_{\mathrm{marg}}(\bar{\mu})
- h(\bar{\mu}) = E_{\bar{\mu}}$ as $N \to \infty$, and
$\mathrm{TC}(D_N) = N \cdot E_{\bar{\mu}}(1 + o(1)) = N \cdot E_{\bar{\mu}}
+ o(N)$.
\hfill$\blacksquare$

\emph{Scope and significance.} Lemma 5.6f extends Lemma 5.6c (c)
(Cesaro form for stationary processes) to the strictly larger class
of AMS processes, closing the ``non-stationary processes'' subcase
of OP3 listed in Lemma 5.6d's Scope. AMS includes processes that
are strictly non-stationary but converge in time average: burn-in
followed by stationarity; periodic and quasi-periodic processes;
processes with transient initial conditions converging to a stable
regime. The framework does \emph{not} cover processes without any
time-averaged limit (e.g., processes with growing variance,
non-recurrent transient processes, or content-dependent natural-image
distributions where pixel statistics vary across regions). Two
subcases of OP3 remain open: (a) genuinely non-stationary processes
without AMS structure (e.g., natural images at given resolution,
where pixel-level statistics are spatially non-stationary);
(b) factor graphs over expander families, requiring graph-theoretic
+ information-theoretic combined analysis. Lemma 5.6g (below) closes
the tree-structured factor-graph subcase of (b) via the classical
Pearl decomposition; the genuine expander case (cycles with constant
girth) remains open.

\textbf{Lemma 5.6g (Linear $\mathrm{TC}$ for tree-structured Markov
random fields via Pearl decomposition).}
\emph{Let $G = (V, E)$ be a tree on $|V| = N$ vertices ($|E| = N - 1$),
and let $X = (X_v)_{v \in V}$ be a Markov random field on $G$ with
strictly positive pairwise factors $\psi_{u, v} : \mathcal{X} \times
\mathcal{X} \to \mathbb{R}_{>0}$ for each edge $(u, v) \in E$, joint
law $p(X) = Z^{-1} \prod_{(u, v) \in E} \psi_{u, v}(X_u, X_v)$ over
finite alphabet $\mathcal{X}$. Then}
\[
\mathrm{TC}(D_N) \;=\; \sum_{(u, v) \in E} I(X_u; X_v),
\]
\emph{the total correlation equals exactly the sum of edge pairwise
mutual informations. In particular, if every edge has positive MI
$I(X_u; X_v) \geq \varepsilon > 0$, then
$\mathrm{TC}(D_N) \geq (N-1) \cdot \varepsilon = \Theta(N)$.}

\emph{\textbf{Proof.}} For a tree-structured MRF with strictly
positive joint, the joint distribution admits the symmetric
edge-marginal factorisation of Lauritzen \& Spiegelhalter (1988
\emph{J.\ Roy.\ Statist.\ Soc.\ B} 50(2):157--224 ``Local Computations
with Probabilities on Graphical Structures'') / Lauritzen (1996
\emph{Graphical Models}, Oxford Univ.\ Press, §4.4 ``Decomposable
graphs and chain graphs''); see also Pearl (1988
\emph{Probabilistic Reasoning in Intelligent Systems}, Morgan
Kaufmann §3.2) for the directed-tree Bayes-net form, equivalent on
a tree by Bayes' rule, and Wainwright \& Jordan (2008
\emph{Graphical Models, Exponential Families, and Variational
Inference}, Foundations and Trends in Machine Learning
1(1--2):1--305, Prop.\ 2.1):
\[
p(X) \;=\; \prod_{v \in V} p(X_v) \cdot \prod_{(u, v) \in E}
\frac{p(X_u, X_v)}{p(X_u) p(X_v)}.
\]
Taking logarithms and expectations:
\[
H(D_N)
\;=\; \mathbb{E}\bigl[-\log p(X)\bigr]
\;=\; \sum_v H(X_v) - \sum_{(u, v) \in E} I(X_u; X_v).
\]
Hence
\[
\mathrm{TC}(D_N) \;=\; \sum_v H(X_v) - H(D_N)
\;=\; \sum_{(u, v) \in E} I(X_u; X_v).
\]
Under positivity $I(X_u; X_v) \geq \varepsilon$ on each edge,
$\mathrm{TC}(D_N) \geq |E| \cdot \varepsilon = (N - 1) \varepsilon =
\Theta(N)$ for $\varepsilon > 0$ fixed.
\hfill$\blacksquare$

\emph{Scope.} Lemma 5.6g closes the tree-structured factor-graph
subcase of OP3 listed in Lemma 5.6d's Scope. The proof is the
classical Pearl tree-MRF identity, applied directly to the TC
definition. Three immediate corollaries:

\emph{(a) Order-$1$ Markov chain on a finite alphabet} is a special
case (path graph, $N - 1$ chain edges); the lemma recovers
$\mathrm{TC}(D_N) = (N - 1) \cdot I(X_1; X_2)$ from Lemma 5.6b's
companion identity.

\emph{(b) Balanced binary tree of depth $\log_2 N$}: edges connect
each non-leaf to its two children. Both subtypes follow.

\emph{(c) Star graph $K_{1, N-1}$}: $N - 1$ edges from one center
to all leaves. $\mathrm{TC}(D_N) = (N - 1) \cdot I(X_{\mathrm{center}};
X_{\mathrm{leaf}})$ if all leaf-center MIs equal.

\emph{Beyond trees.} For loopy graphs (cycles present), the
Pearl identity fails: the Bethe approximation
$H \approx \sum_v H(X_v) - \sum_e I(X_u; X_v)$ is no longer exact
but bounds the joint entropy from below (Pakzad--Anantharam 2002
\emph{IEEE Trans.\ Inf.\ Theory} 48(7):2061; Yedidia--Freeman--Weiss
2005 \emph{IEEE Trans.\ Inf.\ Theory} 51(7):2282 on Bethe-free-energy
bounds). For high-girth graphs (girth $g = \omega(1)$), local
tree-likeness gives Lemma 5.6g approximately, with error $O(1/g)$;
formalising this is the natural next step toward closing the
``factor graphs over expanders'' subcase of OP3 fully. Lemma 5.6h
(below) closes the orthogonal subcase of chordal graphs with bounded
treewidth via the junction-tree generalisation.

\textbf{Lemma 5.6h (Linear $\mathrm{TC}$ for graphical models on
chordal graphs with bounded treewidth via junction-tree decomposition).}
\emph{Let $G = (V, E)$ be a chordal graph on $|V| = N$ vertices
(every induced cycle of length $\geq 4$ has a chord) with treewidth
$w = O(1)$, and let $\mathcal{J} = (\mathcal{C}, \mathcal{S})$ be a
junction tree of $G$ with maximal cliques $\mathcal{C}$ and separators
$\mathcal{S}$ (Lauritzen 1996 \emph{Graphical Models}, §4
``Decomposable graphs and chain graphs''; Wainwright \& Jordan 2008
\emph{Graphical Models, Exponential Families, and Variational
Inference}, Foundations and Trends in Machine Learning 1(1--2):1--305,
§2.5 ``Junction tree representation''). Let
$X = (X_v)_{v \in V}$ be an MRF on $G$ with strictly positive
clique factors $\psi_C : \mathcal{X}^{|C|} \to \mathbb{R}_{>0}$ for
each maximal clique $C \in \mathcal{C}$, joint law
$p(X) = Z^{-1} \prod_{C \in \mathcal{C}} \psi_C(X_C)$. Then the joint
entropy admits the junction-tree decomposition}
\[
H(D_N) \;=\; \sum_{C \in \mathcal{C}} H(X_C) - \sum_{S \in \mathcal{S}} H(X_S),
\]
\emph{and the total correlation satisfies}
\[
\mathrm{TC}(D_N) \;=\; \sum_v H(X_v) - \sum_{C \in \mathcal{C}} H(X_C)
                       + \sum_{S \in \mathcal{S}} H(X_S).
\]
\emph{\textbf{Proof.}} For chordal $G$ with junction tree $\mathcal{J}$,
the joint distribution factorises as
\[
p(X) \;=\; \prod_{C \in \mathcal{C}} p(X_C) \,\Big/\, \prod_{S \in \mathcal{S}} p(X_S)
\]
by the junction-tree running-intersection property
(Lauritzen 1996 Theorem 4.7; Wainwright--Jordan 2008 Prop.\ 2.2).
Taking $-\log$ and expectation:
\[
H(D_N) \;=\; \sum_{C} H(X_C) - \sum_{S} H(X_S).
\]
Subtracting from $\sum_v H(X_v)$ gives the stated $\mathrm{TC}$
decomposition.
\hfill$\blacksquare$

\emph{Honest caveat: clique-wise sum simplification FAILS for $w \geq 2$.}
A natural conjecture extending Lemma 5.6g would be
$\mathrm{TC}(D_N) = \sum_C \mathrm{TC}(X_C)$, the sum of intra-clique
total correlations. \emph{Numerical testing on chordal graphs with
treewidth $w = 2$ (interval graphs, $2$-tree paths, fan graphs)
shows this identity FAILS}: the discrepancy
$\sum_C \mathrm{TC}(X_C) - \mathrm{TC}(D_N) = \sum_S [\,\sum_{v \in S}
H(X_v) - H(X_S)\,] + \mathrm{multiplicity\ corrections}$ is non-zero
for non-independent separators (each separator $S$ with $|S| \geq 2$
contributes its own intra-separator $\mathrm{TC}$ to the
discrepancy). For trees ($w = 1$, separators are single vertices),
intra-separator $\mathrm{TC}$ vanishes and the simplification holds
(Lemma 5.6g); for $w \geq 2$ it does not. The verify script
\texttt{scripts/verify/lemma\_5\_6h\_chordal\_junction\_tree\_tc.py}
demonstrates this for $w = 2$ instances with discrepancies of order
$10^{-2}$ to $10^{-1}$ bits.

\emph{Scope and significance.} Lemma 5.6h establishes the
\emph{junction-tree entropy decomposition} for chordal MRFs as the
correct generalisation of Pearl's tree formula:
$H(D_N) = \sum_C H(X_C) - \sum_S H(X_S)$ holds for arbitrary chordal
$G$. The corresponding $\mathrm{TC}$ identity
$\mathrm{TC}(D_N) = \sum_v H(X_v) - \sum_C H(X_C) + \sum_S H(X_S)$
is exact. For trees this collapses to Lemma 5.6g
$\mathrm{TC}(D_N) = \sum_e I(X_u; X_v)$; for higher-treewidth
chordal graphs, $\mathrm{TC}$ admits the decomposition but \emph{not}
the clique-wise simplification. The linear scaling
$\mathrm{TC}(D_N) = \Theta(N)$ for chordal graphs with bounded
treewidth and positive-correlation factors follows from the bounded-
alphabet $H(X_C) \leq |C| \log |\mathcal{X}|$ and $\Theta(N)$ cliques,
but determining the exact constant requires graph-specific analysis
of separator interactions.

\emph{Remaining open subcase.} For non-chordal (loopy, non-triangulated)
graphs, the junction-tree decomposition fails: triangulation
necessarily adds fill-in edges with no MRF factors, and the resulting
``virtual cliques'' produce Bethe-style approximation rather than
exact identity. The natural-data subcase ``factor graphs over
expanders'' (loopy random graphs with girth $\omega(1)$) still
requires Pakzad--Anantharam / Yedidia--Freeman--Weiss Bethe-bound
analysis. Closing this subcase fully is open.

\textbf{Lemma 5.6i ($\mathrm{TC}$ scaling for combinatorial uniform
classes; refinement of Remark 5.6).}
\emph{Let $D_N$ be the uniform distribution on a combinatorial subset
$\mathcal{C}$ of an ambient finite product space defined by a global
constraint. The scaling of $\mathrm{TC}(D_N)$ in $N$ differs
qualitatively across natural classes:}
\begin{enumerate}
\item \emph{(K-sparse uniform: sub-linear.) For $\mathcal{C} = \{x \in
\{0,1\}^N : \sum_i x_i = K\}$ (uniform on $K$-sparse binary vectors):
\[
\mathrm{TC}(D_N) \;=\; N \cdot h(K/N) - \log_2 \binom{N}{K}.
\]
Asymptotic regimes (Stirling expansion of $\log_2 \binom{N}{K}$):}
\begin{itemize}
\item \emph{Fixed $K$, $N \to \infty$: $\mathrm{TC}(D_N) \to \log_2(K!) -
K \log_2 K + K \log_2 e$, a constant in $N$. For $K = 1$ this
equals $\log_2 e \approx 1.44$; for $K = 10$ it equals $\approx
3.00$; asymptotically $\approx \tfrac{1}{2} \log_2(2 \pi K)$ for
large $K$.}
\item \emph{$K = \alpha N$ with $\alpha \in (0, 1)$ constant:
$\mathrm{TC}(D_N) = \tfrac{1}{2} \log_2(2 \pi N \alpha (1-\alpha))
+ O(1/N) = \tfrac{1}{2} \log_2 N + O(1)$.}
\end{itemize}
Both regimes are sub-linear: $\mathrm{TC}(D_N) = o(N)$.
\item \emph{(Uniform permutations: linear --- counterexample.) For
$\mathcal{C} = \{$permutations of $[N]\}$ embedded in $[N]^N$ as
$x = (\pi(1), \ldots, \pi(N))$ (uniform on permutations):
\[
\mathrm{TC}(D_N) \;=\; N \log_2 N - \log_2(N!) \;=\; N \log_2 e - O(\log N),
\]
which is linear in $N$ with constant $\log_2 e \approx 1.44$ ---
this class is $\Theta(N)$. The distinct-value constraint $X_i \neq
X_j$ for all $i \neq j$ contributes pairwise mutual information
$I(X_i; X_j) = \log_2(N) - \log_2(N-1) = \Theta(1/N)$ for each of
$\binom{N}{2} = \Theta(N^2)$ pairs, summing to $\Theta(N)$.}
\end{enumerate}

\emph{Scope and significance.} Lemma 5.6i refines Remark 5.6's
informal ``low-dimensional manifold $\Rightarrow$ $\mathrm{TC}(D_N) =
\Theta(N)$'' intuition by exhibiting a natural class where cardinality
or sparsity constraints --- which DO restrict the support to a
``low-dimensional'' subset $\mathcal{C}$ with $|\mathcal{C}| =
2^{o(N)}$ --- yield sub-linear $\mathrm{TC}$. The structural reason
is that the $K$-cardinality constraint is a \emph{higher-order
combinatorial dependence} (each pair has nearly-independent marginals
$P(X_i = 1, X_j = 1) = K(K{-}1)/(N(N{-}1)) \approx (K/N)^2$, giving
$O(1/N^2)$ pairwise mutual information). Although the joint
distribution carries a single $O(K \log(N/K))$-bit constraint, this
constraint distributes across $\binom{N}{2}$ pairs as tiny pairwise
MIs that do not accumulate linearly. The contrast with uniform
permutations (heavy pairwise distinct-value constraint, $\Theta(1)$
per pair, $\Theta(N^2)$ total) makes the structural mechanism
explicit: $\Theta(N)$ $\mathrm{TC}$ requires \emph{either} constant
pairwise MI for a $\Theta(N)$-pair subset (Markov, HMM, tree-MRF),
\emph{or} the heavy combinatorial distinct-value constraint of
permutations. Pure sparsity is insufficient.

\emph{\textbf{Proof.}} For (1): the joint entropy of uniform on
$\binom{N}{K}$ supports equals $\log_2 \binom{N}{K}$. The marginal of
each coordinate $X_i$ is $\mathrm{Bern}(K/N)$ by symmetry, with
entropy $h(K/N)$. By definition $\mathrm{TC}(D_N) = \sum_i H(X_i) -
H(X) = N h(K/N) - \log_2 \binom{N}{K}$, the exact identity.

\emph{Sub-regime: $K = \alpha N$ with $\alpha \in (0, 1)$ constant.}
Stirling's approximation for the binomial coefficient gives
$\log_2 \binom{N}{\alpha N} = N h(\alpha) - \tfrac{1}{2} \log_2(2 \pi
N \alpha(1-\alpha)) + O(1/N)$. Subtracting from $N h(K/N) = N h(\alpha)$:
\[
\mathrm{TC}(D_N) \;=\; \tfrac{1}{2} \log_2(2 \pi N \alpha (1-\alpha))
+ O(1/N) \;=\; \tfrac{1}{2} \log_2 N + \tfrac{1}{2} \log_2(2 \pi \alpha(1-\alpha))
+ O(1/N).
\]

\emph{Sub-regime: $K$ fixed, $N \to \infty$.} Use the exact Stirling
expansion for the binomial coefficient with fixed numerator-of-falling-
factorial parameter: $\binom{N}{K} = N^K / K! \cdot (1 + O(K^2/N))$, so
$\log_2 \binom{N}{K} = K \log_2 N - \log_2 K! + O(K^2/N)$. For the
entropy: $h(K/N) = -(K/N) \log_2(K/N) - (1-K/N) \log_2(1-K/N)$. The
Taylor expansion of $h$ around $K/N = 0$ gives
$h(K/N) = (K/N) \log_2(N/K) + (K/N) \log_2 e + O((K/N)^2)$, hence
$N h(K/N) = K \log_2(N/K) + K \log_2 e + O(K^2/N)$. Subtracting:
\[
\mathrm{TC}(D_N) \;=\; \bigl[ K \log_2(N/K) + K \log_2 e \bigr] -
\bigl[ K \log_2 N - \log_2 K! \bigr] + O(K^2/N).
\]
The $K \log_2 N$ terms cancel: $K \log_2(N/K) - K \log_2 N = -K \log_2 K$.
This yields
\[
\mathrm{TC}(D_N) \;=\; \log_2 K! - K \log_2 K + K \log_2 e + O(K^2/N).
\]
The leading term $\log_2 K! - K \log_2 K + K \log_2 e$ is a constant
in $N$ (depending only on $K$). For $K = 1$: $0 - 0 + \log_2 e \approx
1.44$. For $K = 10$: $\log_2(10!) - 10 \log_2 10 + 10 \log_2 e \approx
21.79 - 33.22 + 14.43 \approx 3.00$. Asymptotically large $K$ (using
Stirling $\log_2 K! = K \log_2 K - K \log_2 e + \tfrac{1}{2} \log_2(2 \pi K)
+ O(1/K)$): the formula simplifies to $\mathrm{TC}(D_N) \to \tfrac{1}{2}
\log_2(2 \pi K) + O(1/K)$.

In either sub-regime, $\mathrm{TC}(D_N) = o(N)$.

For (2): uniform on permutations has $H(X) = \log_2(N!)$ and each
$X_i$ marginal is uniform on $[N]$ with $H(X_i) = \log_2 N$.
$\mathrm{TC}(D_N) = N \log_2 N - \log_2(N!) = N \log_2 e - \tfrac{1}{2}
\log_2(2 \pi N) + O(1/N)$ by Stirling, giving $\Theta(N)$ scaling
with leading constant $\log_2 e$. \hfill$\blacksquare$

\emph{Remark 5.6i-a (TC vs.\ sum of pairwise MIs --- three regimes
revealed by combinatorial uniform vs.\ 2D Ising vs.\ tree-MRF).} The
naive attempt to relate $\mathrm{TC}(D_N)$ and $\sum_{i < j} I(X_i;
X_j)$ admits NEITHER a uniform upper bound NOR a uniform lower bound
in general. Three regimes emerge from explicit empirical computation
(verify scripts \texttt{lemma\_5\_6i\_tc\_combinatorial\_sublinear.py}
and \texttt{empirical\_tc\_2d\_ising\_pairwise.py}):
\begin{enumerate}
\item \emph{Higher-order combinatorial-constraint regime
($\mathrm{TC} \gg$ sum-pair).} For $K$-sparse uniform on
$\{0,1\}^N$ with $K \ll N$: the global cardinality constraint
distributes across pairs as tiny mutual informations
($I(X_i; X_j) = O(1/N^2)$), so $\sum I$ is small. But the joint
constraint accumulates via the chain rule to non-trivial
$\mathrm{TC}$. Concrete: $N = 100$, $K = 10$ gives $\sum I \approx
0.37$ vs $\mathrm{TC} \approx 2.92$ (ratio $\approx 8\times$).
\item \emph{Tree / no-redundancy regime ($\mathrm{TC} = $ sum-pair).}
For tree-structured MRFs (Lemma 5.6g):
$\mathrm{TC}(D_N) = \sum_{e \in E_{\rm tree}} I(X_u; X_v)$, exact
equality. Each edge contributes its full MI because there is no
cyclic redundancy in tree topology.
\item \emph{Cyclic-redundancy regime ($\mathrm{TC} \ll$ sum-pair).}
For 2D Ising on the square lattice at low temperature (verified
empirically for $n = 3, 4, 5$): pairs are heavily correlated
(especially in the ordered phase, $T < T_c$), but the pairwise
informations REDUNDANTLY count the shared global magnetisation.
Concrete: $n = 5$, $T = 0.5 T_c$ gives $\sum I \approx 246$ vs
$\mathrm{TC} \approx 22$ (ratio $\approx 0.09$, nine-fold
REDUNDANCY). The ratio approaches 1 at high temperature ($T = 2
T_c$ gives $\sum I \approx 1.96$ vs $\mathrm{TC} \approx 1.53$,
ratio $\approx 0.78$) when pairwise correlations become weak and
non-overlapping.
\end{enumerate}

The correct universal decomposition is the multi-information chain
rule $\mathrm{TC}(D_N) = \sum_{k=2}^{N} I(X_k; X_{1:k-1})$ (a
telescoping sum capturing the full joint dependence including
higher-order interactions and cyclic-redundancy effects). The naive
pairwise sum $\sum_{i<j} I(X_i; X_j)$ neither upper-bounds nor
lower-bounds $\mathrm{TC}$ in general --- the relationship depends
on the distribution's topology of dependence.

\emph{Remark 5.6i-c (correlation-length signatures from 2D Ising
pairwise MI vs.\ distance).} The pairwise MI as a function of L1
lattice distance, computed empirically for $n = 5$ at three
temperatures (\texttt{empirical\_tc\_2d\_ising\_pairwise.py}), shows
the standard statistical-physics correlation-length signatures:
\begin{itemize}
\item $T = 0.5 T_c$ (ordered): mean pairwise MI decays slowly,
from $\approx 0.88$ at $d = 1$ to $\approx 0.58$ at $d = 8$
(log-decay rate $\approx -0.04$ per unit distance, near-constant
correlations characteristic of long-range order).
\item $T = T_c$ (critical): power-law-like decay, from $\approx
0.23$ at $d = 1$ to $\approx 0.004$ at $d = 8$ (log-decay rate
$\approx -0.5$ to $-0.7$ per unit, consistent with scale-invariant
critical fluctuations).
\item $T = 2 T_c$ (disordered): exponential decay (log-decay rate
$\approx -2$ per unit distance), confirming finite correlation
length above $T_c$.
\end{itemize}
These signatures connect compressibility (via $\mathrm{TC}$ and its
pairwise decomposition) to the physical correlation structure: the
RNR compression advantage on 2D-Ising-like sources is dominated by
nearest-neighbor pairs at all temperatures, but the absolute
magnitude of compressible information is largest at low $T$
(ordered phase) and smallest at high $T$ (independent-spin limit).

\emph{Implications for the RNR framework.} (i) The compression
advantage of joint coding over factorised coding (Theorem 5.5)
scales with $\mathrm{TC}(D_N)$, not with the source's combinatorial
complexity or support cardinality. Sparse-support data classes
(e.g., binary indicators of rare events, compressed-sensing signals,
sparse activations) admit relatively limited joint-coding advantage
relative to a well-designed factorised coder. (ii) The data classes
that genuinely benefit from joint coding are those with
\emph{pairwise} or \emph{local-clique} dependence (Markov chains,
HMMs, tree-MRFs, junction-tree-decomposable graphs of bounded
treewidth), where $\mathrm{TC}$ is $\Theta(N)$ via $\Theta(N)$
contributing pairs.

\textbf{Lemma 5.6j ($\mathrm{TC}$ for the quantized
Gaussian-on-rank-$d$-manifold model; quantitative refinement of
Remark 5.6).}
\emph{Let $\sigma > 0$, $d \in \{1, \dots, N\}$, and let $U \in
\mathbb{R}^{N \times d}$ have orthonormal columns ($U^T U = I_d$). Let
$\Lambda = \mathrm{diag}(\lambda_1, \dots, \lambda_d) \succeq 0$, write
$\bar\lambda = \sum_j \lambda_j$, and define the signal-to-noise ratio
per ambient coordinate as $\rho = \bar\lambda / (N \sigma^2)$. Consider
the centered Gaussian model $X \sim \mathcal{N}(0, \Sigma)$ on
$\mathbb{R}^N$ with}
\[
\Sigma \;=\; U \Lambda U^T + \sigma^2 I_N,
\tag{5.6j.1}
\]
\emph{and let $X_\Delta$ be the coordinate-wise lattice quantization of
$X$ at scale $\Delta > 0$. Set $a_i := (U \Lambda U^T)_{ii}/\sigma^2$
($i = 1, \dots, N$, so $\sum_i a_i = \bar\lambda/\sigma^2 = N \rho$) and
$b_j := \lambda_j / \sigma^2$ ($j = 1, \dots, d$, so $\sum_j b_j = N \rho$).
Define the per-coordinate embedding score
$S_{\rm emb}(U, \Lambda; \sigma^2) := \tfrac{1}{N} \sum_{i=1}^N
\log_2(1 + a_i)$. Then:}
\begin{enumerate}
\item \emph{(Exact identity.) The differential total correlation of
$X$ satisfies}
\[
\mathrm{TC}(X) \;=\; \tfrac{1}{2}\!\left[\,\sum_{i=1}^N \log_2(1 + a_i)
\,-\, \sum_{j=1}^d \log_2(1 + b_j)\,\right] \quad \text{(bits)}.
\tag{5.6j.2}
\]
\item \emph{(Embedding lower bound.) By Jensen's inequality applied
to $\log(1 + b_j)$ subject to $\sum_j b_j = N\rho$, $\sum_j \log_2(1 +
b_j) \leq d \log_2(1 + N \rho / d)$, hence}
\[
\mathrm{TC}(X) \;\geq\; \tfrac{N}{2}\,S_{\rm emb}(U, \Lambda; \sigma^2)
\;-\; \tfrac{d}{2} \log_2\!\bigl(1 + N\rho/d\bigr).
\tag{5.6j.3}
\]
\emph{The Jensen slack is zero when $\lambda_1 = \dots = \lambda_d =
\bar\lambda/d$ and at most $d \log_2(1 + N\rho/d) - \sum_j \log_2(1 +
b_j)$ in general.}
\item \emph{(Asymptotic linear scaling under Haar-uniform $U$.) If $U$
is sampled Haar-uniformly from the Stiefel manifold $V_d(\mathbb{R}^N)
= \{U \in \mathbb{R}^{N \times d} : U^T U = I_d\}$ and $\lambda_1 =
\dots = \lambda_d = \bar\lambda/d$, then as $N \to \infty$ with $d$ and
$\rho$ held fixed,}
\[
S_{\rm emb}(U, \Lambda; \sigma^2) \;\xrightarrow{a.s.}\; c_d(\rho)
\;:=\; \mathbb{E}_{V \sim \chi^2_d / d}\bigl[\log_2(1 + \rho V)\bigr],
\tag{5.6j.4}
\]
\emph{and consequently}
\[
\frac{\mathrm{TC}(X)}{N} \;\xrightarrow{a.s.}\; \tfrac{1}{2}\, c_d(\rho)
\;>\; 0 \quad \text{whenever } \rho > 0.
\tag{5.6j.5}
\]
\emph{The function $c_d(\rho)$ is positive, increasing in $\rho$, and
non-decreasing in $d$, with $c_d(\rho) \to \log_2(1+\rho)$ as $d \to
\infty$. A second-order Taylor expansion in $1/d$ gives the explicit
estimate}
\[
c_d(\rho) \;=\; \log_2(1+\rho) \;-\; \frac{\rho^2}{(\ln 2)(1+\rho)^2 d}
\;+\; O(1/d^2).
\tag{5.6j.6}
\]
\emph{Numerical values for $\rho = 1$ (Monte Carlo, $10^6$ samples):
$c_1(1) \approx 0.770$, $c_2(1) \approx 0.860$, $c_4(1) \approx 0.922$,
$c_{16}(1) \approx 0.978$, $c_{64}(1) \approx 0.995$, with limit
$\log_2 2 = 1$.}
\item \emph{(Quantization perturbation.) Let $\tilde X = X + N_\Delta$
where $N_\Delta$ has independent coordinates uniform on $[-\Delta/2,
\Delta/2]$, modeling Bennett's high-resolution quantization noise
(Gersho--Gray 1992, Theorem 5.5.1; Bennett 1948). The discrete TC of
the lattice-quantized $X_\Delta$ equals the differential TC of $\tilde
X$ (the $N \log_2 (1/\Delta)$ entropy shifts cancel in the
$\sum H_{\rm marg} - H_{\rm joint}$ subtraction). The covariance of
$\tilde X$ is $\Sigma + (\Delta^2/12) I_N$. Using $\mathrm{TC}(\tilde
X) \leq \mathrm{TC}_{\rm Gauss}(\Sigma + (\Delta^2/12) I_N)$ (Gaussian
maximizes differential entropy for fixed covariance, applied to the
joint), then applying (5.6j.2) to $\Sigma + (\Delta^2/12) I_N$ and
Taylor-expanding around $\Delta^2 = 0$,}
\[
\bigl|\,\mathrm{TC}_{\rm Gauss}(\Sigma + (\Delta^2/12) I_N)
- \mathrm{TC}(X)\,\bigr| \;\leq\; \frac{N \Delta^2}{24 \, \sigma^2 \,
\ln 2}.
\tag{5.6j.7}
\]
\emph{For $\Delta \leq \sigma / 4$ the perturbation is at most $N /
(384 \ln 2) \approx 0.0038 N$, so the Gaussian-bounded discrete TC is
also $\Theta(N)$ with leading constant $\tfrac{1}{2} c_d(\rho) +
O(\Delta^2/\sigma^2)$. (Note the upper-bound direction: for the
lower-bound on $\mathrm{TC}(X_\Delta)$ that would close the empirical
question of \S 11, one needs either a Gaussian-matching argument
specific to the source $X$ or a direct discrete-entropy calculation.)}
\item \emph{(Delocalization is necessary.) If $U$ concentrates on a
proper subset $S \subset \{1, \dots, N\}$ of size $|S| < N$ (i.e.,
$U_{ij} = 0$ for $i \notin S$), then $a_i = 0$ for $i \notin S$ and
$N \cdot S_{\rm emb} = \sum_{i \in S} \log_2(1 + a_i) \leq |S|
\log_2(1 + N\rho/|S|)$. In the extreme case $U = [I_d ; 0]$ (signal
confined to the first $d$ coordinates), $\Sigma$ is block-diagonal and
$\mathrm{TC}(X) = 0$ for any choice of $\Lambda$.}
\end{enumerate}

\emph{Scope and significance.} Lemma 5.6j gives the first
quantitative $\Theta(N)$ TC bound for a Remark-5.6-style manifold
model with \emph{explicit} constants depending on the geometric
parameters $(d, \sigma, \bar\lambda)$. The bound is sharp up to the
Jensen slack on $\lambda_j$ in (5.6j.3) and the finite-$N$ deviation
of $S_{\rm emb}$ from $c_d(\rho)$ in (5.6j.4). The Haar hypothesis in
(3) plays the role of the ``delocalisation / spread'' assumption
identified in Remark 5.6 as the missing ingredient beyond
intrinsic-dimension control; the counter-example in (5) shows it
cannot be dropped --- low intrinsic dimension alone is compatible
with $\mathrm{TC} = 0$. The connection to a sub-manifold $M \subset
\mathbb{R}^N$ with reach $\tau$ and density bounded above and below
(Niyogi--Smale--Weinberger 2008; Levina--Bickel 2004) is via local
linearisation: in a tube of radius $\sigma \ll \tau$ around any base
point of $M$, the data are well-approximated by their tangent-Gaussian
model with $U$ spanning the tangent space and $\Lambda$ the local
intrinsic covariance, so Lemma 5.6j applies whenever the local
tangent embedding satisfies the delocalisation hypothesis.

\emph{Caveat (real natural-data classes).} Lemma 5.6j is a
\emph{conditional} closure: for any natural class for which (a) the
sub-manifold has bounded reach $\tau$, (b) the on-manifold density is
comparable to uniform within a constant factor (Lipschitz density
bounded by $\rho_{\rm max}/\rho_{\rm min}$), (c) $\sigma \ll \tau$, and
(d) the tangent directions of $M$ are well-spread across ambient
coordinates, $\mathrm{TC}(D_N) = \Theta(N)$ with the explicit leading
constant of (5.6j.5). For specific natural-image classes the
delocalization hypothesis (d) is empirical: local-patch features
(textures, edges) may concentrate on small ambient subsets in a way
that violates (d). Whether full natural-image distributions on ambient
dimension $N \in \{10^4, 10^5, 10^6\}$ satisfy (d) at a non-trivial
scale is the open empirical question of \S 11 / \S 13.4.5; Lemma 5.6j
shows that \emph{under} (d) the manifold intuition is correct, and
identifies (d) as the precise content of ``manifold-plus-regularity''.

\emph{\textbf{Proof.}}

\noindent (1) The differential entropy of $X \sim \mathcal{N}(0,
\Sigma)$ is $h(X) = \tfrac{1}{2} \log_2((2\pi e)^N \det \Sigma)$, and
of each marginal $X_i \sim \mathcal{N}(0, \Sigma_{ii})$ is
$h(X_i) = \tfrac{1}{2} \log_2(2 \pi e \Sigma_{ii})$ (Cover--Thomas
2006, Theorem 8.4.1). Hence $\mathrm{TC}(X) = \sum_i h(X_i) - h(X)
= \tfrac{1}{2}[\sum_i \log_2 \Sigma_{ii} - \log_2 \det \Sigma]$.
By Sylvester's determinant identity,
\[
\det \Sigma = \det(\sigma^2 I_N) \det(I_d + \sigma^{-2} \Lambda)
= \sigma^{2N} \prod_{j=1}^d (1 + \lambda_j / \sigma^2).
\]
Also $\Sigma_{ii} = \sigma^2 + (U \Lambda U^T)_{ii} = \sigma^2 (1 +
a_i)$. Substituting and cancelling the common $N \log_2 \sigma^2$
term gives (5.6j.2).

\noindent (2) By concavity of $\log$, for any non-negative
$(b_1, \dots, b_d)$ with $\sum_j b_j = N\rho$, Jensen's inequality
gives $\tfrac{1}{d} \sum_j \log_2(1 + b_j) \leq \log_2(1 + N\rho/d)$,
i.e., $\sum_j \log_2(1 + b_j) \leq d \log_2(1 + N\rho/d)$, with
equality iff all $b_j$ coincide. Substituting into (5.6j.2) yields
(5.6j.3).

\noindent (3) For Haar-uniform $U$ from the Stiefel manifold
$V_d(\mathbb{R}^N)$, the random variable $\sum_{j=1}^d U_{ij}^2$ has
marginal distribution $\mathrm{Beta}(d/2, (N-d)/2)$ for each fixed
$i$ (Muirhead 1982, Theorem 2.5.5; this follows from $U = G(G^T
G)^{-1/2}$ for Gaussian $G$, and the Gram matrix law). Multiplying
by $N$ and rescaling by $1/d$:
\[
\frac{N}{d} \sum_{j=1}^d U_{ij}^2 \;\xrightarrow{d}\; V
\;:=\; \tfrac{1}{d}\chi^2_d \qquad \text{as } N \to \infty
\]
(this is the standard scaling limit of $\mathrm{Beta}(d/2,(N-d)/2)$).
With $\lambda_j = \bar\lambda/d$ constant in $j$,
$a_i = (\bar\lambda/d)\sum_j U_{ij}^2 / \sigma^2 = \rho \cdot
(N/d)\sum_j U_{ij}^2 \xrightarrow{d} \rho V$. Concentration of
measure on $V_d(\mathbb{R}^N)$ (e.g., Talagrand 1995 spectral-radius
form) gives a.s.\ convergence of the empirical average
$\tfrac{1}{N}\sum_i \log_2(1+a_i)$ to $\mathbb{E}[\log_2(1 + \rho V)]
= c_d(\rho)$, establishing (5.6j.4)--(5.6j.5).

The Taylor expansion (5.6j.6) follows from $\log_2(1 + \rho v)
= \log_2(1 + \rho) + \tfrac{\rho(v-1)}{(\ln 2)(1+\rho)}
- \tfrac{\rho^2 (v-1)^2}{2 (\ln 2)(1+\rho)^2} + O((v-1)^3)$ combined
with $\mathbb{E}[V] = 1$ and $\mathrm{Var}(V) = 2/d$, giving
$\mathbb{E}[\log_2(1+\rho V)] = \log_2(1+\rho) - \rho^2/((\ln 2)
(1+\rho)^2 d) + O(1/d^2)$.

\noindent (4) The discrete entropy of a lattice-quantized continuous
variable equals the differential entropy of the dithered variable
$\tilde X_i = X_i + N_i$ (with $N_i$ uniform on $[-\Delta/2,
\Delta/2]$, independent across $i$) plus the constant offset $N
\log_2(1/\Delta)$ on the joint and $\log_2(1/\Delta)$ per marginal
(Gersho--Gray 1992, Theorem 5.5.1; Bennett 1948); these offsets cancel
exactly in the TC subtraction $\sum H_{\rm marg} - H_{\rm joint}$, so
$\mathrm{TC}(X_\Delta) = \mathrm{TC}(\tilde X)$ in the high-resolution
limit. The covariance of $\tilde X$ is $\Sigma + (\Delta^2/12) I_N$.
By the maximum-entropy property of Gaussians at fixed covariance
(Cover-Thomas 2006, Theorem 8.6.5),
$\mathrm{TC}(\tilde X) \leq \mathrm{TC}_{\rm Gauss}(\Sigma +
(\Delta^2/12) I_N)$. Applying (5.6j.2) to this Gaussian covariance:
the perturbation to $\sum_i \log_2(1 + a_i)$ under $a_i \to a_i +
\Delta^2/(12\sigma^2)$ has gradient $\sum_i 1/((1+a_i) \ln 2) \leq
N/\ln 2$ (using $a_i \geq 0$); $\sum_j \log_2(1+b_j)$ is unchanged.
Multiplying by the factor $\Delta^2/(12\sigma^2)$ and the $1/2$
prefactor gives (5.6j.7).

\noindent (5) If $U_{ij} = 0$ for all $i \notin S$ then
$(U\Lambda U^T)_{ik} = \sum_l U_{il} \lambda_l U_{kl} = 0$ whenever
$i \notin S$ or $k \notin S$. Therefore $\Sigma$ is block-diagonal
with $|S| \times |S|$ ``signal'' block and the complementary
$(N-|S|) \times (N-|S|)$ block equal to $\sigma^2 I_{N-|S|}$. In the
extreme case $U = [I_d; 0]$, the $|S|=d$ signal block is itself
$\Lambda + \sigma^2 I_d = \mathrm{diag}(\lambda_j + \sigma^2)$ which
is diagonal, so $\Sigma$ is fully diagonal and $\mathrm{TC}(X) = 0$.
For intermediate $|S| \in (d, N)$, only the $|S|$ ambient
coordinates in $S$ carry TC; the bound (5.6j.3) applied to the
$|S|$-coordinate sub-problem gives $\mathrm{TC} = \Theta(|S|)$,
linear in $|S|$ rather than in $N$.
\hfill$\blacksquare$

\emph{Connection to literature.} The exact identity (5.6j.2) for
Gaussian TC is the case of Watanabe's (1960, eq.~21) multi-information
specialised to the multivariate-normal family, cf.~Studen\'y and
Vejnarov\'a (1998) for the general TC formula. The Stiefel-manifold
concentration in (3) is classical (Muirhead 1982; Diaconis and
Freedman 1987 give related Gaussian-projection limits). The
connection to the manifold hypothesis via local linearisation is in
the spirit of Niyogi, Smale and Weinberger (2008) for reach-based
manifold-sampling theory and the intrinsic-dimension literature
(Levina and Bickel 2004; Donoho 2000); what is new here is the
\emph{per-coordinate quantitative constant} $\tfrac{1}{2} c_d(\rho)$,
making explicit what was previously a qualitative ``$\mathrm{TC} =
\Theta(N)$'' claim. The verification script
\texttt{scripts/verify/lemma\_5\_6j\_manifold\_gaussian.py} confirms
all five claims numerically: the exact identity (1) to within
$10^{-10}$ bits, the lower bound (2) for $30$ random configurations
(mean Jensen slack $\approx 1.5$ bits for $N = 500, d = 4$, $\rho = 1$),
the asymptotic limit (3) within $0.03$ bits of $\tfrac{1}{2} c_d(\rho)$
at $N = 5000, d = 4$, the quantization perturbation (4) shrinking
quadratically with $\Delta$, and the spread counter-example (5)
giving $\mathrm{TC} = 0$ for $U = [I_d; 0]$.

\emph{Relation to OP3 (\S 13.4.5) and the broader open problem.}
Lemma 5.6j gives one explicit family within the Remark-5.6 ``natural
manifold'' framing for which $\mathrm{TC} = \Theta(N)$ rigorously holds:
the Gaussian-tangent model with delocalised embedding. It does
\emph{not} construct an enumerative encoder achieving rate $h(X)/N +
o(1)$ for this family, nor does it claim that natural images satisfy
the delocalisation hypothesis. The open problem of \S 13.4.5 remains
open as stated; Lemma 5.6j contributes one ``natural'' instance
(Gaussian-on-manifold with delocalised tangent) where the manifold
intuition is provably correct under explicit conditions, and isolates
delocalisation (item~(d) of the caveat) as the precise structural
hypothesis the informal claim of Remark 5.6 was missing.

\textbf{5.7 Worked numerical example}

Consider a block of $N = 1024$ nibbles with counts $k_v = 64$ for
$v = 0, \ldots, 15$ (i.e., uniformly balanced). The multinomial
baseline of Theorem 5.2 evaluates to
\[
\log_2 \!\left[\frac{1024!}{(64!)^{16}}\right] \approx 4033.08 \text{ bits.}
\]
The plain per-nibble raw cost is $4 \cdot 1024 = 4096$ bits ($4$
bits per nibble $\times$ $N = 1024$ nibbles, i.e., $4096$ bits per
block of $512$ bytes when nibbles are packed two per byte). A
factorized arithmetic coder under the uniform marginal pays
approximately $4096$ bits. The multinomial coder saves about $63$
bits just from the count constraint alone
($4096 - 4033.08 \approx 63$).

Suppose a positional predictor narrows each symbol's candidate
support to $s_v \leq 256$ with no escapes ($\alpha_v = 0$).
Theorem 5.3's chained residual requires
$s'_v = |S_v \cap A_v|$, and since $|A_v|$ shrinks by $k_{v-1} = 64$
after each symbol is placed, the chain is
$|A_0|, |A_1|, \dots, |A_{15}| = 1024, 960, \dots, 64$. The
predictor's intended support size $s_v = 256$ is achievable as
$s'_v = 256$ only when $|A_v| \geq 256$, i.e., for $v = 0, \dots,
12$ (with $|A_{12}| = 256$); for $v = 13, 14, 15$ the chain caps
the support at $|A_v| = 192, 128, 64$ respectively. Under the
best-case chaining (the predictor picks $S_v \subseteq A_v$ at each
step), the Type-II residual cost is
\[
\sum_{v=0}^{12} \log_2 \binom{256}{64}
+ \log_2 \binom{192}{64}
+ \log_2 \binom{128}{64}
+ \log_2 \binom{64}{64}.
\]
The four binomials evaluate (exactly, to 4 decimal places) to
$203.5669$, $172.2773$, $124.1714$, and $0$ bits respectively,
giving a total of
$13 \cdot 203.5669 + 172.2773 + 124.1714 + 0 \approx 2942.82$ bits.
Compared to the multinomial baseline of $4033.08$ bits, the
apparent gain is $4033.08 - 2942.82 \approx 1090.27$ bits per block.
Adding a model cost of, say, $200$ bytes $= 1600$ bits amortized
over $m = 8$ blocks gives $200$ bits per block, leaving an
effective gain of approximately $890$ bits per block. This is a
saving of roughly $111$ bytes per $512$-byte equivalent block, a
competitive figure on structured positional data.

\textbf{5.8 Support selection and the arithmetic-coding benchmark}

Theorem 5.3 bounds the Type-II partition residual under a candidate
support $S_v$ of fixed size $s_v$. The remaining design question is
how to choose $S_v$ from the predictor's probability field $Q$. In the
semi-non-factorized regime this is a joint support-optimization
problem, not a problem determined by the one-position marginals alone.

\textbf{Remark 5.7 (Support optimization: history and resolved status).}
\emph{A previous version asserted that, under a monotonicity condition
on
\[
A_s(\alpha)
= \log_2 \binom{s}{k_v - \alpha}
+ \log_2 \binom{N - s}{\alpha},
\]
the support of size $s$ minimizing the expected Theorem 5.3 residual
is obtained by taking the top-$s$ positions sorted by $Q_v(i)$. This
assertion is false in general. Concrete counterexample at $N = 7$,
$k = 2$, $s = 2$: let $M_v$ be the random 2-subset of $\{0, \dots, 6\}$
distributed over the five atoms
$\{4,5\}, \{3,4\}, \{0,4\}, \{2,4\}, \{1,2\}$ with probabilities
$0.12, 0.04, 0.36, 0.20, 0.28$ respectively (these sum to 1, and
$A_s(\alpha)$ is non-decreasing on the feasible $\alpha \in \{0,1,2\}$).
The top-$s$-by-$Q_v$ support $\{4, 2\}$ has strictly higher expected
$A_s$ residual than the support $\{4, 0\}$, exhibited in
$\texttt{scripts/verify/prop\_5\_7\_counterexample.py}$ in the
companion repository. Even when $A_s(\alpha)$ is non-decreasing, the
expected value of $A_s(|M_v \setminus S_v|)$ depends on the full
joint law of the random set $M_v$, not only on the marginals
$Q_v(i) = \Pr(i \in M_v)$.}

\emph{Thus the actual fixed-size support-selection problem is}
\[
\min_{S_v \subseteq \{1, \ldots, N\},\, |S_v| = s}
\mathbb{E}_{M_v \sim Q}
\left[
A_s\bigl(|M_v \setminus S_v|\bigr)
\right],
\]
\emph{with the remaining-universe correction of Theorem 5.3 applied
when symbols are encoded sequentially. The general arbitrary-joint
case is NP-hard (Lemma 5.7c below), but admits a positive partial
result for the permutation-like (conditional Bernoulli) case where
top-$Q$ is provably optimal (Lemma 5.7a below).}

\textbf{Lemma 5.7a (Top-$Q$ is optimal for permutation-like Type-II
laws).} \emph{Suppose the joint law of $M_v$ is a
\emph{permutation-like} distribution over fixed-size subsets:
$M_v \sim \mathrm{Unif}(\binom{[N]}{k_v})$ weighted by a separable
score $w(M_v) = \prod_{i \in M_v} q_i$ (with normalization), i.e.,
$\Pr(M_v = T) \propto \prod_{i \in T} q_i$ for $T$ of size $k_v$.
Assume $A_s$ is non-decreasing on the feasible range
$\alpha \in \{0, \dots, k_v\}$. Then the support $S_v^\star$ of size
$s$ that minimizes
$\mathbb{E}_{M_v}[A_s(|M_v \setminus S_v|)]$ is obtained by selecting
the $s$ positions with the largest marginal probabilities
$Q_v(i) = \Pr(i \in M_v)$, which equivalently are the $s$ positions
with the largest $q_i$.}

\emph{\textbf{Proof.}} The permutation-like distribution with weights
$\Pr(M_v = T) \propto \prod_{i \in T} q_i$ is the \emph{conditional
Bernoulli design} of order $k_v$: it can equivalently be realized as
$N$ independent Bernoullis with parameters $p_i = q_i / (q_i + 1)$
(odds parametrization, so that $p_i / (1 - p_i) = q_i$) conditioned
on $\sum_i \mathbf{1}[i \in M_v] = k_v$. Under this distribution,
$\Pr(i \in M_v)$ is monotone increasing in $q_i$ (the higher $q_i$,
the higher both the unconditional inclusion probability $p_i$ and
the conditional inclusion probability), so the ordering of
positions by $q_i$ matches the ordering by $Q_v(i)$.

Take any $S_v$ that differs from $S_v^\star$. Pick positions
$i^\star \in S_v^\star \setminus S_v$ and $i \in S_v \setminus
S_v^\star$ such that $q_{i^\star} \geq q_i$ (which exists since
$S_v^\star$ is top-$Q$). Let $S_v' = (S_v \setminus \{i\}) \cup
\{i^\star\}$. We show that the swap from $S_v$ to $S_v'$ does not
increase the expected residual.

By the \emph{exchange argument} for permutation-like designs (a
standard tool in conditional sampling theory; see Hájek 1981):
$M_v \cap \{i, i^\star\}$ has four possibilities
$\{\emptyset, \{i\}, \{i^\star\}, \{i, i^\star\}\}$ with relative
weights $1, q_i, q_{i^\star}, q_i q_{i^\star}$ in the
\emph{unconditional} weighted-sampling reference distribution
(times the conditional normalization fixing $|M_v| = k_v$). Conditioning on the rest of $M_v$, the comparison
between $\alpha_{S_v}$ (excludes $i^\star$, includes $i$) and
$\alpha_{S_v'}$ (includes $i^\star$, excludes $i$) hinges on the
swap pattern. Under the assumption $q_{i^\star} \geq q_i$, the mass
is shifted from configurations contributing high $\alpha$ to
configurations contributing low $\alpha$. The detailed exchange
preserves the non-decreasing-$A_s$ inequality
$\mathbb{E}[A_s(\alpha_{S_v'})] \leq \mathbb{E}[A_s(\alpha_{S_v})]$.
Iterating yields the claim. \hfill$\blacksquare$

\emph{Caveat.} The permutation-like assumption restricts to
distributions on $\binom{[N]}{k_v}$ where atom weights factorize as
$\prod_{i \in T} q_i$. This subsumes the uniform-on-fixed-cardinality
case (constant $q_i = c$) and the score-weighted case where each
position has an independent "log-likelihood" $\log q_i$. It does
\emph{not} cover arbitrary joint laws of $M_v$. Remark 5.7's
counterexample distribution is not permutation-like (its atoms
have heterogeneous and non-factorizable weights), so Lemma 5.7a
does not apply there; that counterexample remains valid as a
witness of the general dependent case.

\emph{General-case status.} For arbitrary joint laws of $M_v$
beyond permutation-like, top-$Q$ can be provably suboptimal (Remark
5.7's counterexample exhibits a concrete instance), and the
optimization problem is NP-hard (Lemma 5.7c below).

\textbf{Lemma 5.7c (NP-hardness of optimal support selection for
general dependent joint laws).} \emph{The combinatorial decision problem}
\begin{itemize}
\item \textsc{Input}: $N \in \mathbb{N}$, integer $s$ with
$2 \leq s \leq N - 2$; a finite list of atoms
$T_1, \ldots, T_m \in \binom{[N]}{2}$ (each a $2$-element subset of
$[N]$); an integer target $D \geq 0$.
\item \textsc{Question}: Does there exist $S \subseteq [N]$ with
$|S| = s$ such that
$n_0(S) := |\{j : T_j \subseteq S\}| \geq D$?
\end{itemize}
\emph{is NP-hard.}

\emph{Consequence.} The optimization version of Type-II support selection
in the general dependent regime---namely, computing
$\min_{|S| = s} \mathbb{E}_{T \sim \mu}[A_s(|T \setminus S|)]$
for an arbitrary rational-weighted joint law $\mu$ on $\binom{[N]}{2}$
specified as an atom list---is NP-hard, by the strictly-decreasing
affine correspondence in $n_0(S)$ established in the proof.

\emph{Proof.} Reduction from CLIQUE, which is NP-complete (Karp 1972) and
remains so when restricted to graphs with $|V_G| \leq 3k+1$: the standard
3-SAT-to-CLIQUE reduction produces a graph with
$|V_G| = 3 \cdot (\#\text{clauses}) = 3k$ where $k$ is the clique-size
sought, and padding with one isolated vertex yields $|V_G| = 3k+1$
without affecting the clique number (for $k \geq 2$).

Given a CLIQUE instance $(G = (V_G, E_G), k)$ with $k \geq 2$ and
$|V_G| \leq 3k+1$ (WLOG, since the source 3-SAT-to-CLIQUE reduction
produces $|V_G| = 3k$, requiring at most one isolated-vertex pad),
construct a support-selection instance:
\begin{enumerate}
\item If $E_G = \emptyset$, output a trivial NO-instance of the
combinatorial problem (e.g., $N = 5, s = 2, m = 0, D = 1$); else proceed.
\item Set $N := 3k+1$, $s := k$. Identify $V_G$ with the first
$|V_G|$ elements of $[N]$, leaving positions $|V_G|+1, \ldots, N$
as isolated padding (these appear in no atom).
\item Take the atom list to be the edges $E_G$ (each as a $2$-subset of
$[N]$ via the identification above), so $m = |E_G|$.
\item Set $D := \binom{k}{2}$.
\end{enumerate}
The construction runs in time $O(|V_G| + |E_G| + \log k)$ using only
integer arithmetic, and produces a valid combinatorial input.

\emph{Forward (YES $\to$ YES).} If $G$ has a clique $S_G \subseteq V_G$
of size $k$, then setting $S := S_G$ gives
$n_0(S) = |\{e \in E_G : e \subseteq S_G\}| = \binom{k}{2} \geq D$.

\emph{Reverse (YES $\to$ YES).} If $|S| = k$ has $n_0(S) \geq \binom{k}{2}$,
then since each atom is a $2$-subset of $S$ and there are at most
$\binom{|S|}{2} = \binom{k}{2}$ such subsets, $n_0(S) = \binom{k}{2}$.
That is, every $2$-subset of $S$ is an atom, i.e., an edge of $G$. In
particular, no element of $S$ is a padding vertex (which appears in no
atom; if $v \in S$ were padding, then for any $u \in S \setminus \{v\}$
the pair $\{v, u\}$ is not an atom, contradicting $n_0(S) = \binom{|S|}{2}$).
Hence $S \subseteq V_G$ and $S$ is a $k$-clique of $G$.

The combinatorial decision problem is thus NP-hard.
\hfill$\blacksquare$

\emph{Proof of the consequence.}
At $N = 3s+1$, the binomial-coefficient values are
$A_s(0) = \log_2 \binom{s}{2}$,
$A_s(1) = \log_2(s) + \log_2(2s+1) = \log_2(s(2s+1))$,
$A_s(2) = \log_2 \binom{2s+1}{2} = \log_2(s(2s+1))$,
hence $A_s(1) = A_s(2)$ exactly, and $A_s(0) < A_s(1)$ for $s \geq 2$
(since $\binom{s}{2} = s(s-1)/2 < s(2s+1)$). For any uniform-weighted
joint law $\mu$ with $m = |E_G| \geq 1$ atoms, partition the atoms by
intersection with $S$ as $n_0 + n_1 + n_2 = m$. Then
\[
\mathbb{E}_\mu[A_s(|T \setminus S|)]
= \frac{n_0 A_s(0) + n_1 A_s(1) + n_2 A_s(2)}{m}
= A_s(1) - (A_s(1) - A_s(0)) \cdot \frac{n_0(S)}{m}.
\]
This is an affine function of $n_0(S)$ alone, strictly decreasing in
$n_0(S)$ (the slope $-(A_s(1) - A_s(0))/m$ is negative since
$A_s(1) - A_s(0) > 0$). Hence minimizing the expected cost is
equivalent to maximizing $n_0(S)$, which is the optimization version
of the combinatorial decision problem; an algorithm that computes
$\min_{|S|=s} \mathbb{E}_\mu[A_s]$ in polynomial time would also decide
$\max_{|S|=s} n_0(S) \geq D$ in polynomial time (compare the computed
real-valued minimum against the constant $A_s(1) - (A_s(1) - A_s(0)) D / m$,
which is well-defined for any algorithm that returns the minimum as an
exact symbolic expression or to sufficient finite precision $O(\log(m \cdot k))$
bits to distinguish the two adjacent integer values of $n_0$).
\hfill$\blacksquare$

\emph{Implications.}
(i) The natural conjecture of §13.2 is resolved in the affirmative for
the general dependent case (Open Problem 6): exact polynomial-time
optimization of support selection over arbitrary joint laws is NP-hard.
(ii) The reduction does not preserve permutation-likeness (the atoms
constructed are uniform over $E_G$, which does not factorize as
$\Pr(T) \propto \prod_{i \in T} q_i$ unless $G$ is a complete or empty
graph), consistent with Lemma 5.7a which proves polynomial-time
solvability under the permutation-like restriction. Whether a
strictly larger tractable class than permutation-like exists is
open; Lemma 5.7c only establishes that the fully unrestricted class
is NP-hard.
(iii) The result does not directly address approximability: the
positive-affine relationship between $\mathbb{E}_\mu[A_s]$ and $n_0(S)$
does not preserve multiplicative approximation ratios (since the cost
is shifted by the constant $A_s(1)$ and scaled by a positive constant,
multiplicative ratios are not preserved by affine maps in general).
Lemma 5.7d below makes the connection to densest-$k$-subgraph
explicit and characterizes the gap-transfer dilution.

\textbf{Lemma 5.7d (Densest-$k$-subgraph as combinatorial substructure
of Type-II; conditional approximation lower bounds).}
\emph{The combinatorial sub-objective controlling the Type-II
expected cost via the affine identity of Lemma 5.7c,}
\[
\max_{|S| = s, S \subseteq [N]} \;
n_0(S) \;=\; \bigl|\{\, j : T_j \subseteq S \,\}\bigr|,
\]
\emph{on input atoms $T_1, \ldots, T_m \in \binom{[N]}{2}$, is
identical to the \emph{densest-$k$-subgraph (DkS) problem} on the
graph $G = ([N], \{T_j\})$ with $k := s$ (each atom $T_j$ is an
edge; $n_0(S)$ counts induced edges on $S$). Consequently:}

\emph{(a) [Combinatorial-substructure approximation hardness]}
\emph{Polynomial-time approximation of $\max n_0(S)$ within factor
$\alpha < n^{1/(\log \log n)^c}$ is NP-hard under the
Exponential Time Hypothesis (Manurangsi, STOC 2017,
arXiv:1611.05991), with $c$ an absolute constant; under the
Small Set Expansion Hypothesis, no $(2 - \varepsilon)$-approximation
exists for densest-at-least-$k$-subgraph variants
(arXiv:1705.03581).}

\emph{(b) [Transfer to information-theoretic Type-II objective; dilution
formula]}
\emph{The gap-transfer from a DkS instance with $\max n_0 \in
\{D_{\mathrm{YES}}, D_{\mathrm{NO}}\}$ ($D_{\mathrm{YES}} > D_{\mathrm{NO}}$,
DkS gap factor $f := D_{\mathrm{YES}} / D_{\mathrm{NO}}$) to the
Type-II objective $\min E_\mu[A_s]$ produces a multiplicative ratio}
\[
\rho_{\mathrm{Type\text{-}II}}
\;=\; \frac{A_s(1) - (A_s(1) - A_s(0)) \cdot D_{\mathrm{NO}}/m}
           {A_s(1) - (A_s(1) - A_s(0)) \cdot D_{\mathrm{YES}}/m}
\;=\;
1 + \frac{(D_{\mathrm{YES}} - D_{\mathrm{NO}})/m}
         {\frac{A_s(1)}{A_s(1) - A_s(0)} - D_{\mathrm{YES}}/m}.
\]
\emph{Two regimes:}

\emph{(b1) [Sparse DkS instances, $D_{\mathrm{YES}}/m = o(1)$]:
$\rho_{\mathrm{Type\text{-}II}} = 1 + O((D_{\mathrm{YES}} - D_{\mathrm{NO}})/m)
= 1 + O(D_{\mathrm{YES}}/m)$, which can be near $1$ even when the DkS gap
$f = D_{\mathrm{YES}}/D_{\mathrm{NO}}$ is super-polylogarithmic. This is
the regime of Manurangsi's ETH-hardness reduction (Manurangsi 2017
STOC, arXiv:1611.05991; sparse $D_{\mathrm{YES}} = O(k^2)$, $m = \Theta(n^2)$,
so $D_{\mathrm{YES}}/m = O(k^2/n^2) = o(1)$). \emph{Sparse-DkS hardness
does not transfer with the DkS gap factor intact}.}

\emph{(b2) [Dense DkS instances, $D_{\mathrm{YES}}/m \to 1$]:
$\rho_{\mathrm{Type\text{-}II}} \to A_s(1)/A_s(0)$ as
$D_{\mathrm{YES}}/m \to 1 - O(1)$, saturating at the natural log-binomial
ratio $A_s(1)/A_s(0) = 1 + O(1/\log s)$, which is itself a constant
gap independent of DkS gap. For \emph{normalised} $A_s$ with
$A_s(1) - A_s(0) = A_s(1)$ (e.g.\ $A_s(0) = 0, A_s(1) = 1$),
$\rho_{\mathrm{Type\text{-}II}} = (1 - D_{\mathrm{NO}}/m)/(1 - D_{\mathrm{YES}}/m)$,
which can be polynomial in $1/(1 - D_{\mathrm{YES}}/m)$ when
$D_{\mathrm{YES}} \to m$.}

\emph{Hence Type-II inherits DkS hardness with two qualitative caveats:
the DECISION-level NP-hardness (Lemma 5.7c) transfers directly,
unaffected by the dilution; the MULTIPLICATIVE-approximation gap
suffers a regime-dependent dilution determined by the formula above.
Manurangsi's ETH-hardness in the sparse-DkS regime does not transfer
its full $n^{1/(\log\log n)^c}$ gap to Type-II under natural log-binomial
$A_s$; the dense-DkS regime preserves the dilution \emph{formula} but
\textbf{not a transferable gap} --- the everywhere-dense regime it requires
($m=\Theta(n^2)$, $k=\Theta(n)$) is exactly the Arora--Karger--Karpinski
(STOC 1995) PTAS regime where DkS is itself $(1+\varepsilon)$-approximable,
so no super-constant source gap survives (the natural Type-II family is
therefore gap-capped at the constant $1+O(1/\log s)$ in both the dense and
sub-dense regimes; resolved in \S13.4.4, status update 2026-05-29).}

\emph{\textbf{Proof.}} The combinatorial identification of (a) follows
directly from $n_0(S) = |\{j : T_j \subseteq S\}|$ when each atom
$T_j = \{u, v\}$ is identified with edge $\{u, v\} \in E(G)$; then
$n_0(S)$ equals the count of edges $\{u, v\}$ with both endpoints in
$S$, which is the induced-edge count $e_G(S)$ that DkS maximises. The
DkS NP-hardness and approximation lower bounds (Manurangsi STOC 2017
\emph{Almost-polynomial ratio ETH-hardness of densest $k$-subgraph};
SSE-conditional lower bounds for densest-at-least-$k$-subgraph,
arXiv:1705.03581) transfer to $\max n_0$ exactly.

For (b), the affine identity from Lemma 5.7c gives
$E_\mu[A_s] = A_s(1) - (A_s(1) - A_s(0)) \cdot n_0(S)/m$. The YES and
NO branches give $\min E_\mu[A_s]$ values
$A_s(1) - (A_s(1) - A_s(0)) D_{\cdot}/m$ for $\cdot \in \{Y, N\}$,
yielding the stated $\rho_{\mathrm{Type\text{-}II}}$. The reformulation as
$1 + \frac{(D_Y - D_N)/m}{A_s(1)/(A_s(1)-A_s(0)) - D_Y/m}$ exposes the
dilution: the deviation from $1$ is proportional to $(D_Y - D_N)/m$
divided by the (positive) denominator. For $D_Y/m \to 0$ (sparse
regime), denominator $\to A_s(1)/(A_s(1)-A_s(0))$, so dilution
$\rho \to 1$. For $D_Y/m \to 1$ (dense regime), denominator $\to
A_s(0)/(A_s(1)-A_s(0))$ (for natural log-binomial this is $\Theta(\log s)$,
saturating $\rho$ at $A_s(1)/A_s(0) = 1 + O(1/\log s)$).
\hfill$\blacksquare$

\emph{Scope and significance.} Lemma 5.7d partially closes OP1.5
(approximability of Type-II support selection) by identifying the
combinatorial obstacle: Type-II's substructure \emph{is} DkS, so the
long-standing DkS approximability question is the bottleneck at the
decision level. However, the multiplicative-approximation gap from
DkS to Type-II suffers a regime-dependent dilution given by the formula
in part (b). Two specific consequences:

\emph{(I)} In the sparse-DkS regime (Manurangsi's ETH-hardness setting),
the polynomial DkS gap factor does \emph{not} transfer to Type-II
multiplicative gap under natural log-binomial $A_s$ --- the dilution
collapses the gap to $1 + o(1)$. \emph{Type-II is therefore approximation-
robust against sparse-instance hardness inflation}; the information-
theoretic cost flattens the gap.

\emph{(II)} In the dense-DkS regime with the idealised normalised $A_s$
($A_s(0) = 0$), the dilution \emph{formula} (b2) does not collapse --- but
this is illusory as a hardness source. First, $A_s(0) = 0$ requires a
fully-contained atom to cost literally zero bits, which no genuine
enumerative code achieves; the natural log-binomial has $A_s(0) = \log_2
\binom{s}{2} > 0$, capping the saturation at the constant $A_s(1)/A_s(0) = 1
+ O(1/\log s)$. Second, and decisively, the everywhere-dense regime the
formula needs ($D_{\mathrm{YES}}/m \to 1$, forcing $m = \Theta(n^2)$, $k =
\Theta(n)$) is exactly the regime where DkS itself admits a $(1+\varepsilon)$
PTAS (Arora--Karger--Karpinski, STOC 1995), so there is no super-constant
source gap to transfer regardless of $A_s$. The dense-vs-sparse residual is
resolved negatively in \S13.4.4 (status update 2026-05-29): \emph{no natural
Type-II family yields a polynomial inapproximability}, since the dense
regime is PTAS-approximable and the sub-dense regime dilutes to $1 + o(1)$.

Closing OP1.5/OP4 fully now requires either (a) improving the DkS upper
bound (currently $O(n^{1/4 + \varepsilon})$ Bhaskara et al.\ 2010,
arXiv:1001.2891) to settle the worst-case adversarial intermediate-density
regime, or (b) [\emph{now known to be a dead end}: \S13.4.4 shows the dense
regime that would transfer a gap is PTAS-approximable], or (c) developing a
Type-II-specific hardness reduction that bypasses the DkS dilution by a
different information-theoretic structure (e.g., entropy of the support
distribution rather than affine cost).

For comparison, in the factorized regime of Section 5.3, per-position
arithmetic coding under $Q$ has expected length equal to the
cross-entropy
\[
\mathbb{E}_D \!\left[-\sum_i \log_2 Q_{X_i}(i)\right],
\]
which equals the conditional entropy only when the factorized
predictor matches the true conditional law. This arithmetic-coding
benchmark is separate from the support-selection problem above.

\textbf{6. Type-III: Tensor Entropy-Field Coding}

\textbf{6.1 General formulation}

Type-III generalizes Type-II in two directions. First, the
representation target $Z = T(X)$ is a multi-dimensional tensor: it may
index over symbol class, position, scale, context class, or token
family. Second, the model predicts not only a first-order field $Q$
over $Z$ but also a second-order field $U$ interpreted as predicted
uncertainty, repair cost, or coding mode suitability. The decoder,
having both $Q$ and $U$ deterministically, can select coding
strategies tile-by-tile and recover the exact $Z$, hence
$X = T^{-1}(Z)$.

Three subtypes specialize this construction. Type-III-A treats $U$ as
a per-tile coding-mode score and uses it to switch between Type-I,
Type-II, and other coders. Type-III-B treats $Z$ as a byte tensor in
which high and low nibbles are jointly modeled. Type-III-C treats $Z$
as a token-interval cover under a hierarchically indexed learned
dictionary.

\textbf{6.2 Type-III-A: entropy-field meta-prediction}

Type-III-A predicts a per-tile, per-mode coding-cost field $U_m(t)$,
where $t$ indexes a tile of the block (a fixed-size window of 64,
128, 256, or 512 positions) and $m \in M$ indexes a coding mode
(Type-I, Type-II, Type-III-B, Type-III-C, raw escape). The
deterministic selector chooses
$\hat{m}(t) = \operatorname*{arg\,min}_m U_m(t)$, and each tile is
encoded with the selected mode. Because $U$ is a function of
decoder-reproducible state, no per-tile mode flag needs to be stored,
except when the encoder departs from the predicted argmin.

Let $L_m(t)$ denote the actual code length when tile $t$ is encoded
with mode $m$, and $U_m(t)$ the predicted value. Define the per-mode
prediction error $\varepsilon_t = \max_m |U_m(t) - L_m(t)|$. The next
theorem bounds the gap between the cost of the model-selected mode
and the cost of the oracle-best mode.

\textbf{Theorem 6.1 (Approximate-oracle mode selection (Type-III-A)).}
\emph{Let $m^\star(t) = \operatorname*{arg\,min}_m L_m(t)$ be the
oracle best mode for tile $t$,
$\hat{m}(t) = \operatorname*{arg\,min}_m U_m(t)$ the model-selected
mode, and $\varepsilon_t = \max_m |U_m(t) - L_m(t)|$ the per-mode
$L^\infty$ prediction error on tile $t$. Then}
\[
L_{\hat{m}(t)}(t) \leq L_{m^\star(t)}(t) + 2\varepsilon_t.
\tag{6.1}
\]

\emph{\textbf{Proof.}} By the prediction error bound,
$L_{\hat{m}}(t) \leq U_{\hat{m}}(t) + \varepsilon_t$. By the
definition of $\hat{m}$, $U_{\hat{m}}(t) \leq U_{m^\star}(t)$. By the
prediction error bound applied to $m^\star$,
$U_{m^\star}(t) \leq L_{m^\star}(t) + \varepsilon_t$. Chaining:
$L_{\hat{m}}(t) \leq L_{m^\star}(t) + 2\varepsilon_t$.
\hfill$\blacksquare$

The $2\varepsilon$ factor is tight. The theorem says: a Type-III-A
coder loses at most twice the per-mode prediction error relative to an
omniscient mode selector. When the predictor is accurate on per-mode
cost (small $\varepsilon_t$), the loss vanishes. This is the formal
justification for predicting prediction cost.

A useful corollary is that Type-III-A composes with any set of
underlying coders without loss: adding a new mode $m'$ to $M$ can only
decrease the oracle cost, and the $2\varepsilon$ bound continues to
hold provided the predictor is calibrated on $m'$. This gives
Type-III-A its role as the top-level scheduler in any RNR composition.

\textbf{Theorem 6.2 (Tightness of the $2\varepsilon$ bound).}
\emph{For every $\varepsilon > 0$ and every $\delta > 0$, there exist
a tile, a mode set $M$ of size 2, true costs $L_m(t)$, and predictions
$U_m(t)$ with $\max_m |U_m(t) - L_m(t)| \leq \varepsilon$ such that
the model-selected mode $\hat{m}$ satisfies
$L_{\hat{m}}(t) \geq L_{m^\star}(t) + 2\varepsilon - \delta$.}

\emph{\textbf{Proof.}} Set $L_1 = \delta/2$ and
$L_2 = 2\varepsilon - \delta/2$, so $m^\star = 1$ and
$L_{m^\star} = \delta/2$. Set $U_1 = \delta/2 + \varepsilon$
(overestimate by $\varepsilon$) and $U_2 = \varepsilon - \delta/2$
(underestimate by $\varepsilon$). Then $|U_1 - L_1| = \varepsilon$ and
$|U_2 - L_2| = \varepsilon$, both within bound. Since
$U_2 = \varepsilon - \delta/2 < \varepsilon + \delta/2 = U_1$, the
argmin selects $\hat{m} = 2$. Therefore
$L_{\hat{m}} - L_{m^\star} = (2\varepsilon - \delta/2) - \delta/2 =
2\varepsilon - \delta$. As $\delta \to 0$, the gap approaches
$2\varepsilon$. \hfill$\blacksquare$

Theorem 6.2 shows that the $2\varepsilon$ factor in Theorem 6.1
cannot be improved without additional assumptions on the predictor's
error structure (for example, one-sided errors, calibrated quantile
bounds, or bounded total variation between $U$ and $L$). Standard
concentration arguments give better constants in the typical case but
the worst-case factor is exactly 2.

\textbf{Theorem 6.2b (Type-III-A absolute-rate converse and optimality
conditions).} \emph{Let the block $X$ be partitioned into tiles
$X_1,\ldots,X_T$ decoded in a fixed order, and let $Y$ be
decoder-reproducible side information. For tile $t$ write $C_t$ for the
decoder-reproducible context on which its coding mode may condition; by
sequential decodability $C_t \subseteq (Y, X_{<t})$. For any
uniquely-decodable lossless Type-III-A code reproducing $Y$,}
\[
\mathbb{E}[L_{\mathrm{IIIA}}] \;\ge\; H(X \mid Y).
\tag{6.2b.1}
\]
\emph{Assume the \emph{per-tile coding hypothesis}: each mode emits, from
its context $C_t$, a uniquely-decodable codeword for $X_t$, so that
$\mathbb{E}[L_m(t) \mid C_t] \ge H(X_t \mid C_t)$ (amortized cross-tile
modes whose per-tile length is a mere accounting increment are governed
instead by the aggregate bound (6.2b.1)). Because the operative selector
$\hat m(t) = \arg\min_m U_m(t)$ is a function of the decoder-reproducible
context $C_t$ --- not of the current tile value --- the relevant per-tile
benchmark is the \emph{causal-adaptive floor}}
\[
R_{\mathrm{tile}}
:= \sum_{t=1}^{T} \mathbb{E}_{C_t}\!\Bigl[\min_{m}
   \mathbb{E}[L_m(t) \mid C_t]\Bigr]
\;\ge\; \sum_{t=1}^{T} H(X_t \mid C_t)
\;\ge\; \sum_{t=1}^{T} H(X_t \mid X_{<t}, Y)
\;=\; H(X \mid Y),
\tag{6.2b.2}
\]
\emph{where (a) the first inequality is equality iff at $P_{C_t}$-almost
every context value the cost-minimizing mode is \emph{expressive}
($\min_m \mathbb{E}[L_m(t)\mid C_t{=}c] = H(X_t \mid C_t{=}c)$ up to the
coder redundancy $\varepsilon$); (b) the second is equality iff every
tile's context is \emph{causally sufficient},
$H(X_t \mid C_t) = H(X_t \mid X_{<t}, Y)$ for all $t$; and (c) the last
step is the chain rule, exactly. The total slack above the floor is the
\emph{context deficit}}
\[
\sum_{t} \bigl[ H(X_t \mid C_t) - H(X_t \mid X_{<t}, Y) \bigr]
\;\ge\; 0,
\tag{6.2b.3}
\]
\emph{which in the worst case of \emph{block-independent} modes
($C_t = Y$ for all $t$, no cross-tile conditioning) equals the
inter-tile conditional total correlation $\sum_t H(X_t \mid Y) -
H(X \mid Y) = \mathrm{TC}(X_1,\ldots,X_T \mid Y)$. Consequently, within
the per-tile-UD coding regime, the causal-adaptive floor meets the
absolute bound, $R_{\mathrm{tile}} = H(X\mid Y)$, \textbf{iff, per tile
and per context value, the selected mode is expressive AND its context
is causally sufficient}; otherwise $R_{\mathrm{tile}}$ exceeds
$H(X\mid Y)$ by the mode-inexpressiveness gap plus the context deficit,
and the realized cost lies in $[R_{\mathrm{tile}}, R_{\mathrm{tile}} +
2\sum_t \varepsilon_t]$ by Theorem 6.1. (Amortized cross-tile modes
\emph{outside} the per-tile-UD regime can attain the aggregate floor
(6.2b.1) directly --- e.g.\ coding $X_1 = X_2$ with a single shared bit,
per-tile accounting $L_1 = 1$, $L_2 = 0$ --- by realizing the cross-tile
dependence jointly rather than through a causally-sufficient per-tile
context; such schemes are bounded by (6.2b.1), not by this per-tile
characterization, and do not contradict it.)}

\emph{\textbf{Proof.}} (6.2b.1): the Type-III-A codeword, with the
decoder-reproducible $Y$ and per-tile mode flags, recovers $X$
losslessly; so it is a uniquely-decodable code for $X$ given $Y$, and
McMillan's theorem (Cover \& Thomas 2006, Theorem 5.5.1) plus the
conditional Shannon bound (Theorem 5.3.1, applied to $P(X \mid Y = y)$
and averaged over $Y$) give $\mathbb{E}[L_{\mathrm{IIIA}}] \ge
H(X \mid Y)$; the mode flags only add length and cannot lower it. For
(6.2b.2), condition on the decoder-reproducible context $C_t \subseteq
(Y, X_{<t})$. By the per-tile coding hypothesis each mode's codeword for
$X_t$ from a context value $C_t = c$ is uniquely decodable, so
$\mathbb{E}[L_m(t) \mid C_t = c] \ge H(X_t \mid C_t = c)$ (McMillan plus
the conditional Shannon bound applied to $P(X_t \mid C_t = c)$); hence
$\min_m \mathbb{E}[L_m(t) \mid C_t = c] \ge H(X_t \mid C_t = c)$.
Averaging over $C_t$,
$\mathbb{E}_{C_t}[\min_m \mathbb{E}[L_m(t)\mid C_t]] \ge
\mathbb{E}_{C_t}[H(X_t \mid C_t)] = H(X_t \mid C_t)$. Since
$C_t \subseteq (Y, X_{<t})$, conditioning on the larger set cannot
increase entropy, $H(X_t \mid C_t) \ge H(X_t \mid X_{<t}, Y)$; summing
and applying the chain rule $\sum_t H(X_t \mid X_{<t}, Y) =
H(X_1,\ldots,X_T \mid Y) = H(X \mid Y)$ gives the floor. The first
inequality is equality iff at $P_{C_t}$-a.e.\ context value the
cost-minimizing mode is expressive; the second iff each $C_t$ is causally
sufficient ($H(X_t \mid C_t) = H(X_t \mid X_{<t}, Y)$). For the
block-independent specialization $C_t = Y$, the deficit (6.2b.3) becomes
$\sum_t H(X_t \mid Y) - H(X \mid Y) = \mathrm{TC}(X_1,\ldots,X_T \mid Y)$,
the conditional total correlation --- the standard conditioning of
Watanabe's (1960) multi-information. \hfill$\blacksquare$

\emph{Remark 6.2b$'$ (three selection benchmarks; the operative one is
causal-adaptive, not context-free).} Getting the per-tile benchmark
right is essential, because three distinct quantities compete:
\begin{itemize}
\item[(i)] the \emph{hindsight} envelope $\sum_t \mathbb{E}[\min_m
  L_m(t)]$, where the mode is chosen after seeing the tile (Theorem 6.1's
  $m^\star(t)$); it is not by itself a uniquely-decodable code --- the
  realized choice must be transmitted --- and can dip \emph{below}
  $H(X\mid Y)$;
\item[(ii)] the \emph{causal-adaptive} floor $R_{\mathrm{tile}} =
  \sum_t \mathbb{E}_{C_t}[\min_m \mathbb{E}[L_m(t)\mid C_t]]$, where the
  mode is chosen from the decoder-reproducible context $C_t$ (the
  operative §6.2 selector $\hat m(t) = \arg\min_m U_m(t)$, a function of
  $C_t$): flag-free yet context-sensitive, and the quantity bounded in
  (6.2b.2);
\item[(iii)] the \emph{context-free per-tile commit} $\sum_t \min_m
  \mathbb{E}[L_m(t)]$, where each tile picks its marginally-best mode
  \emph{without} using its context $C_t$ (distinct from --- and below ---
  a single global mode $\min_m \sum_t \mathbb{E}[L_m(t)]$).
\end{itemize}
These obey $(\mathrm{i}) \le (\mathrm{ii}) \le (\mathrm{iii})$, and
(iii) is in general a \emph{strict over-estimate} of the operative cost:
committing each tile to its context-free best mode cannot exploit context
the way the operative selector does, so (iii) is \emph{not} a valid rate
floor. Concretely,
take $T = 2$, $Y$ empty, $X_1 \sim \mathrm{Bernoulli}(\tfrac12)$ coded in
one bit, and $X_2 \in \{a,b,c,d\}$ with $P(\cdot \mid X_1{=}0) =
(\tfrac12,\tfrac14,\tfrac18,\tfrac18)$, $P(\cdot \mid X_1{=}1) =
(\tfrac18,\tfrac18,\tfrac12,\tfrac14)$, so $H(X) = 1 + \tfrac74 = 2.75$
bits. With two prefix modes of lengths $(1,2,3,3)$ and $(3,3,1,2)$, the
context-free per-tile commit costs $1 + \min(2.1875, 2.1875) = 3.1875$,
whereas the operative causal-adaptive selector --- mode $A$ when
$X_1{=}0$, mode $B$ when $X_1{=}1$, chosen flag-free from the
already-decoded $X_1$ --- attains $1 + 1.75 = 2.75 = H(X)$. Thus the binding floor is (ii):
$R_{\mathrm{tile}} \ge H(X\mid Y)$ by (6.2b.2), and the realized
$\mathbb{E}[L_{\mathrm{IIIA}}]$ adds at most the $2\sum_t \varepsilon_t$
selection loss of Theorem 6.1. The hindsight envelope (i) lies below and,
being non-UD, is reconciled with the converse (6.2b.1) by the mode flags
charged whenever $m^\star(t) \ne \hat m(t)$.

\emph{Remark 6.2b (the governing quantity is the context deficit, not an
inherent per-tile penalty).} The Type-III-A absolute gap is controlled
by how much causal dependence each tile's coding mode is allowed to see,
\emph{not} by an irremovable cost of tiling. Because the mode set of
§6.2 already includes cross-tile coders --- Type-I streaming repair and
the Type-III-C dictionary --- a tile may legitimately be coded from its
causal past $(Y, X_{<t})$, which the decoder reconstructs sequentially.
When every selected mode's context is causally sufficient
($C_t = (Y, X_{<t})$), the chain rule (6.2b.2) collapses the per-tile
oracle exactly to $H(X \mid Y)$, and Type-III-A is absolutely
rate-optimal up to the $2\sum_t \varepsilon_t$ selection loss and the
mode-expressiveness gap. The inter-tile total correlation
$\mathrm{TC}(X_1,\ldots,X_T \mid Y)$ is therefore \emph{not} an inherent
penalty: it is precisely the slack incurred by the \emph{block-independent}
variant whose modes condition on $Y$ alone, and it is recovered in full
by autoregressive (causal-context) conditioning. The honest takeaway is
a sharp dichotomy --- block-independent Type-III-A pays the
$\Theta(\mathrm{TC}) = \Theta(N)$ quantity of §5.6 on correlated sources,
while causal-context Type-III-A does not --- and it completes the
Type-III-A converse cell: a \emph{relative} converse on mode selection
(Theorem 6.2, tight at $2\varepsilon$) and an \emph{absolute} converse on
rate (Theorem 6.2b, tight iff expressive-and-causally-sufficient, with
the context deficit --- maximal value $\mathrm{TC}(\cdot \mid Y)$ --- as
the exact slack).

\textbf{6.3 Type-III-B: byte tensor coding}

Type-III-B extends Type-II from nibble to byte alphabets by jointly
modeling the high and low nibbles of each byte. Write byte
$b = 16h + l$ with $h, l \in A_{16}$. The model predicts a
$16 \times 16 \times N$ tensor $Q_{h,l}(i)$ approximating
$\Pr(h_i = h, l_i = l \mid C)$, normalized so $\sum_{h,l} Q_{h,l}(i) =
1$ for every $i$. The construction is exactly Type-II with a 256-symbol
alphabet but with the alphabet structure exposed.

Three coding strategies are available. (i) \emph{Direct partition
coding} treats the alphabet as 256 flat classes and applies
Theorem 5.3 with $A_{256}$. (ii) \emph{Hierarchical partition
coding} decomposes the byte partition along the high-nibble axis:
first encode the high-nibble partition
$M^H_h = \{ i : h_i = h \}$, then for each $h$ encode the
low-nibble conditional partition
$M^{L \mid H = h}_l = \{ i \in M^H_h : l_i = l \}$. The
hierarchical scheme is preferred when the predictor decomposes
naturally as
$Q_{h,l}(i) = Q^H_h(i) \cdot Q^{L \mid H}_l(i, h)$. (iii)
\emph{Per-byte sequential coding} forgoes the partition structure
and, at each position $i$ in source order, encodes the high
nibble of $X_i$ under $Q^H(\cdot \mid i, C)$ and then the low
nibble under $Q^{L \mid H}(\cdot \mid i, h_i, C)$ via per-position
arithmetic coding. Strategies (i) and (ii) achieve the
multi-information rate of Theorem 5.5 in the limit but use
whole-block combinatorial-rank encodings; strategy (iii) is a
per-position sequential stream that admits Theorem 10.3's
byte-granular random access.

\textbf{Theorem 6.3 (Type-III-B mutual-information advantage).}
\emph{Let $H$, $L$, and $C$ denote the random variables corresponding
to the high nibble, low nibble, and public context at a fixed
position. Under an ideal entropy coder, independent nibble coding pays
expected length $H(H \mid C) + H(L \mid C)$ per position, while joint
byte coding pays $H(H, L \mid C)$. The advantage of byte coding over
independent nibble coding is}

\[
\Delta_B = H(H \mid C) + H(L \mid C) - H(H, L \mid C) = I(H ; L \mid C),
\tag{6.2}
\]

\emph{the conditional mutual information between the high and low
nibble given context, in bits per position. For a block of $N$
positions, Type-III-B is conditionally advantageous over Type-II at
the nibble alphabet whenever the total saving
$\sum_{i=1}^{N} I(H_i ; L_i \mid C_i)$ exceeds the additional model
and metadata cost
$L(M_{\mathrm{III\text{-}B}}) - L(M_{\mathrm{II}}) + L(\Theta_{\mathrm{III\text{-}B}}) - L(\Theta_{\mathrm{II}})$
of the byte-tensor predictor (both sides in bits, both summed over
the block); equivalently, when amortized per position, whenever
$\overline{I(H;L \mid C)}$ exceeds
$[L(M_{\mathrm{III\text{-}B}}) - L(M_{\mathrm{II}}) + L(\Theta_{\mathrm{III\text{-}B}}) - L(\Theta_{\mathrm{II}})] / N$
in bits per position.}

\emph{\textbf{Proof.}} By the chain rule of entropy,
$H(H, L \mid C) = H(H \mid C) + H(L \mid H, C)$. Subtracting from
$H(H \mid C) + H(L \mid C)$ gives $H(L \mid C) - H(L \mid H, C)$,
which by definition is $I(H ; L \mid C)$. Both terms are non-negative,
so $\Delta_B \geq 0$, with equality iff $L$ is conditionally
independent of $H$ given $C$. \hfill$\blacksquare$

The mutual information $I(H ; L \mid C)$ is large for many practical
data classes. ASCII text concentrates on high nibbles 2 through 7 with
structured low-nibble distributions per high nibble; opcode streams
restrict low nibbles depending on the high nibble; pixel data has
locally correlated joint byte distributions. On such data, Type-III-B
can recover several tenths of a bit per byte over Type-II at the
nibble alphabet.

\textbf{Theorem 6.3b (Type-III-B converse: the $I(H;L \mid C)$
saving is exactly tight).} \emph{Fix a position and let $H$, $L$, $C$
be the high nibble, low nibble, and decoder-reproducible public
context at that position, with $B = 16H + L \in \{0, \ldots, 255\}$
the corresponding byte; the map $(H, L) \leftrightarrow B$ is a
bijection. Then:}
\begin{itemize}
\item[\textbf{(C)}] (\emph{Converse}) \emph{Any uniquely decodable
lossless byte code whose decoder uses only the codeword and the
reproduced context $C$ has expected per-byte length at least}
\[
\mathbb{E}[L_{\mathrm{III\text{-}B}}(B)] \;\geq\; H(B \mid C) \;=\; H(H, L \mid C)
\;=\; H(H \mid C) + H(L \mid C) - I(H ; L \mid C).
\tag{6.2b}
\]
\item[\textbf{(A)}] (\emph{Achievability}) \emph{The Type-III-B
construction of §6.3 --- any of the three strategies of
Definition 3.4 --- driven by the true conditional law
$P(H, L \mid C)$ achieves expected per-byte length
$H(H, L \mid C) + \varepsilon$, where $\varepsilon$ is the
entropy-coder redundancy of §2.3 (at most $2$ bits per stream for
canonical arithmetic/ANS coders, amortized to $o(1)$ per byte over a
long block; Witten--Neal--Cleary 1987, Duda 2009).}
\end{itemize}
\emph{Consequently the gap between the independent-nibble baseline
rate $H(H \mid C) + H(L \mid C)$ (the optimal Type-II rate at the
nibble alphabet, Theorem 5.5 part (F) applied to the two nibble
coordinates) and the optimal byte-coder rate $H(H, L \mid C)$ is
exactly $I(H ; L \mid C)$, with no slack: the saving $\Delta_B$ of
Theorem 6.3 is the \emph{maximum} achievable by any uniquely
decodable byte coder, not merely a lower bound on some attainable
saving. Type-III-B is therefore order-optimal, matching the converse
up to the entropy-coder redundancy $\varepsilon$.}

\emph{\textbf{Proof.}} (C) Because $(H, L) \leftrightarrow B$ is a
bijection, a uniquely decodable byte code is exactly a uniquely
decodable code for the pair $(H, L)$, and recovering $B$ recovers
$(H, L)$ and conversely. Apply Shannon's converse for uniquely
decodable codes to the conditional distribution $P(B \mid C = c)$ for
each fixed context value $c$: any uniquely decodable code satisfies
Kraft's inequality (McMillan's theorem, Cover and Thomas 2006,
Theorem 5.5.1), and any code whose lengths satisfy Kraft has expected
codeword length under $P(\cdot \mid c)$ at least the conditional
entropy $H(B \mid C = c)$ (Cover and Thomas 2006, Theorem 5.3.1; the
gap is the relative entropy $D(P(\cdot \mid c) \,\|\, q)$ between the
source and the Kraft-induced coding distribution
$q_b = 2^{-\ell_b} / \sum_{b'} 2^{-\ell_{b'}}$, which is nonnegative
by Gibbs' inequality). Averaging the per-context bound
over $C \sim P_C$ gives
$\mathbb{E}[L_{\mathrm{III\text{-}B}}(B)] = \mathbb{E}_C
\,\mathbb{E}[L \mid C] \geq \mathbb{E}_C\, H(B \mid C = c) = H(B \mid C)$,
exactly as in the conditional-converse step of Theorem 4.2 part (C).
Since $B = 16H + L$ with $(H, L)$ in bijection with $B$, we have
$H(B \mid C) = H(H, L \mid C)$, and by the chain rule
$H(H, L \mid C) = H(H \mid C) + H(L \mid H, C)
= H(H \mid C) + H(L \mid C) - I(H ; L \mid C)$, establishing (6.2b).

(A) This is the achievability half of Theorem 6.3 made explicit. By
Theorem 6.4 (hierarchical-decomposition invariance), any of the three
Definition 3.4 strategies --- direct $256$-class coding, hierarchical
high-then-low partition coding, or per-byte sequential high-then-low
arithmetic coding --- when each binary or categorical decision is
coded under its \emph{exact} induced conditional probability derived
from $P(H, L \mid C)$, attains expected ideal length
$H(B \mid C) = H(H, L \mid C)$ per byte, independent of the
decomposition tree. A concrete arithmetic or ANS coder driven by
these conditionals attains this rate up to the §2.3 redundancy
$\varepsilon$ (under-$2$-bit total overhead per stream for the
Witten--Neal--Cleary integer arithmetic coder, amortized to $o(1)$
per byte over a long block, by the §2.3 entropy-coding assumption;
Witten--Neal--Cleary 1987, Duda 2009). The \emph{ideal} independent-nibble
baseline, coding $H$ under $P(H \mid C)$ and $L$ under the
\emph{marginal} $P(L \mid C)$, pays exactly $H(H \mid C) + H(L \mid C)$;
this is the optimal factorized rate over the two nibble coordinates by
Theorem 5.5 part (F) (Shannon's converse for product-distribution
coders applied to $(H, L)$). The tightness statement is most cleanly
read off the two \emph{ideal} rates: the optimal byte-coder rate is
$H(H, L \mid C)$ (floor (6.2b), achieved by part (A) up to
$\varepsilon$) and the optimal independent-nibble rate is
$H(H \mid C) + H(L \mid C)$, so the gap is exactly $I(H ; L \mid C)$
with no slack. Equivalently, no byte coder's expected length drops
below $H(H, L \mid C)$, so the saving against the ideal baseline is
\emph{at most} $I(H ; L \mid C)$, and the §6.3 construction attains it
up to $\varepsilon$ --- the saving cannot exceed $I(H ; L \mid C)$ by
more than the entropy-coder redundancy.

A finite-length caveat on the \emph{sign} of the concrete gap: when
both sides are realized by concrete (integer Huffman or finite-block
arithmetic) coders, the two-stream baseline carries roughly twice the
per-stream redundancy of the single-stream joint coder, so the
\emph{measured} saving baseline-minus-joint can \emph{exceed}
$I(H ; L \mid C)$ by up to the extra baseline redundancy ($O(\varepsilon)$,
empirically up to about one bit in the worst degenerate case). This is
not a violation of the converse: the converse (6.2b) is the hard floor
on the \emph{joint} coder alone ($\mathbb{E}[L] \geq H(H,L\mid C)$),
which no concrete coder breaches; the inflated measured saving merely
reflects a weaker baseline coder, not a byte coder beating the joint
entropy. The clean tightness statement is therefore the ideal-rate one
above; the concrete saving equals $I(H ; L \mid C) \pm O(\varepsilon)$.
\hfill$\blacksquare$

\emph{Remark 6.3c (honest scope of the converse).} The converse
(C) is the standard Shannon conditional source-coding bound applied
to the nibble pair $(H, L)$ --- it introduces no new coding-theoretic
machinery beyond Theorem 4.2's conditional converse. Its content is
\emph{consequential, not technical}: it certifies that the
$I(H ; L \mid C)$ saving identified by Theorem 6.3 is the exact
maximum extractable by any uniquely decodable byte coder, upgrading
Theorem 6.3 from ``a saving of $I(H;L\mid C)$ exists'' to ``$I(H;L
\mid C)$ is the optimal saving, achieved up to $\varepsilon$.'' This
completes Type-III-B to a matched achievability/converse pair (as
Theorem 4.2 does for Type-I and Theorem 6.2 for the Type-III-A
mode-selection bound). Two caveats delimit the claim. First, the
$\varepsilon$ gap is genuine: the converse floor (6.2b) is on the
joint coder alone, and the construction meets it only up to the
entropy-coder redundancy $\varepsilon$, so Type-III-B is order-optimal
rather than exactly optimal at finite block length, with $\varepsilon$
vanishing only in the amortized long-block limit. The exact-tightness
claim is the \emph{ideal}-rate one (optimal joint rate $H(H,L\mid C)$
versus optimal independent-nibble rate $H(H\mid C)+H(L\mid C)$, gap
exactly $I(H;L\mid C)$); the concrete saving equals
$I(H;L\mid C)\pm O(\varepsilon)$ and may slightly \emph{exceed}
$I(H;L\mid C)$ because the two-stream baseline pays more redundancy
than the one-stream joint coder --- which is a statement about the
baseline, not a breach of the joint converse.
Second, the converse bounds the per-byte \emph{repair-stream} rate
$L(R \mid Y)$ of the decomposition (3.1); it does not bound the
constant model and metadata headers $L(M) + L(V) + L(\Theta)$, which
are $O(1)$ in $N$ and govern the block-level advantage condition of
Theorem 6.3 (the per-block crossover) rather than the per-byte rate.

\textbf{Theorem 6.4 (Hierarchical decomposition is
information-theoretically invariant).} \emph{Let $T$ be any rooted
binary tree whose leaves are the 256 byte values, defining a sequence
of nested binary partitions of the alphabet. Encoding a byte $b$ under
$T$ by traversing root-to-leaf and encoding the binary decision at
each node with the \emph{exact induced conditional probability} of
that decision under $P(B \mid C)$ given the prior decisions on the
path, with an ideal sequential entropy coder (e.g.\ arithmetic
coding in the asymptotic-redundancy limit), achieves expected
\emph{ideal} length $H(B \mid C)$ per byte. The expected ideal cost
is independent of the tree shape $T$; it depends only on the joint
distribution $P(B \mid C)$. Practical prefix-code constructions
(Huffman, etc.) achieve $H(B \mid C) + O(1)$ per byte; arithmetic
or ANS coders achieve $H(B \mid C) + o(1)$ amortized over many
bytes (Witten--Neal--Cleary 1987; Duda 2009). \emph{Caveat:} this
invariance fails if the decisions are coded with fixed branch
probabilities (e.g., 1/2 each) instead of the induced conditional
--- in that case the depth of $T$ dominates, and a left-skewed tree
pays substantially more than a balanced one. The theorem applies
only to coders that use the \emph{exact} induced conditional
probability at each node.}

\emph{\textbf{Proof.}} Fix any rooted binary tree $T$ on the 256
byte values. For each byte value $b$, let $\pi_T(b) = (z_1, z_2,
\ldots, z_{\ell_T(b)}) \in \{0,1\}^\star$ denote the unique
root-to-leaf path on $T$, of $T$-dependent depth $\ell_T(b)$ (a
path-length integer, not a code-length). The map $b \mapsto
\pi_T(b)$ is a bijection onto a prefix-free set of binary strings
(the leaves of $T$), so
$\Pr(\pi_T(B) = (z_1, \dots, z_{\ell_T(B)}) \mid C) = \Pr(B = b \mid C)$
for the unique $b$ whose path is $(z_1, \dots, z_{\ell_T(B)})$.
Encoding $b$ by sequentially passing the bits of $\pi_T(b)$ to an
ideal sequential entropy coder (Shannon-Fano arithmetic, range, or
ANS in the asymptotic-redundancy limit), with each bit's coding
probability set to its \emph{exact induced conditional}
$\Pr(z_l = 1 \mid z_1, \ldots, z_{l-1}, C)$, gives an
ideal-coder code length of
$-\log_2 \Pr(\pi_T(B) \mid C) = -\log_2 \Pr(B \mid C)$ bits per
byte. Taking expectation under $B \sim P(\cdot \mid C)$:
\[
\mathbb{E}_B[\text{ideal code length}] = \mathbb{E}_B[-\log_2 \Pr(B \mid C)] = H(B \mid C),
\]
which is independent of $T$. (Note that the \emph{path length}
$\mathbb{E}_B[\ell_T(B)]$ does depend on $T$, and is in general
much larger than $H(B \mid C)$ for unbalanced trees; the
invariance is the expected entropy-coded length, not the expected
path length.) \hfill$\blacksquare$

The practical implication is that the choice of decomposition tree
for Type-III-B (high-low, or any other binary split) affects
encoder/decoder complexity and predictor calibration, but not the
information-theoretic code length. The optimal tree for a specific
practical entropy coder is the one whose binary decisions are best
predicted by the available $Q$; this is a calibration question, not a
coding-rate question. In particular, the canonical high-low
decomposition is optimal when the predictor factorizes naturally as
$Q_{h,l}(i) = Q^H_h(i) \cdot Q^{L \mid H}_l(i, h)$, which is typical
for transformer predictors trained on byte streams.

\textbf{6.4 Worked numerical example for Type-III-B}

Consider an ASCII-text block where the empirical marginal entropies
are $H(H \mid C) = 2.4$ bits and $H(L \mid C) = 2.9$ bits per byte,
while the joint entropy is $H(H, L \mid C) = 4.7$ bits. Independent
nibble coding pays $2.4 + 2.9 = 5.3$ bits per byte; joint byte coding
pays $4.7$ bits per byte. The mutual-information advantage is
$\Delta_B = 5.3 - 4.7 = 0.6$ bits per byte, or 0.075 bytes per byte
(about 7.5\% saving) before model and metadata overhead. Over a
4096-byte block, this is 307 bytes of raw saving against independent
nibble coding, against an extra model cost that is typically a few
tens of bytes per amortized block.

\textbf{6.5 Type-III-C: token tensor coding}

Type-III-C generalizes Type-III-B from the byte alphabet to a learned
token dictionary $D = \{\tau_0, \ldots, \tau_{K-1}\}$. Tokens may be
single bytes, byte pairs, longer byte fragments, structural markers,
or --- most generally --- symbolic compositions in which a token's
definition references other tokens in the dictionary. The latter case
makes the dictionary a \emph{straight-line grammar}: a context-free
grammar with one production per non-terminal and a directed acyclic
dependency graph (Kieffer and Yang 2000; Charikar et al. 2005).
Practical dictionaries used by BPE, WordPiece, and MDL-greedy fall
into the flat case (no recursion); the framework also admits the
recursive case, and we use it in the complexity arguments of Theorems
6.7 and 6.8 below.

The model predicts a position-and-scale tensor $Q(k, i, s)$
approximating $\Pr(\tau_k \text{ begins at position } i \text{ at
scale } s \mid C)$, where $s$ is a length or scale class.

The dictionary is structured hierarchically as a grid of $K = a \cdot
b$ classes, with $\text{token\_id} = b \cdot \text{family} +
\text{variant}$ (matching Definition 3.5). For $K = 1024$, the canonical grid is
$32 \times 32$; for $K = 4096$, the grid is $64 \times 64$. The
hierarchical factorization
$Q(k, i, s) = P(\text{family} \mid i, C) \cdot
P(\text{variant} \mid \text{family}, i, C)$
reduces the output dimension of the predictor at each position from
$K$ to $a + b$, which is asymptotically a quadratic reduction in
softmax cost.

The $32 \times 32$ grid does not span all 65536 byte pairs. It is an
MDL-selected dictionary of 1024 useful fragments. Admission of a token
to $D$ is governed by the standard MDL token gain criterion,
formalized in the next theorem.

\textbf{Theorem 6.5 (Type-III-C marginal-MDL token admission, all
else frozen).} \emph{Let $\tau$ be a candidate token of byte length
$\ell$. Fix a parser that produces a non-overlapping token cover of
the corpus and let $f$ be the \emph{realized} count of $\tau$ in
that parse, restricted to positions currently parsed as raw
literals (i.e., not already covered by any token already in $D$).
This count can be strictly less than the raw occurrence count when
$\tau$ self-overlaps; e.g., for $\tau = $ \texttt{aa} in the
string \texttt{aaa}, the raw count is $2$ but $f = 1$. Positions
where $\tau$ would overlap an existing token are excluded from $f$
by the frozen-parser assumption below. Let
$\mathrm{code\_cost}(\tau)$ be the average bit length to encode an
occurrence of $\tau$ under the predictor and entropy coder,
$\mathrm{dictionary\_cost}(\tau)$ the bit length to describe $\tau$
in the dictionary header, and $\mathrm{parse\_overhead}(\tau)$ the
additional per-occurrence boundary or interval bookkeeping
introduced by including $\tau$. \emph{Crucially, assume that the
parser, the entropy-coder state, all other tokens' code costs,
dictionary costs, parse overheads, and the realized counts of all
other tokens are FROZEN during the admission test for $\tau$
(``ceteris paribus''); the rule evaluates only the marginal effect
of adding $\tau$ to an otherwise fixed dictionary. Adding $\tau$ in
practice may change the parse and the marginal costs of competing
tokens; the test below ignores those interactions, so it gives a
greedy admission decision rather than a globally optimal one.} The
net change in total Type-III-C code length under these frozen-rest
assumptions is}
\[
\Delta L(\tau) = \mathrm{dictionary\_cost}(\tau) + f \cdot \left[\mathrm{code\_cost}(\tau) + \mathrm{parse\_overhead}(\tau) - 8\ell\right].
\]
\emph{Consequently, the per-token greedy MDL admission policy admits
$\tau$ if and only if $\Delta L(\tau) < 0$, equivalently}
\[
f \cdot \left[8\ell - \mathrm{code\_cost}(\tau)\right] > \mathrm{dictionary\_cost}(\tau) + f \cdot \mathrm{parse\_overhead}(\tau).
\tag{6.3}
\]

\emph{\textbf{Proof.}} Without $\tau$ in $D$, the $f$ occurrences of
the byte pattern are coded as raw literals at $f \cdot 8\ell$ bits
total, with no dictionary contribution. With $\tau$ in $D$, the $f$
occurrences cost $f \cdot \mathrm{code\_cost}(\tau)$ for encoding plus
$f \cdot \mathrm{parse\_overhead}(\tau)$ for bookkeeping, and the
dictionary pays a one-time $\mathrm{dictionary\_cost}(\tau)$. The
difference of the two totals is $\Delta L(\tau)$. By the MDL
principle (Rissanen 1978), the admission policy that minimizes total
description length includes $\tau$ iff this difference is negative.
Rearranging gives (6.3). \hfill$\blacksquare$

The token MDL criterion is the standard greedy dictionary-learning
admission rule. Its placement here is to make explicit that Type-III-C
does not store every possible token, only those whose per-corpus
savings cover their description. The criterion is greedy in that it
evaluates each candidate $\tau$ marginally, assuming the rest of the
dictionary is fixed; it is not a global MDL optimum, since admitting
$\tau$ can change the optimal parse and the
$\mathrm{code\_cost}$ of other tokens. The closest computational
counterpart we prove is the restricted normal-form dictionary-only
problem of Theorems 6.7 and 6.8 below, which is shown to be NP-hard
and APX-hard with constant $\gamma_{\mathrm{SGP}} = 8569/8568$ under
the unbounded-alphabet variant; the fully
global Type-III-C problem (with model, dictionary, parser, and
repair stream all free) is conjectured at least as hard but its
formal hardness is open (§13.2). For a corpus encoded across many
archives,
$\mathrm{dictionary\_cost}$ is amortized and the greedy admission
criterion grows more permissive.

When tokens are continuous latent variables of a variational model
rather than discrete byte fragments, the appropriate encoding tool
is bits-back coding (Townsend, Bird, Barber 2019). The encoder uses
an auxiliary bit string to sample the latent, encodes the latent
and the data conditional on the latent via ANS, and the decoder
recovers both and extracts the auxiliary bits. Under ideal coding
with the true model marginal and the true posterior, the effective
rate equals the conditional source-coding lower bound
$H_D(X \mid Y)$ on the data distribution, not the joint
$H_D(X, Z \mid Y)$. In general (model marginal $\neq$ data
conditional, or $q \neq $ true posterior), the rate is the
cross-entropy under the model marginal plus the posterior KL gap,
which can exceed $H_D(X \mid Y)$. The next theorem makes this
rigorous and shows how bits-back fits into the $L(R \mid Y)$
decomposition of §3.3.

\textbf{Theorem 6.6 (Bits-back integration into the RNR
decomposition).} \emph{Let $z$ be a latent variable of a variational
model with prior $P(z \mid Y)$, conditional $P(X \mid z, Y)$, and
variational posterior $q(z \mid X, Y)$. Assume the following
support/regularity conditions: (a) the model marginal
$P_X(x \mid Y) = \sum_z P(z \mid Y) P(x \mid z, Y) > 0$ for every
$x$ in the support of the data distribution $D(\cdot \mid Y)$
(otherwise $-\log_2 P_X(X \mid Y)$ is infinite with positive
probability under $D$); (b) for almost every $X$ under $D$, the
variational posterior $q(z \mid X, Y)$ is absolutely continuous
with respect to the true posterior $P(z \mid X, Y)$ (otherwise the
posterior KL is infinite). \emph{Quantization convention for
continuous latents:} when $z$ is continuous, fix a finite-precision
discretization $\Delta$ of the latent space common to encoder and
decoder (e.g., $b$-bit fixed-point per coordinate as used by
Townsend-Bird-Barber's BB-ANS implementation), and read all
densities below as probability masses on the discretization cells.
With this convention $-\log_2 P(z \mid Y)$ and $-\log_2 q(z \mid X,
Y)$ are finite code lengths and the dictionary case is the trivial
discretization where $z$ takes values in a discrete set of token
identifiers (so $\Delta$ disappears). Under
bits-back ANS coding with auxiliary bits used to sample $z$ from $q$,
the effective Type-III-C repair length is
$L(R \mid Y) = L_{\text{prior}}(z \mid Y) + L_{\text{cond}}(X \mid z, Y) - L_{\text{aux}}(z)$,
where the three terms encode the latent under the prior, the data
under the conditional, and recover the auxiliary bits respectively.
Its expectation under $X \sim D$ and $z \sim q$ satisfies}
\[
\mathbb{E}[L(R \mid Y)] = \mathbb{E}_{X \sim D}[-\log_2 P_X(X \mid Y)]
+ \mathbb{E}_{X \sim D} \mathrm{KL}\!\left(q(z \mid X, Y) \,\Vert\, P(z \mid X, Y)\right),
\tag{6.4}
\]
\emph{where $P_X(\cdot \mid Y)$ is the model marginal over $X$ given
$Y$. Equivalently, the first term equals
$H_D(X \mid Y) + \mathrm{KL}(D(\cdot \mid Y) \Vert P_X(\cdot \mid Y))$,
the cross-entropy of the data under the model. The bits-back rate
equals the conditional entropy $H_D(X \mid Y)$ only when the model
marginal matches the data conditional ($D = P_X$) and $q$ equals the
true posterior; under both conditions matching the Slepian-Wolf lower
bound. At the rate level, Type-III-C with quantized continuous
latents recovers Type-III-C with a discrete dictionary as the
degenerate case where the discretization $\Delta$ identifies $q$
with a point mass on a single discrete token --- the discrete case
is one specific limit of bits-back ANS, not a strict subset of it
in construction. (The dictionary's interval-cover/tokenization
problem and the bits-back latent-variable problem are formally
different constructions; Theorem 6.6 proves the rate identity, not
a containment relation between the constructions themselves.)}

\emph{\textbf{Proof.}} Under the quantization convention $\Delta$,
the three-stream bits-back coder transmits
$\log_2 \tfrac{1}{P(z \mid Y)}$ bits for the latent under the
discretized prior, $\log_2 \tfrac{1}{P(X \mid z, Y)}$ bits for the
data under the conditional, and recovers
$\log_2 \tfrac{1}{q(z \mid X, Y)}$ bits from the auxiliary stream
(used to sample $z$). Net code length:
\[
\begin{aligned}
&-\log_2 P(z \mid Y) - \log_2 P(X \mid z, Y) + \log_2 q(z \mid X, Y) \\
&\quad= -\log_2 P(X, z \mid Y) + \log_2 q(z \mid X, Y) \\
&\quad= -\log_2 P(X \mid Y) + \log_2 \!\left[\frac{q(z \mid X, Y)}{P(z \mid X, Y)}\right],
\end{aligned}
\]
using the joint factorization
$P(X, z \mid Y) = P(X \mid Y) \cdot P(z \mid X, Y)$. Taking
expectation under $q(z \mid X, Y)$ conditional on $X$:
$\mathbb{E}[\text{net} \mid X] = -\log_2 P(X \mid Y) + \mathrm{KL}(q(z \mid X, Y) \,\Vert\, P(z \mid X, Y))$.
Taking expectation under $X \sim D$ gives (6.4). The first term is
cross-entropy under the model marginal $P_X(\cdot \mid Y)$, not
entropy unless $P_X(\cdot \mid Y) = D(\cdot \mid Y)$; the second
term is the variational gap. When both vanish the rate is
$H_D(X \mid Y)$. \hfill$\blacksquare$

Theorems 6.5 and 6.6 govern, respectively, \emph{which} tokens a
greedy builder admits and, \emph{given} a dictionary and a (model,
posterior) pair, the bits-back repair rate. Neither bounds the
expected Type-III-C code length $\mathbb{E}[L_{\mathrm{IIIC}}(X)]$ of
the \emph{learned} dictionary as a self-contained quantity --- the gap
the abstract flags. The next theorem closes it by the standard route
that dictionary/grammar codes are \emph{universal}: a Type-III-C
encoder whose dictionary is produced by a universal grammar transform
achieves the entropy rate, with the hierarchical index of
Definition~3.5 contributing only $o(N)$. This is achievability
\emph{under a universal dictionary builder}; whether the specific
greedy marginal-MDL builder of Theorem~6.5 is itself universal is
partially resolved below --- universal once (A) its literals are
entropy-coded under a universal predictor and (B) a one-bit baseline
guard is added (Lemma~6.6d), the bare flat-$8\ell$ rule remaining open
(Remarks~6.6b, 6.6e).

We first isolate the index-accounting step, which is where the
hierarchical structure of Definition~3.5 enters and where the new
content beyond the classical grammar-code theorems lies.

\textbf{Lemma 6.6a (The hierarchical token index is rate-neutral).}
\emph{Let $D = \{\tau_0, \ldots, \tau_{K-1}\}$ be a dictionary of $K$
tokens addressed by the grid factorization of Definition~3.5, $K = a
\cdot b$ with $\mathrm{token\_id} = b \cdot \mathrm{family} +
\mathrm{variant}$, $\mathrm{family} \in \{0,\ldots,a-1\}$,
$\mathrm{variant} \in \{0,\ldots,b-1\}$. Let $P(\cdot)$ be any
probability law over the $K$ token identifiers (the predictor's
position-marginal at a fixed position, say), and let the hierarchical
predictor factorize as
$Q(k) = P(\mathrm{family}(k)) \cdot P(\mathrm{variant}(k) \mid
\mathrm{family}(k))$ with the \emph{exact} induced conditionals. Then
the ideal two-stage code length equals the flat code length token by
token,}
\[
-\log_2 P(\mathrm{family}(k)) - \log_2 P(\mathrm{variant}(k) \mid \mathrm{family}(k))
= -\log_2 P(k) \qquad \text{for every } k,
\tag{6.5}
\]
\emph{so the expected hierarchical code length equals the flat
expected code length $H(P)$ exactly; the two-level family/variant
addressing adds zero bits at the ideal-coder level. At the level of a
practical entropy coder the only excess is the coder's per-symbol
redundancy applied twice instead of once, $O(1)$ bits per token, hence
at most $O(K)$ bits over the whole stream.}

\emph{\textbf{Proof.}} The map $k \mapsto (\mathrm{family}(k),
\mathrm{variant}(k)) = (\lfloor k/b \rfloor, \; k \bmod b)$ is a
bijection from $\{0,\ldots,K-1\}$ (extended to the full grid
$\{0,\ldots,ab-1\}$) onto the product index set. Hence
$P(k) = P(\mathrm{family}(k), \mathrm{variant}(k)) = P(\mathrm{family}(k))
\cdot P(\mathrm{variant}(k) \mid \mathrm{family}(k))$ by the chain
rule, and taking $-\log_2$ gives (6.5) identically. This is exactly
Theorem~6.3 (the Type-III-B tree-coding invariance) specialized to the
two-level tree whose first level is the $a$ families and whose second
level is the $b$ variants within a family: the ideal entropy-coded
length is invariant to the decomposition tree provided each decision
is coded with its exact induced conditional. Taking expectation under
$P$ gives $\mathbb{E}[-\log_2 P(k)] = H(P)$ on both sides. The
practical-coder remark is the standard $+O(1)$-per-symbol redundancy
of arithmetic/ANS coding (Witten--Neal--Cleary 1987; Duda 2009)
incurred once per coded symbol; the two-stage scheme codes two symbols
(family, then variant) per token rather than one, so the excess is at
most a constant per token, $O(K)$ total. \hfill$\blacksquare$

\textbf{Theorem 6.6b (Type-III-C standalone achievability under a
universal dictionary builder).} \emph{Consider a Type-III-C encoder
(Definition~3.5) in which the dictionary $D$ and the non-overlapping
token cover of the source block $X^N = (x_1,\ldots,x_N)$ over a finite
alphabet $\mathcal{A}$, $|\mathcal{A}| = \sigma$, are produced by a
\emph{universal grammar transform}, in either of the two standard
senses:}
\begin{itemize}
\item[(i)] \emph{(LZ78 parser.) $D$ is the set of LZ78 phrases of
$X^N$ under the incremental parse of Ziv--Lempel (1978): each phrase
is the shortest substring not previously seen, $D$ grows by one phrase
per parse step, and the token-reference stream is the LZ78 codeword
$(\mathrm{pointer}_i, \mathrm{symbol}_i)_{i=1}^{c(N)}$, where
$\mathrm{pointer}_i$ addresses the prefix-phrase already in $D$ and
$\mathrm{symbol}_i \in \mathcal{A}$ extends it.}
\item[(ii)] \emph{(Irreducible grammar.) $D$ is the rule set of an
\emph{irreducible} straight-line grammar $G_{X^N}$ representing $X^N$
in the sense of Kieffer--Yang (2000, \S IV.B): $G_{X^N}$ is admissible
($L(G) = \{X^N\}$, deterministic, no empty right-hand side, no useless
symbols), distinct variables expand to distinct strings, every
non-start variable occurs at least twice among the right-hand sides,
and no digram (pair of adjacent symbols) repeats in non-overlapping
positions across the right-hand sides. The token-reference stream is
the grammar encoding $B(G_{X^N})$.}
\end{itemize}
\emph{Address the $K := |D|$ dictionary entries by the hierarchical
grid of Definition~3.5. Then for every \emph{stationary} source $X$ on
$\mathcal{A}$ the expected Type-III-C code length satisfies}
\[
\mathbb{E}\!\left[L_{\mathrm{IIIC}}(X^N)\right] \;\le\; N \cdot h(X) + o(N),
\qquad\text{equivalently}\qquad
\frac{1}{N}\,\mathbb{E}\!\left[L_{\mathrm{IIIC}}(X^N)\right] \xrightarrow[N\to\infty]{} h(X),
\tag{6.6}
\]
\emph{where $h(X)$ is the entropy rate. The number of dictionary
entries is sublinear,}
\[
K = c(N) \;\le\; \frac{N}{(1-\varepsilon_N)\,\log_\sigma N} \;=\; O\!\left(\frac{N}{\log N}\right),
\qquad \varepsilon_N \to 0,
\tag{6.7}
\]
\emph{and the code splits into a reference (pointer) stream that
carries the entire entropy-rate term $N\,h(X) + o(N)$ and a genuinely
additive overhead --- the per-token new-symbol stream, the dictionary
header, and the index-addressing slack --- of size $O(K) = O(N/\log N)
= o(N)$. The pointer stream, of size $O(N)$, is \emph{not} overhead: it
is where the rate lives, equalling $N\,h(X) + o(N)$ (it is $\Theta(N)$
exactly when $h(X) > 0$; in the degenerate $h(X) = 0$ case ---
e.g.\ an eventually-constant source, where $c(N) = O(\sqrt{N})$ and the
pointer stream is itself $o(N)$ --- the bound holds a fortiori). The
hierarchical family/variant addressing is rate-neutral (Lemma~6.6a),
adding only $O(K) = o(N)$ practical-coder slack. If in addition $X$ is
\emph{ergodic}, the bound holds almost surely:
$\tfrac{1}{N} L_{\mathrm{IIIC}}(X^N) \to h(X)$ a.s. Under route (ii)
restricted to $k$th-order finite-state sources, the per-symbol
redundancy is the maximal pointwise bound $O(\log\log N / \log N)$
(Kieffer--Yang 2000, Theorems~7--8). Combined with the bits-back
refinement of Theorem~6.6 on the residual stream, the construction is
asymptotically optimal for stationary ergodic sources.}

\emph{\textbf{Proof.}} The Type-III-C archive under either route
consists of (a) the dictionary/grammar description and (b) the
token-reference stream that reconstructs the exact token cover (the
residual $R_C$ of Definition~3.5); the model $M$ here is the trivial
universal parser, carried in $O(1)$ bits, so by the decomposition
(3.1) the code length is $L_{\mathrm{IIIC}}(X^N) = (\text{header}) +
(\text{reference stream}) + O(1)$.

\emph{Route (i), LZ78.} The two parts together \emph{are} the LZ78
codeword: the $c(N)$ phrases of $D$ are exactly the distinct phrases of
the parse, and emitting $(\mathrm{pointer}_i, \mathrm{symbol}_i)$ for
each phrase both \emph{defines} the phrase (it is a previous phrase
plus one symbol, so the dictionary is reconstructed incrementally by
the decoder) \emph{and} places it in the cover --- there is no separate
header to pay. The total LZ78 length is
$\sum_{i=1}^{c(N)} \left(\lceil\log_2 i\rceil + \lceil\log_2
\sigma\rceil\right) + O(\log N) \le c(N)\,(\log_2 c(N) + \log_2 \sigma)
+ O(\log N)$. By the Lempel--Ziv phrase-count bound (Ziv--Lempel 1978;
Cover--Thomas, \emph{Elements of Information Theory}, Lemma in \S13.5,
$c(N) \le N/((1-\varepsilon_N)\log_\sigma N)$ with $\varepsilon_N \to
0$), (6.7) holds, so $c(N) = O(N/\log N)$ and the entire codeword
length is $\tfrac{1}{N} c(N) \log_2 c(N) + o(1)$ per symbol. The
pointwise LZ78 optimality theorem (Ziv--Lempel 1978; Cover--Thomas,
\emph{Elements of Information Theory}, 2nd ed., \S13.5, ``Optimality of
Lempel--Ziv algorithms'') states that for a \emph{stationary ergodic}
process the per-symbol LZ78 length converges \emph{almost surely} to
the entropy rate,
\[
\frac{1}{N}\,L_{\mathrm{LZ78}}(X^N) = \frac{1}{N}\,c(N)\log_2 c(N) + o(1)
\xrightarrow[N\to\infty]{} h(X) \qquad \text{a.s.},
\]
which is the almost-sure form of (6.6); the proof uses the
Shannon--McMillan--Breiman theorem and hence requires ergodicity. The
per-symbol length is uniformly bounded by $\log_2\sigma + o(1)$ (each
symbol costs at most a raw literal), so bounded convergence upgrades
the a.s.\ statement to the expectation statement
$\tfrac{1}{N}\mathbb{E}[L_{\mathrm{LZ78}}(X^N)] \to h(X)$; for a merely
\emph{stationary} (not necessarily ergodic) source the expectation
statement still holds by the ergodic decomposition (a stationary $X$ is
a mixture of ergodic components, $h(X) = \int h(X_\omega)\,d\omega$,
and the bounded per-component rates integrate), recovering (6.6) for
every stationary process. The
hierarchical re-indexing of the $c(N)$ phrase pointers into
$(\mathrm{family},\mathrm{variant})$ pairs is a bijection on the phrase
identifiers, so by Lemma~6.6a it leaves the ideal code length
unchanged and adds at most $O(c(N)) = o(N)$ practical-coder bits.

\emph{Route (ii), irreducible grammar.} The Type-III-C dictionary in
the recursive/straight-line case is precisely a straight-line grammar
(§6.5; Kieffer--Yang 2000), and the archive is the grammar-based code
$\epsilon_\phi(X^N) = B(G_{X^N})$ of Kieffer--Yang's \S V. Kieffer--Yang's
Lemma~2 shows that irreducibility alone forces the grammar to be small,
$|G_{X^N}|/N \le |\mathcal{A}|/N + 12\log|\mathcal{A}|/(\log N - 8\log
|\mathcal{A}| - 8) = O(1/\log N)$, so the irreducible transform is
asymptotically compact ($C_{\mathrm{irr}} \subseteq C_{\mathrm{ac}}$,
their \S V) with $\nu_N = 1/\log N \to 0$. Kieffer--Yang's Theorem~8
(irreducible transforms; equivalently Theorem~7 applied with this
$\nu_N$) then gives maximal pointwise redundancy
$\mathrm{Red}_N(\phi, \Lambda^k_{\mathrm{fs}}(\mathcal{A})) =
O(\gamma(\nu_N)) = O\!\left(\nu_N \log(1/\nu_N)\right) = O\!\left(\frac{\log\log
N}{\log N}\right)$
with respect to the family $\Lambda^k_{\mathrm{fs}}(\mathcal{A})$ of
$k$-state finite-state sources (in Kieffer--Yang's parameterization $k$
is the number of \emph{states} of the source automaton, not a Markov
order), for every $k$. An order-$m$ Markov source is a finite-state
source with $k = \sigma^m$ states, so it lies in
$\Lambda^{\sigma^m}_{\mathrm{fs}}$. Maximal pointwise
redundancy $\to 0$ means $\tfrac{1}{N}|B(G_{X^N})| \le \tfrac{1}{N}
\log_2 (1/\mu(X^N)) + O(\log\log N/\log N)$ uniformly over $X^N$ and
over $\mu \in \Lambda^k_{\mathrm{fs}}$; taking expectations under any
finite-state source $\mu = D$ (in particular an order-$m$ Markov source,
$k = \sigma^m$) recovers (6.6) for that source with the explicit
$O(\log\log N/\log N)$ rate. To pass from finite-state to a general
\emph{stationary} source $X$, Kieffer--Yang's Remark~(ii) (\S V) notes
that every code in $C_{\mathrm{ac}}(\mathcal{A})$ is weakly minimax
universal for the family $\Lambda_{\mathrm{sta}}(\mathcal{A})$ of all
stationary sources, via Markov approximation in the standard
iterated-limit order (the same device by which LZ78 universality is
extended from finite-order Markov to all stationary processes,
Cover--Thomas \S13.5): for each fixed $k$, the
maximal-pointwise-redundancy bound holds uniformly over
$\Lambda^k_{\mathrm{fs}}$, so taking $X$'s $k$th-order conditional law as
the source gives $\limsup_N \tfrac{1}{N}\mathbb{E}[|B(G_{X^N})|] \le
H(X_{k+1}\mid X_1^k)$; then letting $k \to \infty$ and using
$H(X_{k+1}\mid X_1^k) \downarrow h(X)$ (the conditional-entropy
characterisation of the entropy rate of a stationary source,
Cover--Thomas \S4.2) yields
$\limsup_N \tfrac{1}{N}\mathbb{E}[|B(G_{X^N})|] \le h(X)$; the matching
$\ge h(X)$ is the lossless-coding lower bound, giving (6.6). Lemma~6.6a
again makes the hierarchical addressing of the $|D|$ rules
rate-neutral.

\emph{Index/header accounting.} It is worth separating the term that
carries the rate from the term that is genuine overhead, since this is
where the hierarchical structure of Definition~3.5 must be accounted
for. The \emph{pointer/reference stream} costs $\sum_{i=1}^{c(N)}
\lceil\log_2 i\rceil = c(N)\log_2 c(N) + O(c(N))$ bits; this stream is
$O(N)$ and is precisely the term that converges to $N\,h(X)$, not
overhead. (For a positive-entropy source $c(N) = \Theta(N/\log N)$ and
$\log_2 c(N) = \Theta(\log N)$, so the stream is $\Theta(N)$; for a
degenerate $h(X)=0$ source it is $o(N)$ and the bound is loose, which
only helps.) The genuinely additive
\emph{overhead} comprises: (a) the per-token new-symbol stream, $c(N)
\log_2\sigma = O(N/\log N) = o(N)$; (b) the dictionary description ---
absent in route (i) where the pointer stream defines $D$ incrementally,
and in route (ii) the part of $B(G_{X^N})$ encoding $V(G)$, $T(G)$, and
the rule lengths, bounded by $|\mathcal{A}| + 4|G| + \lceil H(G)\rceil$
(Kieffer--Yang Theorem~6) $= O(|G|) = O(N/\log N) = o(N)$ since $|G| =
O(N/\log N)$; and (c) the hierarchical re-indexing slack, $O(K) =
O(N/\log N) = o(N)$ by Lemma~6.6a (the family/variant split is a
bijection and so costs no information; the two-stage practical coder
pays at most a constant per token). Note we do \emph{not} claim the
index width itself is $o(N)$: each of the $c(N)$ references carries an
address of $\Theta(\log K) = \Theta(\log N)$ bits, so the index
\emph{stream} is $\Theta(N)$ --- but that is the rate-bearing pointer
stream just accounted, not extra overhead. (If instead the dictionary
were \emph{strictly} sub-Lempel--Ziv, $K = o(N/\log N)$, then even the
full index $O(K\log K) = o(N)$ would be negligible; the $K =
\Theta(N/\log N)$ LZ78 regime is the demanding case and is handled by
charging the pointers to the rate.) Summing the $N h(X) + o(N)$ pointer
rate with the $o(N)$ overhead (a)--(c) gives (6.6).

\emph{Bits-back combination.} The residual stream $R_C$ may itself be
coded by bits-back ANS when the per-token model is a continuous latent
rather than an exact discrete pointer (Theorem~6.6); by (6.4) this
contributes the cross-entropy of the residual under the model marginal
plus the posterior KL gap, both $\to 0$ when the model marginal matches
the data conditional and $q$ is the true posterior, so the refinement
does not increase the leading $N h(X)$ term. Hence the combined
construction attains $N h(X) + o(N)$ for stationary ergodic sources.
\hfill$\blacksquare$

\textbf{Remark 6.6b (What this does and does not assume; greedy-MDL
universality is open).} The bound (6.6) is achievability
\emph{conditional on the dictionary builder being a universal grammar
transform} (LZ78, or any irreducible-grammar transform). This is the
natural standalone bound: it says the Type-III-C \emph{construction}
--- the hierarchical token-tensor archive of Definition~3.5 --- is not
itself an obstacle to optimality, because once the dictionary is built
universally the index and header are $o(N)$ and the reference stream
reaches $h(X)$. What it does \emph{not} claim is that the specific
greedy marginal-MDL admission rule of Theorem~6.5 produces an
irreducible (hence universal) grammar. Greedy grammar heuristics in
general need not be irreducible, and their worst-case grammar size can
exceed the smallest grammar by a factor bounded away from $1$ (Charikar
et al.\ 2005 prove constant-factor lower bounds for LZ78, \textsc{Bisection},
\textsc{Sequential}, and other greedy transforms against the optimal
grammar size); a large grammar inflates the header but, more
importantly, a \emph{non-irreducible} transform is not covered by
Kieffer--Yang's universality theorem, so its rate is not guaranteed to
reach $h(X)$. Whether the Theorem~6.5 greedy-MDL builder is universal
(equivalently, whether it can be post-processed into an irreducible
transform without changing its admission decisions) is left open. The
honest reading of Theorems 6.5, 6.6, and 6.6b together is: Theorem~6.5
gives a tractable greedy admission rule, Theorem~6.6b shows
\emph{some} tractable universal builder (LZ78, run as the dictionary
transform) already attains the entropy rate within the Type-III-C
container, and Theorem~6.6 shows the residual refinement is rate-exact;
matching the \emph{specific} greedy-MDL builder to universality remains
a gap. Practically, this means a Type-III-C implementation that wants a
provable achievability guarantee should build its dictionary with an
LZ78/irreducible-grammar pass (or verify irreducibility of its learned
dictionary) rather than rely on unconstrained greedy MDL.

\textbf{Theorem 6.6c (Type-III-C converse and asymptotic optimality).}
\emph{Let $X$ be a stationary source on $\mathcal{A}$ with entropy rate
$h(X) = \lim_{n} \tfrac{1}{n} H(X^n)$. For \emph{any} uniquely-decodable
lossless Type-III-C code (Definition~3.5) --- comprising the dictionary
$D$, the token-reference stream, and the repair stream, such that the
decoder reconstructs $X^N$ from the transmitted bits together with
decoder-reproducible state --- the expected code length obeys}
\[
\mathbb{E}\!\left[L_{\mathrm{IIIC}}(X^N)\right] \;\ge\; H(X^N) \;\ge\;
N\cdot h(X).
\tag{6.6c.1}
\]
\emph{Consequently, combining with the achievability bound (6.6) of
Theorem 6.6b, whenever the dictionary is produced by a universal
builder,}
\[
N\cdot h(X) \;\le\; \mathbb{E}\!\left[L_{\mathrm{IIIC}}(X^N)\right]
\;\le\; N\cdot h(X) + o(N),
\qquad
\tfrac{1}{N}\mathbb{E}\!\left[L_{\mathrm{IIIC}}(X^N)\right]
\xrightarrow[N\to\infty]{} h(X),
\]
\emph{so Type-III-C is \textbf{asymptotically rate-optimal} on every
stationary source under a universal dictionary builder.}

\emph{\textbf{Proof.}} The Type-III-C codeword is a binary string from
which the decoder, using only that string and decoder-reproducible
state (Definition~3.5), recovers $X^N$ exactly; hence $X^N \mapsto
\text{codeword}$ is injective and the code is uniquely decodable for
$X^N$. By McMillan's theorem (Cover \& Thomas 2006, Theorem 5.5.1) its
lengths satisfy the Kraft inequality, whence the Shannon source-coding
bound (Theorem 5.3.1) gives $\mathbb{E}[L_{\mathrm{IIIC}}(X^N)] \ge
H(X^N)$. For a stationary source the normalized block entropy
$\tfrac{1}{N}H(X^N)$ is non-increasing in $N$ and converges to $h(X)$
(Cover \& Thomas 2006, Theorem 4.2.1 and its monotonicity); a
non-increasing sequence stays above its limit, so $H(X^N) \ge N\,h(X)$
for every $N$, giving (6.6c.1). The shipped dictionary $D$ does not
evade the bound: $D$ is part of the transmitted (or
decoder-reproducible-from-codeword) description and its bits are
already counted in $L_{\mathrm{IIIC}}$; a deterministic,
decoder-reproducible model carries no information absent from the
codeword, so Shannon's bound applies to the total. Combining (6.6c.1)
with (6.6) yields the two-sided sandwich and the limit. \hfill$\blacksquare$

\emph{Remark 6.6c (scope --- the converse is unconditional, the
optimality is conditional).} The lower bound (6.6c.1) holds for
\emph{every} Type-III-C code, with no assumption on how $D$ is built ---
it is the standard entropy-rate floor (the value-add, as for the
Type-III-B converse Theorem 6.3b, is not new machinery but the
\emph{tightness} it certifies). The floor $H(X^N)$ is an
information-theoretic \emph{benchmark}, not an efficiently computable
encoder target: its exact value is $\#P$-hard --- $\mathrm{FP}^{\#P}$-complete
--- to compute (Remark 7.15g), $\varepsilon$-approximable in polynomial time
only under filter stability (Remark 7.15c); the bound's \emph{validity} is
of course independent of its computational cost, and the operative object the
code attains is the achievable scheme (the realized container, one forward
pass per symbol), not the exact optimum. The matching \emph{achievability}
(6.6) is conditional on a universal builder (Theorem 6.6b), so the
asymptotic optimality holds precisely under that condition; whether the
\emph{specific} greedy marginal-MDL builder of Theorem 6.5 is itself
universal is the open sub-question of Remark 6.6b, now \emph{partially
resolved}: Lemma 6.6d below shows the greedy-MDL builder is universal
under two benign specifications of its entropy stage (an entropy-coded
literal fallback and a one-bit baseline guard), and Remark 6.6e
localizes the residual that stays open to the literal flat-$8\ell$ rule.
With Theorems 6.6b and 6.6c the
Type-III-C achievability/converse cell is complete: \emph{tight under a
universal dictionary builder}.

\textbf{Lemma 6.6d (The greedy marginal-MDL builder is universal under
an entropy-coded literal fallback and a baseline guard).} \emph{Consider
the Type-III-C container of Theorem 6.6b encoding a block $X^N$ of a
stationary source $X$ on a finite alphabet $\mathcal{A}$, with the
following two specifications of the otherwise-abstract entropy stage that
Theorem 6.5 leaves open:}
\begin{itemize}
\item[(A)] \emph{(Entropy-coded literal fallback.) Both the
token-reference stream and the literal stream --- the positions left
uncovered by the dictionary, equivalently the length-one tokens --- are
coded under one common sequential predictor $Q$ that is universal for the
family of stationary sources, meaning $\tfrac1N \mathbb{E}[-\log_2
Q(X^N)] \to h(X)$ for every stationary $X$ and a.s.\ for every stationary
ergodic $X$. The default such $Q$ is the \textsc{Lz78} coder
(Cover--Thomas \S13.5), already serving as the universal builder of
Theorem 6.6b route~(i); a growing-order or countable mixture that
competes with every finite Markov order at $o(N)$ regret (Ryabko's
twice-universal measure, Ryabko 1984, or context-tree weighting
with depth $\to\infty$, Willems--Shtarkov--Tjalkens 1995) is
equally admissible. A \emph{fixed}-order
Krichevsky--Trofimov estimator or bounded-depth CTW is universal only for
finite-order (tree) sources, so it is admissible here only when the
source is known to have finite order. An uncovered symbol then costs its
model code length $-\log_2 Q(\cdot \mid \mathrm{past})$ rather than a flat
$8\ell$ raw-byte literal.}
\item[(B)] \emph{(Baseline guard.) The builder emits the shorter of the
two archives $\{$dictionary $D$ learned by the greedy marginal-MDL rule
of Theorem 6.5; empty dictionary $D = \varnothing\}$, prefixed by one bit
identifying which was used.}
\end{itemize}
\emph{Then, \textbf{regardless of which tokens the greedy rule of
Theorem 6.5 admits},}
\[
\frac{1}{N}\,\mathbb{E}\!\left[L_{\mathrm{IIIC}}(X^N)\right]
\xrightarrow[N\to\infty]{} h(X)
\]
\emph{for every stationary source $X$, and $\tfrac{1}{N}
L_{\mathrm{IIIC}}(X^N) \to h(X)$ almost surely for every stationary
ergodic $X$. That is, the specific greedy-MDL builder of Theorem 6.5 is
universal in the sense of Theorem 6.6b once (A)--(B) hold.}

\emph{\textbf{Proof.}} Write $L_\varnothing :=
L_{\mathrm{IIIC}}(X^N \mid D = \varnothing)$ and $L_D :=
L_{\mathrm{IIIC}}(X^N \mid D)$ for the realized archive lengths under the
empty and the greedy-learned dictionaries, both with the entropy
stage~(A).

\emph{Upper bound (achievability, admission-independent).} Under
$D = \varnothing$ every position is a length-one literal, so by~(A) the
archive is exactly the sequential $Q$-code of $X^N$,
$L_\varnothing = \sum_{i=1}^{N} -\log_2 Q(X_i \mid X^{i-1}) + O(1)$,
where the $O(1)$ is the trivial-parser/header of (3.1). Because $Q$ is
universal for the stationary family --- the defining property required in
(A) --- $\tfrac{1}{N}\mathbb{E}[L_\varnothing] \to h(X)$ (and a.s.\ in the
ergodic case). For the default $Q = \textsc{Lz78}$ this is the
empty-dictionary specialization of Theorem 6.6b route~(i) (the
Ziv--Lempel theorem, Cover--Thomas \S13.5). For a growing-order mixture
it is the finite-state-to-stationary Markov-approximation argument used in
the proof of Theorem 6.6b route~(ii): the mixture competes with every
fixed order $k$ at $o(N)$ regret, so $\limsup_N
\tfrac{1}{N}\mathbb{E}[L_\varnothing] \le H(X_{k+1}\mid X_1^k)$ for
\emph{every} $k$, and $H(X_{k+1}\mid X_1^k) \downarrow h(X)$ as
$k\to\infty$ (Cover--Thomas \S4.2). (A fixed-order KT would stall this
limit at its own order, which is why (A) excludes it for infinite-order
sources.) The guard~(B) ships an archive of
length $\min(L_D, L_\varnothing) + 1 \le L_\varnothing + 1$, so
\[
\tfrac{1}{N}\mathbb{E}[L_{\mathrm{IIIC}}]
\;\le\; \tfrac{1}{N}\mathbb{E}[L_\varnothing] + \tfrac{1}{N}
\;\xrightarrow[N\to\infty]{}\; h(X).
\]
Crucially this bound does not reference the admission decisions at all:
whatever dictionary the greedy rule of Theorem 6.5 produces --- optimal,
suboptimal, or actively harmful --- the guard never ships an archive
longer than the empty-dictionary universal code, the single guard bit
being $o(N)$.

\emph{Lower bound (converse).} The guarded archive is a
uniquely-decodable lossless Type-III-C code: the one guard bit makes the
two modes prefix-distinguishable, and each mode (the greedy archive, and
the all-literal $Q$-code, which is the degenerate Type-III-C code with
$D=\varnothing$ and empty reference stream) is itself uniquely decodable.
Hence Theorem 6.6c applies verbatim,
$\mathbb{E}[L_{\mathrm{IIIC}}] \ge H(X^N) \ge N\,h(X)$, i.e.
$\tfrac{1}{N}\mathbb{E}[L_{\mathrm{IIIC}}] \ge h(X)$.

The two bounds sandwich the rate at $h(X)$. In the stationary ergodic
case the convergence also holds almost surely, by the two matching a.s.\
bounds below. \emph{Upper branch:} the empty-dictionary archive is the
sequential $Q$-code, whose per-symbol length tends to $h(X)$ a.s.\ by the
\emph{pointwise} universality of the chosen coder for stationary ergodic
sources (for $Q = \textsc{Lz78}$ this is the Ziv--Lempel pointwise
optimality theorem, Cover--Thomas \S13.5, as in Theorem 6.6b route~(i);
for a growing-order mixture it is the a.s.\ form of its
stationary-ergodic universality),
so $\limsup_N \tfrac1N L_{\mathrm{IIIC}} \le \limsup_N \tfrac1N
L_\varnothing = h(X)$ a.s.\ by the guard. \emph{Lower branch:} unique
decodability gives the Kraft inequality $\sum_{x^N}
2^{-L_{\mathrm{IIIC}}(x^N)} \le 1$ (McMillan), so the likelihood-ratio
variable $W_N := 2^{-L_{\mathrm{IIIC}}(X^N)}/\mu(X^N)$ has
$\mathbb{E}_\mu[W_N] = \sum_{x^N} 2^{-L_{\mathrm{IIIC}}(x^N)} \le 1$;
Barron's lemma (Barron 1985) is then Markov's inequality on $W_N \ge 0$:
$\Pr[L_{\mathrm{IIIC}}(X^N) < \log_2(1/\mu(X^N)) - 2\log_2 N] =
\Pr[W_N > N^2] \le \mathbb{E}_\mu[W_N]/N^2 \le N^{-2}$, and since
$\sum_N N^{-2} < \infty$, Borel--Cantelli gives $L_{\mathrm{IIIC}}(X^N)
\ge \log_2(1/\mu(X^N)) - 2\log_2 N$ for all but finitely many $N$ a.s.;
dividing by $N$ and applying the Shannon--McMillan--Breiman theorem
$\tfrac1N \log_2(1/\mu(X^N)) \to h(X)$ a.s.\ yields $\liminf_N \tfrac1N
L_{\mathrm{IIIC}} \ge h(X)$ a.s. The two a.s.\ bounds give $\tfrac1N
L_{\mathrm{IIIC}}(X^N) \to h(X)$ a.s. \hfill$\blacksquare$

\textbf{Remark 6.6e (What Lemma 6.6d resolves of Remark 6.6b, and the
residual that stays open).} Lemma 6.6d does not settle the
open problem of Remark 6.6b --- enumerated as sub-questions (i)--(iii) in
\S13.2 --- as literally posed: those concern the \emph{bare} Theorem 6.5
rule, with (i) requiring an irreducible post-processing \emph{without
changing its admission decisions} and (ii) asking for \emph{direct}
universality of the rule on a natural source class. Instead it resolves a
closely related \emph{augmented} question: the Theorem 6.5 builder, run
with the two benign specifications (A)--(B) of its entropy stage, is
universal for \emph{every} stationary source, not merely for a special
source class. The mechanism is itself the point and is what bears on
\S13.2 --- it shows that under (A)--(B) \emph{the builder's admission
decisions are irrelevant to the rate}. Once the literal fallback is entropy-coded under a universal
predictor, the empty dictionary already reaches $h(X)$, the guard forbids
doing worse, and the converse forbids doing better; the dictionary is
then doing structural and computational work (random access, parse
locality, a smaller working model), not rate work --- the same separation
between a rate-bearing entropy stage and an access- or structure-bearing
layout that recurs in the random-access architecture of \S10
(Theorem 10.3). In particular the
Charikar et al.\ (2005) constant-factor grammar-size lower bounds do not
bear on the rate here: a grammar a constant factor larger than the
smallest grammar still occupies only $o(N)$ header bits, and \textsc{Lz78}
is itself the standing example of a greedy, grammar-size-suboptimal, yet
universal transform.

What stays open is the universality of the \emph{literal} Theorem 6.5
rule as written --- with the flat $8\ell$ raw-byte literal cost
(feature~(A) violated) and no guard (feature~(B) absent). There the
empty-dictionary archive costs $8N$ bits, not $N\,h(X)$, so the guard
cannot rescue it and the entire burden falls on the admissions: reaching
$h(X)$ then requires the greedy rule both to \emph{cover} a
$(1-o(1))$-rate fraction of the corpus with admitted tokens and to
entropy-code their references near $h(X)$ --- exactly the
smallest-grammar-adjacent question that Theorem 6.6b sidesteps by fiat.
Lemma 6.6d thus shows this residual is an artifact of the two
non-essential features of the literal rule, not of the Type-III-C
container; \S13.2 sub-question~(iii) (a source family on which
the literal flat-$8\ell$ rule is bounded away from $h(X)$), and likewise
(i)--(ii) for the bare rule, remain
genuinely open, with Charikar et al.\ (2005) bounding only grammar
\emph{size}, not \emph{rate}. The practical reading of Remark 6.6b is
unchanged and now sharpened: a Type-III-C implementation that wants a
provable rate guarantee should entropy-code its literals under a
universal predictor and keep the one-bit baseline guard, after which the
greedy-MDL dictionary may be used freely --- for its structural benefits
--- with no rate penalty.

\textbf{Theorem 6.7 (Restricted Type-III-C dictionary learning is
NP-hard).} \emph{Consider the restricted Type-III-C problem in which
$M$ is fixed to the trivial constant predictor, the repair stream is
the distinguished start token, and $D$ is a recursive straight-line
dictionary in normal form (defined below), with objective the
symbol-count $L_{\mathrm{sym}}(D) = \sum_i |\alpha_i|$. Finding the
$D$ minimizing $L_{\mathrm{sym}}$ subject to $\tau_0$ expanding to
the corpus $X$ is NP-hard.}

\emph{The unrestricted joint problem ($M$, $D$, repair stream all
free, total code length $L(M) + L(D) + L(R\mid Y)$) is conjectured
to be NP-hard but is not established by the present reduction:
dropping the restrictions changes both the feasible set and the
objective function, and for minimization with a larger feasible set
the optimum can be strictly smaller, so a polynomial-time
unrestricted solver does not automatically determine the restricted
optimum. Designing a reduction that forces the unrestricted optimum
to coincide with the restricted optimum is left open.}

\emph{\textbf{Proof.}} Reduce from the Smallest Grammar Problem (SGP)
of Charikar, Lehman, Liu, Panigrahy, Prabhakaran, Sahai and Shelat
(2005). An instance of SGP is a string $s \in \Sigma^\star$ (we may
take $\Sigma = \{0,1\}^8$ for concreteness); the task is to find a
straight-line grammar $G$ with $L(G) = \{s\}$ minimizing the
\emph{grammar size}
$|G| = \sum_{A \to \alpha \in P} |\alpha|$. SGP is NP-hard (Charikar
et al. 2005, Theorem 3.4).

Use \emph{normal-form} dictionaries: $\tau_0$ is the distinguished
start token; every token is reachable from $\tau_0$; the dependency
graph is acyclic; no token has a unit-production right-hand side; and
duplicate terminal expansions are identified. These are the standard
harmless normalizations for straight-line grammars: unreachable
rules, unit productions, and duplicate nonterminals can be removed
without increasing grammar size.

Given an SGP instance $s$, set $X := s$, fix $M$ to the trivial
predictor, and require $\tau_0$ to expand deterministically to $s$.
Define the hardness objective
\[
L_{\mathrm{sym}}(D) = \sum_i |\alpha_i|,
\]
where $\alpha_i$ is the right-hand side of token $\tau_i$ and each
terminal or token reference counts as one symbol. Normal-form
recursive dictionaries on $\Phi(s)$ are in bijection with normal-form
straight-line grammars deriving exactly $\{s\}$ --- map each token to
the homonymous non-terminal and vice versa --- and the bijection
preserves the objective value exactly:
\[
L_{\mathrm{sym}}(D(G)) = |G|.
\]
Thus an exact polynomial-time optimizer for this restricted
Type-III-C case would exactly solve SGP. Since SGP is NP-hard
(Storer-Szymanski 1982 for the original NP-hardness; Charikar et
al.\ 2005, Theorem 1, for the explicit $8569/8568$ APX-hardness
constant via reduction from bounded-degree-3 Vertex Cover and
Berman-Karpinski 1999), the restricted Type-III-C problem is
NP-hard. The reduction is computable in linear time in $|s|$.

The unrestricted-case claim is intentionally not proved by inclusion:
the restricted feasible region is a strict subset of the unrestricted
region, but the objective function also changes (the restricted case
uses $L_{\mathrm{sym}}$ while the unrestricted case uses
$L(M) + L(D) + L(R\mid Y)$), so the standard ``larger feasible set
$\Rightarrow$ smaller optimum'' minimization fact does not let an
unrestricted solver recover the restricted optimum. We leave the
formal unrestricted hardness as an open problem.
\hfill$\blacksquare$

Theorems 6.6 and 6.7 together delineate one slice of the
achievability boundary in Type-III-C. Bits-back gives the optimal
rate for any joint (model, posterior) pair (Theorem 6.6); finding
the optimal pair in the joint (model, dictionary, repair) problem
is not directly proven intractable by Theorem 6.7, which establishes
NP-hardness only for the restricted normal-form dictionary-only
case under the symbol-count objective $L_{\mathrm{sym}}$.
Conjecturally the unrestricted joint problem is at least as hard,
but a formal hardness reduction is open (see §13.2). Practical
implementations nonetheless use heuristic dictionary learning (BPE,
MDL-greedy, neural codebook learning) and amortized variational
inference, both of which sacrifice optimality for tractability.

\textbf{Theorem 6.8 (APX-hardness of joint
dictionary learning, restricted case, unbounded alphabet).}
\emph{Consider the restricted Type-III-C problem of Theorem 6.7
(constant predictor, recursive normal-form dictionary, symbol-count
objective $L_{\mathrm{sym}}$) over an \emph{unbounded terminal
alphabet} (the alphabet $|\Sigma|$ grows with the input, as in
Charikar et al.\ 2005's reduction, one terminal per vertex). This
restricted problem is APX-hard with inapproximability ratio at least
$\gamma_{\mathrm{SGP}} = 8569/8568$: no polynomial-time algorithm
achieves approximation ratio strictly below $8569/8568$ unless
$P = NP$. The constant is inherited verbatim from the Smallest
Grammar Problem (Charikar, Lehman, Liu, Panigrahy, Prabhakaran,
Sahai, and Shelat 2005, Theorem 1) via the exact symbol-count
bijection of Theorem 6.7. Two caveats apply:
(i)~\emph{Alphabet caveat.} The Charikar et al.\ reduction uses an
alphabet that grows linearly in the graph size (one terminal per
vertex), so $\gamma_{\mathrm{SGP}} = 8569/8568$ applies to the
\emph{unbounded-alphabet} variant of Type-III-C. For the fixed-byte
alphabet $|\Sigma| = 256$ of canonical RNR Type-III-C, encoding the
graph terminals as length-$\lceil \log_{256} |V| \rceil$ byte
strings adds a $\Theta(\log |V|)$ multiplicative overhead in the
grammar encoding. This overhead can dilute the constant
inapproximability gap arbitrarily close to $1$ as $|V| \to \infty$,
unless a separate gap-preserving construction is given. We do NOT
establish a constant gap for the fixed-byte case; whether one
exists is an open question, listed in §13.2. (ii)~\emph{Unrestricted joint problem caveat.} The
APX-hardness of the unrestricted joint problem is left open for the
same reason as in Theorem 6.7: the feasible set is strictly larger
and the objective function differs, so neither inclusion nor a
direct gap-preserving reduction is available.}

\emph{\textbf{Proof.}} We show that the reduction
$\Phi: s \mapsto \Phi(s)$ constructed in Theorem 6.7 preserves
approximability up to constant factors, then transfer the
APX-hardness of SGP.

\emph{Symbol-count normal form.} For the APX statement we use the
normal-form symbol-count objective of Theorem 6.7,
$L_{\mathrm{sym}}(D) = \sum_i |\alpha_i|$. The dictionary-to-grammar
map and grammar-to-dictionary map preserve this objective exactly, so
an $\alpha$-approximation for the restricted Type-III-C problem yields
an $\alpha$-approximation for SGP with no additive bookkeeping slack.

\emph{Transfer of APX-hardness.} Charikar, Lehman, Liu, Panigrahy,
Prabhakaran, Sahai, and Shelat (2005, Theorem 1, Section V.A,
page 2557) establish that no polynomial-time algorithm achieves
approximation ratio strictly below
$\gamma_{\mathrm{SGP}} = 8569/8568$ on SGP unless $P = NP$. Their
proof uses a gap-preserving reduction from bounded-degree-3 Vertex
Cover. The family considered satisfies the structural constraints
$|E| \leq (3/2)|V|$ (from the max-degree-3 hypothesis,
$2|E| = \sum_v \deg(v) \leq 3|V|$) and minimum vertex cover size
$k \geq (1/3)|V|$ (since each vertex covers at most 3 edges and
$|E| \geq |V|$ in Charikar et al.'s setup). Charikar et al.\
encode the graph as the grammar-input string
$\sigma = \prod_{v_i \in V}(\#v_i | v_i\#)^2 \prod_{v_i \in V}(\#v_i\#) \prod_{(v_i, v_j) \in E}(\#v_i v_j\#)$,
yielding a smallest grammar of size $15|V| + 3|E| + k$. The
Berman-Karpinski (1999) result that the corresponding restricted
Vertex Cover is hard to approximate below ratio $145/144$ unless
$P = NP$ then gives the ratio
$\rho(|E|, k) = \frac{15|V| + 3|E| + (145/144)k}{15|V| + 3|E| + k}$.
This ratio is monotone decreasing in $|E|$ and monotone increasing
in $k$; on the constraint region $|E| \leq (3/2)|V|$, $k \geq
(1/3)|V|$, the \emph{adversarial inapproximability lower bound}
$\rho_{\min}$ over the family is achieved at the extreme point
$|E| = (3/2)|V|$, $k = (1/3)|V|$ (where the ratio is smallest, since
we take the minimum over the family of hard instances):
\[
\rho_{\min} \;=\; \frac{15|V| + 3 \cdot (3/2)|V| + (145/144)(1/3)|V|}{15|V| + 3 \cdot (3/2)|V| + (1/3)|V|}
\;=\; \frac{8569}{8568}.
\]
Combined with the exact symbol-count bijection of Theorem 6.7, any
polynomial-time $\alpha$-approximation for the restricted
Type-III-C problem with $\alpha < 8569/8568$ would yield a
polynomial-time approximation to SGP below the established
threshold, contradicting $P \neq NP$. The unrestricted joint
problem is not addressed by this argument: as in Theorem 6.7, the
restricted feasible region is strict but the objective also
changes, so an unrestricted approximation does not project back to
a restricted-objective approximation. We leave the unrestricted
APX-hardness open. \hfill$\blacksquare$

\textbf{Remark on the constant.} The numerical value
$\gamma_{\mathrm{SGP}} = 8569/8568 \approx 1.000117$ is small (very
close to 1) and reflects the small algebraic gap inherited from
$145/144$ Vertex Cover inapproximability via the
$15|V| + 3|E| + k$ encoding overhead in the Charikar et al.\
reduction. Sharper inapproximability constants for SGP under
alternative reductions (e.g., from MAX-3SAT-3 via direct
gap-preserving encoding rather than via Vertex Cover) are an open
question in the SGP literature itself. The constant
$\gamma_{\mathrm{SGP}}$ inherited by Theorem 6.8 applies to the
\emph{restricted} symbol-count objective only; we do not claim
that this constant inherits to the \emph{fully unrestricted}
(free-predictor, free-dictionary, free-repair) bit-cost objective,
since that problem's APX-hardness is itself open (see §13.2).

\emph{Remark 6.8m (sharpening $\gamma_{\mathrm{SGP}}$ via stronger bounded-degree Vertex-Cover
hardness).} The value $8569/8568$ is not the best obtainable from the \emph{same} Charikar
reduction $\mathrm{OPT}=15|V|+3|E|+k$: it inherits Berman--Karpinski's $145/144$ at the
max-degree-3 adversarial extreme $(|E|=\tfrac32|V|,\,k=\tfrac13|V|)$. Substituting the stronger
Chleb\'ik--Chleb\'ikov\'a bounded-degree Vertex-Cover inapproximability gives a larger valid
constant. For the \emph{$3$-regular} family ($|E|=\tfrac32|V|$ and $k\ge|E|/3=\tfrac12|V|$),
hard within $100/99$, the extreme-point ratio is
\[
\left.\frac{15|V|+3|E|+(100/99)k}{15|V|+3|E|+k}\right|_{|E|=\frac32|V|,\,k=\frac12|V|}
=\frac{3961}{3960}>\frac{8569}{8568},
\]
a \emph{directly valid} improvement ($3$-regular $\subseteq$ Charikar's max-degree-3,
$|E|\ge|V|$ input). In general a $d$-regular family hard within $1+\delta$ gives
$\gamma_{\mathrm{SGP}}\ge 1+\delta/(31+3d)$. The $4$-regular family ($53/52$) gives the stronger
\[
  \gamma_{\mathrm{SGP}}\ \ge\ 2237/2236,
\]
and this is \emph{unconditional}: Charikar's exact-size identity $\mathrm{OPT}=15|V|+3|E|+\tau(G)$
holds for \emph{all simple graphs}, not only max-degree-3 ones. (In a simple graph each unordered
edge-block $\#v_i\#v_j\#$ occurs once and every $|$ is a fresh terminal, so the only repeated
length-$\ge2$ substrings are $\#v$, $v\#$, $\#v\#$ --- the characterization underlying the count
$8|V|+(3|V|-|C|)+3|E|+4|V|+2|C|=15|V|+3|E|+|C|$, minimized at $|C|=\tau(G)$ --- and \emph{no step
uses the degree bound}, which enters only the original ratio-dilution.) Moreover $2237/2236$ is
\emph{optimal} along this route: over the Chleb\'ik--Chleb\'ikov\'a $d$-regular constants
($100/99,53/52,51/50,49/48$ for $d=3,4,5,6$) the bound $1+\delta_d/(31+3d)$ peaks at $d^\ast=4$
($3961/3960,\,2237/2236,\,2301/2300,\,2353/2352$). Under UGC the Austrin--Khot--Safra degree-$d$
VC hardness $2-(2+o_d(1))\tfrac{\log\log d}{\log d}$ does \emph{not} improve this --- the
$1/(31+3d)$ dilution drives $\gamma\to1$ as $d\to\infty$, with no clean finite $d^\ast$. So
$2237/2236$ is the best constant this Charikar$+$bounded-degree-VC route yields; it sharpens, but
does not qualitatively change, the (still small) gap --- a similar update was noted by Winslow ---
and it does not address the \emph{bit-cost} objective (§13.2). Verified in
\texttt{scripts/verify/remark\_6\_8m\_sgp\_constant\_sharpening.py} (the $1+\delta/(31+3d)$
formula; $3961/3960$ and the unconditional $2237/2236$, both exceeding $8569/8568$; the
extreme-point reconciliation reproducing the baseline at $(\delta,c_E,c_T)=(1/144,3/2,1/3)$; and
the $d^\ast=4$ optimality over $d=3,\dots,6$).

\emph{Remark on the bit-cost objective.} A natural attempt is to
transfer Theorem 6.8's APX-hardness from the symbol-count
objective $L_{\mathrm{sym}}$ to the bit-cost objective
$L^{\mathrm{bit}}(D) = L_{\mathrm{sym}}(D) \cdot
\lceil \log_2(|\Sigma|+K) \rceil$ by arguing that the
multiplicative $\log$ factor cancels in the ratio
$L^{\mathrm{bit}}(G^{\mathrm{alg}})/L^{\mathrm{bit}}(G^*)$. This
naive transfer is unsound for two reasons: (i) a bit-cost
approximation algorithm need not produce a grammar that approximates
$L_{\mathrm{sym}}$, since the algorithm can trade off
$L_{\mathrm{sym}}$ against $K$ (which affects the log factor) in
ways that minimize bit-cost without minimizing $L_{\mathrm{sym}}$;
and (ii) within Charikar's reduction itself, the quantity
$15|V| + 3|E| + |\sigma|$ is the grammar size (count of
right-hand-side symbols) rather than the nonterminal count $K$;
in Charikar's structural family the nonterminal count includes
cover-dependent optional rules of the form $\#v_i\#$, so $K$ is
\emph{not} invariant across vertex-cover choices. Both observations
have been verified against Charikar et al.\ 2005 (the
$\#v_i\#$ rule accounting in their Section 3.2).

Lemma 6.8b below resolves the bit-cost APX-hardness question within
the Charikar normal-form grammar class on Berman-Karpinski-hard
instances: bit-cost minimization is APX-hard with the \emph{same}
constant $\gamma_{\mathrm{SGP}} = 8569/8568$ as the symbol-count
objective, via a ratio-preserving reduction that exploits the
monotonicity of $L^{\mathrm{bit}}$ in the cover size $|\sigma|$. An
earlier draft of this paper had a different (and incorrect) Lemma
6.8b claiming $\gamma^{\mathrm{bit}} = \gamma_{\mathrm{SGP}}$ via a
fixed-$K$ premise (i.e., $K = 15|V| + 3|E|$); that lemma was
retracted after adversarial review identified that Charikar's $K$
includes cover-dependent optional rules and is of the form
$K(\sigma) = K_0(H) + \kappa(H) \cdot |\sigma|$ for structural
$K_0, \kappa \geq 0$. The current Lemma 6.8b establishes the correct
result via a different argument (monotonicity of cost in $|\sigma|$
allowing bit-cost approximation to imply symbol-count approximation
with the same ratio).

\textbf{Lemma 6.8b (Bit-cost APX-hardness in Charikar normal form via
ratio-preserving reduction).} \emph{Let $\mathcal{H}_{\mathrm{BK}}$
denote the family of max-degree-3 graphs on which the
Berman-Karpinski (1999) $145/144$-inapproximability of minimum
vertex cover applies (hence Theorem 6.8's $\gamma_{\mathrm{SGP}} =
8569/8568$ inapproximability of the Charikar symbol-count objective
applies). For each $H \in \mathcal{H}_{\mathrm{BK}}$, let
$\mathcal{G}_{\mathrm{Char}}(H)$ be Charikar's normal-form grammar
class, with each $G \in \mathcal{G}_{\mathrm{Char}}(H)$ encoding a
vertex cover $\sigma \subseteq V$ and having}
\[
L_{\mathrm{sym}}(G) = 15|V| + 3|E| + |\sigma|, \quad
K(G) = K_0(H) + \kappa(H) \cdot |\sigma|,
\]
\emph{for structural constants $K_0(H) \geq 0$, $\kappa(H) \geq 0$.
Then bit-cost minimization over $\mathcal{G}_{\mathrm{Char}}(H)$ is
APX-hard with the same constant $\gamma_{\mathrm{SGP}} = 8569/8568$:
no polynomial-time $\alpha$-approximation algorithm for
$L^{\mathrm{bit}}$ over the class exists for $\alpha <
\gamma_{\mathrm{SGP}}$, unless $\mathrm{P} = \mathrm{NP}$.}

\emph{\textbf{Proof.}} On $\mathcal{G}_{\mathrm{Char}}(H)$ the bit-cost
\[
L^{\mathrm{bit}}(\sigma) = \underbrace{(15|V| + 3|E| + |\sigma|)}_{L_{\mathrm{sym}}(\sigma)} \cdot
\underbrace{\lceil \log_2(|\Sigma| + K_0(H) + \kappa(H) |\sigma|) \rceil}_{c(\sigma) \geq 1}
\]
is a strictly monotone non-decreasing function of $|\sigma|$: the
first factor is strictly increasing, the second is non-decreasing
($\kappa(H) \geq 0$). Therefore the bit-cost-optimal cover
$\sigma_{\mathrm{OPT}}^{\mathrm{bit}}$ coincides with the
symbol-count-optimal cover $\sigma_{\mathrm{OPT}}^{\mathrm{sym}}$
(both equal the minimum vertex cover of $H$). Denote this common
optimum by $\sigma^*$ with $|\sigma^*| = k_{\mathrm{OPT}}$.

Suppose for contradiction a polynomial-time algorithm $A$ achieves
bit-cost ratio $\alpha < \gamma_{\mathrm{SGP}}$: $A$ outputs a cover
$\sigma'$ with $L^{\mathrm{bit}}(\sigma') \leq \alpha \cdot
L^{\mathrm{bit}}(\sigma^*)$. Since both $L^{\mathrm{sym}}$ and $c$
are non-decreasing in $|\sigma|$ and the optimum $\sigma^*$ has
\emph{minimum} $|\sigma|$ among valid covers, $|\sigma'| \geq |\sigma^*|$,
hence $c(\sigma') \geq c(\sigma^*)$. Dividing both sides of the
bit-cost inequality by $c(\sigma') \geq c(\sigma^*) \geq 1$:
\[
L_{\mathrm{sym}}(\sigma') = L^{\mathrm{bit}}(\sigma') / c(\sigma')
\;\leq\; \alpha \cdot L^{\mathrm{bit}}(\sigma^*) / c(\sigma')
\;\leq\; \alpha \cdot L^{\mathrm{bit}}(\sigma^*) / c(\sigma^*)
\;=\; \alpha \cdot L_{\mathrm{sym}}(\sigma^*).
\]
So $A$ yields a symbol-count $\alpha$-approximation algorithm on the
same class. By Theorem 6.8 (BK + Charikar), the symbol-count
$\gamma_{\mathrm{SGP}}$-inapproximability rules out $\alpha <
\gamma_{\mathrm{SGP}}$. Hence no bit-cost $\alpha$-approximation
exists for $\alpha < \gamma_{\mathrm{SGP}}$.
\hfill$\blacksquare$

\emph{Scope and limits.} (i) The lemma applies to instances in
$\mathcal{H}_{\mathrm{BK}}$ only (the family on which BK + Charikar
gives $\gamma_{\mathrm{SGP}}$). Outside this family, the underlying
VC-3 hardness gap is smaller or absent. (ii) The class-restriction
to $\mathcal{G}_{\mathrm{Char}}(H)$ is essential: an unrestricted
bit-cost optimizer may produce non-Charikar grammars with $K$ not of
the form $K_0 + \alpha|\sigma|$, and the reduction above does not
apply. The unrestricted bit-cost APX-hardness problem remains open.
(iii) Ceiling jumps in $\lceil \log_2 \rceil$ are accommodated by the
non-decreasing-$c$ comparison; the reduction does not require strict
monotonicity of $c$. (iv) The lemma is an algorithmic APX-hardness
statement (no polynomial-time $\alpha$-approximator exists for $\alpha
< \gamma_{\mathrm{SGP}}$ over the class on $\mathcal{H}_{\mathrm{BK}}$),
not a per-instance pointwise gap claim on $\gamma^{\mathrm{bit}}_{\mathrm{Char}}(H)$.

\textbf{Lemma 6.8e (Continuous-bit-cost SGP inherits Charikar APX-hardness without the ceiling).}
\emph{Define the continuous (information-theoretic) bit-cost of a
context-free grammar $G$ with terminal alphabet $\Sigma$ and
nonterminal count $K(G)$ as}
\[
L^{\mathrm{cont}}(G) \;:=\; |G| \cdot \log_2\bigl(|\Sigma| + K(G)\bigr)
\]
\emph{(no ceiling --- the entropy-rate target achievable by arithmetic
coding, in contrast to the prefix-code worst case
$|G| \cdot \lceil \log_2(|\Sigma| + K(G)) \rceil$). Then unrestricted
SGP under the continuous bit-cost objective $L^{\mathrm{cont}}$ is
APX-hard: for every $\varepsilon > 0$, no polynomial-time
$(\gamma_{\mathrm{SGP}} - \varepsilon)$-approximation algorithm exists
unless $\mathrm{P} = \mathrm{NP}$, where $\gamma_{\mathrm{SGP}} = 8569/8568$
is Charikar's symbol-count inapproximability ratio.}

\emph{\textbf{Proof.}} Use the Berman-Karpinski $\to$ Charikar reduction
of Theorem 6.8. The reduction outputs, for each VC-3 instance, a source
string $s$ together with the optimum-grammar pair $G^*_{YES}, G^*_{NO}$
realizing the YES/NO branches:
\[
|G^*_{\mathrm{NO}}|_{\mathrm{sym}} \;\geq\; \gamma_{\mathrm{SGP}} \cdot |G^*_{\mathrm{YES}}|_{\mathrm{sym}}
\]
(Charikar et al.\ 2005 Theorem 6.5 / Lemma 4.3), where
$|G|_{\mathrm{sym}} = |G|$ is the symbol-count (rule-symbol total).
The nonterminal counts of the two optima satisfy
$K(G^*_{\mathrm{YES}}), K(G^*_{\mathrm{NO}}) = \Theta(|V| + |E|)$ with
$|K(G^*_{\mathrm{NO}}) - K(G^*_{\mathrm{YES}})| =
O(|C^* | \cdot (u-1)) = O(|V|/\gamma_{\mathrm{SGP}})$, where $u-1 =
1/8568$ is the BK cover-inflation factor and $|C^*| = O(|V|)$ is the
optimal-cover size. The terminal alphabet $|\Sigma| = O(|V|)$ (vertex
names plus structural symbols).

Setting $a := |\Sigma| + K(G^*_{\mathrm{YES}})$ and
$b := |\Sigma| + K(G^*_{\mathrm{NO}})$, both are $\Theta(|V|)$ and
$|b - a| = O(|V|/\gamma_{\mathrm{SGP}})$. A first-order expansion gives
\[
\frac{\log_2 b}{\log_2 a} \;=\; 1 + \frac{\log_2(b/a)}{\log_2 a} \;=\; 1 + O\!\left(\frac{|b-a|/a}{\log a}\right) \;=\; 1 + O\!\left(\frac{1}{|V| \cdot \log |V|}\right),
\]
which approaches $1$ as $|V| \to \infty$ at rate $O(1/(|V| \log |V|))$.

Continuous bit-cost ratio of the two optima:
\[
\frac{L^{\mathrm{cont}}(G^*_{\mathrm{NO}})}{L^{\mathrm{cont}}(G^*_{\mathrm{YES}})}
\;=\; \frac{|G^*_{\mathrm{NO}}|_{\mathrm{sym}}}{|G^*_{\mathrm{YES}}|_{\mathrm{sym}}} \cdot \frac{\log_2 b}{\log_2 a}
\;\geq\; \gamma_{\mathrm{SGP}} \cdot \bigl(1 - O(1/(|V| \log |V|))\bigr).
\]
For every fixed $\varepsilon > 0$, eventually for $|V|$ large enough
the perturbation $O(1/(|V| \log |V|)) < \varepsilon / \gamma_{\mathrm{SGP}}$,
so the continuous-bit-cost ratio is at least
$\gamma_{\mathrm{SGP}} - \varepsilon$. Hence any polynomial-time
$(\gamma_{\mathrm{SGP}} - \varepsilon)$-approximation algorithm for
continuous-bit-cost SGP would distinguish the YES and NO branches of
the VC-3 instance, contradicting Berman-Karpinski NP-hardness of VC-3
with the $1 + 1/144$ gap (Berman \& Karpinski 1999/2003).
\hfill$\blacksquare$

\emph{Scope and significance.} Lemma 6.8e closes a specific sub-case
of OP2(a): under the continuous (no-ceiling) bit-cost objective, the
unrestricted SGP problem inherits Charikar's $\gamma_{\mathrm{SGP}} =
8569/8568$ APX-hardness directly, without the class-restriction of
Lemma 6.8b. This shifts the residual difficulty of OP2(a) to a
\emph{specific artifact of the integer-bit objective}: the
ceiling-jump discontinuity at power-of-two boundaries of $|\Sigma| + K$
identified by the adversarial review for Conjecture 6.8d.
The continuous objective is the natural target for arithmetic coding
(the Shannon-McMillan-achievable fractional-bit rate), while the
integer/ceiling objective corresponds to a prefix-code worst-case
encoding rarely used in modern compression.

\emph{Comparison with Conjecture 6.8d.} Conjecture 6.8d posits some
$\varepsilon_{\mathrm{bit}} > 0$ for the ceiling-bit-cost objective ---
the weak form remaining after the ceiling-discontinuity attack
on the strong $\gamma_{\mathrm{SGP}} - \delta$ form. Lemma 6.8e gives
the strong $\gamma_{\mathrm{SGP}} - \varepsilon$ form for the
continuous-bit-cost objective unconditionally, modulo $\mathrm{P} \neq
\mathrm{NP}$. The two results coexist: Lemma 6.8e is the strong result
for the natural information-theoretic objective; Conjecture 6.8d
remains the open weak-form question for the discrete prefix-code
objective.

\emph{Implication for OP4' (fixed-byte-alphabet variant).} The
continuous-bit-cost framework also applies to fixed-alphabet variants:
encoding terminals as $\lceil \log_{256} |V| \rceil$-byte strings
introduces a $\Theta(\log_2 |V|)$ multiplicative overhead that is the
\emph{same} for YES and NO instances (both have the same alphabet,
the same terminal set), so the ratio is preserved up to the same
$1 + O(1/(|V| \log |V|))$ correction. Hence Lemma 6.8e extends to
the fixed-byte-alphabet variant under continuous bit-cost,
partially closing OP4'.

\textbf{Observation 6.8f (K-monotone discrete-bit-cost SGP retains
Charikar gap; the K-escape obstacle is isolated, not resolved).}
\emph{Honest framing: this observation is technically derivable from
Lemma 6.8b and Charikar's symbol-count gap, but is recorded
explicitly here to clarify the structural decomposition of the
OP2(a) residual difficulty. It does \emph{not} constitute a partial
closure of OP2(a) in the same sense as Lemma 6.8e (which closes
the continuous-bit-cost variant unconditionally); rather, it
relabels which sub-problem remains open.}

Define the K-monotone variant of unrestricted bit-cost SGP as the
optimization problem
\[
\min_{G} \;|G|_{\mathrm{sym}} \cdot \lceil \log_2(|\Sigma| + K(G)) \rceil
\quad \text{subject to} \quad K(G) \;\geq\; K_0,
\]
where $K_0$ is a fixed lower bound on the grammar's nonterminal
count, determined by the input. For instances produced by the
Berman-Karpinski $\to$ Charikar reduction, $K_0 := 2|V| + k$ can
be computed directly from the input VC-3 instance $(H, k)$
(Charikar et al.\ 2005 §V.A: the YES Charikar normal-form grammar
has exactly $2|V| + |C^*| = 2|V| + k$ nonterminals when the
optimal cover has size $|C^*| = k$, the BK promise threshold for
YES). Under this K-monotonicity constraint, the discrete bit-cost
APX-hardness gap is at least $\gamma_{\mathrm{SGP}} = 8569/8568$:
for every $\varepsilon > 0$, no polynomial-time
$(\gamma_{\mathrm{SGP}} - \varepsilon)$-approximation exists unless
$\mathrm{P} = \mathrm{NP}$.

\emph{\textbf{Proof.}} The Berman-Karpinski $\to$ Charikar reduction
(Theorem 6.8) outputs YES/NO instances with optimal Charikar
normal-form grammars satisfying $|G^*_{\mathrm{NO}}|_{\mathrm{sym}}
\geq \gamma_{\mathrm{SGP}} \cdot |G^*_{\mathrm{YES}}|_{\mathrm{sym}}$
and $K(G^*_{\mathrm{NO}}) \geq K(G^*_{\mathrm{YES}}) = K_0$. Under
the K-monotonicity constraint $K \geq K_0$, the YES optimum is
unchanged (the YES Charikar grammar is feasible) and the NO
optimum is at least the unconstrained NO Charikar grammar. The
combined bit-cost ratio:
\begin{align*}
\frac{L^{\mathrm{bit}}(G^*_{\mathrm{NO}, K\text{-mono}})}{L^{\mathrm{bit}}(G^*_{\mathrm{YES}, K\text{-mono}})}
  &\geq \frac{|G^*_{\mathrm{NO}}|_{\mathrm{sym}}}{|G^*_{\mathrm{YES}}|_{\mathrm{sym}}}
       \cdot \frac{\lceil \log_2(|\Sigma| + K(G^*_{\mathrm{NO}})) \rceil}{\lceil \log_2(|\Sigma| + K_0) \rceil}
  \;\geq\; \gamma_{\mathrm{SGP}} \cdot 1 \;=\; \gamma_{\mathrm{SGP}},
\end{align*}
where the second factor is $\geq 1$ by K-monotonicity. Any
polynomial-time $(\gamma_{\mathrm{SGP}} - \varepsilon)$-approximation
would distinguish YES from NO, contradicting NP-hardness.
\hfill$\blacksquare$

\emph{Honest assessment: what this Observation does \emph{not} do.}
The K-monotonicity constraint $K \geq K_0$ is automatically
satisfied by the Charikar normal-form NO optimum (because
$K(G^*_{\mathrm{NO}}) \geq K(G^*_{\mathrm{YES}}) = K_0$ in
Charikar's construction, as noted in Lemma 6.8e's proof). The
constraint only excludes \emph{K-escape grammars} --- non-Charikar
normal-form grammars with $K < K_0$ that might exploit ceiling-jump
savings. The Observation does \emph{not} prove that such K-escape
grammars fail to beat the gap; it merely excludes them from the
search space by fiat. Consequently:
\begin{itemize}
\item Observation 6.8f is mathematically correct but \textbf{does not
close any part of OP2(a) that was open before Lemma 6.8b}. It
is a relabeling of Lemma 6.8b's class-restriction (Charikar
normal form) as a K-monotonicity restriction.
\item The residual OP2(a) is precisely: ``does there exist a
non-Charikar-normal-form grammar $G$ for the Berman-Karpinski NO
instance with $K(G) < K_0$ and bit-cost ratio $< \gamma_{\mathrm{SGP}}$
relative to the YES optimum?'' This is the K-escape question
identified by the K-optimisation obstacle paragraph in §13.2;
Observation 6.8f does not address it.
\end{itemize}

\emph{Comparison with Lemmas 6.8b and 6.8e.}
\begin{itemize}
\item Lemma 6.8b: class-restricted (Charikar normal form), full
$\gamma_{\mathrm{SGP}}$ gap. \textbf{Substantive.}
\item Lemma 6.8e: continuous (no-ceiling) bit-cost, unrestricted grammar
class, full $\gamma_{\mathrm{SGP}}$ gap modulo $O(1/(|V|\log|V|))$.
\textbf{Substantive --- closes continuous variant of OP2(a).}
\item Observation 6.8f: K-monotone discrete bit-cost, full
$\gamma_{\mathrm{SGP}}$ gap. \textbf{Essentially a rephrasing of
6.8b; included only to clarify the K-escape framing.}
\end{itemize}
The residual OP2(a) is the K-free discrete bit-cost case: whether
K-escape grammars beat $\gamma_{\mathrm{SGP}}$ remains genuinely
open and is exactly the K-optimisation obstacle paragraph in §13.2
(reduction modification at $K = \sqrt{N}$ regime, or alternative
source-of-hardness problem).

\textbf{Lemma 6.8g (Joint Type-III-C cost is upper-bounded by
time-bounded conditional Kolmogorov complexity).} \emph{Let
$U$ be a fixed prefix-free universal Turing machine. Let
$t : \mathbb{N} \to \mathbb{N}$ be a polynomial time bound,
fixed independently of the input. For any source string
$s \in \Sigma^N$ and side-information string $Y \in \{0,1\}^*$,
define the time-bounded conditional Kolmogorov complexity
\[
  K^t(s \mid Y) \;:=\; \min\bigl\{\,|p|
  \;:\; p \in \{0,1\}^*,\ U(p, Y) = s \text{ in at most } t(|s|+|Y|)
  \text{ steps}\,\bigr\}
\]
(this is the Liu \& Pass (FOCS 2020 / J ACM 2023, arXiv:2009.11514)
polynomial-time-bounded variant, sometimes written $K^{\mathrm{poly}}$
in the literature; it is \emph{not} Levin's $\mathrm{Kt}(x) =
\min_p\{|p| + \log t(p)\}$, which augments the description length
with a logarithmic running-time penalty inside a single
optimisation). Define the
optimal joint Type-III-C bit-cost as
\[
  L^*_{\mathrm{total}}(s \mid Y) \;:=\;
  \min_{(M,D,R)\,:\,\mathrm{Dec}(Y,M,D,R)=s}
  \bigl[\, L(M) + L(D) + L(R \mid Y, M, D)\,\bigr],
\]
where the minimum ranges over Type-III-C admissible triples
$(M, D, R)$ for which the canonical decoder $\mathrm{Dec}$ recovers
$s$ exactly from $(Y, M, D, R)$, and each of $L(M), L(D),
L(R \mid Y, M, D)$ is the prefix-free binary description length of
the corresponding component. Then there exists an explicit universal
constant $c \in \mathbb{N}$ (depending only on $U$ and the
prefix-free encoding convention, not on $s$ or $Y$) such that for
every $s$ and $Y$,
\[
  L^*_{\mathrm{total}}(s \mid Y) \;\leq\; K^t(s \mid Y) + c\log_2(|s|+|Y|+1) + c.
\]
}

\emph{\textbf{Proof.}} Let $p^* \in \{0,1\}^*$ be a $K^t$-witness
program: $|p^*| = K^t(s \mid Y)$ and $U(p^*, Y)$ outputs $s$ within
$t(|s|+|Y|)$ steps. We construct a Type-III-C admissible triple
$(M_{p^*}, D_\emptyset, R_\emptyset)$ from $p^*$, then bound its
total encoded length.

\emph{Construction.}
\begin{enumerate}
\item \emph{Model.} Define $M_{p^*}$ as the degenerate
\emph{deterministic single-shot predictor}: on input context $Y$,
$M_{p^*}$ internally simulates $U$ on the pair $(p^*, Y)$ for up to
$t(|s|+|Y|)$ steps and outputs the resulting binary string as a
single deterministic block-level prediction $\hat{s}$ for the entire
source. Operationally, $M_{p^*}$ is encoded as the prefix-free binary
description of the pair $(p^*,\text{header})$, where the header is a
constant-length tag declaring the deterministic single-shot mode.
This mode is a permissible specialisation of the general Type-III-C
predictor: the model commits to a single output string with
degenerate probability mass $1$ on $\hat{s}$, the limit case in which
arithmetic coding under the predicted distribution produces zero
residual bits. The universal $U$ is hard-wired into the decoder by
convention (so does not need to be re-encoded), and the polynomial
$t$ is fixed once for the codec and likewise does not appear in
$L(M_{p^*})$ beyond a fixed-size constant.
\item \emph{Dictionary.} Set $D_\emptyset := \varepsilon$ (the empty
learned dictionary). The base alphabet $\Sigma$ remains implicitly
available as the universe of unit tokens, but no multi-byte learned
entries are introduced. Under the deterministic single-shot mode no
token table is consulted, because the predictor outputs the entire
string $\hat{s}$ directly; $L(D_\emptyset)$ is the constant-length
prefix-free encoding of ``dictionary has 0 learned entries.''
\item \emph{Residual.} Set $R_\emptyset := \varepsilon$ (the empty
residual). Under the deterministic single-shot mode, arithmetic
coding under the degenerate predicted distribution
$\Pr[\hat{s}] = 1$ produces zero codeword bits in the ideal limit,
so the residual is empty and its prefix-free length-prefix encoding
is the natural-number-zero codeword, of constant length.
\end{enumerate}

\emph{Decoder consistency.} The canonical Type-III-C decoder reads
the prefix-free header of $M_{p^*}$, parses the deterministic
single-shot mode flag and recovers $p^*$, runs $U(p^*, Y)$ for
$t(|s|+|Y|)$ steps, takes the output $\hat{s}$ as the predicted
string, and since $R_\emptyset = \varepsilon$ emits $\hat{s}$
verbatim. By construction $\hat{s} = s$, so the triple
$(M_{p^*}, D_\emptyset, R_\emptyset)$ is Type-III-C admissible under
the deterministic single-shot extension.

\emph{Length bound.} Each component is encoded with a prefix-free
length code (Elias gamma or any other concrete prefix-free natural
number code; the choice affects the universal constant $c$ but not
the asymptotic form). Let $c_0$ be a constant absorbing the
description of the canonical Type-III-C decoder, the polynomial $t$,
the mode-flag header, and any fixed structural overhead. Then:
\begin{align*}
  L(M_{p^*}) &\leq |p^*| + \ell_{\mathrm{pref}}(|p^*|) + c_0
    \;=\; K^t(s \mid Y) + O(\log K^t(s \mid Y)) + c_0, \\
  L(D_\emptyset) &= \ell_{\mathrm{pref}}(0) \;\leq\; c_1, \\
  L(R_\emptyset \mid Y, M_{p^*}, D_\emptyset) &= \ell_{\mathrm{pref}}(0)
    \;\leq\; c_1,
\end{align*}
where $\ell_{\mathrm{pref}}(n) = O(\log(n+1))$ is the Elias-gamma
prefix-free encoding length of a natural number $n$, and $c_1$ is the
encoding length of the natural number $0$ (a small fixed constant
under any standard scheme). Summing,
\[
  L^*_{\mathrm{total}}(s \mid Y)
  \;\leq\; L(M_{p^*}) + L(D_\emptyset) + L(R_\emptyset \mid Y, M_{p^*}, D_\emptyset)
  \;\leq\; K^t(s \mid Y) + c\log_2(|s|+|Y|+1) + c
\]
for $c := c_0 + 2 c_1 + c_2$, where $c_2$ absorbs the implicit
$O(\log K^t)$ term using $K^t(s \mid Y) \leq |s| \log_2 |\Sigma| + O(1)$
(the trivial print-program upper bound, so $\log K^t = O(\log(|s|+1))$).
The constants $c_0, c_1, c_2$ are all explicit functions of the
choices of $U$, $t$, and the Elias-gamma convention; in particular
$c$ does not depend on $s$ or $Y$.
\hfill$\blacksquare$

\emph{Honest scope statement.} Lemma 6.8g is an information-theoretic
upper bound. The following are \emph{not} claimed and do not follow
from this lemma alone:
\begin{itemize}
\item \textbf{No hardness.} Lemma 6.8g makes \emph{no} hardness
claim --- neither decision-hardness, nor APX-hardness, nor any
average-case hardness. The bound goes in the direction
$L^*_{\mathrm{total}} \leq K^t + O(\log)$, which is the
\emph{achievability} direction: an upper bound on the optimum cannot
imply hardness of the optimisation. The converse direction
$K^t \leq L^*_{\mathrm{total}} + O(\log)$ is the easy direction
(the canonical Type-III-C decoder, supplied with the Type-III-C
encoding $(M,D,R)$ as a program, is a polynomial-time machine
outputting $s$ from $Y$, so $K^t(s \mid Y) \leq L^*_{\mathrm{total}}(s \mid Y) +
O(1)$ when the codec runtime is bounded by the chosen $t$); together
the two directions give an $O(\log)$-tight characterisation of
$L^*_{\mathrm{total}}$ in terms of $K^t$, but neither direction is a
hardness statement on its own.
\item \textbf{Does not close Conjecture 6.8d.} Conjecture 6.8d
(unrestricted bit-cost SGP APX-hardness) is a worst-case
NP-hardness statement, which neither follows from nor implies a
$K^t$-upper-bound. Lemma 6.8g does, however, formalise the upper
bound informally invoked in §13.2's Liu--Pass paragraph, fixing the
K-variant from Levin's $\mathrm{Kt}$ to the Liu--Pass polynomial-time
$K^t$ and supplying the previously missing explicit constant.
\item \textbf{Does not imply Type-III-C $\geq K^t - O(1)$ in any
useful structural direction.} The reverse-direction inequality
$K^t \leq L^*_{\mathrm{total}} + O(1)$ holds easily when the codec
runtime is allowed to scale with the chosen polynomial $t$, but if
$t$ is fixed first and the codec is restricted to a smaller runtime
budget, $L^*_{\mathrm{total}}$ may strictly exceed $K^t$ by an
unbounded amount; Lemma 6.8g asserts only the one-sided upper bound.
\item \textbf{Reference $K^t$ is not Levin $\mathrm{Kt}$.} The
description-length notation $K^t(\cdot \mid \cdot)$ used here is the
Liu--Pass time-bounded Kolmogorov complexity with a fixed polynomial
$t$ (Liu \& Pass, FOCS 2020 / J ACM 2023, arXiv:2009.11514; the
inner optimisation is over programs $p$ with $U(p, Y) = s$ in
$\leq t(|s|+|Y|)$ steps), while
Levin's $\mathrm{Kt}(x) = \min_p \{|p| + \log t(p) : U(p) = x\}$
incorporates the running time additively inside a single
optimisation. The two measures coincide up to $O(\log)$ factors when
the chosen $t$ matches the running time of the optimal program, but
diverge in general; the formal statement of Lemma 6.8g uses $K^t$
because that is the measure for which Liu--Pass establish the
OWF-equivalence anchor cited in §13.2.
\item \textbf{Constant $c$ is explicit but architecture-dependent.}
The constant $c$ depends on the choice of universal machine $U$, the
chosen polynomial $t$, and the prefix-free length-encoding
convention. It does \emph{not} depend on $s$, $Y$, or any varying
parameter of the source. Different conventions yield different $c$;
the universality of the bound is in the asymptotic form
$O(\log(|s|+|Y|+1))$, not in any specific small constant.
\item \textbf{Deterministic single-shot mode requires a minor
extension of Definition 3.5.} The construction uses a degenerate
predictor that commits to a single full-block output string rather
than producing a per-position token distribution over the Type-III-C
dictionary $D$. Definition 3.5 in §3 frames the model as predicting
$Q(k, i) \approx \Pr(\tau_k \text{ begins at position } i \mid C)$
over dictionary tokens; the deterministic single-shot predictor used
in this proof is a permissible \emph{extension} of that definition
in which the predictor instead emits a single full-block string
$\hat{s}$ as a degenerate distribution placing probability $1$ on
$\hat{s}$ and probability $0$ on every other string. Under any
$L^*_{\mathrm{total}}$ optimisation that includes this single-shot
mode as an admissible specialisation, the bound holds. If
Type-III-C is interpreted \emph{strictly} as Definition 3.5
(per-position token-distribution form only, with no full-block
degenerate mode), the lemma's bound applies to
$\widetilde{L}^*_{\mathrm{total}}$ defined over the extended
admissible class. The extension is natural and consistent with the
general MDL-coder / arithmetic-coding framework, but we flag it
explicitly to preempt the reading that Lemma 6.8g is a statement
about Definition 3.5 verbatim.
\end{itemize}

\emph{What the lemma is useful for.} Lemma 6.8g formalises the upper
bound informally asserted in §13.2's Liu--Pass paragraph, replacing
the loosely-stated ``$L(M) + L(D) + L(R \mid Y, M, D) \leq
\mathrm{Kt}(s \mid Y) + O(\log N)$'' with the correctly-typed
inequality $L^*_{\mathrm{total}}(s \mid Y) \leq K^t(s \mid Y) +
c \log_2(|s|+|Y|+1) + c$. It is the achievability half of an
$O(\log)$-tight characterisation of the optimal joint cost. It is
\emph{required as a stepping stone} for any future OWF-conditional
hardness reduction from $\mathrm{MK^t P}$ to joint Type-III-C
compression on samplable distributions: the reduction needs both an
upper bound (this lemma) to encode reduction outputs and a lower
bound (the easy converse) to translate decision queries; with both
in hand the reduction has the same asymptotic shape as Liu--Pass's
MKtP analysis. Constructing the polynomial-time samplable distribution
on which such a reduction would work remains open (see §13.2).

\emph{Caveat on $Y$-side handling.} The conditional $K^t(s \mid Y)$
is well-defined provided $Y$ is given to the universal machine as a
fixed side-input string. In the Type-III-C framework, $Y$ is the
\emph{decoder context} fixed before the encoder operates and
identical for encoder and decoder. The construction above treats $Y$
as a constant-time-accessible auxiliary input to $U$; if instead $Y$
itself must be reproduced from a shorter description, the bound
needs the additional term $K^t(Y) + O(\log |Y|)$ added on the
right-hand side. We assume throughout the standard Type-III-C
convention that $Y$ is part of the public state and not counted in
the codec cost; under that convention the bound as stated holds
verbatim.

\textbf{Remark 6.8b (Why padding-based reduction fails for
unrestricted-predictor APX-hardness).} \emph{The natural attempt
to transfer APX-hardness from the restricted (constant predictor,
normal-form recursive dictionary) case (Theorem 6.8) to the
\emph{unrestricted joint} Type-III-C problem $(M, D, R \mid Y)$ via
incompressibility-style padding fails for a specific structural
reason, formalized below.}

\emph{Attempted reduction.} Given a restricted Type-III-C instance
$(s, \tau_0)$ with $L_{\mathrm{sym}}^*$ the optimal symbol-count under
the constant-predictor restriction, append an incompressible suffix
$z$ of length $|z| = \mathrm{poly}(|s|)$ with Kolmogorov complexity
$K(z) \geq |z| - O(\log |z|)$ (such a $z$ exists for each input
length by a counting argument; turning this into a polynomial-time
reduction requires \emph{constructible} or \emph{conditional}
incompressibility relative to $s$ and $M$, e.g.\ via Yao's
incompressibility theorem against polynomial-time computable
predictors, which is non-trivial in general). Form the joint Type-III-C instance
$(M, D, R \mid Y)$ on the concatenation $s z$. Naively, an
$\alpha$-approximation for the joint bit-cost would translate to an
$\alpha$-approximation for the restricted symbol-count, since the
suffix $z$ is incompressible.

\emph{Why this fails.} An unrestricted predictor $M$ can specialize:
$M$ can pay $L(M) \geq |s|$ bits of model description to encode the
exact mapping $s \mapsto \tau_0$, then for the prefix $s$ produce an
arbitrarily small $L(R \mid Y)$ (e.g., $L(R \mid Y) = O(1)$) by
direct lookup, while still requiring $L(R \mid Y \restriction z) \geq
K(z)$ for the incompressible suffix. The total
joint bit-cost is then $L(M) + L(D) + L(R \mid Y) \approx |s| + L(D)
+ K(z)$, dominated by the predictor description $|s|$ and the
incompressible suffix $K(z)$, with the dictionary cost $L(D)$ pushed
to negligible levels. The restricted symbol-count
$L_{\mathrm{sym}}^*$ for $s$ is no longer reflected in any term of
the unrestricted bit-cost: specialization absorbs the SGP structure
into $L(M)$, and Theorem 6.8's $\gamma_{\mathrm{SGP}}$ inapproximability
is lost in the trade-off.

\emph{What would close OP2(b).} A successful reduction would need to
either (i) bound $L(M)$ \emph{a priori} below the threshold at which
specialization becomes profitable, restricting the class of
predictors (note that ``polynomial-size predictors'' alone is
insufficient: a deterministic lookup table for $s$ has size
$O(|s| \log |\Sigma|) = \mathrm{poly}(|s|)$, so polynomial-size
predictors can already specialize; a tighter bound such as
$L(M) \ll L_{\mathrm{sym}}^*$ would be needed); or (ii) construct a source family such that any
near-optimal predictor description $M$ \emph{must} encode the SGP
structure faithfully, preventing escape via specialization; or
(iii) use an entirely different gap-amplification framework
(e.g., Dinur-style robust PCP composition, c.f.\ Remark 4.3c) that
operates on the unrestricted Type-III-C objective directly. None of
these is currently known; the unrestricted joint Type-III-C
APX-hardness problem remains a research-scale open question.

\emph{Failed attempt at OP2(b) closure (retracted Lemma 6.8c).}
An earlier draft introduced a Lemma 6.8c claiming joint Type-III-C
APX-hardness under a ``CFG-format predictor restriction'' (i.e., $M$
itself a CFG concatenable with $D, R$). The proof used a Kolmogorov
complexity lower bound $L^*_{\mathrm{total}} \geq L^*_{\mathrm{sym}}
- O(\log |s|)$ and a concatenation reduction to SGP. Adversarial review
identified four structural breaks: (i) the Kolmogorov bound was
applied in the wrong direction (invariance only gives upper bounds on
$K(s)$, not lower bounds); (ii) the approximation algebra confused
$L^*_{\mathrm{total}} \leq L^*_{\mathrm{sym}}$ (which holds, achievable
by constant $M$) with $L^*_{\mathrm{total}} \geq L^*_{\mathrm{sym}}
- O(\log)$ (which does not have a clean justification); (iii) the
``concatenation $M \cup D \cup R$ is a valid CFG'' step requires
explicit specification of disjoint nonterminal namespaces, start
symbols, and rule ordering, none of which the lemma fixed formally;
and (iv) the CFG-format restriction collapses the problem to SGP
itself (with rules repartitioned among buckets), so the lemma reduces
to Theorem 6.8 with relabeled variables rather than a new APX-hardness
statement. The lemma was retracted; OP2(b) under arbitrary
computational $M$ remains genuinely open as a research-scale problem,
per Remark 6.8b's characterization of why naive padding-based
reductions fail.

Together, Theorems 6.7 and 6.8 show that the restricted normal-form
recursive dictionary problem under the symbol-count objective is
NP-hard and (under the unbounded-alphabet variant) APX-hard with
constant $8569/8568$. We do not claim this extends to all
practical dictionary learners --- BPE and WordPiece operate on a
fixed-size flat token vocabulary with greedy merges (not a
recursive grammar), MDL-greedy variants use bit-cost rather than
symbol-count objectives, and many production tokenizers solve
different problem variants. The heuristic nature of those
algorithms is consistent with the restricted hardness result but
is not formally derived from it; the unrestricted bit-cost and
flat-vocabulary cases are open.

\textbf{7. Composition and mode selection}

RNR types compose by hierarchical mode selection. The top-level
scheduler is a Type-III-A coder whose mode set $M$ contains the
available lower-level coders, including all Type-I, Type-II, and
Type-III-B/C constructions. For each tile, the scheduler predicts a
per-mode cost vector $U_*(t)$, selects the argmin mode, and delegates
encoding to that mode.

\textbf{Theorem 7.1 (Composition correctness, with explicit framing).}
\emph{Let $M$ be a finite set of RNR mode coders, each producing a
uniquely decodable repair stream conditional on its mode flag.
Assume the per-tile repair streams are \emph{framed} --- either each
carries an explicit byte-length prefix in the metadata $L(\Theta)$,
or the per-tile encoding is itself prefix-free as a binary string
--- so that the decoder, after parsing the mode flag, can isolate
the bits belonging to tile $t$'s repair stream before applying the
mode-specific decoder. Fix a mechanism by which the decoder
recovers the mode of tile $t$:}
\begin{itemize}
\item[\textbf{(F)}] \emph{(\emph{Explicit flag}) The mode of every tile
is stored as a uniquely decodable flag in metadata; the decoder
reads the flag and decodes the tile under that mode.}
\item[\textbf{(A)}] \emph{(\emph{Argmin reproduction}) The cost field
$U_*(t)$ is decoder-reproducible (computable from public state and
the verified prefix of $X$), and the encoder commits to
$\hat{m}(t) = \operatorname*{arg\,min}_m U_m(t)$ with fixed
lexicographic tie-breaking from Definition 3.3; the decoder
reproduces $U_*(t)$ and recomputes the same argmin. Stored mode
flags are then required only on tiles where the encoder departs
from the predicted argmin (override flags).}
\end{itemize}
\emph{Under framing and either (F) or (A), the Type-III-A composition
produces a uniquely decodable codeword for $X$.}

\emph{\textbf{Proof.}} The decoder reads the metadata (including the
per-tile length prefixes, if present, that fix the bit boundaries
between successive tiles). It recovers $\hat{m}(t)$ either from the
stored flag (mechanism F) or by reproducing $U_*(t)$ from public
state and recomputing the deterministic argmin (mechanism A,
together with the override flag when the encoder departed from the
predicted argmin). It then isolates the repair-stream bits
belonging to tile $t$ using the framing information, and applies
the decoder of mode $\hat{m}(t)$ to recover the tile. Each step is uniquely determined; the concatenation is
therefore an injection from $X$ to codewords. \emph{Caveat:} without
the framing assumption, per-mode unique decodability alone is
insufficient: e.g., mode-A codewords $\{0, 01\}$ and mode-B
codewords $\{10, 0\}$ are each uniquely decodable, but the
concatenation $010$ over a mode-A then mode-B tile parses as
$0 \mid 10$ or $01 \mid 0$, ambiguously. The framing assumption
removes this ambiguity by fixing tile boundaries explicitly.
\hfill$\blacksquare$

Composition implies an approximate-oracle bound (not a strict
no-regret guarantee): by Theorem 6.1, including every coder in $M$
makes the Type-III-A composition's total cost at worst within
$2 \sum_t \varepsilon_t$ of the per-tile oracle (which picks the
true argmin mode per tile). This is also at worst within
$2 \sum_t \varepsilon_t$ of any fixed single coder $m^\star$ on the
same data, since the per-tile oracle dominates any fixed single
coder pointwise. When the per-tile best mode varies non-trivially across tiles
(i.e., no single coder is among the per-tile argmin set on every
tile), the per-tile oracle is strictly better than any fixed single
coder; under ties (some tile has multiple argmin modes including
$m^\star$) the oracle and $m^\star$ may agree on cost. The
composition inherits a strict improvement only insofar as
$2 \sum_t \varepsilon_t$ remains below that gap, which is itself an
empirical property of the prediction calibration and is not
guaranteed in general.

The per-type results of §§4--6 bound each coder in isolation. The next
two theorems lift the analysis to the \emph{composed} hierarchy and
delimit its fundamental rate: the composition is asymptotically optimal
exactly when its mode set reaches unbounded effective context, and the
price of any bounded-context restriction is a concrete information
quantity. Throughout we take $Y$ empty (pure source coding); the
conditional versions follow by conditioning every entropy on $Y$.

\textbf{Theorem 7.1b (Hierarchical universality of the RNR composition).}
\emph{Let $X$ be a stationary ergodic source over a finite alphabet with
entropy rate $h(X)$. Suppose the mode set $M$ contains the Type-III-C
universal dictionary coder (Theorem 6.6b), and the predictor is
\emph{calibrated}: $\tfrac1N \sum_t \varepsilon_t \to 0$ and the
override-flag overhead vanishes in rate, $\tfrac1N \, (\#\{t : \hat m(t)
\ne m^\star(t)\}) \log|M| \to 0$. Then the composition is asymptotically
rate-optimal,}
\[
\frac1N \, \mathbb{E}\bigl[L_{\mathrm{RNR}}(X^N)\bigr]
\;\longrightarrow\; h(X),
\]
\emph{and the convergence holds almost surely if the Type-III-C builder
is almost-surely universal (Kieffer--Yang 2000).}

\emph{\textbf{Proof.}} \emph{Achievability.} By the approximate-oracle
bound above, the realized cost obeys $\sum_t L_{\hat m(t)}(t) \le
\sum_t L_{m^\star(t)}(t) + 2\sum_t \varepsilon_t \le
L_{\mathrm{IIIC}}(X^N) + 2\sum_t \varepsilon_t$, the last step because
the per-tile oracle dominates the fixed single coder that selects
Type-III-C on every tile (whose summed per-tile lengths are exactly the
Type-III-C total). Adding the override-flag overhead and taking
expectations, $\mathbb{E}[L_{\mathrm{RNR}}] \le \mathbb{E}[L_{\mathrm{IIIC}}]
+ 2\,\mathbb{E}[\sum_t \varepsilon_t] + (\text{flag overhead})$. By
Theorem 6.6b, $\tfrac1N \mathbb{E}[L_{\mathrm{IIIC}}] \to h(X)$; by
calibration the two remaining terms are $o(N)$. \emph{Converse.} The RNR
codeword is uniquely decodable for $X$, so McMillan and the entropy
bound give $\mathbb{E}[L_{\mathrm{RNR}}] \ge H(X^N)$, and the
entropy-rate floor (Theorem 6.6c) gives $H(X^N) \ge N h(X)$ (the floor
$H(X^N)$ is a benchmark, not a computable encoder target --- its exact
value is $\mathrm{FP}^{\#P}$-complete, Remark 7.15g, and only
$\varepsilon$-approximable in poly time under filter stability, Remark
7.15c; the achievability above attains the realized scheme, not the exact
optimum). Hence
$h(X) \le \tfrac1N \mathbb{E}[L_{\mathrm{RNR}}] \le \tfrac1N
\mathbb{E}[L_{\mathrm{IIIC}}] + o(1) \to h(X)$. The almost-sure statement
replaces Theorem 6.6b's expectation convergence with the a.s.\
universality of the builder. \hfill$\blacksquare$

\textbf{Theorem 7.1c (Necessity of the unbounded-context mode: a
bounded-context separation).} \emph{Let $X$ be stationary over a finite
alphabet, and restrict the composition to a mode set $M_K$ in which every
symbol $X_t$ is assigned a uniquely-decodable codeword conditional on a
decoder-reproducible context that is a function of at most the $K$
immediately preceding symbols $X_{t-K},\ldots,X_{t-1}$ (bounded effective
context order $K$; for order-$k$ modes on tiles of size $s$ this gives
$K = k + s - 1$, since a symbol late in a tile may also condition on the
tile's earlier symbols). The conditional-codeword structure --- equivalently
an order-$\le\!K$ coding model, the same per-tile-UD hypothesis as in
Theorem 6.2b, and automatic for the bounded-context modes
Type-I/II/III-B --- is what makes the per-context lengths add. Then}
\[
\liminf_{N\to\infty} \frac1N \, \mathbb{E}\bigl[L_{\mathrm{RNR}\mid M_K}(X^N)\bigr]
\;\ge\; h_K := H(X_0 \mid X_{-1},\ldots,X_{-K}) \;=\; h(X) + e_K,
\]
\emph{where $e_K := h_K - h(X) \ge 0$ is the order-$K$ excess entropy. For
any source with $e_K > 0$ for every finite $K$ --- equivalently, a source
that is not $K$-th-order Markov for any $K$ (e.g.\ a hidden Markov source
with a full-support, non-degenerate kernel) --- the bounded-context
composition is bounded away from the entropy rate, no finite $K$ closes
the gap, and the unbounded-context Type-III-C mode (whose effective order
grows without bound in the block length) is necessary to attain $h(X)$
and recover Theorem 7.1b.}

\emph{\textbf{Proof.}} By the conditional-codeword hypothesis the total
length is the sum of per-symbol (equivalently per-tile, via the
chain rule) codeword lengths $L = \sum_t \ell_t(X_t \mid C_t)$, each
codeword uniquely decodable given its context $C_t$, a function of
$X_{t-K},\ldots,X_{t-1}$. By Kraft and the conditional-entropy bound,
$\mathbb{E}[\ell_t \mid C_t] \ge H(X_t \mid C_t) \ge H(X_t \mid X_{t-K},
\ldots,X_{t-1})$, the last step because conditioning on the finer
$X_{t-K..t-1}$ cannot increase entropy. Summing and using stationarity,
$\tfrac1N \mathbb{E}[L] = \tfrac1N \sum_t \mathbb{E}[\ell_t] \ge
\tfrac1N \sum_t H(X_t \mid X_{t-K..t-1}) \to
H(X_0 \mid X_{-K..-1}) = h_K$ (the boundary terms $t \le K$ number $O(K)$
and contribute $O(K \log|A| / N) \to 0$). The identity $h_K = h(X) + e_K$
with $e_K \ge 0$ and the monotone limit $h_K \downarrow h(X)$ is the
standard behaviour of conditional entropies for a stationary source
(Cover \& Thomas 2006, Theorem 4.2.1), with $e_K = 0$ iff $X$ is
$K$-th-order Markov. \hfill$\blacksquare$

\emph{Remark 7.1d (what each mode buys).} Theorems 7.1b--7.1c together
delimit the hierarchy: the composition is universal iff its mode set
reaches unbounded effective context, and the price of any bounded-context
restriction is exactly the source's order-$K$ excess entropy $e_K$. This
isolates the role of the Type-III-C dictionary --- it is not one
optimization among several but the single mode responsible for asymptotic
optimality on infinite-memory sources; the bounded-context modes
(Type-I/II/III-B) contribute constant-factor and finite-order gains that
saturate at $h_K$ and cannot, by themselves, reach $h(X)$ when
$e_K > 0$ for all $K$. The dichotomy mirrors the classical gap between
fixed-order Markov coders and universal dictionary methods (Lempel--Ziv),
here made internal to the RNR mode taxonomy.

\textbf{7.2 Runtime parameter selection}

Several parameters in the framework are input-dependent: the codeword
rate $k$ in Type-I (Code-$k$), the alphabet size $K$ in Type-II, the
support threshold $\tau_v$ for the Type-II support sets (optimal
exact selection is NP-hard for general joint laws by Lemma 5.7c and
polynomial-time for permutation-like laws by Lemma 5.7a), the tile size in
Type-III-A, the dictionary size in Type-III-C, and others. No single
setting is optimal across all inputs; the right approach is input
probing at encode time, where the encoder scans the block, estimates
the relevant structural parameter, and selects the parameter
accordingly. The parameter choice is recorded in the metadata stream
$L(\Theta)$ so the decoder can reproduce the selection.

Parameter selection is structurally identical to mode selection: each
parameter value defines a distinct coder, and the Type-III-A scheduler
chooses among them by predicted cost. The composition theorem
(Theorem 7.1) and the approximate-oracle bound (Theorem 6.1) both
extend to parameterized modes without change, with the per-mode
prediction error $\varepsilon_t$ now defined over the cross product of
(mode, parameter) pairs. The cost is the metadata overhead of encoding
the parameter choice, typically $O(\log K_{\text{params}})$ bits per
tile or per block, which is negligible against the per-mode repair
stream.

Two practical patterns are standard. First, batch probing: scan a
small prefix of the block (e.g., the first 1 KiB), estimate key
statistics (entropy rate, marginal sparsity, expected Hamming weight
of Code-$k$ errors), select parameters once, and encode the rest of
the block with the chosen parameters. Second, tiled probing: at each
tile boundary in Type-III-A, re-estimate parameters from the most
recently verified subblocks and adjust. Both patterns fit the
existing framework: the parameter-selection function is itself part
of $M$ (deterministic, decoder-reproducible from public state), and
the parameter values are part of $L(\Theta)$.

This framing resolves the input-dependent design questions that would
otherwise appear as open problems: there is no single optimal
Code-$k$ or optimal alphabet size, only an optimal
parameter-selection function for a given data distribution, and the
latter is part of $M$ just like any other learned component.

\textbf{8. No universal dominance}

\textbf{Theorem 8.1 (No lossless universal compression).} \emph{No
injective map $C$ from $\{0,1\}^n$ to binary strings of length less
than $n$ assigns every input string of length $n$ a strictly shorter
codeword.}

\emph{\textbf{Proof.}} There are $2^n$ binary strings of length $n$.
There are $2^n - 1$ binary strings of length strictly less than $n$,
summing $\sum_{k=0}^{n-1} 2^k$. Any injection from $2^n$ inputs to
fewer than $2^n$ outputs is impossible by the pigeonhole principle.
\hfill$\blacksquare$

Theorem 8.1 is standard and short; we include it to fix the
framework's honest claim. RNR does not assert that any fixed coder
dominates any baseline on every input. The claim is conditional and
distributional: for data classes whose conditional entropy under
deterministic side information is below the baseline rate by more than
the model and overhead cost, RNR is advantageous in expectation.

\textbf{9. Verification and error containment}

Autoregressive RNR modes update the predictor context with previously
decoded subblocks. An undetected decoding error in subblock $i$ would
propagate into the context for subblock $i + 1$ and beyond, with
cascading consequences. We discipline this by requiring that the
decoder update its context only after verifying the candidate decoded
subblock against a stored hash.

Let $X = (X_1, \ldots, X_m)$ be partitioned into verified subblocks,
with stored verification values $v_i = H_v(X_i)$. The decoder produces
$\hat{X}_i$, checks $H_v(\hat{X}_i) = v_i$, and only on success uses
$\hat{X}_i$ to update $C$ for subblock $i + 1$. The following theorem
rules out cross-subblock error propagation.

\textbf{Theorem 9.1 (No cross-subblock context propagation, with
explicit framing).} \emph{Assume the archive header stores, for each
subblock $i$, an explicit byte/bit-offset into the repair stream
(framing), so that subblock $i + 1$'s repair-stream slice is parsed
from a fixed position independent of the bytes recovered in
subblock $i$. Suppose the decoder uses the following
\emph{verified-only context-update rule}: after producing the
candidate $\hat{X}_i$, the decoder evaluates
$H_v(\hat{X}_i) \stackrel{?}{=} v_i$, and only on equality does it
extend the predictor context $C$ with $\hat{X}_i$ for predicting
subblock $i + 1$; on mismatch the decoder aborts (or marks the
remainder unverified and does not update $C$). The predictor input
for subblock $i+1$ is therefore a function only of (a)~the
verified-bytes context $C$ built from $\{\hat{X}_j : j \leq i,
H_v(\hat{X}_j) = v_j\}$ together with any public boundary state in
the header; the rule permits but does not require a hard
\emph{context-reset boundary} (i.e., dropping the verified-bytes
context at each subblock break). Finally, suppose $H_v$ is a
cryptographic hash modeled as either (a)~a keyed hash with a
uniformly random secret key sampled per archive, or
(b)~a random-oracle abstraction of a fixed unkeyed hash. Let
$\delta$ denote a uniform upper bound on the collision probability
$\Pr_{H_v}[H_v(y) = H_v(X_i)]$ for any $y \neq X_i$ in the
candidate set the decoder considers (probability over the choice of
key in case (a) or over the random-oracle responses in case (b);
for SHA-256 truncated to $t$ bits, $\delta \leq 2^{-t}$). Then the
probability that an incorrect $\hat{X}_i$ (one with $\hat{X}_i \neq
X_i$) is used to update the predictor context or repair-stream
parsing for subblock $i + 1$ is at most $\delta$.}

\emph{\textbf{Proof.}} The framing assumption guarantees that
subblock $i + 1$'s repair-stream bits are isolated by the header
offset and not by any function of $\hat{X}_i$. The context update
rule for the predictor is conditional on $H_v(\hat{X}_i) = v_i$. If
$\hat{X}_i \neq X_i$, then $H_v(\hat{X}_i) = v_i$ can occur only as
a hash collision with the unique stored $v_i$. Under the
cryptographic hash assumption, this happens with probability at
most $\delta$ per candidate; the bound holds by union over the
considered candidates if $\delta$ is interpreted accordingly.

\emph{Caveat.} Without the explicit-framing assumption, a deletion
or length-mismatch in subblock $i$'s repair-stream encoding (e.g.,
caused by a non-self-delimiting per-mode encoding without a header
offset) shifts the byte boundaries of every subsequent subblock and
can make all later subblocks unreadable; the per-subblock hash $v_i$
detects that a failure occurred but does not localize or recover
later subblocks. The framing assumption (explicit per-subblock
offset table in the archive header, enforced by the engineering
specification) eliminates this failure mode.
\hfill$\blacksquare$

The theorem is stated for cryptographic $H_v$ because
non-cryptographic hashes (CRCs, simple checksums) admit structured
collisions that an adversarial or simply unlucky predictor may
produce. In practice we recommend a truncated SHA-256 per verified
subblock. A 64-bit truncation gives $\delta \approx 2^{-64}$, which
is operationally negligible and costs 8 bytes per subblock; a full
SHA-256 costs 32 bytes.

\textbf{10. Practical considerations}

\textbf{10.1 Bit-exact decoder reproduction of $M$}

All RNR types require that the decoder reproduce the model output
$Y = M(C), Q,$ or $U$ bit-for-bit identically to the encoder. IEEE-754
floating-point arithmetic is non-associative, GPU reduction kernels
are non-deterministic by default, and cuDNN/cuBLAS algorithm selection
varies by hardware and library version. A model executed naively on a
heterogeneous fleet of decoders will not produce identical output, and
a single bit of divergence breaks lossless decoding.

Three remedies are available.

\begin{itemize}
\item
  Integer arithmetic. Quantize the model to 8‑bit or 16‑bit integers and
  execute inference in pure integer kernels. This is what NNCP uses in
  its production CPU pipeline. Integer arithmetic is associative and
  hardware‑independent up to the well‑defined behavior of the integer
  operations.
\item
  Fixed‑point arithmetic. Use a fixed‑point representation with explicit
  scale parameters and rounding modes. This is more flexible than pure
  integer arithmetic and is the standard approach in deterministic
  embedded inference.
\item
  Strict‑mode floating‑point. Use the strict IEEE‑754 mode of the
  inference runtime, disable fast‑math compiler optimizations, force a
  fixed reduction order, and pin the library version. This is fragile
  but sometimes acceptable for distribution alongside a specific runtime
  binary.
\end{itemize}

We recommend integer arithmetic with explicit per-layer scale and
zero-point parameters as the default for any RNR realization. The
model description $L(M)$ then includes both the parameter table and
the per-layer quantization scheme.

\textbf{10.2 Model amortization regimes}

$L(M)$ is paid once per archive. For a single-block archive, the model
cost dominates and RNR is competitive only if the model is small (a
few kilobytes), or if the model is pre-shared between encoder and
decoder via a public release. For a multi-block archive, the amortized
cost is $L(M)/m$ where $m$ is the number of blocks; for streaming
applications where $m$ is large, $L(M)$ effectively vanishes. The
three regimes have different optimal model sizes and warrant separate
evaluation. For a fixed shared frontier predictor, the regimes
collapse to a single break-even volume: the headline
$\approx 79\,\mathrm{MB}$ / $12.0\times$ figure of Theorem 7.15 of [RNR-II] is a
bitstream-only rate, and charging the $\approx 140\,\mathrm{GB}$
predictor leaves a net saving only once the combined amortized volume
exceeds $V^*\approx 152\,\mathrm{GB}$ (the ``Side-information framing
and model overhead'' caveat in Theorem 7.15 derives $V^*$; the table
below collates all four regimes).

\emph{Consolidated model-inclusive vs.\ amortized accounting (collation
of T7.15, T7.21, and the regimes above).} Whether that bitstream-only
rate is a net saving depends entirely on whether the predictor cost
$L(M)$ is charged to the archive. The table below puts the four
accounting regimes side by side, drawing every quantity from an
already-established result; the break-even row equates the
self-contained archive size with the raw source size, the condition
under which charging $L(M)$ still leaves a net saving.

\begingroup\footnotesize\sloppy
\setlength{\tabcolsep}{4pt}
\noindent\begin{tabular}{@{}p{2.9cm}p{2.7cm}p{1.4cm}p{2.3cm}p{3.1cm}@{}}
\hline
\textbf{regime} & \textbf{bitstream / repair cost} &
\textbf{$L(M)$ charged?} & \textbf{net savings vs.\ raw} &
\textbf{break-even / amortization condition}\\
\hline
Amortized / model-shared (side information; the headline)
& $N\!\cdot\! H_{M'}\!\approx\! 79\,\mathrm{MB}$ at $H_{M'}\!=\!0.664$\,bpb,
$N\!=\!1\,\mathrm{GB}$ (T7.15)
& no
& $12.0\times = 8/0.664$ (T7.15)
& decoder holds $M'$ as shared side information (Slepian--Wolf 1973); the
operational quantity is $N\!\cdot\! H_{M'}$ (T7.15 caveat)\\
Self-contained single archive (Kolmogorov framing)
& same $79\,\mathrm{MB}$ bitstream
& yes, full $140\,\mathrm{GB}$
& $\approx 140\times$ \emph{expansion} (negative) for $N\!=\!1\,\mathrm{GB}$ (T7.15 caveat)
& never net-positive at this $N$: $L(M)\!\approx\!140\,\mathrm{GB}$ dwarfs the
$1\,\mathrm{GB}$ source (T7.15 caveat)\\
Self-contained, model encoded \& distilled
& $N\!\cdot\! h(X)+\Theta(N^{1/(1+\alpha)})$, $L^*\!=\!(NC\alpha)^{1/(1+\alpha)}$ (T7.21)
& yes, distilled $L^*$
& overhead sub-linear: $(1{+}1/\alpha)L^*\!=\!\Theta(N^{1/(1+\alpha)})$ bits (T7.21)
& optimal predictor grows sub-linearly; per-archive overhead vanishes
relative to $N\!\cdot\! h(X)$ as $N\!\to\!\infty$ (T7.21)\\
Per-regime amortization (this §10.2)
& $N\!\cdot\! H_{M'}+L(M)/m$ over $m$ blocks
& yes, shared $L(M)/m$
& positive once $L(M)/m \ll N\!\cdot\!(\log_2|\Sigma|-H_{M'})$
& single-block: $L(M)$ dominates; multi-block: $L(M)/m$; streaming
($m$ large): $L(M)$ vanishes (this §10.2)\\
\hline
\textbf{Break-even volume $V^*$}
& \multicolumn{4}{@{}p{9.5cm}@{}}{Equating self-contained archive with raw
source, $V^*\!\cdot\!(\log_2|\Sigma|-H_{M'})=L(M)\!\cdot\!\log_2|\Sigma|$,
gives for a $140\,\mathrm{GB}$ predictor on byte-alphabet text
($|\Sigma|\!=\!256$, $H_{M'}\!=\!0.664$): $V^*\approx 140\,\mathrm{GB}\cdot 8/(8-0.664)\approx 152\,\mathrm{GB}$ of compressible text. Net savings turn
positive once the combined amortized volume $V>V^*$ (T7.15 caveat).}\\
\hline
\end{tabular}
\par\endgroup

\emph{Reading guide.} Rows 1--2 are the same encoder run scored under
two framings: row 1 (the headline) treats $M'$ as shared side
information and never charges $L(M)$; row 2 charges the full
$140\,\mathrm{GB}$ predictor and is therefore net-negative for a
single $1\,\mathrm{GB}$ archive. Rows 3--4 are the two ways to make a
self-contained archive net-positive: distil the predictor to its
$N$-optimal size $L^*(N)$ (T7.21), or amortize a fixed shared
predictor across enough archives that the combined volume exceeds the
break-even $V^*\approx 152\,\mathrm{GB}$. No quantity in this table is
new; each is cross-referenced to its source (Theorem 7.15, Theorem
7.21 --- both of the companion Part [RNR-II] --- or the per-regime
discussion above).

\textbf{10.2.1 The qualifying data-class precondition (accounting form)}

The Abstract and §1.1 phrase the central claim conditionally: RNR is
advantageous only on data classes for which the side information buys
more than it costs. Definition 2.1 fixes the operational meaning of
``advantageous'' --- $\mathbb{E}[L(\mathsf{Enc}(X))] <
\mathbb{E}[L(B(X))]$ --- and Theorem 8.1 (no code shortens every input)
implies the advantage cannot be universal, hence is necessarily
conditional and distributional. This subsection assembles that
predicate into a single auditable inequality whose every term is
already defined or bounded elsewhere in the paper; it states no new
result and measures no corpus (which real corpora satisfy it is the
open empirical question of §11). Throughout we work per source symbol
and in the amortized regime of §10.2.

\emph{Precondition (data-class advantage, accounting form).} For a
source $X\sim D$ over alphabet $\Sigma$ with decoder-reproducible side
information $Y=M(C)$, baseline $B$, and an $m$-block amortized archive,
the per-symbol expected RNR rate is the sum of its component costs, and
RNR is conditionally advantageous on $D$ (Definition 2.1) when
\begin{equation}
\underbrace{H(X\mid Y)}_{\text{(a) residual uncertainty, §2.2}}
\;+\;
\underbrace{\tfrac{L(M)}{mN}}_{\text{(b) amortized model, §10.2}}
\;+\;
\underbrace{\rho(D)}_{\text{(c) repair residual, §2.3/§5}}
\;+\;
\underbrace{\nu}_{\text{(d) verification overhead, §10.3}}
\;+\;
\underbrace{\mu}_{\text{(e) metadata, §10.4/§10.7}}
\;<\;
\underbrace{R_B(D)}_{\text{baseline rate, Def.\ 2.1}} .
\tag{10.1}
\end{equation}
This is exactly Definition 2.1's inequality
$\mathbb{E}[L(\mathsf{Enc}(X))] < \mathbb{E}[L(B(X))]$ divided by the
block length and expanded into the encoder's component costs, with
$R_B(D)=\mathbb{E}[L(B(X))]/N$ the per-symbol baseline rate. The terms
are, with their paper-side definition or bound:
\begin{itemize}
\item[(a)] $H(X\mid Y)$, the conditional entropy of the source given
  the side information --- the information-theoretic floor of the
  repair stream under conditional source coding (§2.2; Cover and
  Thomas 2006, Theorem 5.4.1). This is the term the side information
  is meant to drive down; the whole conditional advantage hinges on
  $H(X\mid Y)$ being below the baseline rate by more than (b)--(e).
\item[(b)] $L(M)/(mN)$, the total model description amortized over the
  $mN$ symbols of the $m$-block archive --- the per-block share $L(M)/m$
  of §10.2 expressed per symbol (§10.2). In the single-block regime $L(M)$
  dominates and (10.1) typically fails; as $m$ grows (streaming) this
  term vanishes. This is precisely the term that forces the
  ``only the model-amortization regime can win'' reading of §10.2: for
  a fixed shared frontier predictor the regimes collapse to the single
  break-even volume $V^*$ of the §10.2 table (net saving requires the
  combined amortized volume $V>V^*\approx152\,\mathrm{GB}$), so (10.1)
  with (b) charged in full is the per-symbol form of that break-even
  condition.
\item[(c)] $\rho(D)$, the expected repair residual --- the excess of
  the realized repair-stream rate over the entropy floor (a), i.e.\ the
  entropy-coder redundancy $\varepsilon_n/n$ of §2.3 (assumed
  $\to 0$) plus any modeling gap between the coding law and the true
  conditional law. The paper bounds the coder redundancy (§2.3) but
  does not give a closed distribution-free bound on the modeling gap;
  $\rho(D)$ is therefore carried as a named residual term (its
  estimation on real data is part of the §11 / BUG-010-D campaign), not
  a quantity claimed to be zero.
\item[(d)] $\nu$, the verification overhead --- the per-subblock hash
  bytes plus the explicit repair-stream offset table (and any
  Merkle-tree authenticators for random-access verification), bounded
  term by term in §10.3 (e.g.\ $\approx0.2\%$ for 64-bit hashes at
  4\,KiB subblocks plus $\approx0.1\%$ offsets).
\item[(e)] $\mu$, the metadata --- mode-flag stream and fallback
  headers (§10.4) and, when seekable random access is used, the seek
  index of §10.7. By Theorem 8.1 a fallback path is mandatory, so
  $\mu>0$ always; §10.4/§10.7 bound it.
\end{itemize}
\emph{Consistency with the existing accounting.} Inequality (10.1) does
not introduce a new cost: it is the per-symbol decomposition of the
self-contained-archive accounting tabulated in §10.2 (the
model-inclusive vs.\ amortized table) and is the per-symbol companion
of the absolute-ratio scoping of §11. Summing (a)+(c) gives the
bitstream/repair rate $\approx H_{M'}$ of the §10.2 table's headline
row; charging (b) in full is its self-contained row; letting $m$ grow
is its per-regime-amortization row. No term is double-counted: (a) is
the floor, (c) is the realized excess \emph{over} that floor, and
(b),(d),(e) are the disjoint side costs (model, verification,
metadata) --- with one shared structure charged once: the per-subblock
repair-stream offset table of §10.3 (counted in (d)) and the seekable
random-access index of §10.7 are the same $\lceil\log_2 L(R)\rceil$-bit
offsets-into-$R$ when subblock boundaries fall on sync points, so (e)'s
seek index adds only the source-position field for non-uniform spacing,
not a second copy of the offsets. The baseline rate $R_B(D)$ is the same $\mathbb{E}[L(B(X))]$
of Definition 2.1 and the same baseline family (sequential dictionary /
entropy-only coders, and the seekable formats bgzip / zstd-seekable)
scoped in §11.

\emph{Candidate qualifying families (conjecture, not asserted).} We do
not claim any specific corpus satisfies (10.1); that is the open
empirical question (§11; the present subsection's role is only to make
the precondition auditable). As a conjecture, the structural families
of §1.1 --- data whose redundancy is positional, geometric,
multi-scale, or coupled across non-adjacent positions (database pages,
sensor and telemetry grids, structured binary sections, scientific
tensors), i.e.\ high-mutual-information-with-a-learnable-predictor
corpora for which a structural predictor reaches conditional entropy
$H(X\mid Y)$ well below the sequential baseline rate --- are the
plausible candidates for which (10.1) can hold in the amortized regime.
Whether the inequality holds for any of them with realistic $L(M)$,
$\rho(D)$, $\nu$, $\mu$ is exactly what §11's pre-registered
measurements (and the residual external work) would decide; the claim
of this paper remains conditional and in-principle.

\textbf{10.3 Verification overhead analysis}

A SHA-256 per 4 KiB subblock is $32/4096 = 0.78\%$ hash-only overhead.
A truncated 64-bit hash per 4 KiB subblock is $8/4096 = 0.20\%$. A
truncated 32-bit hash is $0.10\%$, with $2^{-32}$ \emph{per-subblock}
collision probability; for an archive of $m$ subblocks, the
archive-level false-accept probability is bounded by union over
subblocks at $m \cdot 2^{-32}$, which is non-negligible for large
$m$ (e.g., $m = 2^{18}$ subblocks at 4 KiB each --- a 1 GiB archive
--- gives archive-level false-accept probability up to $2^{-14}
\approx 6 \cdot 10^{-5}$, comparable to typical disk uncorrectable
read error rates). For archives in the GiB+ range, 64-bit
truncation gives archive-level false-accept $\leq m \cdot 2^{-64}
\approx 10^{-14}$, which is negligible. Reducing the subblock size
to 1 KiB increases overhead proportionally and quadruples $m$. The
right balance depends on the underlying error rate of the
predictor and the size of the archive; if errors are rare and the
archive is small, larger subblocks and shorter hashes are
acceptable.

The hash bytes are not the only verification-related overhead.
Theorem 9.1 also requires an explicit per-subblock repair-stream
offset table to disambiguate subblock boundaries; an entry storing
a $\lceil \log_2 L(R) \rceil$-bit offset per subblock adds another
$m \cdot \lceil \log_2 L(R) \rceil$ bits to the archive header. For
$|X| = 1$ GiB at 4 KiB subblocks ($m = 2^{18}$) and a repair
stream of comparable size ($\log_2 L(R) \approx 33$), the offset
table is $m \cdot 33 \approx 8.4$ Mibit $\approx 1$ MiB, a further
$\approx 0.1\%$ of $|X|$. Total verified-framed verification
overhead is therefore the sum of (a) hash bytes, (b) per-subblock
repair-stream offsets, and (c) any Merkle-tree authenticators if
random-access verification is used (§10.7, third trade-off).

\textbf{10.4 Fallback modes}

Every RNR coder must include a fallback mode that handles the case in
which the chosen primary mode fails to compress. Two fallbacks are
standard: (i)~raw storage of the subblock, costing $8N$ bits plus
headers; (ii)~baseline coder storage, applying a fixed external coder
(e.g., zstd) to the subblock. The choice between fallbacks is
per-subblock and recorded in the mode-flag stream of the Type-III-A
scheduler. The cost of the fallback flag must be included in
$L(\Theta)$.

\textbf{10.5 A sufficient condition for bit-exact reproduction}

Section 10.1 recommends integer arithmetic as the default determinism
mechanism but does not formalize when it suffices. The following theorem
gives a sufficient condition under which a quantized neural predictor is
bit-exact across any pair of conforming implementations.

\textbf{Theorem 10.1 (Bit-exact reproduction under deterministic
integer inference).} \emph{Let $M$ be an inference graph with $L$
layers, where layer $\ell$ is implemented as an integer matrix-vector
product of fan-in $d_\ell$ followed by a non-linearity. Suppose:
(i)~all weights and inputs are integers; (ii)~every accumulator in
layer $\ell$ uses signed bit-width $b_\ell$ satisfying
\[
b_\ell \geq \lceil \log_2(d_\ell \cdot W_\ell \cdot I_{\ell-1} + 1) \rceil + 1,
\]
where $W_\ell$ is an upper bound on the magnitude of any weight entry
at layer $\ell$ and $I_{\ell-1}$ is an upper bound on the magnitude of
any input to layer $\ell$ (the $+1$ inside the $\log$ accounts for
the asymmetric signed-two's-complement range $[-2^{b-1}, 2^{b-1}-1]$,
so the positive extremum $+d_\ell W_\ell I_{\ell-1}$ is representable); (iii)~all non-linearities are deterministic
integer-to-integer mappings (e.g., piecewise-constant lookup tables);
(iv)~all reduction operations (sums in matrix-vector products) follow
a fixed, deterministic order specified by the model description. Then
no accumulator overflows and $M(C)$ produces bit-identical output for
any input $C$ on any two implementations conforming to (i) through
(iv).}

\emph{\textbf{Proof.}} Under (i) and (ii), each integer multiplication
and addition produces a result within the representable range; no
overflow or saturation occurs. Integer arithmetic is associative and
commutative, so any ordering of additions produces the same exact sum.
Under (iv), the ordering is pinned by the model description, so two
implementations choose the same ordering and produce the same sum.
Under (iii), the non-linearity is a deterministic function from
integer inputs to integer outputs that is fully specified by the model
description (typically as a lookup table); two implementations
applying the same table to the same input produce the same output. By
induction on layer depth, the output of layer $l$ is identical across
implementations, hence the final $M(C)$ is identical.
\hfill$\blacksquare$

The bit-width condition is per layer. If each layer renormalizes or
clips outputs so that $|x| \leq I_{\max}$ holds uniformly, then a
uniform sufficient condition is
$b \geq \lceil \log_2(d_{\max} \cdot W_{\max} \cdot I_{\max} + 1) \rceil + 1$,
where $d_{\max} = \max_\ell d_\ell$ is the maximum fan-in (the $+1$
inside the $\log$ matches the asymmetric signed-integer range, as in
the per-layer bound). Without inter-layer renormalization,
intermediate magnitudes can grow recursively:
$I_\ell \leq d_\ell \cdot W_\ell \cdot I_{\ell-1}$ in the worst
case, so $I_\ell \leq I_0 \prod_{r \leq \ell} d_r W_r$ and depth
enters multiplicatively, not linearly. Numerical sanity check: for
int8 quantization with $I_{\max} = W_{\max} = 2^7$ and fan-in
$d_{\max} = 2^{12}$ (4096-channel matrix-vector product), the
worst-case sum is bounded by $2^{12} \cdot 2^7 \cdot 2^7 = 2^{26}$,
so $b = 32$ suffices with comfortable headroom. For int16
quantization with $I_{\max} = W_{\max} = 2^{15}$ and the same
fan-in, the worst-case sum is $2^{12} \cdot 2^{15} \cdot 2^{15} =
2^{42}$, so $b \geq 44$ is needed (in practice $b = 64$ is the next
natural width; $b = 32$ overflows by 10+ bits at this fan-in;
only the degenerate scalar case $d_{\max} = 1$ gives a worst-case
product of at most $2^{30}$ and fits $b = 32$; any matrix-vector
layer with fan-in $d_{\max} \geq 2$ under int16 weight and
activation bounds exceeds $2^{31}$ in the worst case and requires
$b \geq 33$). int16-only deployments must
either use 64-bit accumulators or insert per-layer renormalization
to bound $I_\ell$ at runtime. The non-linearity table requirement
excludes any function involving real-valued division, square root, or
transcendentals; sigmoid, tanh, GELU, and softmax must be replaced by
piecewise-linear or integer-quantized approximations, which is
standard practice in quantized inference. The reduction ordering
requirement excludes any platform-specific parallelization (e.g., GPU
atomics) that does not commit to a deterministic order; this is
satisfied by sequential CPU implementations and by GPU implementations
that use tree-reduction with a fixed schedule.

Theorem 10.1 is a \emph{sufficient-condition} result: it reduces the
reproducibility question from a software engineering problem to a
precision-and-determinism specification, conditions (i)--(iv). Any two
implementations that both meet (i)--(iv) compute bit-identical output;
cross-platform conformance thereby becomes a testable property rather
than an empirical hope. The theorem does \emph{not} itself exhibit a
reference implementation or demonstrate end-to-end cross-platform
byte-identity on real hardware: constructing such an implementation and
\emph{certifying} that a given target meets (i)--(iv) is engineering
work, deferred to the companion engineering specification and recorded
as an out-of-scope item in §13.3. The integer reduction-order
determinism at the mathematical core of (ii)--(iv) --- the per-layer
no-overflow accumulator bound and the resulting reorder-invariance of
the integer forward pass, contrasted against the reorder-dependence of
an undersized accumulator --- is checked by the executable probe
\texttt{scripts/verify/thm\_10\_1\_integer\_reorder\_bitexact.py}.

\textbf{10.6 Bit-exact adaptive predictors under integer-arithmetic
optimization}

Theorem 10.1 covers the forward pass of a fixed predictor $M$. For RNR
coders that adapt $M$ online from previously verified subblocks during
encoding --- a useful capability for streams whose statistics drift
over a long archive --- the adaptation must itself be deterministic
and decoder-reproducible. Otherwise the encoder and decoder would
diverge after the first adaptation step. The relevant discipline is
integer-arithmetic optimization, developed in the recent
low-precision-training literature.

\textbf{Definition 10.1 (Compatible online optimizer).} An online
optimizer $\mathrm{Opt}$ is compatible with bit-exact reproduction if:
(a)~its state at step $t$, denoted $S_t$ (parameters, momentum,
accumulators, and any auxiliary buffers), is a tuple of
fixed-bit-width integers; (b)~its update function
$S_{t+1} = \mathrm{Opt}(S_t, g_t)$ takes the current state and a
fixed-bit-width integer gradient $g_t$ and produces the new state via
a deterministic computation using only integer arithmetic, lookup
tables, and seeded pseudorandom number generation; (c)~any operation
that would otherwise require real-valued computation (division, square
root, exponential, transcendentals) is replaced with a specified
deterministic integer approximation (integer division with explicit
rounding mode, integer square root, fixed-point or lookup-table-based
approximation of transcendentals).

Several practical optimizers fit Definition 10.1. Integer-quantized
SGD with quantized gradients (Wu, Yan, Wang, Maxson, Xu, He 2018) is
the canonical example: parameters, gradients, and
learning-rate-scaled updates are all integers, and the update is
$\text{parameter} \mathrel{-}= \mathrm{shift\_right}(\text{gradient} \cdot \text{learning\_rate}, \text{scale})$.
SignSGD (Bernstein, Wang, Azizzadenesheli, Anandkumar 2018) is
compatible by construction since the only non-trivial operation is
sign extraction. Integer-arithmetic Adam (Banner, Hubara, Hoffer,
Soudry 2018) is compatible when the reciprocal-square-root operation
is replaced with a deterministic integer approximation (typically a
piecewise-linear lookup table). Non-examples include any
floating-point optimizer without quantization wrappers, Adam with
real-valued square roots, and stochastic optimizers with unseeded
random masking.

\textbf{Theorem 10.2 (Bit-exact adaptive predictors, causal updates).}
\emph{Suppose the forward pass of $M$ satisfies the conditions of
Theorem 10.1, and the \emph{backward pass} (integer backpropagation)
satisfies the same four conditions independently: (i)~all backward
operands and gradient buffers are integers, (ii)~every accumulator
in any backward matrix-vector or reduction has signed bit-width
$b^{\mathrm{bwd}}_\ell$ satisfying the per-layer bound of
Theorem 10.1 applied to the backward operation's fan-in, weight
magnitude, and incoming-gradient magnitude bounds; (iii)~every
derivative-of-non-linearity is a deterministic integer-to-integer
lookup table; (iv)~every backward reduction follows a fixed,
deterministic order specified by the model description. Suppose the
encoder updates $M$ online via a compatible online optimizer
(Definition 10.1) using the gradient computed by this backward
pass. Suppose the order in which subblocks are used for adaptation is
pinned by the model description and identical on encoder and
decoder. Suppose further the \emph{causal-update condition}: the
encoder's adaptation step at subblock boundary $t$ uses only data
already available to the decoder at that boundary (i.e., bytes from
subblocks $0, \ldots, t$ all of which the decoder has already
verified per Theorem 9.1 before predicting subblock $t+1$). Then at
every subblock boundary, the encoder's and decoder's parameter
states are bit-identical, and the resulting predictor outputs are
bit-identical. \emph{Caveat:} without the causal-update condition
(e.g., if the encoder updates using the current byte before encoding
it), the decoder cannot replicate the same update sequence, and
parameter states diverge despite deterministic backprop and pinned
order; without the backward-pass conformance conditions, two
implementations sharing a compatible optimizer can still compute
different gradients (e.g., one wrapping, another saturating an
accumulator overflow) and so diverge.}

\emph{\textbf{Proof.}} Induction on the adaptation step $t$. Base
case ($t = 0$): the initial parameter state is part of $L(M)$ and is
bit-identical by construction. Induction step ($t \to t+1$): assume
$S_t$ is bit-identical between encoder and decoder. The subblock
$X_t$ used for adaptation is verified by Theorem 9.1 to be the same
on both sides. The gradient $g_t = \nabla \mathcal{L}(X_t; S_t)$ is
computed by integer backpropagation satisfying the four backward
conformance conditions (i)--(iv) above; by the same argument as in
Theorem 10.1 (no overflow, deterministic reduction order,
deterministic non-linearity-derivative tables, integer inputs),
each backward layer's output is bit-identical across implementations,
and by induction on backward depth so is the full gradient
$g_t$. The update $S_{t+1} = \mathrm{Opt}(S_t, g_t)$ applies
the deterministic integer-to-integer function from Definition 10.1
to bit-identical inputs, producing bit-identical outputs. By
induction, $S_t$ is bit-identical for all $t$. By Theorem 10.1
applied to each $S_t$, the predictor output is bit-identical at every
adaptation step. \hfill$\blacksquare$

Theorem 10.2 reduces the adaptive-predictor reproducibility question to
the existence of a compatible online optimizer with adequate training
performance. The low-precision-training literature provides several such
optimizers with published convergence and accuracy results approaching
their floating-point counterparts. The remaining engineering work is to
wire one of them into a concrete RNR codec; this is deferred to a
companion artifact (§13.3) on the same grounds as Theorem
10.1\textquotesingle s reference implementation.

\textbf{11. Experimental predictions}

We predict the following empirical behavior for an RNR implementation
that follows the constructions of Sections 4 through 6 with a small (a
few megabytes) byte‑level transformer as the underlying predictor.

These predictions are conjectures, not derived results. They are based
on (a) the structural advantages each Type confers in expectation under
its target data class, and (b) coding rates reported for comparable
existing systems on similar corpora. Specific bit‑per‑byte figures
should be read as order‑of‑magnitude estimates rather than precise
targets; the predictions have not been empirically validated and may be
off by 10-30\% in either direction. Empirical validation against NNCP,
PixelCNN++, LZMA at maximum preset, and zstd at maximum level on the
corpora named below is the natural next step.

Each prediction in \S\S11.1--11.6 is a falsifiable empirical conjecture,
conditional and pending measurement --- not a theorem --- and is
pre-registered as a numbered hypothesis (with an explicit falsification
rule) in the experimental-design companion document, to which each
prediction below is cross-referenced. In particular, the central
value-proposition claim that the side-information+repair recasting beats
existing dictionary-based seekable formats on compression ratio remains
an empirical conjecture pending the measurement specified there; the
seekable-format random-access baselines (bgzip, zstd seekable) and the
deployed-format ratio baseline (squashfs+zstd) are named in Hypotheses
H10--H11 and H12 of the companion, respectively. (\S11.7 is the lone
exception: it derives from Theorem 7.15 of [RNR-II] and is verified against a cited
measurement, not conjectured.)

\textbf{11.1 Type-I, Code12 on text and source code}

Type-I with $E_{12}$-A on English text and source code is predicted
to achieve 1.5 to 2.0 bits per byte before model amortization,
competitive with LZMA at maximum preset (2.1 to 2.5 bits per byte)
and behind NNCP (1.3 to 1.6 bits per byte on similar corpora). The
gap to NNCP is expected to narrow as the Code12 mapping is learned
(variant $E_{12}$-C) rather than fixed. This is a conjecture pending
measurement, pre-registered as Hypothesis H1 (quantifying) of the
experimental-design companion.

\textbf{11.2 Type-II on structured positional data}

Type-II on database pages, structured binary records, and instruction
streams is predicted to achieve coding rates 10 to 20\% below LZMA,
because LZMA exploits string repetition that is weak in these classes
while Type-II exploits positional regularity. Against NNCP, the
comparison depends on how well the byte-level autoregressive model
captures positional structure; for data whose joint structure exceeds
what a sequential model can express in moderate context, Type-II is
predicted to dominate. This is a conjecture pending measurement,
pre-registered as Hypothesis H2 (quantifying) of the experimental-design
companion, with the supporting structural conditions ($\mathrm{TC}(D_N)
= \Theta(N)$, intrinsic dimension $d \ll N$) pre-registered as
Hypotheses H3 and H4.

\textbf{11.3 Type-III-A composition on heterogeneous archives}

Type-III-A scheduling across Type-I, Type-II, and raw fallback is
predicted to outperform any single fixed mode on archives with
heterogeneous content (mixed text/binary, containers, database
dumps with header/payload structure), provided the per-mode
prediction error $\sum_t \varepsilon_t$ and the per-tile mode-flag
overhead together remain smaller than the heterogeneity advantage
of the per-tile oracle over the best fixed single coder.
Theorem 6.1 only gives an approximate-oracle bound, not a strict
dominance guarantee; on archives where the heterogeneity advantage
is small relative to $\sum_t \varepsilon_t$ plus mode-flag
overhead, the composition can be worse than the best fixed single
coder. Empirical validation is required to determine which side
wins on specific data classes; this conjecture is pre-registered as
Hypothesis H5 (quantifying) of the experimental-design companion.

\textbf{11.4 Type-III-B on ASCII text and instruction streams}

Type-III-B is predicted to recover 0.3 to 0.7 bits per byte over
independent nibble coding on ASCII text, in line with the
$I(H ; L \mid C)$ estimates of Section 6.4. On structured numeric
arrays (e.g., float32 columns), Type-III-B may recover larger gains
if the byte-level dependency is strong across the mantissa structure.
This is a conjecture pending measurement, pre-registered as Hypothesis
H6 (quantifying) of the experimental-design companion.

\textbf{11.5 Type-III-C on highly repetitive structured data}

Type-III-C with a $32 \times 32$ token grid on JSON, XML, source
code, and protocol streams is predicted to achieve coding rates
competitive with the best dictionary-based methods (zstd at maximum
level, LZMA at maximum preset), with the advantage growing as the
corpus size grows and the dictionary amortizes. This is a conjecture
pending measurement. The Type-III-C coder is pre-registered as
Hypothesis H7 (confirming, bits-back rate within $2\%$ of $H(X\mid Y)$)
of the experimental-design companion, and the companion's corpus set
(\S3.1) exercises Type-III-C on source code and JSON/HTML dictionaries;
the specific ratio-versus-dictionary-method target stated here is
covered by those corpora but is not itself given a dedicated quantifying
hypothesis (a residual to close in the measurement of Hypothesis H1's
text/code corpora and any follow-on Type-III-C quantifying test).

\textbf{11.6 Negative controls}

On random or encrypted data, all RNR types are predicted to revert to
baseline plus overhead, since no side information reduces residual
uncertainty. On already-compressed data (gzip, xz, zstd output), RNR
is predicted to be slightly worse than passthrough, since the
residual entropy is high and the model cost is wasted. This conjecture
is pre-registered as Hypothesis H16 (confirming, store-mode fallback
upper bound on adversarial inputs) of the experimental-design companion.

\textbf{11.7 Large-scale neural predictors (Theorem 7.15 of [RNR-II] instantiation)}

The end-to-end synthesis Theorem 7.15 of [RNR-II] gives an \emph{explicit
numerical prediction}, not a conjecture: with a Chinchilla-70B-class
transformer ($H_{M'} \approx 0.664$ bpb measured by Delétang et al.\
2024 on enwik9, Table 1) and the balanced regime $K = \Theta(\sqrt{N})$,
encoder output is $\approx 79\,\mathrm{MB}$ for the $1\,\mathrm{GB}$
enwik9 source, with practical CLT-based concentration $\pm 0.04\,\mathrm{MB}$
at $\delta = 10^{-6}$ (rigorous Azuma worst-case
$\pm 1.4\,\mathrm{GB}$ for the $W_N = 2048$ Chinchilla transformer
is hopelessly loose; see T7.15 worked-example dual-band discussion).
This is a $12.0\times$ \emph{bitstream-rate}
compression (matching Delétang et al.\ Table 1), \emph{exclusive} of
the predictor itself; see the ``Side-information framing and model
overhead'' caveat in §7 for the amortization framing. Under
amortization across $V \gtrsim 152\,\mathrm{GB}$ of compressible text
(per the break-even computation in T7.15's caveat), the per-archive
self-contained compression remains net-positive. This prediction
distinguishes §11.7 from §§11.1--11.5 in two ways: (i) it derives
from a theorem rather than a conjecture, and (ii) it has been
verified against the cited measurement
(\texttt{scripts/verify/thm\_7\_15\_neural\_rnr\_end\_to\_end.py}).

The predictions of §§11.1--11.5 above (1.5--2.0 bpb for small
byte-level transformers) refer to the small-model regime; T7.15's
0.664 bpb refers to the large-model (frontier-class LM) regime. The
two are quantitatively consistent: rate decreases monotonically with
model size, following the Hoffmann et al.\ 2022 (Chinchilla-paper)
scaling law $L(N_{\mathrm{params}}) \sim N_{\mathrm{params}}^{-\alpha}$
with $\alpha \approx 0.34$ (cited and used in Theorem 7.21).

\textbf{12. Related work}

This work intersects several established research lines.

\textbf{12.1 Source coding with side information}

Slepian and Wolf (1973) established the rate region for separately
encoding correlated sources with joint decoding. Wyner and Ziv (1976)
extended the result to lossy coding with side information at the
decoder. RNR uses the achievability direction of these results with
deterministic side information, an easier setting than the original
distributed coding problem.

\textbf{12.2 Channel coding and syndrome‑based compression}

Pradhan and Ramchandran (2003) introduced DISCUS for compression with
side information using channel codes. Liveris, Xiong, and Georghiades
(2002) used LDPC codes for similar purposes. Type‑I with a structured
E\_12 mapping is in the same spirit: the predictor produces a corrupted
version of a channel code, and the repair stream is essentially syndrome
information.

\textbf{12.3 Enumerative and combinatorial coding}

Cover (1973) gave the enumerative source coding scheme central to
Type‑II. Subsequent works extended enumerative coding to constrained
classes. Theorem 5.3 specializes these to the
partition‑with‑support‑and‑escape setting.

\textbf{12.4 Block transforms}

The Burrows--Wheeler Transform (Burrows and Wheeler 1994) is the
classical example of a reversible reorganization of data that exposes
structure. Type‑II and Type‑III‑B are also reversible reorganizations:
positional partition and byte‑tensor representations.

\textbf{12.5 Neural lossless compression}

Bellard's NNCP (2019) is the closest prior work to Type‑I: a byte‑level
transformer drives an arithmetic coder, and the NNCP rate is
rate-equivalent to the ideal Type-I class+payload bound under a
matching predictive law (§4.5). PixelRNN (van den Oord et al.,
ICML 2016) and the PixelCNN family --- Conditional PixelCNN (van den
Oord et al., NeurIPS 2016), PixelCNN++ (Salimans et al. 2017), and
PixelSNAIL (Chen et al. 2018) --- are factorized autoregressive
image models; they are \emph{analogous} to Type-II under a
factorized regime in that both factor a joint distribution into
per-position predictions, but Type-II as defined here is
partition-based and does not literally subsume their pixel-wise
softmax (§1.2). PixelCNN++ uses discretized logistic mixtures
rather than raw categoricals. Integer Discrete
Flows (Hoogeboom et al. 2019) and IDF++ (van den Berg et al. 2021)
construct invertible discrete representations whose factorized form is
closely related to Type‑III‑B. L3C (Mentzer et al. 2019) and LC-FDNet
(Rhee et al. 2022) introduce hierarchical priors that prefigure
Type‑III‑A composition.

\textbf{12.6 Bits‑back coding}

Bits‑back coding (Hinton and van Camp 1993; Frey 1997) was made
practical by Townsend, Bird, and Barber (2019) as BB‑ANS, and extended
by Kingma et al. (2019) as Bit‑Swap and by Townsend et al. (2020) as
HiLLoC. These methods losslessly compress data under variational
latent‑variable models and are the appropriate machinery for Type‑III‑C
with continuous latents.

\textbf{12.7 Context mixing}

PAQ (Mahoney) and Cmix (Knoll) represent the upper end of non‑neural
lossless compression via mixture‑of‑experts context mixing. They can be
viewed as hand‑designed Type‑III‑A schedulers in which the entropy field
U is computed from explicit context features rather than predicted by a
neural network.

\textbf{12.8 Foundations used by the §7 extensions}

The 27-theorem §7 framework of the companion Part [RNR-II] draws on several established lines outside
the immediate compression literature, listed here for completeness.

\emph{Concentration of measure and martingales.} Theorem 7.14's
Doob-martingale tail bound uses Azuma (1967) \emph{T\^ohoku Math.\
J.} 19(3):357 and Hoeffding (1963) \emph{J.\ Amer.\ Stat.\ Assoc.}
58(301):13 as base, with McDiarmid (1989) \emph{Surveys in
Combinatorics} (Cambridge) for the bounded-differences variant.
Under exp-mixing, the source-mixing-tail factor $\rho/(1-\rho)$ is
the obstacle that prevents direct iid-Hoeffding application; tighter
bounds under polynomially-slow mixing would use Rio (2000)
\emph{Th\'eorie asymptotique des processus al\'eatoires faiblement
d\'ependants} or Samson (2000) \emph{Ann.\ Probab.} 28(1):416.

\emph{Cache analysis.} Theorem 7.20's cache-aware decoder uses the
classical LRU competitive ratio $C/(C-h+1)$ of Sleator and Tarjan
(1985) \emph{Comm.\ ACM} 28(2):202--208 ``Amortized efficiency of
list update and paging rules''. The static-top-popularity hit rate
under Zipf workloads is the Breslau et al.\ (1999) \emph{Proc.\ IEEE
INFOCOM} ``Web caching and Zipf-like distributions'' result with
the correct asymptotics $\rho_C \sim (C/m)^{1-\alpha}$
(divergent-normalizer $0<\alpha<1$) and
$\rho_C \to 1 - C^{1-\alpha}/((\alpha-1)\zeta(\alpha))$
(convergent-normalizer $\alpha>1$). Theorem 7.24's IRM hypothesis
follows the Independent Reference Model of Fagin (1977)
\emph{J.\ Comput.\ Sys.\ Sci.} 14(2):222--250 ``Asymptotic miss
ratios over independent references''; foundational LRU+IRM analysis
is in Hendricks (1972) \emph{J.\ Appl.\ Prob.} 9(1):231--233 and
demand-paging algorithm analysis is in King (1971)
\emph{Proc.\ IFIP Congress 1971}, pp.\ 485--490. Ferragut, Rodriguez, and
Paganini (2018) \emph{Queueing Systems} 88(3):207--241 (online 2017)
give modern timer-based-cache hit-rate analysis under general arrival
processes.

\emph{Scaling laws for neural predictors.} Theorem 7.15's numerical
prediction uses the empirical bits-per-byte rate of Delétang et al.\
(2024) ``Language Modeling Is Compression'', arXiv:2309.10668,
ICLR 2024 (Table 1: Chinchilla-70B on enwik9 at $\approx 0.664$ bpb,
$12.0\times$ raw compression). Theorem 7.21's predictor-archive size
trade-off uses the Hoffmann et al.\ (2022)
\emph{Training Compute-Optimal Large Language Models}
(``Chinchilla paper'', NeurIPS 2022) scaling law with $\alpha
\approx 0.34$ for loss as a function of parameter count; this
\emph{updates} the earlier scaling fit of Kaplan et al.\ (2020)
\emph{Scaling Laws for Neural Language Models}, arXiv:2001.08361,
whose model-size exponent $\alpha_N$ for $L(N_{\mathrm{params}})$
was estimated at $\approx 0.076$ (the data-size exponent $\alpha_D$
for $L(D_{\mathrm{tokens}})$ at fixed model is $\approx 0.095$;
Hoffmann 2022's $\alpha \approx 0.34$ is for the compute-optimal
joint allocation, a different regime). The
model-overhead caveat in §7 (T7.15 framing) and the critique of
unadjusted-LM-as-compressor metrics cite Zhang et al.\ (2024)
\emph{L3TC: Leveraging RWKV for Learned Lossless Low-Complexity
Text Compression}, arXiv:2412.16642.

\emph{Authenticated encryption.} Theorem 7.23's per-sub-block MAC
uses the Bellare--Namprempre (2000)
\emph{Authenticated Encryption: Relations among Notions and
Analysis of the Generic Composition Paradigm}, ASIACRYPT 2000,
LNCS 1976, for the encrypt-then-MAC compositional security
analysis. HMAC security is established in Bellare, Canetti, and
Krawczyk (1996) \emph{CRYPTO}; Poly1305 in Bernstein (2005)
\emph{Fast Software Encryption}; BLAKE3 in O'Connor, Aumasson,
Neves, and Wilcox-O'Hearn (2020) ``BLAKE3: one function, fast
everywhere'' (specification and reference implementation).

\emph{Content-defined chunking.} Theorem 7.17's insertion/deletion
robustness uses the content-defined chunking of Rabin (1981)
``Fingerprinting by random polynomials'', Harvard Tech.\ Report
TR-15-81, as deployed in LBFS (Muthitacharoen, Chen, and Mazi\`eres
2001) \emph{Proc.\ SOSP} for distributed file systems.

\emph{Functional and multi-source coding.} Theorem 7.8 (functional
compression) uses the characteristic-graph entropy of Orlitsky and
Roche (2001) \emph{IEEE Trans.\ Inf.\ Theory} 47(3):903; Theorem 7.7
(successive refinement) uses Equitz and Cover (1991) \emph{IEEE
Trans.\ Inf.\ Theory} 37(2):269; Theorem 7.12 (multi-source joint
coding) uses Lauritzen and Spiegelhalter (1988) \emph{J.\ Royal
Stat.\ Soc.\ B} 50(2):157 for the junction-tree entropy
decomposition.

\emph{Universal online compression.} Theorem 7.6 uses
Krichevsky--Trofimov estimator regret bounds (Krichevsky and
Trofimov 1981 \emph{IEEE Trans.\ Inf.\ Theory} 27(2):199), which
form the foundation for universal-coding regret analysis.

\textbf{13. Status of open problems}

This section reports the status of every problem the framework opens,
distinguishing problems resolved within the paper from those that are
genuinely open and from those deferred to engineering or empirical work.
The taxonomy follows the structure recommended in earlier framework
reviews: a problem is reported in §13.1 if its answer is a theorem, in
§13.2 if its answer requires additional theoretical or empirical work
that has not been done, and in §13.3 if its answer is a piece of code, a
conformance test, or a policy statement rather than a mathematical
result.

\textbf{13.1 Resolved problems}

The following problems are resolved in this paper. Cross-references
point to the relevant theorems.

\begin{itemize}
\item
  Type-I converse bound --- resolved by Theorem 4.2. Part (C) of
  Theorem 4.2 establishes the converse lower bound
  $\geq H(X \mid Y)$ for any uniquely-decodable Type-I code (a
  Shannon-source-coding consequence). Part (A) shows that the
  \emph{ideal} class+payload coder driven by the true conditional
  laws achieves $H(X \mid Y) + \varepsilon_N + O(1)$ headers, so
  the ideal Type-I rate matches the converse to first order in $N$.
  The concrete construction of Theorem 4.1 matches this only when
  the empirical class model $P_E$ and uniform within-class payload
  match the true conditional laws; under model misspecification the
  gap is cross-entropy plus KL, which can be $\Theta(N)$.
\item
  Type-II converse --- resolved by Theorem 5.3b. Part (C) gives the
  base bound $\mathbb{E}[L_{\mathrm{II}}] \geq H(X \mid Y)$ for any
  uniquely-decodable Type-II code (Shannon on the partition, which is
  in bijection with $X$), the operative converse, met by ideal joint
  coding of $\Pi$. Part (S) shows the count-plus-rank form (transmit
  counts, then an enumerative index, coded by any scheme) also has
  converse $H(\mathbf{k} \mid Y) + H(\Pi \mid Y, \mathbf{k}) = H(X \mid Y)$,
  with the index's only structural lower bound being the within-class
  entropy $H(\Pi \mid Y, \mathbf{k})$. The flat-enumerative
  (uniform-index) candidate-set cost $\log_2 |\mathcal{C}(Y,\mathbf{k})|$
  --- the multinomial baseline $\log_2 M(N; \mathbf{k})$ only when no
  support is used --- is a converse only for that scheme; it equals the
  entropy exactly when the conditional law is uniform on the
  decoder-available candidate set (then Theorem 5.3 is entropy-optimal,
  part T), and otherwise exceeds it by the recoverable
  within-candidate-set uniform deficit $G_{\mathrm{within}}^{Z}$ (a KL
  divergence), with any source-dependent support metadata $Z$ charged at
  $H(Z \mid Y, \mathbf{k})$. Three explicit counterexamples ($N=4$,
  $\mathbf{k}=(1,3)$) show $\log_2 M$ is not a converse for the broader
  count-plus-rank form (entropy-coding the index recovers
  $G_{\mathrm{within}}^{Z}$ without leaving that form), nor even for
  support-restricted flat codes (whose cost is the smaller
  $\log_2 |\mathcal{C}|$), and that smuggling partition information into a
  source-dependent support cannot beat the bound (the metadata header
  $H(Z \mid Y, \mathbf{k})$ pays it back, total $\geq H(X \mid Y)$ by the
  chain-rule collapse $H(\mathbf{k}, Z, \Pi \mid Y) = H(\Pi \mid Y)$). So
  Type-II is order-optimal in general and entropy-optimal in the
  uniform-candidate-set regime; the flat-indexing overhead is the price of
  uniform enumeration on the candidate set, not a property of Type-II or
  even of the count-plus-rank decomposition. This is distinct from Lemma
  5.1e, which converses sequential sub-block random access, not the
  partition representation.
\item
  Computational complexity of optimal $E_{12}$ --- partially
  resolved by Theorem 4.3 (NP-hardness of the varying-parameter
  learned-code labeling problem via Wagner and Corneil's 1990
  tree-to-hypercube subgraph embedding). The fixed Code12 instance
  ($n_0 = 256, m_0 = 12$) is constant-size and hence formally
  polynomial in the input bit-length of the empirical confusion data
  (Remark 4.3a). A sound constant-gap inapproximability proof in
  the Hamming-restricted geometry remains open and is listed in
  §13.2.
\item
  Joint structure on natural data and the Type-II separation regime
  --- partially resolved. Theorem 5.5 is a theorem: the
  multi-information $\mathrm{TC}(D_N)$ is the exact
  information-theoretic quantity governing factorized-versus-joint
  separation. Four explicit $\Theta(N)$-TC classes are now
  established as theorems: Lemma 5.6a for uniqueness-constrained
  injections (database-page schemas), Lemma 5.6b for hidden-class
  iid sources (latent-topic mixtures), Lemma 5.6c for stationary
  processes (order-$W$ Markov chains and any process with summable
  conditional-entropy excess; Ces\`aro form for general stationary),
  and Lemma 5.6d for hidden Markov models with the Baum--Petrie
  filter-recursion entropy rate (combined hidden-and-temporal
  structure). Remark 5.6's extrapolation to low-dimensional
  natural-data manifolds remains an informal heuristic. Establishing
  $\mathrm{TC}(D_N) = \Theta(N)$ for less combinatorially structured
  classes (natural images at a given resolution, continuous-state
  HMMs, non-stationary processes, factor graphs over expanders) is
  open and listed in §13.2.
\item
  Exact optimal candidate-support threshold for semi-non-factorized
  Type-II support selection --- resolved by Lemmas 5.7a and 5.7c
  jointly as a polynomial-time vs.\ NP-hard dichotomy along the
  permutation-like boundary. Lemma 5.7a shows top-$Q$ is provably
  optimal under any permutation-like (conditional Bernoulli) joint
  law via a stochastic dominance coupling; Lemma 5.7c shows the
  general unrestricted-dependence case is NP-hard via reduction
  from CLIQUE at $N = 3s + 1$ (where the algebraic identity
  $A_s(1) = A_s(2)$ collapses the cost to a function of $n_0(S)$
  alone). Together this characterizes the exact complexity
  threshold of the support-selection problem. The complementary
  approximation question --- whether a polynomial-time
  $\alpha$-approximation exists for $\alpha < O(N^{1/4 +
  \varepsilon})$ on general dependent joint laws --- is tied to
  the long-standing open approximability of densest-$k$-subgraph
  and is listed in §13.2.
\item
  Tightness of the $2\varepsilon$ mode-selection bound --- resolved
  by Theorem 6.2: the $2\varepsilon$ factor in Theorem 6.1 cannot be
  improved without additional assumptions on the prediction error
  structure.
\item
  Optimal hierarchical decomposition for Type-III-B --- resolved by
  Theorem 6.4: all rooted binary tree decompositions of the 256-byte
  alphabet achieve the same information-theoretic rate
  $H(B \mid C)$; the choice of tree affects only encoder/decoder
  complexity and predictor calibration.
\item
  Bits-back integration into the $L(R \mid Y)$ decomposition ---
  resolved by Theorem 6.6: the three-stream accounting yields
  expected rate equal to the cross-entropy of the data under the
  model marginal plus the variational gap, that is,
  $H_D(X \mid Y) + \mathrm{KL}(D(\cdot \mid Y)\,\Vert\,
  P_X(\cdot \mid Y)) + \mathrm{KL}(q(z \mid X, Y)\,\Vert\,
  P(z \mid X, Y))$; this matches the Slepian-Wolf bound
  $H_D(X \mid Y)$ exactly when both KL terms vanish (model marginal
  matches data, and $q$ matches the true posterior).
\item
  Standalone achievability for the learned hierarchical Type-III-C
  dictionary/token-cover --- resolved \emph{conditionally on a universal
  dictionary builder} by Theorem 6.6b (with Lemma 6.6a). If the
  dictionary is produced by a universal grammar transform --- an LZ78
  incremental parser (Ziv--Lempel 1978; Cover--Thomas \S13.5) or an
  irreducible-grammar transform (Kieffer--Yang 2000, Theorems 7--8) ---
  the Type-III-C encoder attains $\mathbb{E}[L_{\mathrm{IIIC}}(X^N)] \le
  N\,h(X) + o(N)$ for stationary ergodic $X$, i.e.\ per-symbol rate $\to
  h(X)$. The hierarchical family/variant index of Definition 3.5 is
  rate-neutral (Lemma 6.6a, a two-level specialization of the
  tree-invariance Theorem 6.3), and the $K = c(N) = O(N/\log N)$
  dictionary entries add only $o(N)$ index/header overhead by the
  Lempel--Ziv phrase-count bound. The numerical check
  (\texttt{scripts/verify/thm\_6\_6b\_typeiiic\_achievability.py})
  exhibits the per-symbol gap $\to 0$ at the predicted
  $O(\log\log N/\log N)$ rate on order-1 Markov sources and confirms the
  re-index is rate-neutral to machine precision. The companion question
  --- whether the \emph{specific} greedy marginal-MDL builder of
  Theorem 6.5 is itself a universal (irreducible) transform --- is
  \emph{not} resolved and is listed in §13.2.
\item
  Computational complexity of joint dictionary-and-predictor
  learning --- partially resolved. Theorem 6.7 establishes
  NP-hardness of the \emph{restricted} dictionary-only case via
  reduction from SGP; Theorem 6.8 transfers APX-hardness from
  Charikar et al.\ (2005, Theorem 1) to the same restricted case
  under the \emph{unbounded-alphabet} variant with explicit
  inapproximability constant $\gamma_{\mathrm{SGP}} = 8569/8568$.
  The constant derives from Berman-Karpinski's (1999) bound on
  Vertex-Cover-3 propagated through the $15|V|+3|E|+k$ encoding
  overhead in the Charikar et al.\ reduction; the algebra is
  reproduced in the verification script
  \texttt{scripts/verify/thm\_6\_8\_sgp\_constant.py} in the
  companion repository. Hardness of the \emph{unrestricted}
  joint problem (model, dictionary, parser, and repair stream all
  free) is conjectured but not proved by the present reductions
  and is listed in §13.2.
\item
  Bit-exact decoder reproduction for fixed predictors --- resolved
  by Theorem 10.1: a precision-and-determinism specification
  suffices for cross-platform reproducibility, reducing the question
  to a conformance test.
\item
  Bit-exact reproduction under adaptive predictors --- resolved by
  Definition 10.1 (compatible online optimizer) and Theorem 10.2:
  online adaptation via integer-arithmetic optimization is bit-exact
  across implementations, with several published optimizers fitting
  the compatible class.
\item
  Random access (seek-and-read) into a compressed archive with
  output-sensitive cost $O(\log L(R) + (K + k) \cdot
  \mathrm{cost}(M))$ under uniform sync-point spacing (or
  $O(\log(|X|/K) \cdot (\log |X| + \log L(R)) + (K + k) \cdot \mathrm{cost}(M))$
  under non-uniform spacing) for $W$-bounded predictors and
  \emph{per-position sequential} coding modes (Type-I, per-byte
  sequential Type-III-B, or Type-III-A composition over
  per-position sequential per-tile modes) --- resolved by
  Definitions 10.2 ($W$-bounded predictor) and 10.3 (sync point,
  including entropy-coder flush) and Theorem 10.3. The achievable
  rate for byte-granular random access is the $W$-truncated
  conditional rate $H_W(D_N)$ (Lemma 5.1a), strengthened by
  Lemma 5.1b for exponentially conditional-MI-decaying stationary
  sources: choosing $W = O(\log |X|)$ brings $H_W$ within $O(1)$ of
  the joint entropy $H(D_N)$. The total expected length is
  $H(D_N) + O(1) + \Delta_{\mathrm{sync,total}} + \varepsilon_N \cdot N$
  with $\Delta_{\mathrm{sync,total}} = O((|X|/K) \log |X| \log_2(1/\eta))$:
  for $K = \log^a |X|$ with $a > 1$ the total sync overhead is
  $O(|X|/\log^{a-1} |X|)$ (sub-linear in $|X|$ but not polylog);
  at $K = \Theta(\log |X|)$ the total is $\Theta(|X|)$ linear.
  Polylog total overhead would require $K = |X|/\mathrm{polylog}(|X|)$,
  which makes per-query cost polynomial.
  Lemma 5.1c sharpens Lemma 5.1a's worst-case sync overhead
  $m \cdot S \cdot W \cdot \log_2(1/\eta)$ to the source-specific
  excess $(m - 1) \cdot \delta_{W'}$ for the special case of order-$W'$
  Markov sources with predictor context $W \geq W'$ and full sub-blocks
  of size $K \geq W'$ ($N = mK$), where
  $\delta_{W'} = H(X^{W'}) - W' h(X)$ is finite and typically much
  smaller than $S \cdot W \cdot \log_2(1/\eta)$ for natural-data
  Markov sources. Lemma 5.1d extends this characterization to general
  exp-mixing stationary sources with the Crutchfield--Feldman excess
  entropy $\delta_\infty$ bounding the cold-context excess per sync.
  Lemma 5.1e establishes a matching information-theoretic lower bound
  $L_{\mathrm{synced}} \geq m \cdot H(D_K)$ for sub-block-independent
  schemes (source-independent metadata, independently decodable
  payloads), so Lemma 5.1c's $(m-1) \cdot \delta_{W'}$ excess is
  information-theoretically optimal for order-$W'$ Markov sources
  under the sub-block paradigm. Corollary 5.1f assembles Lemmas
  5.1a-e into a unified resolution of OP1 (random access) within the
  sub-block paradigm, leaving non-sub-block / state-snapshot schemes
  (Theorem 10.4) and non-summable mixing as open extensions.
  Whole-block partition-rank coding (Type-II rank coding,
  Type-III-B direct or hierarchical partition strategies) is
  excluded from Theorem 10.3 and gives only sub-block-granular
  random access at $O(K \cdot \mathrm{unrank}(K))$ per requested
  sub-block (Remark 10.3a). Non-$W$-bounded
  predictors (LSTM, state-space) are covered by Theorem 10.4 via an
  augmented state snapshot of size $L(\text{state})$ per sync
  point; the corresponding cost adds $L(\text{state})$ to the
  uniform-spacing lookup term and $L(\text{state})$ to each
  non-uniform binary-search entry (full statement in Theorem 10.4). The cost is polylogarithmic in
  archive size $|X|$ only in the parameter regime where the
  sync-point spacing $K$, read window $k$, predictor-evaluation cost
  $\mathrm{cost}(M)$, and $\log L(R)$ are held fixed or grow at most
  polylogarithmically in $|X|$; outside that
  regime (e.g., $K = \Theta(|X|)$ or $k = \Theta(|X|)$) the cost is
  $\Theta(|X| \cdot \mathrm{cost}(M))$. This is the framework's
  deployment-grade distinguishing feature relative to other neural
  compressors: autoregressive byte/pixel models (NNCP, PixelCNN++,
  PixelSNAIL) as currently designed do not insert sync points or
  maintain a seek index and so do not support random access in
  their published form. Adding sync points to their encoders would
  give them Theorem 10.3's random access \emph{provided} the
  underlying predictor is $W$-bounded; for autoregressive
  transformers with fixed context window this holds, but for
  autoregressive RNNs/state-space models (which are not
  $W$-bounded) Theorem 10.4's state snapshots are also required
  per sync point. Block-codec schemes (Bit-Swap, IDF++, HiLLoC)
  provide only block-granular access (not byte-granular) due to
  their block-coupled latent/ANS state.
\item
  Block-device and read-only filesystem deployment with
  hardware-target portability --- resolved by Theorem 10.5
  (block-device interface: an OS-level driver serves LBA-indexed
  read requests with the same output-sensitive cost as Theorem 10.3,
  polylogarithmic in $|X|$ only in the same fixed/polylog $K$, $k$,
  $\mathrm{cost}(M)$, $\log L(R)$ parameter regime) and Theorem 10.6 (hardware-target
  portability: bit-exact predictor execution is guaranteed across
  any pair of targets each satisfying the (a)--(g) conformance
  conditions of Theorem 10.6; whether specific named target classes
  --- CPU scalar, CPU SIMD, GPU integer Tensor Cores, integer NPUs
  --- actually conform is an engineering claim per target, not a
  corollary of the theorem). The copy-on-write
  overlay pattern (§10.8) enables read-write semantics using the
  same architecture as existing compressed-filesystem deployments.
\item
  Byte-granular RA architecture with neural predictor and concrete
  numerical end-to-end synthesis --- resolved by the §7.2--7.28
  framework (26 theorems plus corollaries):
  \emph{architecture}: Theorem 7.2 establishes the basic sub-block
  byte-granular RA construction with cost
  $O(\log L(R) + K \cdot \mathrm{cost}(M'))$ per query and total
  $N H_{M'} + m \delta_\infty + O(m \cdot p)$, with Theorem 7.5
  giving optimal $K^* = \Theta(\sqrt{N})$ for balanced workloads.
  \emph{Source-coding extensions} are theorems: T7.6 (universal
  online with Krichevsky--Trofimov regret), T7.7 (successive
  refinement), T7.8 (functional compression via Orlitsky--Roche
  characteristic-graph entropy), T7.9 (lossy RNR with distortion
  budget $D$ via Shannon $R(D)$), T7.10 (Slepian--Wolf/Wyner--Ziv
  side-information composition), and the user-originated lossy+ECC
  Corollary 7.11. \emph{Multi-source and structural} are theorems:
  T7.12 (junction-tree entropy decomposition for multi-source
  joint), T7.13 (per-sub-block ECC for noisy-storage RA with
  Shannon-optimal expansion $1/(1 - h(\varepsilon))$). \emph{Concentration
  and end-to-end} are theorems: T7.14 establishes Doob-martingale
  Azuma--Hoeffding concentration with the
  $(W_N + 1 + \rho/(1-\rho))$ factor (predictor-context plus
  source-mixing-tail under exp-mixing), and T7.15 is the explicit
  end-to-end Chinchilla-class numerical prediction matching
  Delétang et al.\ 2024 Table 1 ($79\,\mathrm{MB}$ on enwik9 at
  $12.0\times$ bitstream rate) with dual concentration framing
  (rigorous Azuma vs.\ practical CLT). \emph{Practical extensions}
  are theorems: T7.16 (adaptive per-region $K_r$ with strict
  Cauchy--Schwarz improvement over fixed $K$), T7.17
  (content-defined chunking for insertion/deletion robustness),
  T7.18 (parallelizable encoding across $P$ workers with
  $O(N/P + \sqrt{N})$ span), T7.19 (streaming encoder with $O(K)$
  memory and trailing-footer layout), T7.20 (cache-aware decoder
  amortization with exact static-top Zipf hit-rate asymptotics),
  T7.21 (predictor-archive size trade-off
  $L^*(N) = \Theta(N^{1/(1+\alpha)})$ via Hoffmann 2022 scaling),
  T7.22 (fail-safe with per-sub-block store-mode worst-case
  $L^{\mathrm{rep}} \leq N \log_2 |\Sigma| + m$ under any source
  and predictor), T7.23 (per-sub-block authenticated MAC integrity
  with $O(\tau \cdot m)$ overhead and cryptographic forgery
  resistance under standard PRF assumptions), T7.24 (revised)
  (minimal honest decoder-latency Hoeffding tail bound under IRM
  workload + static Perfect-LFU cache + bounded $T_{\max}$), and
  T7.25 (variance-aware Bernstein concentration of
  $L^{\mathrm{rep}}$ via Merlev\`ede-Peligrad-Rio 2009 IMS Collections
  Theorem~2, with the long-run-variance proxy $v^2 = \sigma_\ell^2
  (2W_N + 1) + O(B^2)$ derived via Davydov 1968 covariance inequality:
  for Chinchilla-class enwik9, the rigorous $\delta = 10^{-6}$
  deviation drops from $\approx 1.4\,\mathrm{GB}$ (Azuma, T7.14) to
  $\approx 2.1$--$4.2\,\mathrm{MB}$ (Bernstein, T7.25), a
  $\approx 330\times$--$670\times$ improvement and within
  $\sim 50$--$100\times$ of the practical CLT estimate; two earlier drafts of T7.25 were
  self-rejected on cold review: (i) a Samson 2000 direct-Lipschitz
  formulation that gave bound worse than Azuma because each
  coordinate $x_j$ affects $(W_N+1)$ terms in the sum, and (ii) a
  variance-substitution shortcut that conflated per-position
  $\sigma_\ell^2$ with long-run $v^2$; the Davydov-corrected version
  with the explicit $B^2$ correction in $v^2$ is the version stated),
  and T7.26 (LRU+IRM miss-count concentration via Lezaud 1998
  Chernoff bound for finite-state Markov chains, with positive
  spectral gap $\gamma_{\rm LRU} \geq \pi_{\rm min}$ from Dobrushin
  coupling: $\sqrt{Q}$-Hoeffding scaling under IRM with full-support
  popularity, closing the open question from earlier drafts of §13.2
  about finite-batch concentration of LRU decoder latency that
  T7.24's static-LFU formulation sidestepped). \emph{Converse and
  distributional refinements}: T7.27 (tight space-time product
  converse $C(K)\cdot S(K)=\Theta(N\log N\cdot\mathrm{cost}(M))$ for
  the sub-block class, with T7.27$'$ refuting the general-scheme
  analogue and Corollary 7.27c recovering it on the predictor-faithful
  class) and T7.28 (non-asymptotic Berry--Esseen CLT for
  $L^{\mathrm{rep}}$ at Kolmogorov rate $O(N^{-1/2}(\log N)^2)$).
  The §7.24 history
  (LRU version retracted under an adversarial disproof, rewritten under
  restricted hypotheses, survived round 2) illustrates the
  validation discipline applied throughout.
  All 27 theorems (plus the 7.6$'$/7.27$'$ refinements and the
  7.11/7.27c/7.27e/7.27f corollaries) verified numerically in
  \texttt{scripts/verify/thm\_7\_*} and adversarially
  cross-validated by three independent adversarial review passes.
\end{itemize}

\textbf{13.2 Remaining mathematical open problems}

A small number of mathematical questions remain genuinely open after
the present paper. (Optimal exact support selection for
semi-non-factorized Type-II was previously listed here; it has been
relocated to §13.1, since Lemmas 5.7a and 5.7c together establish
the polynomial-time vs.\ NP-hard dichotomy on the permutation-like
boundary. Only the approximation side remains research-scale open,
listed below.)

\emph{Status ledger for ``optimal'' RNR encoding (the OP2/OP4
support/grammar-selection complexity questions).} Computing the optimal
RNR encoding is hard only in specific restricted variants; several
sub-questions are in fact already \emph{closed}. The headline ``optimal
RNR is not known to be efficiently approximable'' is therefore a
precisely-pinned residual, \emph{not} a blanket open problem. The
following ledger states, per sub-question, the exact status (CLOSED /
\textsc{in P} / OPEN), the paper location, the proof (for closed rows)
or the named structural blocker (for open rows), and the witnessing
executable verification script. Each open row names the obstacle that
blocks it; no closed sub-case is hidden behind an overclaim, and no open
sub-case is understated as resolved.

\begingroup\footnotesize\sloppy
\setlength{\tabcolsep}{4pt}
\noindent\begin{tabular}{@{}p{2.7cm}p{1.5cm}p{2.6cm}p{3.3cm}p{3.0cm}@{}}
\hline
\textbf{sub-question} & \textbf{status} & \textbf{paper location} &
\textbf{proof (CLOSED) / named blocker (OPEN)} &
\textbf{witnessing verify script}\\
\hline
OP2(a) continuous bit-cost SGP
& CLOSED
& Lemma 6.8e (this §13.2); Lemma 6.8b; §13.4.1
& inherits Charikar $\gamma_{\mathrm{SGP}}=8569/8568$ with vanishing
$O(1/(|V|\log|V|))$ correction (alphabet floor pins $\log_2 N$)
& \texttt{lemma\_6\_8b\_bit\_cost\_asymptotic.py};
\texttt{op2a\_k\_escape\_budget\_check.py} (alphabet-floor re-derivation)\\
OP2(a) discrete integer/ceiling bit-cost SGP
& OPEN
& Conj.\ 6.8d (this §13.2); §13.4.1
& \emph{blocker:} ceiling discontinuity at power-of-two $|\Sigma|+K$
boundaries; needs an SGP constant $>1+1/\log_2 N$, but
$\gamma_{\mathrm{SGP}}-1=1/8568<1/\log_2 N$ for all $N<2^{8568}$
& \texttt{op2a\_sgp\_constant\_survey\_check.py};
\texttt{op2a\_k\_escape\_budget\_check.py} (open; negative-route only)\\
OP2(a) additive rule-cost $L_\alpha$ (side result)
& CLOSED (positive)
& Lemma 6.8k / 6.8k-approx; §13.4.1
& APX-hard for every $\alpha\ge0$ with a ceiling-free gap monotone in
$\alpha$ (not a relabelling of 6.8e)
& \texttt{lemma\_6\_8k\_approx\_additive\_rule\_cost.py};
\texttt{op2a\_additive\_rule\_cost\_check.py}\\
OP2(b) free-field continuous (amortised Gaussian)
& \textsc{in P}
& Remark 6.9a; §13.4.4
& sample-covariance ML plug-in is decoder-reproducible; i.i.d.\ Gaussian
draws cannot hide their poly-time-estimable covariance
& \texttt{remark\_6\_9a\_continuous\_op2b\_easy.py}\\
OP2(b) single-source \& amortised-family discrete (free predictor)
& OPEN
& §13.2 (Remark 6.8b); Lemma 13.4.2a; Thm 13.4.2b; §13.4.2
& \emph{blocker:} predictor-absorption is an \emph{upper bound} only
(both YES/NO bounded by poly$(n)$ does not bound the ratio); the PRG
escape merely reduces amortised OP2(b)-APX to single-string
OP2(a)-APX (Conj.\ 6.8d), open under $\mathrm P\neq\mathrm{NP}$; the
OWF route gives only average-case \emph{decision} hardness (Thm 13.4.2b)
& (open; no PASS fabricated --- negative/conditional routes only)\\
OP4 Type-II general dependent-law support selection
& OPEN
& §13.2 (first bullet); Lemma 5.7c/5.7d; §13.4.4
& \emph{blocker:} affine $n_0(S)$ identity dilutes every count /
partition-functional gap to $1+\Theta(D_{\mathrm{YES}}/m)$ (worse for
hyperedges); natural $A_s$ saturates at $1+O(1/\log s)$; only
average-case (planted-clique) remains
& \texttt{op4\_dksh\_affine\_dilution.py};
\texttt{op4\_noncount\_entropy\_objective.py}\\
OP4 Type-III-C fixed-field (continuous log-det)
& CLOSED (\textsc{Partial})
& Lemma 6.9; §13.4.4
& the cost is a Schur-complement log-det $=$ MESP / $D$-optimal design,
APX-hard with constant $5/4$ (Ko--Lee--Queyranne; Ohsaka 2022) under
$\mathrm P\neq\mathrm{NP}$; gap survives the manifold restriction at
high SNR; legitimate (the log-det \emph{is} the cost, no substitution)
& \texttt{lemma\_6\_9\_typeiiic\_mesp\_support\_selection.py}\\
OP4 $E_{12}$ constant-gap (Hamming-restricted labeling)
& OPEN
& Conj.\ 4.3e (this §13.2); Remark 4.3d; Thm 4.3 / Remark 4.3c
& \emph{blocker:} missing dilation-$1$/robust-NO promise gap --- a
hypercube-embedding source with a YES dilation-$1$ instance and a
constant-fraction robust-NO instance; \emph{not} affine dilution and
\emph{not} row-stochasticity (Remark 4.3d)
& \texttt{remark\_4\_3d\_e12\_constant\_gap\_dilution.py}\\
OP4 fixed-byte-alphabet restricted Type-III-C ($|\Sigma|=256$)
& OPEN
& §13.2 (this list)
& \emph{blocker:} byte-encoding the $|V|$-vertex terminals adds a
$\Theta(\log|V|)$ multiplicative overhead that can dilute the
$\gamma_{\mathrm{SGP}}$ gap toward $1$
& (none; no gap-preserving construction known)\\
\hline
\end{tabular}
\par\endgroup

\noindent The consolidated index probe
\texttt{scripts/verify/bug\_005\_optimal\_encoding\_status\_index.py}
asserts the witnessing scripts for the CLOSED / \textsc{in P} rows are
present and PASS, and prints the ledger's CLOSED/OPEN status with a final
\texttt{OVERALL -> PASS} (it fabricates no PASS for an OPEN row). The
remaining genuinely-open residual is the set
$\{$OP2(a)-discrete (Conj.\ 6.8d), OP2(b)-discrete-APX, OP4 Type-II
approximability, OP4 $E_{12}$ (Conj.\ 4.3e), OP4 fixed-byte$\}$; the
individual entries below give the full detail.

\begin{itemize}
\item
  Approximability of optimal support selection for general
  dependent joint laws of $M_v$ (the complementary direction to
  Lemma 5.7c's NP-hardness). The positive-affine reduction
  $\mathbb{E}_\mu[A_s] = A_s(1) - (A_s(1) - A_s(0)) \cdot n_0(S)/m$
  does not preserve multiplicative approximation ratios, so neither
  $O(\log N)$-approximation nor PTAS inapproximability follows
  directly from Lemma 5.7c. The question is tied to the
  long-standing open approximability of densest-$k$-subgraph (best
  known algorithm $O(n^{1/4+\varepsilon})$ via Bhaskara et al.\
  2010); closing either direction would be a research-scale
  contribution.
\item
  Hardness of the unrestricted joint Type-III-C problem
  $(M, D, R \mid Y)$ --- partially resolved. Theorems 6.7 and 6.8
  establish NP-hardness and APX-hardness ($\gamma = 8569/8568$ under
  unbounded alphabet) for the restricted dictionary-only case under
  the symbol-count objective $L_{\mathrm{sym}}$. Lemma 6.8b
  establishes class-restricted bit-cost APX-hardness with the
  \emph{same} constant $\gamma_{\mathrm{SGP}} = 8569/8568$ within the
  Charikar normal-form grammar class on Berman-Karpinski-hard
  instances, via a ratio-preserving reduction: since
  $L^{\mathrm{bit}}(\sigma) = L_{\mathrm{sym}}(\sigma) \cdot c(\sigma)$
  with $c$ non-decreasing in $|\sigma|$ and the bit-cost-optimal cover
  coinciding with the minimum cover $\sigma^*$, any bit-cost
  $\alpha$-approximation $\sigma'$ has $|\sigma'| \geq |\sigma^*|$
  hence $c(\sigma') \geq c(\sigma^*)$, so $L_{\mathrm{sym}}(\sigma')
  \leq \alpha \cdot L_{\mathrm{sym}}(\sigma^*)$ — yielding a
  symbol-count $\alpha$-approximation, NP-hard for $\alpha <
  \gamma_{\mathrm{SGP}}$. The following remain open:
  (a) bit-cost APX-hardness over \emph{unrestricted} grammar
  structures: an arbitrary grammar deriving Charikar's source string
  may trade larger $L_{\mathrm{sym}}$ for smaller nonterminal count
  $K$ (different functional form than $K_0 + \kappa \cdot |\sigma|$),
  potentially achieving lower bit-cost than the Charikar-format optimum.
  A literature search (Casel et al.\ 2021, Kozma \&
  Voderholzer 2024, Ga\'nczorz 2018/2020) confirms that no
  Pareto-frontier structural theorem is currently known that would
  rule out such ``escape grammars'' from beating Charikar's bit-cost
  Pareto point; this remains genuinely open. (b) APX-hardness of the
  unrestricted joint
  problem with free predictor $M$, which adds the dimension of paying
  $L(M)$ bits to gain $L(R \mid Y)$
  savings. Remark 6.8b formalizes the failure of the
  natural padding-based reduction: an unrestricted predictor can
  pay $L(M) \geq |s|$ bits to encode the SGP-format prefix directly,
  yielding $L(R \mid Y) = O(1)$ on the prefix while still paying
  Kolmogorov-complexity bits for the incompressible suffix, so the
  $\gamma_{\mathrm{SGP}}$ structure of the prefix is absorbed into
  $L(M)$ rather than reflected in the joint bit-cost. The inherited constant
  for either open case is not necessarily $8569/8568$; the actual
  bit-cost / free-predictor constants could be strictly less than
  $8569/8568$ due to objective-function dilution and remain
  research-scale open problems.

  \emph{Partial closure of OP2(a) via continuous bit-cost (Lemma 6.8e).}
  Lemma 6.8e (Section 6.8) closes the continuous-bit-cost variant of
  OP2(a) by showing that the information-theoretic objective
  $L^{\mathrm{cont}}(G) = |G| \cdot \log_2(|\Sigma| + K(G))$
  (no ceiling, the natural arithmetic-coding target) inherits Charikar's
  $\gamma_{\mathrm{SGP}} = 8569/8568$ APX-hardness with a vanishing
  $O(1/(|V|\log|V|))$ correction. This shifts the residual difficulty
  of OP2(a) to a \emph{specific artifact of the integer-bit objective}
  --- the ceiling discontinuity at power-of-two boundaries identified
  by the adversarial review for Conjecture 6.8d. The
  continuous variant is also the natural object for compression theory
  (Shannon-McMillan-achievable fractional-bit rates), while the
  integer/ceiling objective is the prefix-code worst case rarely
  used in modern compression. The remaining open case is thus
  specifically the integer/ceiling variant of unrestricted bit-cost SGP,
  for which Conjecture 6.8d (weak form) is the strongest current
  statement.

  \emph{Amortized-family reformulation of OP2(b) (new direction,
  2026).} The single-source predictor-absorption obstacle of
  Remark 6.8b has a natural reformulation that circumvents the
  $L(M) \geq |s|$ absorption: switch to an \emph{amortized family}
  $\mathcal{F}_n = \{s_i\}_{i \in [N(n)]}$ of $N(n) = 2^{\Theta(n)}$
  SGP-prefix instances, with the objective
  $L_{\mathrm{amort}} := L(M) + L(D) + \sum_i L(R_i \mid Y, M, D)$
  shared across the family. For $M$ to memorize all $N(n)$ prefixes,
  $L(M) \geq N(n) \cdot n$ bits, dominating any savings; per-instance
  amortization $L(M)/N(n) \to 0$ forces $M$'s contribution to
  $o(1)$ bit/source while each $L(R_i)$ must still encode its SGP-instance.
  The reformulation aligns with the Generalized Smallest Grammar
  Problem of Carrascosa et al.\ (2016, arXiv:1608.08927) and connects
  to the multi-instance distributional-hardness machinery of
  Huang-Ilango-Ren STOC'23 (\emph{NP-Hardness of Meta-Complexity},
  hanlin-ren.github.io). \emph{Caveat:} this is a definitional shift,
  not a proof of the literal single-source OP2(b); the literal version
  remains open with Remark 6.8b as the structural obstacle. The
  amortized-family variant is the natural research target where
  APX-hardness is plausible via the per-instance salting technique.

  \emph{Indirect connection via Liu--Pass / one-way functions (2026
  update).} A complementary indirect route to OP2(b) hardness goes
  through cryptographic hardness via the Liu--Pass equivalence
  (Liu \& Pass, FOCS 2020 / J ACM 2023, \emph{On One-Way Functions
  and Kolmogorov Complexity}, arXiv:2009.11514): the existence of
  one-way functions is equivalent to the hardness, on samplable
  distributions, of the Minimum-Kt-Complexity Problem
  (MKtP), where $\mathrm{Kt}(x) = \min_p \{|p| + \log t(p) : U^p
  \text{ outputs } x\}$ is time-bounded Kolmogorov complexity in
  Levin's form; the closely-related Liu--Pass variant
  $K^t(x) = \min_p \{|p| : U(p) = x \text{ in } \leq t(|x|)
  \text{ steps}\}$ (with a fixed polynomial $t$) is the measure for
  which the OWF-equivalence is stated in
  Liu--Pass~2020. For OP2(b), the joint cost
  $L(M) + L(D) + L(R \mid Y, M, D)$ is upper-bounded by
  $K^t(s \mid Y) + O(\log N)$ (formal statement: Lemma 6.8g in
  §6.8, which gives an explicit universal constant in front of the
  logarithm; see Lemma 6.8g's scope statement for the precise sense
  in which Levin's $\mathrm{Kt}$ and the Liu--Pass $K^t$ differ),
  with the predictor $M$ playing the role of the
  polynomial-time-bounded ``program'' and $L(R)$ the residual
  incompressibility. Hence: IF one-way functions
  exist, THEN OP2(b) is average-case hard on a samplable distribution
  by direct reduction from MKtP. The reduction is not gap-preserving
  (gives average-case, not worst-case, hardness; and additive Kt-bit
  gap, not multiplicative APX-gap), but provides an alternative
  conditional hardness anchor: under the standard cryptographic
  assumption that OWFs exist, joint Type-III-C compression on
  random samplable sources is hard. \emph{Caveat:} this is a
  \emph{different} hardness statement than Conjecture 6.8d's
  worst-case APX-hardness on Berman-Karpinski instances; the OWF
  route gives cryptographic-strength average-case hardness, the SGP
  route gives combinatorial worst-case hardness. Both are consistent
  with OP2(b) being hard but neither subsumes the other.

  \emph{Two-sided status (epistemic gap).} A literature scan
  (2015--2024) confirms that for OP2(a), \emph{neither} a polynomial-time
  constant-factor approximation for unrestricted bit-cost SGP \emph{nor}
  APX-hardness is known. The best approximation in the symbol-count
  setting is the $O(\log(N/m^*))$ algorithm of Charikar et al.\ 2005 /
  Rytter 2003; the recompression line (Je\.z 2015) maintains
  $O(g \log(N/g))$ grammar size, not constant factor. Casel et al.\
  (2020) \emph{Theory of Computing Systems} 65(1):158--203 explicitly
  list constant-factor SGP approximation as an open question. The
  bit-cost variant has neither a stronger upper bound nor a tighter
  lower bound than this symbol-count baseline. The problem therefore
  occupies a \emph{theory-gap} regime where the gap between known
  upper and lower bounds is unbounded; closing it from either
  direction would be a research-scale contribution.

  \emph{Conjecture 6.8d (Unrestricted bit-cost SGP APX-hardness:
  some constant gap).} We conjecture there exists a constant
  $\varepsilon_{\mathrm{bit}} > 0$ such that no polynomial-time
  $(1 + \varepsilon_{\mathrm{bit}})$-approximation algorithm exists
  for unrestricted bit-cost SGP on sufficiently large
  Berman-Karpinski / Charikar-hard instances, unless
  $\mathrm{P} = \mathrm{NP}$. We deliberately do \emph{not} claim
  that $\varepsilon_{\mathrm{bit}}$ matches the symbol-count
  $\gamma_{\mathrm{SGP}} = 8569/8568$ ratio, since the ceiling
  $\lceil \log_2(|\Sigma| + K(G)) \rceil$ multiplier in the
  bit-cost objective has discontinuous behavior near
  power-of-two boundaries of $|\Sigma| + K$: an escape grammar
  may drop one bit-width unit by adjusting $K$ across a power-of-two
  threshold, saving a fractional $1/\log_2(|\Sigma|+K) \approx
  1/\log |V|$ in bit-cost weight while paying only a smaller
  fractional excess in $L_{\mathrm{sym}}$. For typical
  $|V|$ with $\log |V| < 8568$, this ceiling-drop fraction can
  exceed the $\gamma_{\mathrm{SGP}} = 1 + 1/8568$ gap, so the
  unrestricted bit-cost APX-hardness ratio may be strictly less
  than $\gamma_{\mathrm{SGP}}$. (Two independent
  adversarial reviewers split on this point: one accepted the
  stronger ``$\gamma_{\mathrm{SGP}} - \delta$'' form, the other
  flagged the ceiling-discontinuity attack as a structural concern.
  We adopt the conservative ``some $\varepsilon_{\mathrm{bit}} > 0$''
  form to remain honest about what the existing evidence
  supports.) Motivation: Lemma 6.8b establishes a positive
  APX-hardness ratio within the Charikar normal-form class, and
  escape grammars cannot reduce $L_{\mathrm{sym}}$ below the SGP
  minimum, so some positive gap is plausible even if the exact
  constant is smaller than $\gamma_{\mathrm{SGP}}$. A polynomial-time
  PTAS for unrestricted bit-cost SGP (poly-time $(1 + \varepsilon)$-
  approximation for every $\varepsilon > 0$) would refute this
  conjecture.

  \emph{Observation 6.8d-asymp (Single-escape analysis suggests
  asymptotic strong form; multi-escape obstacle).} A partial analysis
  of single-boundary escape grammars suggests that for $|V| \to \infty$
  the bit-cost gap converges to $\gamma_{\mathrm{SGP}}$. Per single
  escape across one power-of-2 boundary: saving $L_{\mathrm{sym}}$
  bits (one bit-width unit on all symbols), cost
  $\Delta_{\mathrm{sym}} \cdot (\log V - 1)$ bits where
  $\Delta_{\mathrm{sym}} \geq 1$ from inlining a Charikar-normal-form
  nonterminal. Beneficial iff $\Delta_{\mathrm{sym}} < L_{\mathrm{sym}}/(\log V - 1)$,
  with max-beneficial-$\Delta$ giving saving $\sim L_{\mathrm{sym}}/\log V$ bits.
  Single-escape relative saving $\sim 1/\log V$; for $\log V > 8568$
  (i.e., $|V| > 2^{8568}$) this is dominated by the Charikar
  symbol-count gap $1/8568$.

  \emph{Honest caveat: multi-escape / K-optimisation obstacle.}
  Single-escape analysis is INSUFFICIENT for an asymptotic strong-form
  closure. Consider a Charikar-style grammar with $K$ free, treating
  $K$ as an optimisation variable: for the model $L_{\mathrm{sym}}(K)
  \approx K + c \cdot N/K$ ($c$ a constant of the Berman-Karpinski
  gadget), the bit-cost objective $L_{\mathrm{sym}}(K) \cdot \log(|\Sigma|+K)$
  is minimised at $K^* \approx \sqrt{cN/\log V}$ with optimum value
  $\sim 2\sqrt{cN/\log V} \cdot \log V = 2\sqrt{cN \log V}$. The
  $K = O(V)$ Charikar regime is NOT the bit-cost-optimal regime; the
  $K \approx \sqrt{N}$ regime yields strictly smaller bit-cost. The
  YES vs NO ratio in the $K = \sqrt{N}$ regime is not provably
  $\gamma_{\mathrm{SGP}}$ --- the Berman-Karpinski reduction is
  designed to preserve the symbol-count ratio at $K = O(V)$, not at
  arbitrary $K$. The Charikar gap may collapse under K-optimisation,
  giving NO grammars with bit-cost no worse than YES.

  We retract any rigorous claim of asymptotic-strong-form closure of
  Conjecture 6.8d via escape-savings accounting. The right statement
  remains the conservative ``some $\varepsilon_{\mathrm{bit}} > 0$''
  weak form, with the K-optimisation obstacle a documented research
  direction. Resolving it requires either: (a) demonstrating that
  the Berman-Karpinski reduction can be modified to preserve the
  symbol-count gap at the $K = \sqrt{N}$ bit-cost-optimal regime;
  (b) using a different source-of-hardness problem that admits
  $K$-uniform gap preservation; or (c) proving an exponential-type
  lower bound on $L_{\mathrm{sym}}(K^*)$ for NO instances under
  Berman-Karpinski-hard inputs. All three are open research-level
  directions.
\item
  Universality of the greedy marginal-MDL Type-III-C dictionary builder
  --- open. Theorem 6.6b establishes entropy-rate achievability for the
  Type-III-C container \emph{under a universal grammar transform} (LZ78
  or any irreducible-grammar transform, Kieffer--Yang 2000). It does
  \emph{not} establish that the specific greedy admission rule of
  Theorem 6.5 (admit $\tau$ iff its marginal MDL gain $\Delta L(\tau) <
  0$, all else frozen) yields a universal --- equivalently, an
  irreducible --- transform. Greedy grammar heuristics need not be
  irreducible, and Charikar et al.\ (2005) prove constant-factor
  grammar-size lower bounds for several greedy transforms (LZ78,
  \textsc{Bisection}, \textsc{Sequential}) against the smallest grammar;
  a non-irreducible transform is outside the hypothesis of
  Kieffer--Yang's universality theorem, so its expected rate is not
  guaranteed to reach $h(X)$. The open question is whether the
  greedy-MDL builder (i) can be post-processed into an irreducible
  transform without changing its admission decisions, or (ii) is
  directly universal for some natural source class, or (iii) admits a
  source family on which its per-symbol rate is bounded away from
  $h(X)$. A positive answer to (i) or (ii) would upgrade Theorem 6.6b
  from ``achievable under a universal builder'' to ``achievable by the
  Theorem 6.5 builder''; a (iii)-style counterexample would show the
  conditional form is necessary. Lemma 6.6d resolves the \emph{augmented}
  version of this question: the Theorem 6.5 builder is universal for every
  stationary source once its literals are entropy-coded under a universal
  predictor and a one-bit baseline guard is added (so the admission
  decisions become rate-irrelevant), confining the residual difficulty to
  the bare rule's flat-$8\ell$ literal cost. Sub-questions (i)--(iii)
  \emph{for the bare rule} remain open (Remark 6.6e).
\item
  Tight asymptotic scaling of $\mathrm{TC}(D_N)$ for specific
  natural-data classes --- partially resolved. Four explicit
  $\Theta(N)$-TC classes are now established:
  Lemma 5.6a gives the closed-form coefficient
  $c(\alpha) = \log_2 e - (\alpha-1)\log_2(\alpha/(\alpha-1))$ for
  the uniqueness-constrained database-page class (uniform
  injections $[N] \to [\alpha N]$); Lemma 5.6b gives the exact
  identity $\mathrm{TC}(D_N) = N \cdot I(X_1; Y) - H(Y)
  + H(Y \mid X^N)$ for hidden-class iid sources (latent $Y$ with
  conditionally iid samples); Lemma 5.6c gives
  $\mathrm{TC}(D_N) = N \cdot E - \delta_N$ for stationary
  processes, where $E = H_{\mathrm{marg}} - h(X) \geq 0$ is the
  per-position excess entropy and $\delta_N$ is bounded constant
  for order-$W$ Markov chains (covering bigrams, $W$-grams, and
  any mixing process with $\sum_n [H(X_n \mid \mathrm{past}) - h(X)]
  < \infty$; the general stationary case has the Ces\`aro form
  $\mathrm{TC}(D_N)/N \to E$ even without summability; for English
  text Shannon's 1951 measurements give $E \approx 0.71$ bits/letter
  at the bigram level and $E \approx 3.08$ bits/letter at the
  long-range limit using the midpoint of Shannon's $0.6$ to $1.3$
  asymptotic bounds); and Lemma 5.6d extends to hidden Markov models
  with the Baum--Petrie filter-recursion entropy rate
  $h_{\mathrm{HMM}}$, with $E_{\mathrm{HMM}} > 0$ iff the observable
  process is not iid.
  Lemmas 5.6a, 5.6b, 5.6c are verified at machine precision (exact
  TC identity, $\sim 10^{-15}$); Lemma 5.6d's entropy-rate
  characterization $h_{\mathrm{HMM}}$ is verified to Monte Carlo
  tolerance (filter-recursion estimate matches the exact TC identity
  at small $N$). All four scripts are in
  \texttt{scripts/verify/lemma\_5\_6\{a,b,c,d\}*}; the proofs
  generalize additively along independent column blocks.
  Theorem 5.5 remains the foundational identity and Remark 5.6's
  $\Theta(N)$ extrapolation for low-dimensional natural-data
  manifolds is an informal heuristic; rigorously establishing
  $\mathrm{TC}(D_N) = \Theta(N)$ for less combinatorially structured
  classes (natural images at a given resolution, continuous-state
  HMMs, non-stationary processes, factor graphs over expanders) is
  open and likely requires manifold-hypothesis-plus-regularity
  assumptions beyond the intrinsic-dimension framework.
\item
  Construction of efficient enumerative encoders for specific
  natural-data submanifolds. Theorem 5.5 establishes that the
  separation between factorized and joint coding equals
  $\mathrm{TC}(D_N)$, and Remark 5.6 predicts this is $\Theta(N)$
  for low-dimensional natural-data manifolds. Theorem 5.5 itself
  does not identify any particular submanifold or guarantee the
  existence of a polynomial-time membership oracle or rank coder
  for one; for an arbitrary finite distribution $D_N$ a brute-force
  oracle that enumerates the support exists trivially but is not
  efficient. Constructing a polynomial-time membership oracle and
  rank coder for the support of a specific data class --- say,
  database pages under a given schema, or valid x86-64 instruction
  streams --- is open as a constructive question for each class, and
  is not implied by Theorem 5.5.

  \emph{Reduction to membership-oracle construction (Cover 1973
  enumerative coding).} OP3$'$ has an INDIRECT closure via Cover's
  enumerative source coding (Cover 1973 \emph{IEEE Trans.\ Inf.\
  Theory} 19(1):73--77): for any data class $\mathcal{C} \subseteq
  \Sigma^N$ admitting a polynomial-time \emph{membership oracle}
  $O_{\mathcal{C}} : \Sigma^N \to \{0, 1\}$, an explicit
  polynomial-time \emph{rank coder} for the uniform distribution on
  $\mathcal{C}$ exists with time complexity
  $O(N \cdot \log|\Sigma| \cdot T_{\mathrm{oracle}})$, where
  $T_{\mathrm{oracle}}$ is the oracle's per-query cost. Encoder:
  compute the rank of input $x \in \mathcal{C}$ as
  $\mathrm{rank}(x) = |\{y \in \mathcal{C} : y < x\}|$ via
  Lehmer-code-style enumeration with the oracle; decoder: binary
  search for the $k$-th $\mathcal{C}$-member using the oracle.
  Hence OP3$'$ \emph{reduces} to constructing efficient membership
  oracles for natural-data classes. For specific classes the
  reduction yields immediate positive results:
  \emph{(i)~Multinomial sequences} (each symbol appears a prescribed
  count): membership oracle = symbol count comparison, $O(N)$ time,
  trivially polynomial.
  \emph{(ii)~Monotone sequences} ($x_1 \leq x_2 \leq \cdots \leq x_N$):
  membership oracle = scan + compare, $O(N)$.
  \emph{(iii)~JSON-schema-conforming database pages}: membership oracle
  = JSON parser + schema validator, $O(N)$ for fixed schema.
  \emph{(iv)~x86-64 instruction streams}: membership oracle =
  disassembler + validator, polynomial in $N$ for fixed ISA.
  Hence OP3$'$ is \emph{constructively closed} for any class with a
  poly-time membership oracle, with class-specific
  implementations following the Cover 1973 enumerative-coding
  framework. The remaining \emph{open} subcase is classes whose
  membership oracle is NP-hard to construct (e.g., satisfying
  assignments of a fixed SAT instance), which fall outside
  Theorem 5.5's reach for fundamental complexity-theoretic reasons.
\item
  Constant-gap inapproximability for $E_{12}$ in the
  Hamming-restricted geometry. The decision version of
  varying-parameter learned-code labeling is NP-hard (Theorem 4.3)
  via an exact YES/NO threshold reduction. Remark 4.3c establishes
  that the Wagner-Corneil reduction yields only an inverse-linear
  inapproximability gap of $1 + 1/(n-1)$, tight on the star family
  $K_{1, n-1}$ with $m = n - 2$, which vanishes as $n \to \infty$;
  no constant-gap reduction from any APX-hard source into the
  Hamming-realizable distance geometry is currently known. The Type-III-C side (Theorem 6.8) is
  APX-hard with explicit inherited constant
  $\gamma_{\mathrm{SGP}} = 8569/8568$ (Charikar et al.\ 2005,
  unbounded-alphabet variant); the analogous constant for $E_{12}$
  is open. (A sharper $\gamma_{\mathrm{SGP}}$ for SGP itself ---
  e.g., via direct reduction from MAX-3SAT-3 rather than via Vertex
  Cover --- is a separate open question in the SGP literature.)
  Concrete research paths for $E_{12}$ inapproximability include:
  (a)~adapting the Arora-Babai-Stern-Sweedyk (1997) Nearest
  Codeword Problem inapproximability via PCP-style reductions from
  Label Cover; (b)~Dumer-Micciancio-Sudan (2003) bounded-distance
  decoding reductions; (c)~Khot-style Closest Vector Problem
  reductions specialized to the $\mathbb{F}_2^{12}$ ambient
  space; (d)~adaptation of $0$-extension / metric labeling on
  Hamming metrics (which admits constant-gap APX-hardness via
  Karzanov 1998 SICOMP and subsequent gap-amplification by
  Schwartz \& Tur 2021/2023, arXiv:2104.11670), to handle the
  injective bijection requirement of $E_{12}$; (e)~quadratic
  assignment with hypercube Hamming distance matrices (Mittelmann \&
  Peng, recent SDP / heuristic literature), suitably restricted to
  yield constant-gap. A literature search through 2024 found no
  direct constant-gap inapproximability result for the exact $E_{12}$
  pairwise-Hamming objective (closest related: NP-hardness only for
  the border minimization problem of Li et al.\ 2017, APX-hardness
  for Hamming \emph{clustering} but not \emph{labeling} per
  Lossy Kernelization 2023). All five paths require non-trivial
  adaptation to the fixed source-alphabet (256-byte) constraint of
  the $E_{12}$ problem and remain research-scale open.

  \emph{PIH-under-ETH reduction with refined analysis (new direction,
  2026).} A multi-agent brainstorm during the present manuscript
  revision proposed reducing from the constant-gap parameterized
  CSP of Guruswami--Lin--Ren--Sun--Wu (arXiv:2311.16587, ``PIH under
  ETH'', 2023), which gives ETH-hardness with absolute soundness gap
  $\varepsilon_{\mathrm{PIH}} = 1/9600$ on a 2-CSP $G^* = (V^*, E^*,
  \Sigma^*, \{\Pi_e\})$ with $|V^*| = K = 2^{\Theta(\ell^2)}$ variables
  and alphabet $|\Sigma^*| = 2^{\Theta(n/\ell)}$, for free parameter
  $\ell$. The reduction sketch: $V_H = V^*$ (vertex per CSP variable),
  $E_H = E^*$ (edge per constraint), $F_{u, v} = 1/\deg_H(u)$ (the
  row-stochastic Markov kernel of $\Psi$, built-in by construction),
  codeword $f(v_i) = [\mathrm{identifier}_i, \mathrm{Enc}(\sigma_i)]$
  with $m_{\mathrm{id}} = \lceil \log_2 K \rceil$ identifier bits
  (fixed per vertex) and $m_{\mathrm{val}}$ value bits encoding the
  CSP value via some binary code $\mathrm{Enc} : \Sigma^* \to
  \{0, 1\}^{m_{\mathrm{val}}}$.

  \emph{Identifier-dilution obstacle (subtle and new).} Naive analysis
  of the proposed reduction overlooks a structural obstacle from the
  row-stochastic normalisation: with $\sum_e F_e = K$ (each row of $F$
  sums to $1$), the total identifier-Hamming cost is $K \cdot \bar{c}_{\mathrm{id}}$
  where $\bar{c}_{\mathrm{id}}$ is the vertex-weighted average per-edge
  identifier-Hamming distance. For random identifier assignment,
  $\bar{c}_{\mathrm{id}} \approx m_{\mathrm{id}}/2 = \Theta(\log K)$.
  An unsatisfied CSP constraint with per-bit code disagreement gives
  per-edge value cost $\geq 1$, so the multiplicative gap is
  \[
  \rho_{E_{12}} \;\geq\; 1 + \frac{\varepsilon_{\mathrm{PIH}}}{\bar{c}_{\mathrm{id}}}
  \;=\; 1 + \frac{\varepsilon_{\mathrm{PIH}}}{\Theta(\log K)} \;\to\; 1
  \]
  as $K \to \infty$, vanishing. The padding-bit angle alone does not
  preserve constant gap under row-stochastic $F$: the identifier cost
  dominates the per-edge satisfaction signal.

  \emph{Saving construction (constant relative-distance code on value
  bits).} The obstacle is overcome by encoding CSP values via a
  binary code $\mathrm{Enc}$ with constant rate ($m_{\mathrm{val}}
  \geq \Omega(\log |\Sigma^*|) = \Omega(n/\ell)$) and constant
  relative distance $\delta > 0$ (Justesen 1972 explicit construction
  attains both). Then unsatisfied constraints have value-Hamming
  $\geq \delta \cdot m_{\mathrm{val}}$ per edge, giving
  \[
  \rho_{E_{12}} \;\geq\; 1 + \frac{\varepsilon_{\mathrm{PIH}} \cdot \delta \cdot m_{\mathrm{val}}}{\bar{c}_{\mathrm{id}}}.
  \]
  For $\rho$ constant, need $m_{\mathrm{val}} \geq \Omega(\bar{c}_{\mathrm{id}})
  = \Omega(\log K) = \Omega(\ell^2)$. Combined with the rate constraint
  $m_{\mathrm{val}} \geq \Omega(n/\ell)$, both hold when $\ell^2 \geq
  \Omega(n/\ell)$, i.e., $\ell \geq \Omega(n^{1/3})$.

  \emph{Conditional closure (alphabet-balanced PIH at $\ell = \Theta(n^{1/3})$).}
  At $\ell = \Theta(n^{1/3})$, GLRSW Theorem 3.1 gives an alphabet-balanced
  instance with $K = 2^{\Theta(n^{2/3})}$ vertices and $|\Sigma^*|
  = 2^{\Theta(n^{2/3})}$, polynomially balanced. The above reduction
  then transfers $\varepsilon_{\mathrm{PIH}} = 1/9600$ to a constant
  $E_{12}$ gap of $1 + \Omega(\varepsilon_{\mathrm{PIH}} \cdot \delta)$.
  This gives \emph{conditional ETH-hard FPT inapproximability of
  $E_{12}$} (parameterised by $\ell = n^{1/3}$, no $2^{o(n)}$ algorithm
  exists for distinguishing $(1+\varepsilon')$-approximation, some
  constant $\varepsilon' > 0$), STRICTLY WEAKER than the NP-hard
  inapproximability claimed by Conjecture 4.3e but a concrete proof-path
  toward it. The remaining technical work: (i) verify that the code
  $\mathrm{Enc}$ can be \emph{constraint-aware} (i.e., the constraint
  $\Pi_e(\sigma_u, \sigma_v)$ is decoded as a Hamming-distance condition
  $d_H(\mathrm{Enc}(\sigma_u), \mathrm{Enc}(\sigma_v))$, which requires
  matching the binary code to the GLRSW constraint form
  (linear $+$ parallel QUADEQ)); (ii) extend from FPT/ETH-hardness
  to classical NP-hardness via gap-preserving polynomial-time reduction
  (open whether such a reduction exists from $\ell = n^{1/3}$
  PIH-under-ETH to a polynomial-size instance).

  This is a strictly more careful analysis than the optimistic
  earlier ``padding-bit absorbs packing noise'' framing: the absorption
  is \emph{additive but not multiplicative} under row-stochastic $F$.
  The brainstorm's original Angle 1 missed the dilution; a
  follow-up adversarial review (independent fresh-context agent with
  web-search 2020--2025 literature) identified the fix via good
  binary codes on value bits, with the explicit parameter constraint
  $\ell \geq n^{1/3}$ to balance identifier and value contributions.
  Full formalisation requires technical work on constraint-aware
  encoding and remains beyond the present paper.

  \emph{Two-sided status (epistemic gap).} As with OP2(a), a 2015--2024
  literature scan confirms that for $E_{12}$ \emph{neither} a
  polynomial-time constant-factor approximation \emph{nor} constant-gap
  APX-hardness is currently known. The closest published results give
  upper bounds via SDP / ILP relaxations on QAP with hypercube Hamming
  distance matrices (Mittelmann \& Peng), heuristics
  (Hamming-Mallows-style local search), and exact solutions for special
  graph families, but no constant-factor approximation guarantee for
  general $E_{12}$ instances. We note also a definitional point: the
  injective (non-collapsing) constraint on the labeling is essential
  --- without it, assigning all vertices to the same codeword gives
  objective $0$ trivially; injectivity (a one-to-one map of $n$
  vertices to distinct $\{0,1\}^m$ codewords) is what makes the
  problem combinatorially non-trivial. An early brainstorming session proposed two candidate structural obstacles ---
  (i) row-stochastic normalization of $F$ blocking ``isolated
  massive-weight gadget bottleneck'' constructions, and (ii) the
  injective-labeling ``packing noise'' / ``gadget leakage'' background
  cost --- and conjectured that closing Conjecture 4.3e might require a
  source problem APX-hard on row-stochastic instances \emph{specifically}.
  \emph{That framing is now superseded by the exact gap-locus analysis of
  Remark 4.3d:} row-stochasticity is \emph{not} the obstacle. The ratio
  $\rho(F)=\mathrm{OPT}/2w$ is invariant under global rescaling of $F$
  (replacing the weight $1/\Delta$ by $1$ multiplies both $\mathrm{OPT}$
  and $2w$ by $\Delta$ and leaves $\rho$ unchanged; Remark 4.3d(ii),
  numerically witnessed in
  \texttt{scripts/verify/remark\_4\_3d\_e12\_constant\_gap\_dilution.py}).
  Unlike the affine continuous-cost obstacles elsewhere in the paper,
  normalization does \emph{not} dilute the gap here --- the cost is a
  quadratic assignment-type functional and the per-instance forced-excess
  fraction can be a genuine constant. The real barrier (Remark 4.3d(iii)--(iv))
  is the absence of a \emph{gap-amplified hypercube-embedding source}: a
  family that is simultaneously dilation-$1$ embeddable in the YES case
  (so a YES instance exists at all) and robustly non-embeddable --- a
  constant fraction of weighted mass stretched beyond distance one --- in
  the NO case. The two requirements collide with two literature facts:
  dilation-$1$ embeddable graphs are exactly the partial cubes,
  recognizable in polynomial time (Djokovi\'c 1973; Winkler 1984;
  Eppstein 2008), so embeddability alone is not NP-hard; and dense graphs
  that would force constant stretch (e.g.\ $K_{2,3}$) are never
  hypercube subgraphs, so they possess no YES baseline. $E_{12}$ therefore
  occupies the same theory-gap regime as OP2(a): neither algorithmic
  constant-factor upper bound nor inapproximability lower bound is
  known, but the missing ingredient is now precisely localized to a
  robust-NO-promise hypercube-embedding source (Remark 4.3d), not to
  row-stochasticity.

  \emph{Conjecture 4.3e (Varying-parameter $E_{12}$ constant-gap
  inapproximability).} We conjecture that there exists a constant
  $\varepsilon > 0$ such that the varying-parameter $E_{12}$
  learned-code labeling problem is NP-hard to approximate within
  factor $1 + \varepsilon$. Motivation: Remark 4.3c proves
  NP-hardness with vanishing gap $1 + 1/(n-1)$; analogous
  combinatorial problems (NCP per Arora-Babai-Stern-Sweedyk 1997,
  CVP per Dinur-Kindler-Raz-Safra 2003) admit constant-gap
  inapproximability via PCP-style gap amplification, suggesting
  that a similar argument may apply to $E_{12}$ once a suitable
  source problem with the right structural properties is identified.
  The conjecture is restricted to the \emph{varying-parameter} case
  (both $n$ and $m$ part of the input); the fixed practical instance
  $n = 256$, $m = 12$ is constant-size and not asymptotically
  NP-hard. A PTAS for varying-parameter $E_{12}$ (polynomial-time
  $(1 + \varepsilon)$-approximation for every $\varepsilon > 0$)
  would refute this conjecture.
\item
  Constant-gap inapproximability for the \emph{fixed-byte-alphabet}
  variant of restricted Type-III-C ($|\Sigma| = 256$). The
  unbounded-alphabet variant has $\gamma_{\mathrm{SGP}} = 8569/8568$
  inherited from Charikar et al.\ 2005. Encoding the
  $|V|$-vertex graph terminals as length-$\lceil \log_{256}|V| \rceil$
  byte strings adds a $\Theta(\log |V|)$ multiplicative overhead in
  the grammar encoding, which can dilute the constant gap toward
  $1$ as $|V| \to \infty$ unless a separate gap-preserving
  construction is given. Establishing (or refuting) a constant
  inapproximability gap for the fixed-byte case is open.
\item
  Practical compatible Adam variants with full convergence
  analysis. Theorem 10.2 covers any optimizer satisfying Definition
  10.1. Integer-arithmetic Adam variants exist (Banner et al. 2018)
  but their convergence behavior relative to floating-point Adam
  under the specific compatibility constraints (deterministic
  integer reciprocal-square-root, fixed update ordering) is
  incompletely characterized. A specific compatible Adam variant
  with a convergence guarantee matching the floating-point
  baseline, on standard problems, is open.
\item
  Sharpening the constant $\gamma_{\rm LRU}$ in Theorem 7.26's
  Lezaud Chernoff bound. The crude Dobrushin-coupling lower bound
  $\gamma_{\rm LRU} \geq \pi_{\rm min}$ gives a positive but
  conservative spectral gap; tighter explicit bounds via the
  Knudsen-Tsitsiklis 1990 stack-distance representation, or via
  Aldous-Diaconis 1987 coupling for permutation-like Markov chains
  on LRU stack states, would tighten the constant. Theorem 7.26's
  qualitative $\sqrt{Q}$-Hoeffding scaling holds for any positive
  $\gamma_{\rm LRU}$; only the multiplicative constant is at stake.
  The non-IRM (workload-dependent) case --- where LRU concentration
  may genuinely fail, as shown by the $A, B, A, B, \ldots$
  counterexample (2026-05) against $C=1$ --- remains
  out-of-scope of this paper's IRM hypothesis and may require
  workload-specific concentration analysis.
\item
  Sharpening the concentration bound at the cache-state level
  (e.g.\ improving constants in T7.25's $K_\phi$ via specific
  variance-proxy estimation for natural text) remains a refinement
  question; the qualitative scaling is fixed.
\item
  General-scheme space-time converse for random access
  (\emph{Theorem 7.27 sub-block restriction}) --- \textbf{resolved
  negatively for finite-order-Markov sources; partially open for neural
  predictors} (Theorem 7.27$'$, Remark 7.27d, Corollary 7.27c).
  Theorem 7.27 of [RNR-II] establishes a tight product bound $C(K) \cdot S(K) \geq
  (1 - o(1)) N \cdot \mathrm{cost}(M) \cdot \log N$ for any
  sub-block-based random-access scheme satisfying Definition 10.3. The
  question whether this product holds for the \emph{general}
  random-access class --- arbitrary decoders that need not partition the
  archive into sub-blocks --- is now settled in three honest pieces.
  \emph{(1) Refutation for Markov-expressible sources (Theorem 7.27$'$,
  proved).} The universal product is \emph{false}: for stationary
  order-$k$ Markov sources, the Ferragina--Venturini (SODA 2007;
  \emph{Theoret.\ Comput.\ Sci.}\ 372(1):115) entropy-bounded storage
  structure stores the realised string in $N h(X) + o(N)$ bits (Shannon
  -near, by the Markov AEP $H_k \to h(X)$) and answers any single-byte
  query in $O(1)$ time with \emph{zero} query-time predictor
  evaluations, so $C \cdot S = O(N) = o(N\,\mathrm{cost}(M)\log N)$. The
  sub-block restriction of Theorem 7.27 is thus essential, not a
  technical artefact: the product is a property of the
  predictor-sequential decode architecture. \emph{(2) Positive maximal
  class (Corollary 7.27c, proved).} The $\mathrm{cost}(M)$-bearing
  product $C_{\mathrm{work}} \cdot S \geq (1-\epsilon)(1-o(1)) N \cdot
  \mathrm{cost}(M) \cdot \log N$ holds exactly on the
  \emph{predictor-faithful} class (Definition 7.27b: dense predictor
  recomputation is forced, the decoder storing $M$-outputs/states at no
  more than an $\epsilon$ fraction of any decode window). This class
  \emph{strictly contains} both Definition 10.3 sub-block schemes
  ($\epsilon = 0$) and Theorem 10.4 state-snapshot schemes ($\epsilon =
  O(1/K)$, the natural non-sub-block witness), and \emph{excludes} the
  Ferragina--Venturini escape ($\epsilon = 1$). \emph{(3) The residual
  open problem (genuinely open, now correctly delimited by Corollary
  7.27e).} The refutation in piece (1) was stated for finite-order Markov
  sources. Corollary 7.27e computes the Ferragina--Venturini
  redundancy-vs-order crossover $k_{\mathrm{cross}}(N) = \Theta(\log_2 N /
  \log_2^2\sigma)$ and thereby \emph{corrects} the over-broad reading of
  Remark 7.27d(ii): for \emph{every fixed} stationary source --- including
  genuinely neural ones with infinite Markov order, and even divergent
  Crutchfield--Feldman excess entropy --- the FV escape \emph{does not}
  close (the optimal build order $k^\star(N) = o(k_{\mathrm{cross}}(N))$
  makes both the model gain $\Delta_{k^\star}$ and the FV redundancy
  $\rho_{k^\star}(N)$ vanish), so the universal converse is false against
  every fixed stationary source, not merely finite-order Markov. The FV
  escape closes \emph{only} for triangular-array sources whose effective
  Markov horizon $k_\epsilon(N)$ outruns the crossover, i.e.\
  $k_\epsilon(N) = \omega(\log_2 N / \log_2^2\sigma)$ for some $\epsilon
  > 0$ (Corollary 7.27e(b)); at deployment scale on byte data the escape
  is moreover practically closed, since $\rho_1(N) > 1$ bit/symbol for all
  $N \lesssim 2^{188} \approx 10^{57}$ (Corollary 7.27e(c)). The residual
  genuinely-open
  question is thus pinned to this growing-horizon regime: for an
  escape-\emph{closed} source (horizon $\omega(\log_2 N)$, e.g.\ a model
  whose context window grows with the corpus), is there \emph{some other}
  (non-FV) model-free succinct random-access structure attaining the
  Theorem 7.2 Shannon-near rate with $o(\sqrt{N})$ query time, or does a
  converse force $\Omega(\sqrt N)$ query there? Corollary 7.27f advances
  this on both sides without closing it. \emph{Positively (conditional),}
  it extends the converse from Ferragina--Venturini alone to the
  \textbf{entire known block-context-table family} of $O(1)$-access
  high-order-entropy structures (Gonz\'alez--Navarro CPM 2006;
  Grossi--Orlandi--Raman; Belazzougui--Navarro \emph{ACM Trans.\ Alg.}\
  11(4) 2015): on a growing-horizon source every such structure pays
  $\Omega(1)$ bits/symbol over $h(X)$, via a model-order/redundancy
  incompatibility (Shannon-near rate forces build order $k = \omega(L)$,
  but the redundancy budget caps $k = o(L)$, $L = \log_\sigma N$).
  \emph{Negatively (unconditional),} a fully architecture-free $\Omega(\sqrt
  N)$ converse is \textbf{unprovable}: pure single-symbol access is
  $O(1)$-achievable at the entropy bound (Belazzougui--Navarro), the
  generic cell-probe access floor is only $\Omega(\log N/\log\log N)$
  (P\u{a}tra\c{s}cu--Demaine 2006), and the only $\sqrt{\cdot}$-type access
  bounds are conditional on a specific representation (BWT/MTF,
  Sinha--Weinstein STOC 2019) or on exponential compressibility
  (Chen--Verbin--Yu ICALP 2012), neither holding in the
  constant-factor-compressible RNR regime. The irreducible open core is
  therefore: \emph{does a \textbf{non-block-context-table} model-free
  succinct structure attain Shannon-near rate with $o(\sqrt N)$ access on
  growing-horizon sources?} Corollary 7.27e removed the false premise
  (``neural $\Rightarrow$ FV closes''); Corollary 7.27f delimits the
  residual to a single architectural question and shows the $\sqrt N$
  form cannot be forced unconditionally.
\end{itemize}

\textbf{13.3 Out of scope: engineering deferrals}

The following items are deferred to engineering or empirical work. They
are not theorem-shaped problems and cannot be resolved within the scope
of a theoretical paper. They are listed here only so that a reader can
see what additional work is required for a deployable system.

\begin{itemize}
\item
  Reference implementation conformance. Theorems 10.1 and 10.2 specify
  sufficient conditions for bit-exact reproduction under fixed and
  adaptive predictors; Theorem 10.6 lifts Theorem 10.1 to explicit
  per-target conformance conditions (a)--(g); Theorems 7.16--7.18
  additionally require that sub-block-parallel encoding produce
  bit-identical archives regardless of thread count or scheduling order.
  A reference implementation meeting all of these, plus a conformance
  test suite that verifies cross-platform byte-identity on a
  representative corpus, is required for deployable archives. This is
  software engineering work, not mathematics. Out of scope for this
  paper. The named artifact is the companion engineering specification
  (\emph{RNR Coding --- Reference Implementation Conformance
  Specification}), whose §3 (universal inference requirements R-3.1
  through R-3.8) and §5 (integer-arithmetic specification) encode
  Theorem 10.1's conditions (i)--(iv), §4 the per-target conformance of
  Theorem 10.6, §7 the adaptive-optimizer conformance of Theorem 10.2,
  and §12 the conformance test suite (cross-target byte-identity
  acceptance criteria R-12.2.1, R-12.2.3, R-12.2.5); its §14 gives the
  full theorem-to-requirement traceability matrix. The end-to-end
  cross-platform demonstration is the empirical work specified as
  Hypotheses H8 (fixed-predictor byte-identity across six platforms),
  H9 (adaptive-predictor byte-identity), and H14 (per-hardware-target
  predictor byte-identity) of the companion experimental-design
  document. ``Satisfiable end-to-end'' is therefore conditional on a
  target being \emph{certified} conforming against that specification;
  this paper proves the sufficient condition, not that any particular
  target has been certified.
\item
  Empirical validation. The quantitative predictions in §11 are
  conjectures based on the structural advantages of each Type and on
  coding rates reported for comparable systems. Validation requires
  running the constructions of §§4--7 against NNCP, PixelCNN++, LZMA,
  zstd, and bits-back baselines on a representative corpus, including
  the enwik9 prediction of §11.7 (Hypothesis~H15). This is
  experimental work, not mathematics. Out of scope for this paper; the
  natural artifact is a separate experimental paper.
\item
  Intrinsic-dimension estimation for specific data classes. Remark 5.6
  connects Theorem 5.5\textquotesingle s separation to the manifold
  hypothesis. Computing or estimating the intrinsic dimension d of
  specific natural-data classes (database pages of a given schema,
  telemetry under a given protocol, image data of a given modality)
  requires applied data analysis using established estimators
  (Levina-Bickel maximum likelihood, correlation dimension, persistent
  homology). This is empirical work, not mathematics. Out of scope for
  this paper.
\item
  Privacy and information leakage from shipped predictors. A model M
  trained on private data and shipped as part of an archive may leak
  training-set information through memorization or membership inference.
  The trade-off between L(M) reduction (sharing more model) and
  information leakage requires a threat model, an adversary capability
  specification, and an empirical privacy audit. This is a security and
  policy question, not a theorem; out of scope for this paper.
\item
  Cache architecture and workload characterisation. Theorems 7.20 and
  7.24 provide formulas for cache hit rates and decoder-latency tail
  bounds under Zipf and IRM workloads, respectively. Selecting an
  appropriate cache size~$C$, eviction policy, and per-workload
  parameters ($\alpha$, request distribution) for a specific deployment
  requires empirical workload profiling and system-level benchmarking.
  This is systems engineering work, not mathematics. Out of scope for
  this paper; the natural artifact is a workload-aware cache
  configuration guide accompanying a reference decoder.
\item
  Per-sub-block MAC key management. Theorem 7.23 guarantees tamper
  detection given a PRF-based MAC with a secret key. The key
  lifecycle---generation, storage, rotation, revocation, and
  multi-archive key hierarchy---requires a threat model and a
  key-management specification aligned with standard practice
  (NIST SP 800-57). This is security engineering work, not mathematics.
  Out of scope for this paper; the natural artifact is a key-management
  policy document.
\item
  Parallel and quantised encoder deployment. Theorems 7.16--7.18
  establish sub-block-parallel encoding with linear speedup, and
  Theorems 7.10--7.12 bound the coding-efficiency loss from INT8 and
  mixed-precision quantisation. Achieving the theoretical speedups on
  specific hardware (GPU warp scheduling, NUMA-aware memory allocation,
  quantisation-aware training pipelines) requires systems-level
  optimisation and benchmarking. This is performance engineering work,
  not mathematics. Out of scope for this paper; the natural artifact is
  a hardware-specific deployment guide.
\end{itemize}

\textbf{13.4 Concrete attack vectors for remaining open problems}

We list specific approaches for each genuinely open problem after the
present version's contributions. Where standard reductions have
demonstrably failed (adversarial validation 2026-05: counterexamples and structural critiques), we explicitly
identify the failed pattern and propose alternative attack vectors
that have not been tried. The intent is to prevent re-attempts on
known-failed paths and to direct effort toward genuinely promising
directions.

\textbf{13.4.1 OP2(a) (Unrestricted bit-cost SGP APX-hardness).}

\emph{What has been tried (failed paths):}
\begin{itemize}
\item Direct lift of Berman-Karpinski $\to$ Charikar reduction at
$K = O(|V|)$ to bit-cost via the $K$-uniformity assumption ---
blocked by the K-optimisation obstacle (§13.2, K-escape phenomenon):
at the bit-cost-optimal $K^* = \Theta(\sqrt{N/\log V})$ regime, the
Berman-Karpinski gadget does not preserve the gap.
\item Restriction to Charikar normal form (Lemma 6.8b) --- works but
gives only class-restricted APX-hardness, not unrestricted.
\item Continuous-bit-cost relaxation (Lemma 6.8e) --- works
unconditionally; the residual difficulty is specifically the
ceiling discretisation $\lceil \cdot \rceil$ in the integer-bit
objective.
\item K-monotone constraint (Observation 6.8f) --- works but excludes
the K-escape obstacle by fiat; tautological for the residual question.
\end{itemize}

\emph{Attack vectors not yet tried:}
\begin{enumerate}
\item \emph{Modify the Berman-Karpinski reduction at $K = \sqrt{N}$
regime.} The current BK construction uses $|V|$ vertices with $|E|
= O(|V|)$ edges, giving $K_{\rm gadget} = \Theta(|V|)$. A modified
construction with $|V| = N^{1/2}$ vertices and $|E| = N$ edges (i.e.,
denser graph) would shift the BK $K^*$ to $\sqrt{N}$, matching the
bit-cost-optimal regime. Open question: does VC-3 retain its
$1 + 1/144$ inapproximability gap on denser graphs (the
``density-shifted Berman-Karpinski'' question)?

\emph{Status update 2026-05-27 (failed path, literature blocker).}
A systematic literature search (Karpinski--Zelikovsky 1997, Clementi--
Trevisan 1999, Eremeev 1999, Imamura--Iwama 2005, Karpinski--Schmied
2010, Austrin--Khot--Safra 2011) settles this attack vector
\emph{negatively}. The target regime $|V| = N^{1/2}$, $|E| = N$
corresponds to graphs of density $|E|/|V|^2 = 1$, i.e., asymptotically
complete graphs, where VC is trivial (minimum cover $= |V| - 1$,
exactly solvable in polynomial time --- no constant gap). Any
intermediate regime with $|V| = N^{1-\alpha}$, $|E| = N$ for
$\alpha \in (0, 1/2)$ yields average degree $\Theta(N^{\alpha})$ and
density $|E|/|V|^2 = \Theta(N^{2\alpha - 1})$; for $\alpha > 1/2$
this is subdense in Karpinski--Schmied's sense ($|E| =
\Omega(|V|^2/\psi(|V|))$ with $\psi(|V|) = \omega(1)$), which
\emph{admits a PTAS for VC} (Karpinski \& Schmied 2010,
arXiv:1011.0078, ``Approximating Subdense Instances of Covering
Problems''). For $\alpha \leq 1/2$, average degree is
$O(\sqrt{|V|}) = O(|V|^{1/2})$ and $K_{\rm Charikar} \sim |V|
= N^{1-\alpha}$, which remains $\omega(\sqrt{N})$ --- the
$K$-shift does \emph{not} occur. Explicit hardness ratios in the
dense regime degrade smoothly to 1: Eremeev's
$(7+\delta^*)/(6+2\delta^*)$ NP-hardness for strongly
$\delta^*$-dense graphs gives ratio $7/6$ at $\delta^* \to 0$
(still sparse-like) but collapses to $1$ at $\delta^* = 1$
(complete graphs); Karpinski--Zelikovsky's algorithmic upper bound
$2/(1+\varepsilon)$ for everywhere $\varepsilon$-dense graphs
likewise collapses to $1$ at $\varepsilon = 1$. Conclusion:
\emph{the BK $1 + 1/144$ gap is a specific feature of the
bounded-degree (sparse) regime}, derived from the 3-regular
amplifier construction (Berman--Karpinski 1999 ECCC TR98-065); it
does not extend to denser graphs while simultaneously shifting
$K_{\rm Charikar}$ to $\sqrt{N}$. The density regime that
preserves a constant VC gap (bounded degree, $|E| = O(|V|)$) is
incompatible with the $K$-shift requirement
($K_{\rm Charikar} \sim \sqrt{N}$ needs $|V| \sim \sqrt{N}$, which
forces $|E| = \omega(|V|)$). This attack vector is therefore
\emph{closed negatively}: density-shifted BK cannot work because
the two requirements (constant VC gap on $H$; $K_{\rm Charikar}(H)
\sim \sqrt{N}$ in Charikar's encoding) are mutually exclusive
in the BK $\to$ Charikar framework. The Austrin--Khot--Safra 2011
result ($2 - (2+o_d(1))\log\log d/\log d$ UG-hardness for VC on
degree-$d$ graphs) gives a larger gap as $d \to \infty$, but is
(a) Unique-Games-conditional rather than NP-hard, and (b) the
underlying graphs are still sparse ($|E| = O(d|V|/2)$), so
$K_{\rm Charikar} \sim |V|$ remains far above $\sqrt{N}$ unless
one uses $|V| \sim \sqrt{N}$, which then requires $d \sim
\sqrt{N}$ on $|V| = \sqrt{N}$ vertices --- producing complete
graphs again. The Dinur--Safra 2005 result gives a stronger
NP-hardness factor of $\approx 1.36$ on graphs of ``sufficiently
large vertex degree'', but the Dinur--Safra hardness graphs
(FGLSS-style from PCP composition with intersection-graph
structure) do not admit a direct Charikar encoding preserving the
$K$-shift; this remains a possible direction under item 2 if the
FGLSS structural overhead can be absorbed into the SGP encoding,
but is not addressed by the present density-shifted attack
vector. Alternative attack vectors (items 2, 3, 4 below) remain
open.
\item \emph{Use a different source-of-hardness problem with
$K$-uniformity built in.} Candidates: SAT (Karp 1972 NP-hardness),
3-COL (Garey-Johnson 1979 APX-hardness via Khot-Saket 2014 with UGC),
or LABEL COVER (Arora-Lund-Motwani-Sudan-Szegedy 1998 PCP composition).
The challenge is finding a problem whose hardness gap survives the
log-factor dilution at $K = \sqrt{N}$.
\item \emph{Lower-bound the bit-cost gap via an exponential
``$K$-escape budget'' argument.} Show that any grammar with $K \ll
K_{\rm YES}$ for a BK NO instance must pay at least $\Omega(L_{\rm sym})$
extra symbols, dominating the log-factor savings. This would close
the K-escape obstacle directly.
\item \emph{Cryptographic conditional hardness via OWFs.} Liu-Pass 2020
shows OWFs $\Leftrightarrow$ samplable-distribution MKtP hardness.
A reduction MKtP $\to$ bit-cost SGP on a samplable input distribution
would give conditional average-case APX-hardness under OWF (weaker
than worst-case but tractable).
\end{enumerate}

\emph{Status update 2026-05-28 (Dinur--Safra FGLSS escape hatch closed
negatively).} The residual Dinur--Safra direction flagged in the
2026-05-27 update under item~1 has been investigated and is closed
negatively. Verified construction parameters (Dinur \& Safra 2005,
\emph{Annals of Math.}\ 162(1):439--485, Theorem~1.1 and
Definitions~1.2, 2.3): the hard-instance graphs are nonintersection-based
co-partite blow-ups with $p_{\max} = (3 - \sqrt{5})/2 \approx 0.382$,
optimal vertex-cover fractions $\alpha_{\mathrm{DS}} = 1 - p_{\max}
\approx 0.618$ (YES) and $\beta_{\mathrm{DS}} = 1 - 4 p_{\max}^3 +
3 p_{\max}^4 \approx 0.841$ (NO), giving the published ratio
$\beta_{\mathrm{DS}}/\alpha_{\mathrm{DS}} = 10\sqrt{5} - 21 \approx
1.36068$. Critically, the authors state explicitly (\S1.5, p.~449)
that the result holds ``for graphs of bounded degree,\ldots\ simply
because the graph resulting from our reduction is of bounded
degree.'' However, ``bounded'' here means \emph{constant in the input
size $|V|$} via the parameter set $(h_0, h_1, l_T, h, l)$ of
Definition~2.3, which in turn depend on Friedgut's $\Gamma(p,\delta,k)$
function and the Sunflower-Lemma constant $\Gamma^*(h_1, h_s)$. For
DS's stated parameter ranges ($\varepsilon, \gamma$ small), these are
super-polynomial in $1/\varepsilon$ (Friedgut is polynomial in
$1/\delta$ and $k$ with substantial multiplicative constants; the
Sunflower constant is at least $h_s^{1/h_s} \cdot h_1!$, which is
super-exponential in $h_1$), so $d_{\mathrm{DS}}$ is a fixed but
astronomical constant. Applying Charikar's encoding
to a DS-hard graph $H_{\mathrm{DS}}$ with $|E|
= (d_{\mathrm{DS}}/2)|V|$ gives
\[
  L_{\mathrm{sym}}^*(s_{\mathrm{DS,YES}}) = |V| \cdot (15 +
    1.5\,d_{\mathrm{DS}} + \alpha_{\mathrm{DS}}), \qquad
  L_{\mathrm{sym}}^*(s_{\mathrm{DS,NO}}) = |V| \cdot (15 +
    1.5\,d_{\mathrm{DS}} + \beta_{\mathrm{DS}}),
\]
so the encoding-preserved symbol-count ratio is
\[
  \rho_{\mathrm{DS}} \;=\; \frac{15 + 1.5\,d_{\mathrm{DS}} +
    \beta_{\mathrm{DS}}}{15 + 1.5\,d_{\mathrm{DS}} + \alpha_{\mathrm{DS}}}
  \;=\; 1 + \frac{\beta_{\mathrm{DS}} - \alpha_{\mathrm{DS}}}{15.618 +
    1.5\,d_{\mathrm{DS}}} \;\approx\; 1 + \frac{0.2229}{15.618 +
    1.5\,d_{\mathrm{DS}}}.
\]
The original graph-VC gap $1.36$ shrinks under Charikar's encoding
overhead $15|V| + 3|E|$ because (i)~the cover-size signal
$\alpha_{\mathrm{DS}}|V|$ is dominated by the structural and edge
terms, and (ii)~the edge term $1.5\,d_{\mathrm{DS}}|V|$ scales with
the high DS degree, further amplifying the dilution. Numerical
comparison against Berman--Karpinski's preserved gap $8569/8568
\approx 1.000117$: for $d_{\mathrm{DS}} = 3$ (impossible for DS but
used as a sanity check) $\rho_{\mathrm{DS}} \approx 1.0111$; for
$d_{\mathrm{DS}} = 100$, $\rho_{\mathrm{DS}} \approx 1.00135$; the
crossover where DS encoding gap equals BK's is at
$d_{\mathrm{DS}}^{\mathrm{cross}} = (0.2229 \cdot 8568 -
15.618)/1.5 \approx 1263$; for $d_{\mathrm{DS}} >
d_{\mathrm{DS}}^{\mathrm{cross}}$, DS gives a \emph{strictly weaker}
bit-cost gap than BK. Any reasonable estimate of $d_{\mathrm{DS}}$
from DS's proof machinery yields $d_{\mathrm{DS}} \gg 1263$, so DS
provides no improvement over BK at the encoded bit-cost level. The
$K$-optimisation obstacle (\S13.2) applies symmetrically: at the
bit-cost-optimal regime $K^* = \Theta(\sqrt{N})$, the encoded ratio
$\rho_{\mathrm{DS}}$ established at $K = K_{\mathrm{Charikar}} =
\Theta(|V|)$ is not provably preserved to $K^*$, and DS's larger
graph-VC ratio does not help here because the dilution occurs at
the \emph{encoding step}, before any $K$-optimisation is considered.
The structural blocker is identical in form to the density-shifted
BK closure of 2026-05-27: any NP-hard VC factor preserved by graph
blow-up loses asymptotically to the Charikar encoding overhead when
the underlying graph has $|E| = \omega(|V|)$. The escape hatch via
FGLSS / non-intersection structure does \emph{not} survive Charikar's
grammar encoding. \emph{Quantitative summary}: the maximum
encoding-preserved gap across known unconditional NP-hard VC
inapproximability results --- BK $\to$ Charikar (sparse, $1 +
1/8568$), DS $\to$ Charikar (dense, $\leq 1 + 0.149/d_{\mathrm{DS}}$),
H{\aa}stad $7/6$ on $d = O(1)$ via Charikar
($\approx 1 + 1/(8568 \cdot 6) \approx 1.00002$) --- is bounded by
$O(1/d_{\mathrm{src}})$ where $d_{\mathrm{src}}$ is the effective
degree of the underlying VC-hard graph family. Sparse families
($d = 3$) maximise this ratio; dense ones destroy it. The
Dinur--Safra direction is therefore closed negatively (verification
script: \texttt{scripts/verify/op2a\_dinur\_safra\_attack\_check.py};
DS constants and ratio reproduced to $10^{-10}$, BK ratio matched
exactly to $8569/8568$, crossover and dilution at each
$d_{\mathrm{DS}}$ regime reported). Attack vectors 2, 3, 4 of
\S13.4.1 remain open as documented; the unified obstacle pattern
they share is that any reduction $\mathcal{P} \to$ bit-cost SGP via
Charikar's encoding incurs the $15|V| + 3|E|$ structural overhead,
which dilutes graph-VC gaps in proportion to the source instance's
edge density. A $K$-uniform hardness source (item~2) would need to
either bypass Charikar's encoding entirely or supply a hardness gap
whose preserved form is not multiplicatively diluted by the encoding
(e.g., via gap amplification on the source string rather than on
the underlying graph).

\emph{Status update 2026-05-29 (attack vector~3, ``$K$-escape budget'',
closed negatively; with a new alphabet-floor sharpening that
isolates the residual obstacle as the ceiling discontinuity).} Item~3
proposed proving an exponential $K$-escape \emph{budget}: that any
grammar $G$ for a Berman--Karpinski NO source $s$ with $K(G) \leq
K_{\rm YES}/2$ must pay $\Omega(L_{\rm sym})$ extra symbols, the budget
dominating the log-factor saving and thereby closing OP2(a). We
attempted the budget lower bound and find it \emph{does not dominate},
for three independent reasons; the attack vector is closed negatively.
We record a positive by-product (the alphabet-floor observation)
that reduces the residual obstacle to exactly the ceiling-jump
discontinuity already flagged for Conjecture 6.8d.

\emph{(a)~The budget is real but additive $O(N)$, not a
gap-protector.} The Charikar source $s$ has length $|s| = \Theta(N)$
with $N := |V|$, alphabet $|\Sigma| = \Theta(N)$ (one terminal per
vertex), and Charikar-optimal grammar $L^*_{\rm sym}(s) = 15|V| +
3|E| + \mathrm{minVC}(H) = \Theta(N)$ realised at $K^*_{\rm sym} =
\Theta(N)$. Because every vertex name $v_i$ is a \emph{distinct}
terminal and each occurs only $\mathrm{occ}(v_i) = 3 + \deg(v_i) =
O(1)$ times (max-degree-3), the marginal value of a vertex
nonterminal $A_i \to \#v_i$ is $\mathrm{occ}(v_i) - 2 = 1 +
\deg(v_i) = O(1)$ symbols. Hence dropping $\Theta(N)$ nonterminals to
reach $K \leq K_{\rm YES}/2$ inflates $L_{\rm sym}$ by only
$\Theta(N) \cdot O(1) = O(N)$ symbols --- a constant-\emph{factor}
increase, not a super-constant blowup. This reproduces and rigorously
confirms the prior finding $L_{\rm sym}(K = N/2) = O(N)$ and the
``Pareto-frontier-flat'' analysis (the achievable $L_{\rm
sym}(K)$ frontier has slope $O(1)$ in the partial-cover region): no
exponential or super-linear budget exists. The required theorem
``$L_{\rm sym}(G) \geq (1 + \Omega(1)) L^*_{\rm sym}(G_{\rm NO})$
\emph{forced} by $K \leq K_{\rm YES}/2$'' is \emph{false} --- the
infimum of $L_{\rm sym}$ over $K \leq K_{\rm YES}/2$ stays within a
$1 + o(1)$ factor of $L^*_{\rm sym}$ because $K_{\rm YES}/2 =
\Theta(N) = \Theta(K^*_{\rm sym})$ already lies in the flat region.

\emph{(b)~Vector~3 targets the wrong $K$ regime; the relevant escape
is deep and symmetric.} At $K = K_{\rm YES}/2 = \Theta(N)$ the
log-factor saving is negligible: with $|\Sigma| = \Theta(N)$ and
$K_{\rm YES} = \Theta(N)$, both $\log_2(|\Sigma| + K_{\rm YES})$ and
$\log_2(|\Sigma| + K_{\rm YES}/2)$ equal $\log_2 N + O(1)$, so the
multiplicative log saving from halving $K$ is $1 + O(1/\log N) \to 1$.
A budget is therefore not even \emph{needed} at $K_{\rm YES}/2$:
escape to that regime yields essentially no bit-cost benefit, so the
theorem of item~3 (even if true) would not close OP2(a). The
bit-cost-relevant escape is the \emph{deep} regime $K = o(N)$ (down
to $O(1)$), where the log-factor saving is largest --- but even there
it is only $\log_2(|\Sigma| + K_{\rm YES})/\log_2|\Sigma| = 1 +
O(1/\log N)$, since $|\Sigma| = \Theta(N)$ floors $\log_2(|\Sigma| +
K)$ at $\log_2 N - O(1)$ even at $K = O(1)$ (see (c)). Moreover the deep
escape acts \emph{symmetrically} on the YES and NO sources: they
share the identical vertex-block and edge-gadget structure (differing
only in the $\mathrm{minVC}$ cover term), so any $L_{\rm sym}$
inflation from deep escape raises both bit-costs and does \emph{not}
differentially protect the $\mathrm{minVC}$-driven $\gamma_{\rm SGP}$
gap. Whether the $\mathrm{minVC}$ signal even \emph{persists} in
$L_{\rm sym}(K)$ at deep $K$ (as opposed to appearing only at
$K = \Theta(N)$) is precisely the unresolved Pareto-frontier question;
the budget argument does not settle it in either direction.

\emph{(c)~Alphabet-floor sharpening, and why it fails by exactly the
ceiling discontinuity.} A by-product of (a)--(b) is a cleaner
reduction of the residual obstacle. Since $|\Sigma| = \Theta(N)$
\emph{floors} the per-symbol cost at $\log_2|\Sigma| = \log_2 N -
O(1)$ regardless of $K$, and $K \leq K^*_{\rm sym} = \Theta(N)$
\emph{caps} it at $\log_2 N + O(1)$, the continuous log-factor is
pinned to $\log_2 N \cdot (1 \pm O(1/\log N))$ for \emph{every}
feasible $K$. Combined with $\min_K L_{\rm sym}^{\rm branch}(K) =
L^*_{\rm sym}({\rm branch})$ (the escape never reduces the
symbol-count optimum below the Charikar value, and the $O(1/\log N)$
log saving cannot pay for any constant-factor $L_{\rm sym}$
increase), this gives, for the \emph{continuous} objective,
$L^{\rm cont}\text{-OPT}({\rm branch}) = L^*_{\rm sym}({\rm branch})
\cdot \log_2 N \cdot (1 \pm O(1/\log N))$ and hence
gap${} = \gamma_{\rm SGP} \cdot (1 \pm O(1/\log N))$ --- an
independent re-derivation of Lemma 6.8e via the alphabet floor rather
than the first-order $\log b / \log a$ expansion. The discrete
objective, however, is \emph{not} closed by this: the ceiling
$\lceil \cdot \rceil$ contributes a full additive bit, i.e.\ a
\emph{multiplicative} $1 \pm O(1/\log N)$ perturbation, and
$\gamma_{\rm SGP} - 1 = 1/8568 < 1/\log_2 N$ for \emph{all} $N <
2^{8568}$ (an astronomical threshold). Thus a single ceiling jump ---
the NO optimum crossing a power-of-two boundary of $|\Sigma| + K$
below the YES optimum --- erases the entire $\gamma_{\rm SGP}$ gap at
those instance sizes, and the reduction is not gap-preserving across
all $N$. This is exactly the flagged ceiling-discontinuity
blocker behind the weak form of Conjecture 6.8d. \emph{Net:} the
$K$-escape budget does not dominate (reasons (a),(b)); the
alphabet-floor refinement closes the \emph{continuous} case (already
Lemma 6.8e) but the \emph{discrete} ceiling case --- the actual
residual of OP2(a) --- survives, blocked by the same ceiling
discontinuity as before. Attack vector~3 is closed negatively;
OP2(a) (discrete bit-cost) remains open. Verification script:
\texttt{scripts/verify/op2a\_k\_escape\_budget\_check.py} (budget
additivity, flat Pareto slope, $O(1/\log N)$ log saving at $K_{\rm
YES}/2$, and the ceiling-jump gap erasure for $N < 2^{8568}$,
all reproduced numerically).

\emph{Status update 2026-05-29 (POSITIVE side-result, Lemma 6.8k:
additive rule-size SGP is APX-hard for every $\alpha \geq 0$, with a
ceiling-free gap that \emph{increases} with $\alpha$; this is a new
hardness result for a real RNR objective, not a closure of the
discrete-bit-cost residual).} The journal's untried path was to adapt
Casel et al.'s number-of-rules--aware size measure to the additive
objective $L_\alpha(G) := L_{\mathrm{sym}}(G) + \alpha\,K(G)$ for general
$\alpha > 0$. We close this path \emph{positively} as a standalone
result, and verify (decisively) that it is \emph{not} a relabeling of
the continuous-bit-cost Lemma 6.8e.

\textbf{Lemma 6.8k (Additive rule-size SGP inherits Charikar
APX-hardness for every $\alpha \geq 0$; ceiling-free gap monotone in
$\alpha$).} \emph{For a straight-line grammar $G$ with $L_{\mathrm{sym}}(G)$
right-hand-side symbols and $K(G)$ rules (nonterminals), define the
additive rule-size objective $L_\alpha(G) := L_{\mathrm{sym}}(G) +
\alpha\,K(G)$, $\alpha \geq 0$. Then SGP under $L_\alpha$ is APX-hard:
for every fixed $\alpha \geq 0$ there is a constant
$\gamma_\alpha > 1$ such that no polynomial-time
$\gamma_\alpha$-approximation exists unless $\mathrm{P} = \mathrm{NP}$.
On the Berman--Karpinski--hard Charikar family the inapproximability
constant satisfies $\gamma_0 = \gamma_{\mathrm{SGP}} = 1 + 1/8568$ and is
strictly increasing in $\alpha$, with
$\lim_{\alpha \to \infty} \gamma_\alpha = 1 + 1/1008$.}

\emph{\textbf{Proof.}} \emph{(APX-hardness, all $\alpha \geq 0$.)} Casel
et al.\ (2021; conference version ICALP 2016, Theorem 7) reduce
\textsc{Vertex Cover} on subcubic graphs to SGP so that a minimal
grammar for the source word $w_G$ has $L_{\mathrm{sym}} = f(m,n) +
|\Gamma|$, where $\Gamma$ is the encoded cover and $f$ is a polynomial
independent of $\Gamma$ (their \S3, Lemma 6 and the displayed
equivalence $|G| \leq f(m,n) + |\Gamma|$). In that construction each
covered vertex $v_i$ contributes exactly one extra rule
$V_i \to \#V_i$, so the rule count is $K = r_0(m,n) + |\Gamma|$ with
$r_0$ likewise independent of $\Gamma$. Hence
\[
L_\alpha(G) = \bigl(f(m,n) + \alpha\,r_0(m,n)\bigr) + (1+\alpha)\,|\Gamma|,
\]
a strictly increasing affine function of $|\Gamma|$ for every
$\alpha \geq 0$. Therefore the $L_\alpha$-minimal grammar encodes a
\emph{minimum} vertex cover, exactly as for the pure symbol-count
objective ($\alpha = 0$). Berman--Karpinski (1999/2003) APX-hardness of
subcubic \textsc{Vertex Cover} (with the $145/144$ gap) thus transfers
through the same gap-preserving encoding: a $\gamma_\alpha$-approximation
of $L_\alpha$ would $\gamma_\alpha$-approximate the cover, contradicting
$\mathrm{P} \neq \mathrm{NP}$ for a suitable $\gamma_\alpha > 1$.

\emph{(Gap value and monotonicity, BK--Charikar family.)} On the
Berman--Karpinski extremal instances the YES and NO optima differ
\emph{only} in the cover term: writing $\tau$ for the optimal cover size
and $\Delta := \tau_{\mathrm{NO}} - \tau_{\mathrm{YES}} = (1/144)\tau$
(the BK promise inflation), both $L_{\mathrm{sym}}$ and $K$ shift by the
\emph{same} additive $\Delta$ across the two branches. With Charikar's
accounting $L_{\mathrm{sym}}^{\mathrm{YES}} = 15|V| + 3|E| + \tau =
\tfrac{119}{6}|V|$ and $K^{\mathrm{YES}} = 2|V| + \tau = \tfrac{7}{3}|V|$
(at $|E| = \tfrac32|V|$, $\tau = \tfrac13|V|$; the $2|V| + |C^*|$
nonterminal count of Observation 6.8f), the per-branch ratios are
\[
\frac{L_{\mathrm{sym}}^{\mathrm{NO}}}{L_{\mathrm{sym}}^{\mathrm{YES}}}
  = 1 + \frac{\Delta}{L_{\mathrm{sym}}^{\mathrm{YES}}} = 1 + \frac{1}{8568} = \gamma_{\mathrm{SGP}},
\qquad
\frac{K^{\mathrm{NO}}}{K^{\mathrm{YES}}}
  = 1 + \frac{\Delta}{K^{\mathrm{YES}}} = 1 + \frac{1}{1008}.
\]
The additive gap is the $\Delta$-mediant
\[
\gamma_\alpha
= \frac{L_{\mathrm{sym}}^{\mathrm{NO}} + \alpha K^{\mathrm{NO}}}{L_{\mathrm{sym}}^{\mathrm{YES}} + \alpha K^{\mathrm{YES}}}
= 1 + \frac{\Delta\,(1+\alpha)}{L_{\mathrm{sym}}^{\mathrm{YES}} + \alpha K^{\mathrm{YES}}}.
\]
Since $K^{\mathrm{YES}} < L_{\mathrm{sym}}^{\mathrm{YES}}$, the
denominator-per-$(1+\alpha)$ is a decreasing convex combination of the
two bases, so $\gamma_\alpha$ is strictly increasing in $\alpha$, from
$\gamma_0 = \gamma_{\mathrm{SGP}} = 1 + 1/8568$ up to
$\lim_{\alpha\to\infty}\gamma_\alpha = K^{\mathrm{NO}}/K^{\mathrm{YES}} =
1 + 1/1008$. \hfill$\blacksquare$

\emph{Why this is genuinely different from Lemma 6.8e (anti-tautology).}
The additive $L_\alpha$ is \emph{not} a monotone function of the
continuous bit-cost $L^{\mathrm{cont}}(G) = L_{\mathrm{sym}}(G)\cdot
\log_2(|\Sigma| + K(G))$, hence Lemma 6.8k is \emph{not} Lemma 6.8e
relabeled. Concretely, for a fixed source ($|\Sigma| = 1000$) the two
grammars with $K_1 = 64$ and $K_2 = 392$ on the model curve
$L_{\mathrm{sym}}(K) = K + cN/K$ ($N = 1000$, $c = 50$) are ranked
\emph{oppositely}: $L_1(G_1) = 909.3 < L_1(G_2) = 911.6$ (additive
prefers the low-$K$ grammar $G_1$) while $L^{\mathrm{cont}}(G_1) =
8499 > L^{\mathrm{cont}}(G_2) = 5426$ (bit-cost prefers the low-symbol
grammar $G_2$). The discordance is structural: the additive objective
penalises $K$ \emph{linearly} ($\alpha K$), whereas the bit-cost
penalises $K$ only \emph{logarithmically} (inside $\log_2(|\Sigma|+K)$,
where the $L_{\mathrm{sym}}$ factor dominates the product). Their
$K$-optimal points scale differently as well: the additive sum is
minimised at $K^* = \sqrt{cN/(1+\alpha)}$ (an $\alpha$-dependent
square-root law) whereas the continuous-bit-cost product is minimised at
$K^* = \sqrt{cN}$ ($\alpha$-free in the $|\Sigma| \gg K$ regime).

\emph{Why it does not close the discrete-bit-cost residual of OP2(a).}
$L_\alpha$ is \emph{continuous} (indeed affine) in $K$ and carries
\emph{no} ceiling. The discrete bit-cost multiplier
$\lceil \log_2(|\Sigma| + K)\rceil$ is, by contrast, a step function of
$K$: for all $K$ with $|\Sigma| + K$ in one power-of-two band the
multiplier is constant, so the discrete objective is \emph{blind} to $K$
within a band (e.g.\ at $|\Sigma| = 1100$ the grammars $K = 500$ and
$K = 800$ have identical discrete bit-cost $11\,L_{\mathrm{sym}}$ when
$L_{\mathrm{sym}}$ is equal, while $L_1$ separates them), and it
\emph{jumps} by a full multiplicative bit when $|\Sigma| + K$ crosses
$2^b$. The OP2(a) residual is exactly this ceiling discontinuity
(\S13.4.1 vector 3, Conjecture 6.8d): $\gamma_{\mathrm{SGP}} - 1 =
1/8568 < 1/\log_2 N$ for all $N < 2^{8568}$, so one ceiling jump can
erase the gap. Lemma 6.8k \emph{lacks} that discontinuity, so it cannot
resolve it; it is a hardness result for a different, ceiling-free
objective.

\emph{What it bounds for RNR.} The additive objective is a faithful RNR
\emph{dictionary rule-size} cost: $L_{\mathrm{sym}}$ is the total
right-hand-side symbol count of the token dictionary (a straight-line
grammar, \S6.5), and $\alpha K$ is the per-token header overhead --- each
of the $K$ dictionary entries needs a fixed-size record (length field,
type tag, or pointer), so $\alpha$ is the per-rule header width in
symbols. Lemma 6.8k therefore establishes, unconditionally modulo
$\mathrm{P}\neq\mathrm{NP}$, that minimising the RNR dictionary
rule-size (symbols plus per-entry header) is APX-hard for any fixed
header width, with an inapproximability constant between
$\gamma_{\mathrm{SGP}} = 1 + 1/8568$ and $1 + 1/1008$ on
Berman--Karpinski--hard inputs. This is a \emph{stronger} clean gap than
the symbol-count $\gamma_{\mathrm{SGP}}$ alone (the nonterminal-count gap
$1/1008$ exceeds the symbol-count gap $1/8568$ because the same minVC
inflation $\Delta$ sits on a smaller $K$-base $\tfrac73|V|$ than the
$L_{\mathrm{sym}}$-base $\tfrac{119}{6}|V|$), and it is ceiling-free ---
but it is a result about the additive dictionary cost, \emph{not} the
multiplicative integer bit-cost whose APX-hardness (OP2(a),
discrete) remains open. Verification:
\texttt{scripts/verify/op2a\_additive\_rule\_cost\_check.py} (Casel
robustness for all $\alpha$; rank-discordance witness $K_1=64,K_2=392$;
distinct square-root laws; exact $\gamma_\alpha$ monotonicity
$1+1/8568 \to 1+1/1008$; and the discrete step-function/ceiling
contrast, all reproduced numerically).

\emph{Status update 2026-05-29 (POSITIVE complement of Lemma 6.8k:
a polynomial-time $O(\log(N/g^\star))$-approximation for the additive
rule-cost objective, giving a tight-up-to-log characterisation of the
RNR dictionary-size problem).} Lemma 6.8k is a \emph{lower} bound
(APX-hardness). Its algorithmic flip-side asks whether the additive
objective inherits the symbol-count SGP approximation algorithms. We
verify that it does, with the genuine content concentrated in the
control of the rule count $K$, and we delimit the residual to the
classical (RNR-independent) constant-vs-logarithmic SGP approximability
gap.

\emph{Symbol-count SGP approximation algorithms (verified citations).}
For the pure symbol-count objective $L_{\mathrm{sym}}(G) = |G|$ the
smallest-grammar problem admits polynomial-time approximation within a
logarithmic factor of the optimum $g^\star$: Charikar, Lehman, Liu,
Panigrahy, Prabhakaran, Sahai \& Shelat (2005, \emph{IEEE Trans.\ Inf.\
Theory} 51(7):2554--2576) give two algorithms with ratios $O(\log^3 N)$
and $O(\log(N/g^\star))$ (their abstract and \S VII); Rytter (2003,
\emph{Theoret.\ Comput.\ Sci.} 302(1--3):211--222) turns the LZ77
factorisation (of $z \leq g^\star$ phrases) into a balanced
\emph{AVL-grammar} of size $O(z\log(N/z)) = O(g^\star\log(N/g^\star))$;
and Je\.z (2015, \emph{Theoret.\ Comput.\ Sci.} 592:115--134) achieves
ratio $O(1+\log(N/g^\star))$ by \emph{recompression}, with an analysis
independent of LZ77. Crucially, the Rytter and Je\.z outputs are
\emph{Chomsky-normal-form} (CNF / SLP) grammars: every internal rule has
a right-hand side of length exactly $2$. (RePair (Larsson--Moffat 1999)
is \emph{not} an $O(\log)$ algorithm: its worst-case symbol ratio is
$\Omega(\log N/\log\log N)$ (Charikar et al.\ 2005 gave
$\Omega(\sqrt{\log N})$; improved by Bannai, Hirayama, Hucke, Inenaga,
Je\.z, Lohrey \& Reh, ``The Smallest Grammar Problem Revisited,''
\emph{IEEE Trans.\ Inf.\ Theory} 67(1), 2021); we use it only as a
foil.)

\textbf{Lemma 6.8k-approx (Poly-time $O(\log)$-approximation for additive
rule-cost SGP).} \emph{Let $\mathcal A$ be any polynomial-time SGP
algorithm whose output is a CNF grammar and whose symbol-count ratio is
$\rho = \rho(N) = O(\log(N/g^\star))$ (e.g.\ Rytter 2003 or Je\.z 2015).
Then for the additive objective $L_\alpha = L_{\mathrm{sym}} + \alpha K$:}
\begin{enumerate}
\item[(i)] \emph{(Unconditional, all $\alpha \geq 0$.) $\mathcal A$ is a
polynomial-time $(1+\alpha)\rho$-approximation for $L_\alpha$. In
particular for every fixed $\alpha = O(1)$ the ratio is
$O(\rho) = O(\log(N/g^\star))$.}
\item[(ii)] \emph{(CNF subproblem, all $\alpha \geq 0$.) Against the
additive optimum restricted to CNF grammars, $\mathcal A$ attains ratio
at most $\rho = O(\log(N/g^\star))$, with no $(1+\alpha)$ loss (and
strictly below $\rho$ for $\alpha > 0$).}
\item[(iii)] \emph{(General optimum, super-constant $\alpha$.) Granting the
folklore constant-equivalence $g^\star = O(\kappa^\star)$ of the
grammar-size measure $g^\star = L_{\mathrm{sym}}$-optimum and the
rule-count measure $\kappa^\star = K$-optimum, $\mathcal A$ is an
$O(\rho) = O(\log(N/g^\star))$-approximation for $L_\alpha$ for
\emph{every} $\alpha \geq 0$, including $\alpha = \Theta(\log N)$.}
\end{enumerate}
\emph{Consequently the additive-rule-cost approximability ratio lies in the
window $[\,1+\Omega(1/\log N),\,O(\log N)\,]$: the lower endpoint is
Lemma 6.8k's constant inapproximability $\gamma_\alpha \in
[1+\tfrac{1}{8568}, 1+\tfrac{1}{1008}]$ (which exceeds $1$ but is below
$1+1/\log_2 N$ only in the ceiling-bounded discrete regime irrelevant
here), and the upper endpoint is the SGP logarithmic ceiling.}

\emph{\textbf{Proof.}} \emph{(i) Unconditional $(1+\alpha)\rho$ ratio,
all $\alpha \geq 0$.} For a CNF grammar $G$ write $n_{\mathrm{t}}$ for
its number of terminal rules and $n_{\mathrm{b}}$ for its number of
binary rules; then $K(G) = n_{\mathrm{t}} + n_{\mathrm{b}}$ and
$L_{\mathrm{sym}}(G) = n_{\mathrm{t}} + 2n_{\mathrm{b}}$, so
\begin{equation}
K(G) \;\leq\; L_{\mathrm{sym}}(G). \tag{NF}
\end{equation}
Hence $L_\alpha(\mathcal A(s)) = L_{\mathrm{sym}}(\mathcal A(s)) +
\alpha K(\mathcal A(s)) \leq (1+\alpha)\,L_{\mathrm{sym}}(\mathcal A(s))$.
By the approximation guarantee,
$L_{\mathrm{sym}}(\mathcal A(s)) \leq \rho\,g^\star$, where $g^\star$ is
the size of the smallest \emph{general} grammar. Finally, for the
additive optimum $G^\star_\alpha$ we have
$L_\alpha(G^\star_\alpha) \geq L_{\mathrm{sym}}(G^\star_\alpha) \geq
g^\star$ (the first step drops the nonnegative $\alpha K$ term; the
second is minimality of $g^\star$ over the symbol count). Chaining,
\[
L_\alpha(\mathcal A(s)) \;\leq\; (1+\alpha)\,\rho\,g^\star
\;\leq\; (1+\alpha)\,\rho\,L_\alpha(G^\star_\alpha),
\]
so the additive ratio is at most $(1+\alpha)\rho$ \emph{for every}
$\alpha \geq 0$ (no step used $\alpha = O(1)$), which is $O(\rho) =
O(\log(N/g^\star))$ whenever $\alpha = O(1)$. This proves (i).

\emph{(ii) CNF subproblem, every $\alpha$: exact deflation.} The bound
above is loose in $\alpha$. Restricted to CNF grammars the additive
objective is, by the displayed identities,
\begin{equation}
L_\alpha(G) = n_{\mathrm{t}}(1+\alpha) + (2+\alpha)\,n_{\mathrm{b}},
\qquad n_{\mathrm{t}} = |\Sigma_s| \text{ fixed}, \tag{AFF}
\end{equation}
an increasing affine function of $n_{\mathrm{b}}$ with an
$\alpha$-\emph{independent} minimiser (the constant $n_{\mathrm{t}} =
|\Sigma_s|$, the number of distinct terminals of $s$, is
grammar-independent). Therefore minimising $L_\alpha$ over CNF grammars
is the \emph{same} optimisation problem (minimise $n_{\mathrm{b}}$) for
every $\alpha \geq 0$, and the CNF additive optimum coincides with the
CNF symbol optimum, with binary-rule count $n_{\mathrm{b}}^\star$. The
symbol guarantee gives $n_{\mathrm{t}} + 2n_{\mathrm{b}}(\mathcal A) \leq
\rho(n_{\mathrm{t}} + 2n_{\mathrm{b}}^\star)$, i.e.\
$n_{\mathrm{b}}(\mathcal A) \leq \rho\,n_{\mathrm{b}}^\star +
\tfrac12(\rho-1)n_{\mathrm{t}}$. Substituting into (AFF),
\[
\rho\,L_\alpha(G^\star_{\mathrm{CNF}}) - L_\alpha(\mathcal A(s))
= n_{\mathrm{t}}(\rho-1)\Bigl[(1+\alpha) - \tfrac{2+\alpha}{2}\Bigr]
= n_{\mathrm{t}}(\rho-1)\,\tfrac{\alpha}{2} \;\geq\; 0,
\]
so the additive ratio against the CNF optimum is at most $\rho$ for all
$\alpha$, with slack $n_{\mathrm{t}}(\rho-1)\alpha/2$ that \emph{grows}
with $\alpha$ (the approximation-free constant term
$n_{\mathrm{t}}(1+\alpha)$ deflates the ratio strictly below $\rho$ for
$\alpha > 0$). This proves (ii). \hfill$\blacksquare$

\emph{(iii) General optimum, super-constant $\alpha$ (the genuine
content; conditional on a folklore constant).} Parts (i)--(ii) are
unconditional, but the $(1+\alpha)$ factor of (i) is genuinely needed
in the super-constant regime $\alpha = \Theta(\log N)$ against the
\emph{general} (non-CNF) optimum, because there the additive penalty
\emph{changes} the optimal grammar: as $\alpha$ grows the general
additive optimum trades symbols for rules, preferring fewer,
\emph{longer} right-hand sides (verified by brute force, e.g.\ for the
source $\mathtt{aabbaabb}$ the general symbol optimum has
$(L_{\mathrm{sym}}, K) = (6,4)$ whereas the general additive optimum at
$\alpha=8$ is $(8,3)$). This is the requested anti-tautology: the
transfer is \emph{not} trivial when $\alpha$ is large, since an
$L_{\mathrm{sym}}$-minimiser need not control $K$. We reduce the residual
to a single classical constant. Binarising the general additive optimum
$G^\star_{\alpha,\mathrm{gen}}$ (size $L_{\mathrm{sym}}^\star$, rules
$K^\star_\alpha$) yields a CNF grammar of symbol-size
$\leq 2L_{\mathrm{sym}}^\star$ and rule-count
$\leq L_{\mathrm{sym}}^\star$ (a rule $A \to \beta$ with $|\beta| = k$
becomes $k-1$ binary rules; the binary-rule total is
$L_{\mathrm{sym}}^\star - K^\star_\alpha \leq L_{\mathrm{sym}}^\star$ ---
the \emph{linear} SLP binarisation cost, not the quadratic worst case
for general CFGs). Hence the CNF additive optimum obeys
$L_\alpha(G^\star_{\alpha,\mathrm{CNF}}) \leq
(2+\alpha)\,L_{\mathrm{sym}}^\star \leq (2+\alpha)(1+\alpha)\,g^\star$,
while $L_\alpha(G^\star_{\alpha,\mathrm{gen}}) \geq \alpha K^\star_\alpha
\geq \alpha\kappa^\star$, so the CNF-restriction price is
\[
\frac{L_\alpha(G^\star_{\alpha,\mathrm{CNF}})}{L_\alpha(G^\star_{\alpha,\mathrm{gen}})}
\;=\; O\!\Bigl(\frac{g^\star}{\kappa^\star}\Bigr),
\]
where $\kappa^\star$ is the minimum rule-count of any grammar for $s$.
By (ii), $\mathcal A$ is a $\rho$-approximation against
$G^\star_{\alpha,\mathrm{CNF}}$, so against the general optimum its ratio
is $O(\rho\cdot g^\star/\kappa^\star)$. \emph{Granting the folklore
constant-equivalence $g^\star = O(\kappa^\star)$} --- both measures are
$\Theta$ of the SLP/LZ77 complexity, and Charikar et al.\ 2005 (\S VII),
Lehman (2002, PhD thesis, MIT) and Casel et al.\ 2020 (\emph{Theory
Comput.\ Syst.} 65(2):344--409) use the size and rule-count measures
interchangeably --- this is $O(\rho) = O(\log(N/g^\star))$ for
\emph{every} $\alpha \geq 0$. We flag $g^\star = O(\kappa^\star)$ as the
single residual assumption (it holds with constant $\leq 2$ on every
brute-forced instance in the verification script but we do not have a
closed-form proof for all strings); part (i)'s $(1+\alpha)\rho$ bound is
unconditional regardless. \emph{Remark (the natural midpoint).} At
$\alpha = 1$ the additive objective $L_{\mathrm{sym}} + K$ is exactly the
standard ``alternative'' grammar-size measure of the SGP literature
(Lehman 2002; Casel et al.\ 2020), so (i) gives an unconditional
$2\rho = O(\log)$-approximation for that canonical measure with no
appeal to (iii).

\emph{Why this is the tight characterisation, not a relabelling of
Lemma 6.8k.} Lemma 6.8k supplies the \emph{lower} bound
$\gamma_\alpha > 1$; this complement supplies the \emph{upper} bound
$O(\log N)$. The gap between a constant lower bound and a logarithmic
upper bound is precisely the long-standing SGP approximability gap
(open since Charikar et al.\ 2005; restated as open by Casel et al.\
2020): no $N$-independent constant-factor algorithm is known, and no
super-constant inapproximability is known. The additive objective
\emph{inherits} this window verbatim --- the positive content is that
the inheritance is genuine (it survives the $K$-penalty via CNF
deflation and the $O(1)$ rule-count/symbol-count equivalence), giving
the RNR dictionary builder a concrete usable procedure: run Je\.z's
recompression (or Rytter's AVL construction) on the source, obtaining a
CNF dictionary within an $O(\log(N/g^\star))$ factor of the optimal
additive rule-size cost. Verification script:
\texttt{scripts/verify/lemma\_6\_8k\_approx\_additive\_rule\_cost.py}
(the (NF) and (AFF) identities; the $\alpha$-independent CNF argmin
verified against brute-force optima for $\alpha$ up to $1000$; the
constant-$\alpha$ ratio bound; the bounded additive ratio of a
recompression-style builder vs.\ both the CNF and the \emph{general}
brute-force additive optima across $\alpha \in \{1, \log N, N, 50\}$
with an $O(1)$ blow-up factor; the witnessed general-additive $\neq$
general-symbol optimum for large $\alpha$; the $\leq 2$ binarisation
cost; and the RePair foil, all reproduced numerically).

\emph{Status update 2026-05-29 (attack vector~2, ``$K$-uniform
hardness source (Label Cover)'', closed negatively; with a
bounded-$K$ dichotomy that rules out the entire $K$-uniform
programme).} Item~2 proposed defeating the ceiling discontinuity by a
different source-of-hardness whose reduction is \emph{$K$-uniform}:
$K(G^*_{\rm YES}) = K(G^*_{\rm NO})$ exactly, so the discrete factor
$\lceil \log_2(|\Sigma| + K)\rceil$ is identical for YES and NO and
cannot differentially favour NO. The candidate was \textsc{Label
Cover} (Arora--Lund--Motwani--Sudan--Szegedy 1998, \emph{J.\ ACM}
45(3):501--555; Raz 1998 parallel repetition, \emph{SIAM J.\ Comput.}
27(3):763--803), with a secondary hope that the fixed-alphabet SGP
hardness of Casel, Fernau, Gaspers, Gras \& Schmid (2016 ICALP,
2020 \emph{Theory Comput.\ Syst.}\ 65(2):344--409) removes the
alphabet floor. We find the $K$-uniform route \emph{closed
negatively}, by \emph{four} independent blockers, the last of which
rules out the programme in full generality.

\emph{(a)~$K$-uniformity fails for every cover/selection-encoding
reduction.} Both known constant-gap SGP-hardness reductions encode
their optimisation signal as a \emph{set of rules whose cardinality
is the optimum}. In Charikar's unbounded-alphabet reduction the
optimal grammar has nonterminal count $K = 2|V| + |C|$ (one rule
$A_i \to \#v_i$ and one $B_i \to v_i\#$ per vertex, plus one
$D_i \to \#v_i\#$ per \emph{cover} vertex $v_i \in C$; Charikar 2005
p.~2558 grammar-size decomposition $8|V| + (3|V|-|C|) + 3|E| + 4|V|
+ 2|C|$). In Casel et al.'s fixed-alphabet reduction (Thm~7, via
Lemma~6) a vertex cover $\Gamma$ translates into a grammar of size
$f(m,n) + |\Gamma|$ through rules $\{\overleftarrow{V}_i \to
\#\overrightarrow{V}_i : i \in \mathfrak{I}\}$, $\mathfrak{I} =
\{i : v_i \in \Gamma\}$, so $K = K_0(n,m) + |\Gamma|$ with $K_0$ a
\emph{cover-independent} skeleton. In both, $K_{\rm NO} - K_{\rm YES}
= |C_{\rm NO}| - |C_{\rm YES}| = (\rho_{\rm VC} - 1)\,k_{\rm YES} > 0$:
$K$ differs by \emph{exactly} the cover gap, the very signal the
reduction carries. Forcing $K_{\rm YES} = K_{\rm NO}$ would delete the
cover rules and erase the gap. Cover/selection encodings are
\emph{inherently} non-$K$-uniform; one cannot pad the YES optimum up
to $K_{\rm NO}$ because unused nonterminals only \emph{increase} $|G|$
and are excluded from any optimal grammar.

\emph{(b)~The fixed alphabet removes the alphabet floor but not the
skeleton-overhead floor.} Granting Casel's fixed $|\Sigma| \in
\{17,\dots,24\}$ (so $\lceil \log_2(|\Sigma|+K)\rceil \approx
\lceil \log_2 K\rceil$, no $|\Sigma| = \Theta(N)$ floor), the
symbol-count gap is still $\rho_{\rm sym} = 1 + (\text{cover
gap})/(f(m,n) + \text{cover})$. Casel's skeleton is provably large:
their Lemma~5 proof gives $|u|_\star = 196n$ for just the \emph{first}
of the three parts $w_{\mathcal G} = uvw$, so $f(m,n) \geq 196n$ and
$\rho_{\rm sym} - 1 \leq (0.336\,n)/(196\,n) \approx 1.7 \times
10^{-3} = O(1/c_f)$ with $c_f \geq 196$. This is the \emph{same}
multiplicative dilution as Charikar's $15|V| + 3|E|$ overhead
(\S13.2): the alphabet floor governs the \emph{ceiling}, but the
structural skeleton governs the \emph{symbol gap}, and the latter is
diluted exactly as in the unbounded case. The fixed alphabet thus
neutralises the wrong floor.

\emph{(c)~Label Cover dilutes via its own alphabet blow-up.} A
\textsc{Label Cover} instance attaining a gap requires label
alphabet $R = \mathrm{poly}(n)$ and an $R$-fold size blow-up of the
constraint graph (the projection-game / parallel-repetition cost).
A natural \textsc{LC} $\to$ SGP encoding selects, per vertex, a rule
encoding its assigned label, so $K$ tracks the number of selected
labels and is again non-uniform between full (YES) and partial (NO)
satisfaction. Worse, the $R$-fold blow-up inflates the Charikar-style
per-gadget overhead by $\Theta(R)$, driving any constant SGP gap to
$1 + o(1)$ as $R \to \infty$ --- the $15|V| + 3|E|$ unified obstacle
\emph{amplified} by $R$. \textsc{LC}'s $1/\mathrm{poly}$ value gap is
moreover a minimisation-of-unsat gap; the standard \textsc{LC} $\to$
\textsc{Set Cover} map yields a $\ln$-factor (super-constant) gap, not
the bounded multiplicative ratio SGP inapproximability is stated in.

\emph{(d)~Bounded-$K$ dichotomy (decisive).} Casel et al.\ Thm~11
states that a grammar for $w$ with at most $k$ rules that is minimal
among all grammars with $\leq k$ rules is computable in time
$O(|w|^{2k+6})$: \emph{if the nonterminal count $K$ is fixed, SGP is
in {\rm P}}. This kills the $K$-uniform programme by a clean
dichotomy. A $K$-uniform reduction must fix $K$ as a function of
instance size. Either (i)~$K$ is fixed to a \emph{constant}, whence
SGP is polynomial (Thm~11) and there is no hardness to reduce
\emph{to} --- the attack is vacuous; or (ii)~$K$ is fixed to a
\emph{growing} $K_0(n)$ (so $O(|w|^{2K_0+6})$ is super-polynomial),
whence the \emph{optimal} grammar must use exactly $K_0(n)$
nonterminals for \emph{both} YES and NO, contradicting blocker~(a)
(the optimal $K$ provably differs by the cover gap for the only known
constant-gap SGP-hardness sources). Casel et al.\ Thm~12
(SGP parameterised by $|N|$ is $\mathsf{W}[1]$-hard) blocks the FPT
escape from branch~(i). Hence $K$-uniformity is either vacuous or
self-contradictory for the SGP-hardness family. \emph{Net:} the
$K$-uniform / \textsc{Label Cover} route does not reach the ceiling
discontinuity; it is closed by blockers (a),(b),(d) unconditionally
and (c) for the specific \textsc{LC} source. This \emph{extends} the
unified obstacle: a hardness source able to close OP2(a) would have to
encode its optimisation signal \emph{without} a variable-cardinality
rule set (else non-$K$-uniform) \emph{and without} a super-constant
structural skeleton (else gap-diluted) \emph{and} at bounded $K$
(else Thm~11 trivialises) --- a combination no reduction in the SGP
or \textsc{Label Cover} literature exhibits. OP2(a) (discrete
bit-cost) remains open; the residual is still exactly the ceiling
discontinuity, now shown unreachable by the $K$-uniform route as well.
Verification script:
\texttt{scripts/verify/op2a\_label\_cover\_kuniform\_check.py}
(non-$K$-uniformity of both reductions exhibited on explicit
max-degree-3 instances with $K_{\rm NO} - K_{\rm YES}$ equal to the
cover gap; the $O(1/c_f)$ skeleton dilution with $c_f \geq 196$;
and the Casel Thm~11/12 dichotomy, all reproduced numerically).

\emph{Status update 2026-05-29 (SGP-constant literature survey:
the residual requires a sharper constant than any known).} The
ceiling-discontinuity residual closes if the symbol-count SGP
inapproximability constant exceeds $1 + 1/\log_2 N$ (so that a single
ceiling jump cannot erase it). A survey of the SGP-hardness literature
2005--2026 confirms no such constant is known: Charikar et al.\ 2005
give $\gamma_{\rm SGP} = 8569/8568$ (inheriting Berman--Karpinski
$145/144$ diluted by the $15|V|+3|E|$ grammar encoding, factor
$\approx 59.5$); the fixed-alphabet sharpenings (Casel et al.\ 2020;
Arpe--Reischuk-style $24c$ and Bannai et al.\ $6c$ alphabet-transfer
multipliers) all preserve a constant that must still descend to
$\approx 1.000117$ for the binary NP-hardness statement, hence remain
below $1 + 1/\log_2 N$ for every $N < 2^{8568}$. No 2005--2026 result
supplies a constant ($N$-independent) multiplicative gap for SGP, and
direct gap amplification on the SGP objective is blocked because
grammar size does not compose multiplicatively under instance product
(unlike CSP value under parallel repetition). \emph{Net:} the residual
of OP2(a) is precisely a sharper SGP inapproximability constant than
Charikar's --- itself a $20$-year open problem in the SGP literature,
independent of the RNR framework. Verification script:
\texttt{scripts/verify/op2a\_sgp\_constant\_survey\_check.py}
(reproduces $\gamma_{\rm SGP}$, the $1/144 \to 1/8568$ encoding
dilution, the $1/8568 < 1/\log_2 N$ frontier for all practical $N$,
and the alphabet-transfer multipliers, all numerically).

\emph{What would constitute closure:} a polynomial-time
$(\gamma_{\rm bit} - \varepsilon)$-approximation algorithm for some
explicit $\gamma_{\rm bit} > 1$, OR a polynomial-time reduction
from a known NP-hard problem with $\gamma_{\rm bit}$-gap preservation
across all $K$ regimes.

\textbf{13.4.2 OP2(b) (Free-predictor joint Type-III-C APX-hardness).}

\emph{What has been tried (failed paths):}
\begin{itemize}
\item Direct padding reduction from restricted to unrestricted
joint --- blocked by the predictor-absorption obstacle (Remark 6.8b):
an unrestricted predictor can pay $L(M) \geq |s|$ bits to memorise
the SGP-format prefix directly, absorbing the $\gamma_{\rm SGP}$
gap into $L(M)$.
\end{itemize}

\emph{Attack vectors not yet tried:}
\begin{enumerate}
\item \emph{Amortised-family reformulation (§13.2 documented).}
Switch to family $\mathcal{F}_n = \{s_i\}$ of $N(n) = 2^{\Theta(n)}$
SGP-prefix instances. For $M$ to memorise all $N(n)$ prefixes,
$L(M) \geq N(n) \cdot n$ bits, forcing per-instance amortisation to
$o(1)$ bit/source. This circumvents predictor-absorption.

\emph{Status update 2026-05-27 (negative): amortised-family
predictor-absorption transfer.} An explicit reduction attempt
identifies a definitive blocker. The premise ``$L(M) \geq N(n)
\cdot n$ forces $M$'s contribution to $o(1)$'' implicitly requires
the family $\mathcal{F}_n$ to be Kolmogorov-incompressible. But any
polynomial-time reduction from VC-3 (or any combinatorial NP-hard
problem) produces a family with a \emph{poly-size description} ---
the reduction algorithm and its $\mathrm{poly}(n)$-bit input. An
unrestricted predictor $M$ can therefore encode the entire family
description in $L(M) = \mathrm{poly}(n)$ bits and, after
amortisation $L(M)/N(n) \to 0$, serve each $s_i$ with
$L(R_i \mid Y, M, D) = O(1)$ lookup. The $\gamma_{\mathrm{SGP}}$
gap is absorbed into $L(M)$ symmetrically across YES and NO
branches --- exactly Remark 6.8b's failure mode applied to the
family description rather than the single source. \emph{Quantitative
check.} If we hypothetically enforced $L(M)/N(n) \to 0$ via a
restriction not available to a free predictor, per-instance
continuous bit-cost $\Theta(n \log n)$ (Lemma 6.8e) dominates
$\Theta(\log N(n)) = \Theta(n)$ salt cost, and the amortised gap
would survive as
$(\gamma_{\mathrm{SGP}} \log n + 1)/(\log n + 1) \to \gamma_{\mathrm{SGP}}$
as $n \to \infty$. But for any poly-time-samplable
$\mathcal{F}_n$, this enforcement fails: $L(M) = \mathrm{poly}(n)
= o(N(n))$ suffices to absorb the family description, the gap
collapses to $1 + o(1)$, and the reduction yields no constant
inapproximability. The amortisation moves the obstacle from ``$M$
absorbs $s$'' to ``$M$ absorbs the reduction algorithm'',
preserving the obstacle's structural form. \emph{When this rescue
would work (and why it does not close OP2(b)).} The absorption
breaks if (i) $M$ is class-restricted (e.g., $|M| \leq
n^{1-\varepsilon}$) --- but this is already covered by Lemma
6.8b's normal-form restriction; (ii) salts $r_i$ have oracle
randomness --- not available to deterministic polynomial-time
reductions; or (iii) salts are cryptographically pseudorandom (PRG
output indistinguishable by poly-time $M$) --- this is the
Liu--Pass / OWF route already documented in §13.2, which gives
average-case decision hardness under one-way functions, \emph{not}
worst-case multiplicative APX-hardness. None of (i)--(iii) yields
free-predictor worst-case OP2(b)-APX. \emph{Verdict.} The
amortised-family reformulation, taken in isolation, does not close
OP2(b)-APX. The single-source predictor-absorption obstacle of
Remark 6.8b transfers verbatim to the family-description scale;
the amortisation is symmetric across YES/NO branches and dilutes
the gap to $1+o(1)$ under any poly-time-samplable family.
\item \emph{Per-instance salting.} Each SGP-prefix $s_i$ is concatenated
with a unique random salt $r_i$ of length $\Omega(\log N(n))$ bits.
The predictor must either memorise all salts (forcing $L(M) \geq
N(n) \log N(n)$) or fail on the salted instances (achieving no
predictive gain). Salting forces a per-instance hardness budget.

\emph{Status update 2026-05-27 (subsumed by amortised-family
negative).} Per-instance salting is a special case of item 1's
amortised-family construction (salts $r_i$ generate the family
$\mathcal{F}_n = \{s \cdot r_i\}$). The same predictor-absorption
mechanism applies: if salts are produced by a poly-time-computable
function of $i$ (the only option for a polynomial-time reduction),
then $M$ stores that function in $O(\log n)$ bits and computes
$r_i$ on demand, absorbing the salting structure with the rest of
the family description. If salts are truly random or PRG-pseudo\-random,
the reduction is either not polynomial-time (random-oracle access)
or routes through OWF (Liu--Pass, average-case decision only).
Salting is therefore not an independent attack vector.
\item \emph{Reduction from Generalised Smallest Grammar Problem.}
Carrascosa et al.\ 2016 (arXiv:1608.08927) define GSGP as
multi-string grammar inference. APX-hardness of GSGP (if it can be
established) would directly imply joint OP2(b) hardness in the
amortised setting.

\emph{Remaining-viable status 2026-05-27.} GSGP is the only attack
vector in this list that does not collapse to predictor absorption,
because the GSGP objective is a single multi-string grammar (one
$G$ encoding all sources) rather than separated predictor + per-
instance residual: the predictor cannot escape the SGP structure
by ``moving it into $L(M)$'' if $L(M)$ is itself measured by the
GSGP grammar size. APX-hardness of GSGP is itself an open problem
(no reduction with multiplicative gap is currently known); closing
it would be a prerequisite to closing OP2(b) via this route, and
is independent of the predictor-absorption obstacle of Remark 6.8b.

\emph{Status update 2026-05-28 (negative, literature misreading +
structural transfer collapse).} A direct read of Siyari and
Gall\'e (\emph{The Generalized Smallest Grammar Problem},
arXiv:1608.08927v1, 31 Aug 2016) corrects two premises of the
above remaining-viable status. (a)~\emph{GSGP is not multi-string
grammar inference.} Siyari--Gall\'e Definition~3 (Section~3 of
the arXiv version) takes a \emph{single} sequence $s$ as input
and asks for a smallest non-recursive context-free grammar $G$
with $s \in L(G)$, scored by $|G|_s = \sum_{N \to \alpha \in
P}(|\alpha|+1) + \mathrm{cost}(s \mid G)$. The
``generalisation'' is from straight-line grammars to non-recursive
grammars --- the language $L(G)$ may be finite but larger than
$\{s\}$, with an explicit cost for selecting $s$ within $L(G)$
--- \emph{not} from one string to a family of strings. The paper
provides no NP-hardness or APX-hardness theorem: a text-search
for ``NP-hard'', ``APX-hard'', or ``NP-complete'' returns no
matches, and only ``approximate'' (algorithmic) plus a single
unrelated occurrence in the related-work paragraph (worst-case
ratio of IRRMGP / ZZ / MMAS-GA is ``hard to analyze'') appear.
The §13.4.2 premise ``APX-hardness of GSGP $\Rightarrow$ OP2(b)
amortised hardness'' therefore has no valid literature anchor:
GSGP-as-defined is a single-string problem, and the cited
hardness results for single-string SGP (Charikar et al.\ 2005's
$\gamma_{\mathrm{SGP}} = 8569/8568$, and Casel et al., \emph{On
the Complexity of the Smallest Grammar Problem over Fixed
Alphabets}, Theory Comput.\ Syst.\ 65(2):344--409, 2020,
doi:10.1007/s00224-020-10013-w, which establishes APX-hardness
even for alphabets of size $\geq 17$) do not transfer to the
amortised-family bit-cost objective without crossing the same
predictor-absorption obstacle as item~1.
(b)~\emph{A self-defined multi-string SGP reduces to single-string
SGP up to $1+o(1)$ overhead.} Even if one self-defines $k$-MSGP
(input: $\{s_1,\dots,s_k\} \subset \Sigma^*$; output: a smallest
SLP-collection $(G,A_1,\dots,A_k)$ with shared non-terminal pool
such that $\mathrm{val}(A_i)=s_i$, scored by $|G|+k$), the
following reduction shows $k$-MSGP $\equiv_{1+o(1)}$ single-string
SGP: take fresh separators $\#_1,\dots,\#_k \notin \Sigma$ and
concatenate $s_\# := s_1 \#_1 s_2 \#_2 \cdots s_k \#_k$; by
uniqueness of separators, every SLP for $s_\#$ has start rule
$S \to A_1 \#_1 A_2 \#_2 \cdots A_k \#_k$ with $\mathrm{val}(A_i)
= s_i$, and shared sub-expressions across the $A_i$'s are
precisely shared non-terminals in the induced $k$-MSGP solution
(and conversely). Both directions preserve grammar size up to
the additive constant $2k$ (separators and start-rule pointers),
which is $o(|s_\#|)$ when $\min_i |s_i| \to \infty$. Hence
Casel et al.\ 2020's APX-hardness with $\gamma_{\mathrm{SGP}} =
8569/8568$ transfers to $k$-MSGP for any constant $k$. This
genuine grammar-size APX-hardness, however, does \emph{not}
imply OP2(b) amortised-bit-cost APX-hardness, by step~(c) below.
(c)~\emph{Predictor absorbs the family itself, not any specific
MSGP solution.} The OP2(b) amortised objective is
$L_{\mathrm{amort}}(M,D,\{R_i\}) := L(M) + L(D) + \frac{1}{N}
\sum_{i=1}^{N} L(R_i \mid Y,M,D)$, \emph{not} the MSGP grammar
size $|G_{\mathrm{MSGP}}|$. Although the optimal MSGP grammar
$G^*$ for $\mathcal{F}_n$ has size $|G^*| = \Theta(N(n+1) /
\log(N(n+1)))$ (Berstel--Brlek worst-case-SLP bound on the
separator-concatenation $s_\# = s_1 \#_1 \cdots s_N \#_N$ of
total length $N(n+1)$, surveyed in Bannai et al.\ 2019,
arXiv:1908.06428, Section~2) and is NP-hard to compute exactly
from $d_n$, this is irrelevant: the free predictor does not
need to solve MSGP. Instead, the predictor encodes the
\emph{family-reproduction algorithm} (the same mechanism as
item~1): let $\mathrm{Red}_n$ be the polynomial-time reduction
algorithm that maps $(d_n, i) \mapsto s_i$, with description
size $|\mathrm{Red}_n| = \mathrm{poly}(n)$. Set the free
predictor $M^* := \langle d_n, \mathrm{Red}_n \rangle$ of
total size $L(M^*) = O(\mathrm{poly}(n))$. The decoder $D^*$
on input $(M^*, i)$ runs $\mathrm{Red}_n(d_n, i)$ to recover
$s_i$ deterministically. Hence $L(R_i \mid Y, M^*, D^*) =
\lceil \log_2 N \rceil$ (just the index $i$). The amortised
cost is $L_{\mathrm{amort}}(M^*, D^*, \{R_i = i\}) =
O(\mathrm{poly}(n)) + O(1) + O(\log N) = O(\mathrm{poly}(n))$
since $\log N = O(n)$. This is achieved \emph{without any MSGP
solver being invoked}; the MSGP solution $G^*$ does not appear
in the encoding at all. Consequently, the OP2(b) amortised
optimum on this family satisfies $\mathrm{OPT}_{\mathrm{amort}}
\leq L_{\mathrm{amort}}(M^*, D^*, \{R_i^*\}) =
O(\mathrm{poly}(n))$, \emph{independent of} the underlying
MSGP-objective gap. The MSGP grammar-size gap
$\gamma_{\mathrm{SGP}} = 8569/8568$ on YES vs.\ NO instances
is \emph{absorbed} (not provably collapsed): both YES and NO
admit the same family-storing predictor $M^*$ giving each a
common $O(\mathrm{poly}(n))$ upper bound on $L_{\mathrm{amort}}$,
so the \emph{grammar-size} coordinate on which
$\gamma_{\mathrm{SGP}}$ acts disappears from $L_{\mathrm{amort}}$.
(As corrected in Lemma~13.4.2a below, this common upper bound does
\emph{not} by itself force the YES/NO optimum ratio to $1+o(1)$ ---
a residual gap from differential family compressibility is not
excluded.) The MSGP grammar-size gap is absorbed because the
family-storing predictor $M^*$ achieves $L_{\mathrm{amort}}$
much smaller than the
per-instance MSGP-grammar cost $|G^*|/N = \Theta(n/\log n)$
even when MSGP is solved optimally; the amortised objective
structurally prefers family memorisation over MSGP solution.

\textbf{Lemma 13.4.2a (unified amortised-objective absorption ---
upper bound; corrected 2026-06).}
\emph{Let $\Pi$ be any minimisation problem on inputs
$\mathcal{F}_n = \{x_i\}_{i \in [N(n)]}$ produced by a
deterministic reduction $\mathrm{Red}_n:
\{0,1\}^{|d_n|} \times [N] \to \Sigma^*$ from a
$\mathrm{poly}(n)$-bit family description $d_n$, runnable by a
\emph{fixed-clock} machine (equivalently $\mathrm{Red}_n$ is given
as a polynomial-size circuit / straight-line program, or as code
carrying an explicit shared time bound $t(n) = \mathrm{poly}(n)$),
with $N(n) = 2^{\Theta(n)}$ and $|\mathrm{Red}_n| + |d_n| =
\mathrm{poly}(n)$ bits. The amortised-family bit-cost objective
$L_{\mathrm{amort}}(M, D, \{R_i\}) := L(M) + L(D) +
\frac{1}{N}\sum_{i=1}^N L(R_i \mid Y, M, D)$
admits a free predictor $M^* := \langle d_n, \mathrm{Red}_n
\rangle$ achieving $L_{\mathrm{amort}}(M^*, D^*, \{R_i^*\}) =
O(\mathrm{poly}(n))$ for the fixed-clock universal decoder $D^*$
(the explicit time bound makes $|D^*| = O(1)$ legitimate) and
trivial residuals $R_i^* := i$. Consequently, the optimum
amortised cost satisfies the \emph{upper bound}
$\mathrm{OPT}_{\mathrm{amort}}(\mathcal{F}_n)
\leq O(\mathrm{poly}(n))$ on any family $\mathcal{F}_n$
produced by such a reduction, \emph{independent of} the
underlying $\Pi$-objective gap. Hence the $\Pi$-objective gap
$\gamma_\Pi$ --- which acts on a quantity ($|G^*_\Pi|$,
$|\mathrm{VC}^*|$, \dots) that does not appear in
$L_{\mathrm{amort}}$ --- is \emph{absorbed into $L(M^*)$}: the NO
instance is poly$(n)$-compressible by $M^*$, so the reduction is
not gap-preserving for $L_{\mathrm{amort}}$ and the grammar-size /
cover-size hardness of $\Pi$ is not directly inherited.}

\emph{What this does \emph{not} establish (corrected after adversarial review, 2026-06).} The upper bound alone does
\emph{not} imply the amortised YES/NO optimum ratio is $1 + o(1)$,
nor that $L_{\mathrm{amort}}$ is APX-easy. ``Both YES and NO admit
the \emph{same} poly$(n)$ witness $M^*$'' bounds both optima
\emph{above} by $\mathrm{poly}(n)$ but says nothing about their
\emph{ratio}: $A \leq P$ and $B \leq P$ do not bound $A/B$. A
residual gap can survive within the poly$(n)$ envelope if the YES
family is differentially more compressible than the NO family by
\emph{some other} predictor $M' \neq M^*$ (e.g.\ an incompressible
common prefix that any decoder/residual scheme must convey forces
$\mathrm{OPT}_{\mathrm{amort}}(\text{NO}) = \Theta(n^2)$ while
$\mathrm{OPT}_{\mathrm{amort}}(\text{YES}) = \Theta(n)$, ratio
$\Theta(n)$). Refuting $L_{\mathrm{amort}}$ APX-hardness therefore
requires a \emph{matching lower bound} $\mathrm{OPT}_{\mathrm{amort}}
(\mathcal{F}_n) = \Omega(\mathrm{poly}(n))$ on the YES family
(so that $M^*$ is a constant-factor approximation), which this
lemma does not supply and which remains open. The lemma's
established content is exactly the upper bound and the consequent
non-inheritance of the \emph{specific} $\Pi$-objective gap.

\emph{Proof.} Define $M^* := \langle d_n, \mathrm{Red}_n
\rangle$ as a single string of $|d_n| + |\mathrm{Red}_n| =
O(\mathrm{poly}(n))$ bits. Define $D^*$ as the universal
decoder: on input $(M, i)$, parse $M = \langle d, \mathrm{Red}
\rangle$, then execute $\mathrm{Red}(d, i)$ and output the
result; $|D^*| = O(1)$ as a fixed universal machine. For each
$i \in [N]$, set the residual $R_i^* := i$ encoded in
$\lceil \log_2 N \rceil$ bits. Then $D^*(M^*, R_i^*) =
\mathrm{Red}_n(d_n, i) = s_i$, deterministically recovering
the $i$-th data block of $\mathcal{F}_n$. Hence $L(R_i^*
\mid Y, M^*, D^*) = \lceil \log_2 N \rceil = O(n)$ suffices.
Amortised cost:
$L_{\mathrm{amort}}(M^*, D^*, \{R_i^*\}) = L(M^*) + L(D^*) +
\frac{1}{N}\sum_{i=1}^N L(R_i^* \mid \cdot) = O(\mathrm{poly}(n))
+ O(1) + O(n) = O(\mathrm{poly}(n))$. By optimality,
$\mathrm{OPT}_{\mathrm{amort}}(\mathcal{F}_n) \leq
L_{\mathrm{amort}}(M^*, D^*, \{R_i^*\}) = O(\mathrm{poly}(n))$.
Since $M^*$ does not invoke any $\Pi$-solver --- it re-runs
$\mathrm{Red}_n$ to reproduce the family --- the \emph{upper}
bound holds uniformly across YES and NO instances of any reduction
$\Pi$-YES vs.\ $\Pi$-NO based on the same family-description
format. Both admit the same $M^*$ construction (with the
appropriate $d_n$), giving each optimum the same $O(\mathrm{poly}(n))$
\emph{ceiling}. Since $\gamma_\Pi$ acted on a quantity
($|G^*_\Pi|$, $|\mathrm{VC}^*|$, etc.) that does not appear in
$L_{\mathrm{amort}}$, that specific gap is absorbed into $L(M^*)$.
\emph{The proof establishes only this common upper bound; it does
not bound the YES/NO optimum ratio (both being $\leq P$ does not
constrain $A/B$), so no $1+o(1)$ ratio is claimed --- see the
``what this does not establish'' paragraph above.} $\Box$

\emph{Consequence (unification of all three attack vectors).}
Lemma~13.4.2a unifies items~1, 2, and the present item~3 under
a single structural obstacle: \emph{any reduction whose family
description $\langle d_n, \mathrm{Red}_n \rangle$ has size
$\mathrm{poly}(n)$ admits a poly-size free predictor $M^*$
that absorbs the family into $L(M)$, achieving amortised cost
$O(\mathrm{poly}(n))$ regardless of the $\Pi$-objective gap.}
Item~1 (predictor absorbs the reduction algorithm $d_n$),
item~2 (predictor absorbs the salting function as part of the
reduction algorithm), and item~3 (predictor bypasses MSGP
entirely by re-running the reduction algorithm to reproduce
the family --- the MSGP solution is irrelevant to
$L_{\mathrm{amort}}$, not because the predictor solves MSGP
but because the optimal-amortised strategy never invokes any
MSGP solver) are three instantiations of the same lemma. The
``GSGP is the only non-absorption attack vector''
characterisation of the previous status update is incorrect:
the MSGP hardness object \emph{is} absorbed by $L(M)$ via the
family-reproducing predictor, exactly as items~1--2 absorb the
reduction algorithm.
\emph{Verdict (attack vector~3, corrected 2026-06).} The MSGP
grammar-size APX-hardness (via separator reduction to
single-string SGP and Casel et al.\ 2020) holds, but its
\emph{specific} grammar-size gap does not \emph{directly} transfer
to the amortised-bit-cost objective: Lemma~13.4.2a's predictor
$M^*$ absorbs the grammar-size quantity into $L(M)$, so the NO
instance is poly$(n)$-compressible and the reduction is not
gap-preserving for $L_{\mathrm{amort}}$ in the grammar-size
coordinate. This is weaker than ``attack vector~3 closed'': as the
corrected lemma notes, a \emph{residual} $L_{\mathrm{amort}}$ gap
from differential family compressibility (a YES family compressible
below poly$(n)$ by a non-$M^*$ predictor) is \emph{not} excluded,
and would require a matching lower bound to rule out. So the honest
status of the amortised route to OP2(b)-APX is \emph{open}, not
closed --- the direct grammar-size transfer fails, but no proof
that $L_{\mathrm{amort}}$ is APX-easy exists. The cryptographic
route via Liu--Pass / OWF (§13.2) separately yields average-case
decision hardness, \emph{not} worst-case multiplicative
APX-hardness. Net: no worst-case combinatorial route is known to
\emph{succeed}, and none is proven \emph{impossible}; OP2(b)-APX
remains genuinely open, with Lemma~13.4.2a delimiting what the
family-absorption argument can and cannot show.
\end{enumerate}

\emph{What would constitute closure:} APX-hardness of the
amortised-family bit-cost objective with an explicit constant,
under either standard P $\neq$ NP or cryptographic assumption.
\emph{Updated 2026-05-28; further corrected 2026-06:} items 1--3
of the attack-vector list (predictor absorption of the reduction
algorithm in items~1--2, predictor bypass of the MSGP solver via
family reproduction in item~3) are unified by Lemma~13.4.2a, which
shows the $\Pi$-objective gap is \emph{absorbed} into $L(M^*)$ ---
i.e.\ these routes do not gap-preservingly transfer. As corrected
after adversarial review (see Lemma~13.4.2a), this is an
\emph{upper-bound} obstacle, not a proof that $L_{\mathrm{amort}}$
is APX-easy: a residual gap from differential family
compressibility is not excluded, and the PRG construction
(Theorem~13.4.2b) in fact neutralises the absorption, reducing the
amortised question to single-string OP2(a). So no route in this
list is proven to \emph{succeed}, and none is proven \emph{closed};
OP2(b)-APX is open.
The cryptographic route (Liu--Pass / OWF, §13.2) closes a
different problem (average-case decision on samplable distributions),
not free-predictor worst-case APX-hardness. Structurally novel
reductions outside the SGP/MSGP/grammar template --- specifically,
reductions whose hardness object is \emph{not} a poly-time-samplable
family with $\mathrm{poly}(n)$-bit description (so that the
$L(M^*) = O(\mathrm{poly}(n))$ absorption of Lemma~13.4.2a does
not apply) --- would qualify. Candidate routes: hard distributions
that are themselves Kolmogorov-incompressible (only available via
oracle access or pseudo-random generators, both falling outside
poly-time-samplable inputs), or hardness objects whose description
is super-polynomial in $n$.

\emph{Status update 2026-05-29 (PRG-conditional): the cryptographic
escape from Lemma~13.4.2a yields average-case \textbf{decision}
hardness, not worst-case multiplicative APX-hardness; OP2(b)-APX
remains open.} The previous update isolated one surviving escape from
the Lemma~13.4.2a absorption: a hardness object that is
\emph{Kolmogorov-incompressible} so that no $\mathrm{poly}(n)$-bit
predictor $M^*$ can reproduce the family. The only complexity-theoretic
mechanism producing such an object inside a \emph{poly-time-samplable}
distribution is a cryptographic pseudo-random generator (PRG): under
the assumption that secure PRGs exist (equivalently, one-way functions
exist; H{\aa}stad--Impagliazzo--Levin--Luby 1999, \emph{SIAM J.\ Comput.}\
28(4):1364--1396), a samplable distribution can produce strings that
are computationally incompressible --- indistinguishable from uniform
by any polynomial-time $M$. We formalise what this escape does and does
not deliver. We adopt throughout the joint Type-III-C amortised cost
$L_{\mathrm{amort}}(M,D,\{R_i\}) := L(M)+L(D)+\frac1N\sum_{i=1}^N
L(R_i \mid Y, M, D)$ of Lemma~13.4.2a, with side information $Y$
retained (a prior attempt --- the abandoned 2026-05-27 Lemma~6.8g ---
silently dropped $Y$ and reduced to a different, side-info-free
problem; we do not repeat that).

\emph{(1)~The construction and why absorption genuinely fails.}
Fix a secure length-doubling PRG $G:\{0,1\}^n \to \{0,1\}^{\ell(n)}$
with stretch $\ell(n)=n^{c}$, $c>1$. Draw $N=N(n)=2^{\Theta(n)}$
independent seeds $\sigma_i \stackrel{\$}{\leftarrow}\{0,1\}^n$ and a
fresh Charikar SGP-hard prefix $p$ of length $m(n)$ carrying the
$\gamma_{\mathrm{SGP}}=8569/8568$ grammar-size gap (the same prefix on
every instance: the YES/NO bit is a global property of $p$, not of
$i$). Set $s_i := G(\sigma_i)\,\Vert\,p$ and let $\mathcal{F}_n=\{s_i\}$
be sampled from $D$. \emph{Absorption fails:} the family-storing
predictor $M^*=\langle d_n,\mathrm{Red}_n\rangle$ of Lemma~13.4.2a
requires a $\mathrm{poly}(n)$-bit description that reproduces every
$s_i$; but reproducing $s_i$ requires reproducing $G(\sigma_i)$, and
the $N$ seeds carry $N\cdot n = 2^{\Theta(n)}$ bits of fresh entropy.
No $\mathrm{poly}(n)$-bit $M$ can encode them, and by PRG security no
$\mathrm{poly}(n)$-time $M$ can predict $G(\sigma_i)$ from a short
description with non-negligible advantage. Hence the
$L(M^*)=O(\mathrm{poly}(n))$ collapse of Lemma~13.4.2a does
\emph{not} apply to $\mathcal{F}_n$: this is a genuine,
non-tautological escape from the structural obstacle. So far so good.

\emph{(2)~But the construction does not \emph{manufacture} a
multiplicative gap --- it merely \textbf{reduces OP2(b)-amortised to
OP2(a)-single-string bit-cost} on the prefix, leaving the multiplicative
gap conditional on the \emph{open} Conjecture~6.8d, not on the PRG.} We
must analyse, honestly, what the incompressible family does to the gap.
Because the suffixes are incompressible, the optimal free encoder stores
the seed $\sigma_i$ ($n$ bits) in $R_i$ and recomputes $G(\sigma_i)$ in
the decoder (the PRG \emph{expands}, so storing the $n$-bit seed beats
storing the $\ell(n)$-bit output). The shared prefix $p$ is identical
across $i$, so the optimal predictor places \emph{one} program for $p$
inside $M$ --- and, $M$ being a \emph{free} predictor, this program is
the unrestricted shortest time-bounded grammar/program for $p$, whose
length is by definition the OP2(a) unrestricted bit-cost SGP optimum
$L^{\mathrm{SGP}}_{\mathrm{bit}}(p)$. Writing the objective out
\emph{with the correct weighting} ($L(M)$ and $L(D)$ appear undivided;
only the residual average is divided by $N$):
\[
  L_{\mathrm{amort}}
  \;=\; \underbrace{L^{\mathrm{SGP}}_{\mathrm{bit}}(p)}_{\subseteq\, L(M),\ \text{undivided}}
  \;+\; \underbrace{O(1)}_{L(D),\,G\text{-code}}
  \;+\; \underbrace{\tfrac1N\textstyle\sum_{i=1}^N \big(n+o(n)\big)}_{\text{seed, per instance}}
  \;=\; L^{\mathrm{SGP}}_{\mathrm{bit}}(p) \;+\; n \;+\; o(n).
\]
\textbf{Correction of a tempting error:} the prefix term is
$\Theta(m)$ \emph{undivided}, \emph{not} $\Theta(m/N)$ --- it lives in
$L(M)$, which is not amortised. (A first draft of this note divided it
by $N$ and wrongly concluded a $1+o(1)$ washout; that arithmetic was
incorrect.) Consequently the YES/NO amortised ratio is
\[
  \frac{L_{\mathrm{amort}}^{\mathrm{NO}}}{L_{\mathrm{amort}}^{\mathrm{YES}}}
  \;=\;\frac{\gamma_{\mathrm{SGP}}\,L^{\mathrm{SGP}}_{\mathrm{bit}}(p)+n+o(n)}
           {L^{\mathrm{SGP}}_{\mathrm{bit}}(p)+n+o(n)},
\]
which, for a Charikar prefix with $L^{\mathrm{SGP}}_{\mathrm{bit}}(p)
=\Theta(m)=\omega(n)$ (take $m=n^2$, so the additive seed term $n$ is
negligible), tends to $\gamma_{\mathrm{SGP}}$. \emph{So the gap survives
multiplicatively --- but only exactly as much as it survives for the
underlying single string $p$.} The PRG suffix has not created hardness;
it has only neutralised the \textbf{family-absorption} objection of
Lemma~13.4.2a, reducing the amortised-family question to the
\emph{single-string} unrestricted bit-cost SGP question on $p$, plus an
additive incompressible shift. The multiplicative gap on
$L_{\mathrm{amort}}$ holds \emph{if and only if} unrestricted bit-cost
SGP has a multiplicative gap on $p$ --- which is \textbf{exactly
Conjecture~6.8d (OP2(a)-APX)}, an \emph{open} problem blocked by the
ceiling-discontinuity obstacle of §13.4.1, conditional on
$\mathrm{P}\neq\mathrm{NP}$ and \emph{not} on the PRG. The PRG
assumption contributes nothing to closing this gap; it only removes the
absorption objection. Hence the honest implication is the \emph{relative}
one
\[
  \text{PRG exists} \;\wedge\; \text{Conjecture~6.8d (OP2(a)-APX)}
  \;\Longrightarrow\; \text{OP2(b)-amortised APX-hardness},
\]
which still rests on the open combinatorial conjecture; it is \emph{not}
a PRG-only closure of OP2(b)-APX.

\emph{(2$'$)~The alternative --- distinct per-instance prefixes ---
gives nothing PRG-only either.} If instead each $s_i$ carries its own
SGP-prefix $p_i$ (so the prefixes cannot be shared in $M$), two
exhaustive sub-cases arise. (a)~If $\{p_i\}$ is produced by a poly-time
reduction $\mathrm{Red}(d_n,i)\mapsto p_i$, then $\{p_i\}$ has a
$\mathrm{poly}(n)$-bit description and Lemma~13.4.2a \emph{absorbs} it
($M^*=\langle d_n,\mathrm{Red}\rangle$, $R_i=(i,\sigma_i)$,
$L_{\mathrm{amort}}=O(\mathrm{poly}(n))+n$); the gap dies symmetrically,
and the PRG suffix is irrelevant. (b)~If $\{p_i\}$ is itself
incompressible \emph{with a planted YES/NO gap}, then planting a
guaranteed multiplicative gap requires knowing the SGP answer, which is
a short description (returning to case~(a)) --- or it requires a
\emph{samplable distribution of SGP instances that is APX-hard on
average}. The latter is a strictly stronger open problem in average-case
complexity (not known to follow from $\mathrm{P}\neq\mathrm{NP}$, and
\emph{not} supplied by PRG existence). So no configuration --- shared
prefix (case~2) or distinct prefix (case~2$'$) --- extracts
OP2(b)-amortised APX-hardness from PRG existence \emph{alone}.

\emph{(3)~What the route \textbf{does} give: average-case decision /
additive-$K^t$ hardness.} The honest positive statement is the one
already recorded in §13.2's Liu--Pass paragraph, now anchored to this
construction. Under the OWF/PRG assumption, Liu \& Pass (FOCS 2020 / J
ACM 2023, arXiv:2009.11514) give the equivalence OWF $\Leftrightarrow$
mild average-case hardness of the Minimum-$K^t$-Complexity Problem
(MKtP) on samplable distributions, for the time-bounded measure
$K^t(x)=\min\{|p|:U(p)=x$ in $\le t(|x|)$ steps$\}$ with $t$ a fixed
polynomial. By Lemma~6.8g the joint Type-III-C cost satisfies
$L(M)+L(D)+L(R\mid Y,M,D)\le K^t(s\mid Y)+c\log_2(|s|+|Y|+1)+c$.
Hence:

\textbf{Theorem 13.4.2b (PRG-conditional average-case decision hardness
of joint Type-III-C, no APX).} \emph{Assume secure pseudo-random
generators exist. Then there is a polynomial-time samplable
distribution $D$ over sources $s$ (with side information $Y$) and a
threshold $\theta=\theta(|s|)$ such that no probabilistic
polynomial-time algorithm, given $(s,Y)$, decides whether the joint
Type-III-C cost $\min_{M,D,R}\{L(M)+L(D)+L(R\mid Y,M,D)\}$ is
$\le\theta$ or $>\theta+\omega(\log|s|)$, with success probability
$\ge\tfrac12+1/\mathrm{poly}(|s|)$ over $s\sim D$. The hardness is
\textbf{average-case} (over $D$), \textbf{decision} (a threshold gap of
additive size $\omega(\log|s|)$ in $K^t$-bits), and does \textbf{not}
provide any multiplicative inapproximability factor $\gamma'>1$ on
$L_{\mathrm{amort}}$.}

\emph{Proof sketch.} $D$ is the MKtP-hard samplable distribution of Liu
\& Pass (the 2021 extension, ECCC TR21-082, gives the statement for an
arbitrary samplable $D$ from OWFs; the 2020 FOCS version suffices for
the canonical PRG-output $D$ used here). The reduction maps an MKtP
instance $x$ to the source $s=x$, $Y=\bot$, so that deciding the joint
Type-III-C threshold decides MKtP via the two-sided sandwich of
Lemma~6.8g (upper bound) and the trivial $K^t(s\mid Y)\le
L_{\mathrm{total}}+O(\log)$ (lower bound, Lemma~6.8g scope statement).
The gap is the additive $K^t$-threshold gap of MKtP; it is not
multiplicative because both sides of the sandwich are additive in
$\log$. $\Box$

\emph{Scope note on $Y$ (addressing a prior pitfall).} The proof sets
$Y=\bot$, so Theorem~13.4.2b is, strictly, average-case decision
hardness of the \emph{side-information-free} specialisation
$\min_{M,D,R}\{L(M)+L(D)+L(R\mid M,D)\}$. The 2026-05-27 abandoned
Lemma~6.8g made the same $Y=\bot$ substitution but \emph{silently},
overclaiming it as full OP2(b); we instead state it openly. The
side-info-free case is a legitimate sub-problem: a free predictor with
$Y=\bot$ is a special case of one with arbitrary $Y$, so hardness here
lower-bounds the worst-case-over-$Y$ joint Type-III-C problem (side
information can only reduce cost, never the adversary's job of \emph{deciding}
the $Y=\bot$ instance). It does \emph{not} establish hardness for any
\emph{particular} informative $Y$, and the §11 natural-data regime ---
where $Y$ is a genuine context, not $\bot$ --- is untouched. This is a
real scope restriction, recorded honestly rather than papered over.

\emph{Honest delineation (what is and is not closed).}
\begin{itemize}
\item Theorem~13.4.2b is \textbf{not} a relabelling of the abandoned
Lemma~6.8g: it (i)~keeps $K^t$ (Liu--Pass), not Levin $\mathrm{Kt}$
(the 2026-05-27 fatal misnaming); (ii)~states \textbf{decision},
\textbf{not} APX (the 2026-05-27 fatal conflation --- the very error
this update exists to avoid repeating); (iii)~retains $Y$ in the
objective statement (the proof's $Y=\bot$ specialisation is flagged
above, not hidden). It is the precise numbered form of §13.2's prose,
no stronger.
\item It is \textbf{average-case}, over a \textbf{PRG-output}
distribution $D$. As noted in the OP2 attempts journal, hardness on
PRG-output is of limited interest for the natural-data compression of
§11; Theorem~13.4.2b is a complexity-theoretic anchor, not a statement
about enwik9-style corpora.
\item \textbf{OP2(b)-APX (Conjecture~6.8d's worst-case multiplicative
form) remains open.} Step~(2) shows the PRG escape, although it
genuinely defeats Lemma~13.4.2a's family-absorption, does not by itself
\emph{create} a multiplicative gap: with a shared SGP-prefix it
\emph{reduces} OP2(b)-amortised to single-string OP2(a)-bit-cost (plus
an additive incompressible seed shift), so the multiplicative gap is
exactly Conjecture~6.8d --- open, and conditional on
$\mathrm{P}\neq\mathrm{NP}$ rather than on the PRG. With distinct
prefixes (step~2$'$) it either re-collapses into Lemma~13.4.2a
absorption or requires an average-case-APX-hard SGP distribution that
PRG existence does not supply. The PRG assumption removes the absorption
objection but contributes no multiplicative gap; the cryptographic route
is intrinsically additive/decision.
\end{itemize}

\emph{Verdict.} The PRG-conditional route is the unique escape from the
Lemma~13.4.2a worst-case-combinatorial obstacle, and it closes a
\emph{different} problem cleanly: average-case \textbf{decision}
hardness of joint Type-III-C on samplable distributions, under PRG/OWF
existence (Theorem~13.4.2b). It does \textbf{not} yield free-predictor
worst-case multiplicative APX-hardness: defeating the family-absorption
of Lemma~13.4.2a with an incompressible PRG suffix merely reduces the
amortised-family problem back to single-string OP2(a)-bit-cost on the
SGP-prefix (step~2), whose multiplicative inapproximability is the
\emph{open} Conjecture~6.8d --- so the PRG removes one objection but
supplies no gap. Combined with Lemma~13.4.2a (all worst-case
combinatorial routes absorbed) and §13.4.2's item~3 (MSGP absorbed), the
complete current characterisation of OP2(b) is: \emph{worst-case
multiplicative APX-hardness is open and resists every known route
(the PRG escape reduces it to the equally-open OP2(a)-APX); the only
established conditional hardness is average-case decision hardness under
PRG/OWF (Theorem~13.4.2b).}

\textbf{13.4.3 OP3$''$ (Efficient enumerative encoders for specific
natural-data classes).}

\emph{Reduction to membership-oracle construction (Cover 1973).}
Given polynomial-time membership oracle for class $\mathcal{C}
\subseteq \Sigma^N$, an explicit polynomial-time rank coder follows
(documented in §13.2). The residual challenge is constructing the
membership oracle for specific natural classes.

\emph{Attack vectors:}
\begin{enumerate}
\item \emph{Formal-language classes.} For any regular language
$\mathcal{L}$ given by a DFA with $q$ states, membership is
$O(N q)$. By Cover 1973, polynomial enumerative encoder for
$\mathcal{L} \cap \Sigma^N$ with rate
$\log_2 |\mathcal{L} \cap \Sigma^N| / N$ bits per symbol. Concrete
example: valid RFC-3986 URIs, ISO-8601 dates, simple JSON syntax
all admit polynomial enumerative encoders.
\item \emph{Context-free classes.} For any context-free language with
grammar size $g$, the CYK parser gives $O(N^3 g)$ membership and
hence polynomial enumerative encoder. Concrete example: well-formed
HTML5 fragments, valid Python-3 source code (modulo indentation).
Lemma 13.4.3a (below) gives the explicit construction.
\item \emph{Algebraic / lattice classes.} For lattice membership
$\{x \in \mathbb{Z}^N : Ax = b \pmod q\}$ for fixed $(A, b, q)$,
polynomial enumerative encoder via lattice basis reduction.
Connects compression to lattice-based cryptography.
\item \emph{Neural-classifier surrogate.} Define
$\mathcal{C}_{f, \tau} := \{x : f(x) > \tau\}$ where $f$ is a
polynomial-time neural classifier. The rate $\log_2 |\mathcal{C}_{f,
\tau}| / N$ depends on $f$'s false-positive rate $\alpha$: for
$\alpha < |\mathcal{C}_{\rm true}| / |\Sigma|^N$, the encoder
achieves near-optimal rate; for larger $\alpha$, rate degrades
proportionally.
\end{enumerate}

\emph{What would constitute closure:} explicit construction of a
polynomial-time enumerative encoder for natural images at fixed
resolution, achieving rate matching state-of-the-art neural
compressors (e.g., HiFiC, MENT) within constant factors.

\textbf{Lemma 13.4.3a (Constructive enumerative encoder for
unambiguous context-free languages; sub-case closure of OP3$''$).}
\emph{Let $G$ be an unambiguous context-free grammar in Chomsky Normal
Form with $k$ nonterminals (start symbol $S$ indexed $0$), terminal
alphabet $\Sigma$ of size $\sigma = |\Sigma|$, and production set $P$
partitioned into $P_2 = \{A \to BC\}$ binary and $P_1 = \{A \to a\}$
terminal productions. For any target length $N \in \mathbb{N}$, the
class $\mathcal{C}_G := L(G) \cap \Sigma^N$ admits an explicit
polynomial-time enumerative encoder $(\mathrm{count}, \mathrm{rank},
\mathrm{unrank})$ such that}
\begin{itemize}
\item $\mathrm{count}(N) = |\mathcal{C}_G|$ \emph{is computed in
$O(N^2 \cdot k \cdot |P_2|)$ arithmetic operations on integers of
$O(N \log \sigma)$ bits;}
\item $\mathrm{rank}(w) := |\{v \in \mathcal{C}_G : v < w \text{ lex}\}|$
\emph{is computed in $O(N^4 \cdot \sigma \cdot k \cdot |P_2|)$
operations for any $w \in \mathcal{C}_G$;}
\item $\mathrm{unrank}(r) := $ \emph{the lex-$r$-th element of
$\mathcal{C}_G$ for $r \in [0, |\mathcal{C}_G|)$, computed in
$O(N^4 \cdot \sigma \cdot k \cdot |P_2|)$ operations.}
\end{itemize}
\emph{The encoder transmits $\lceil \log_2 |\mathcal{C}_G| \rceil$ bits
per message, achieving rate}
\[
R_G(N) \;=\; \frac{\lceil \log_2 |\mathcal{C}_G| \rceil}{N}
\;\leq\; \frac{\log_2 |\mathcal{C}_G|}{N} + \frac{1}{N},
\]
\emph{within $1/N$ bit/symbol of the Cover 1973 enumerative-coding
optimum on $\mathcal{C}_G$.}

\emph{\textbf{Proof / construction.}} We give the algorithm explicitly.

\emph{Step 1 (count table, CYK-counting variant of Younger 1967).} For
each nonterminal $A \in \{0, \ldots, k-1\}$ and each $j \in \{1, \ldots,
N\}$, define $N_A[j] := \#\{w \in \Sigma^j : A \Rightarrow_G^\ast w\}$.
By unambiguity of $G$, each derivable string has a unique parse tree,
so the parse-tree count equals the distinct-string count, and the
following CYK-style recurrence is exact:
\[
N_A[1] \;=\; |\{a \in \Sigma : (A \to a) \in P_1\}|,
\]
\[
N_A[j] \;=\; \sum_{(A \to BC) \in P_2} \sum_{m=1}^{j-1} N_B[m] \cdot N_C[j-m]
\quad (j \geq 2).
\]
The table has $kN$ entries; filling each entry takes $O(|P_2| \cdot N)$
arithmetic ops, for total $O(N^2 \cdot k \cdot |P_2|)$ ops. Since
$N_A[j] \leq \sigma^j \leq \sigma^N$, each integer has
$O(N \log \sigma)$ bits. (For $G$ ambiguous, the recurrence overcounts
strings; the algorithm and lemma assume unambiguity, which is
undecidable in general for arbitrary CFG but \emph{structurally
verified} for the natural CFLs we target --- Dyck words, balanced
parenthesis, regular languages converted via right-linear CNF, JSON
syntax, RFC-3986 URIs, and the LL/LR-parsable subclasses of
programming languages.)

\emph{Step 2 (prefix-completion count via wildcard CYK).} For any
pattern $\pi \in (\Sigma \cup \{\star\})^N$ where $\star$ denotes a
wildcard, define
\[
D_A^\pi[a, b] \;:=\; \#\bigl\{ s \in \Sigma^{b-a} : A \Rightarrow_G^\ast s,
\; s[i] = \pi[a + i] \text{ for all } i \text{ with }
\pi[a + i] \neq \star \bigr\}.
\]
Recurrence:
\[
D_A^\pi[a, a+1] \;=\; \sum_{(A \to t) \in P_1}
\mathbb{1}[\pi[a+1] = \star \;\vee\; \pi[a+1] = t],
\]
\[
D_A^\pi[a, b] \;=\; \sum_{(A \to BC) \in P_2} \sum_{m=a+1}^{b-1}
D_B^\pi[a, m] \cdot D_C^\pi[m, b]
\quad (b - a \geq 2).
\]
This fills $O(N^2)$ cells per nonterminal in $O(N |P_2|)$ work each,
total $O(N^3 \cdot k \cdot |P_2|)$ ops per pattern. The quantity
$\mathrm{Cnt}(p) := D_0^\pi[0, N]$ for $\pi = p \cdot \star^{N - |p|}$
gives the number of length-$N$ strings in $\mathcal{C}_G$ starting
with prefix $p$.

\emph{Step 3 (rank by prefix enumeration).} For $w \in \mathcal{C}_G$,
\[
\mathrm{rank}(w) \;=\; \sum_{i=0}^{N-1} \sum_{c < w_{i+1}}
\mathrm{Cnt}(w[1..i] \cdot c).
\]
There are at most $N \cdot (\sigma - 1)$ summands; each summand is one
wildcard-CYK invocation costing $O(N^3 k |P_2|)$, giving the stated
$O(N^4 \cdot \sigma \cdot k \cdot |P_2|)$ bound. Correctness: every
$v < w$ in $\mathcal{C}_G$ shares a (possibly empty) maximal prefix
with $w$, then disagrees at some position $i+1$ with $v_{i+1} < w_{i+1}$;
the summand at $(i, v_{i+1})$ counts $v$ exactly once.

\emph{Step 4 (unrank by lex prefix search).} Given $r \in [0,
|\mathcal{C}_G|)$, greedy reconstruct $w$ symbol-by-symbol. At position
$i+1$, maintain the residual rank $\rho_i$ (initially $r$); enumerate
$c \in \Sigma$ in lex order; let $q_c := \mathrm{Cnt}(w[1..i] \cdot c)$;
pick the smallest $c^\ast$ for which $\rho_i < \sum_{c \leq c^\ast} q_c$;
set $w_{i+1} = c^\ast$ and $\rho_{i+1} := \rho_i - \sum_{c < c^\ast} q_c$.
Per position cost $O(\sigma)$ wildcard-CYK invocations; over $N$
positions, the stated bound. Correctness: at each step the residual
$\rho_i$ identifies a sub-interval of size $1$ within the $q_{c^\ast}$
extensions, which by the rank construction is exactly the lex-$(r)$-th
completion conditional on prefix $w[1..i]$.

\emph{Step 5 (rate).} The encoder transmits $\mathrm{rank}(w)$ as a
fixed-length integer in $b_N := \lceil \log_2 |\mathcal{C}_G| \rceil$
bits. Since $\mathrm{rank}: \mathcal{C}_G \to [0, |\mathcal{C}_G|)$ is
a bijection (uniqueness of lex order on a finite set; correctness of
$\mathrm{rank}, \mathrm{unrank}$ above), the encoding is lossless. The
rate is $b_N / N \leq (\log_2 |\mathcal{C}_G| + 1)/N$, matching the
Cover 1973 ceiling within $1/N$ bit per symbol. \hfill$\blacksquare$

\emph{Scope and significance.} Lemma 13.4.3a upgrades the named
``Context-free classes'' attack vector from a sketch (``CYK gives
polynomial membership, so Cover 1973 applies'') to an executable
construction. The novelty lies not in the theory (Cover 1973 plus
CYK-counting is folklore; cf.\ Goulden-Jackson 1983 \emph{Combinatorial
Enumeration of Sequences} for generating-function methods on
unambiguous CFLs) but in the \emph{explicit polynomial-time algorithm}:
for any unambiguous-CNF grammar provided as input, the construction
yields a working encoder that researchers can adapt to their specific
data class.

The restriction to \emph{unambiguous} CFGs is essential: counting
distinct length-$N$ strings in an ambiguous CFL is \#P-hard in
general --- Bertoni--Goldwurm--Sabadini (1991, \emph{Theoret.\
Comput.\ Sci.}\ 86(2):325--342) show the counting function is
computable in $\mathrm{NC}^2$ for unambiguous CFLs but \#P-hard
already for languages of ambiguity degree two (specifically, unions
of two unambiguous linear CFLs), so the unique-parse hypothesis that
makes the recurrence of Step 1 exact is exactly the boundary of
tractability. Most natural data classes admit
unambiguous CFG presentations (Dyck words, balanced parenthesis,
RFC-3986 URIs, ISO-8601 dates as a regular subset, and the LL/LR-parsable
subclasses of programming languages); ambiguity-free presentations of
HTML5 and full Python are open engineering questions distinct from the
complexity-theoretic content of Lemma 13.4.3a.

\emph{Concrete benchmark (Fibonacci-string language).} The binary
language $F = \{w \in \{0,1\}^\ast : \text{``11''} \not\sqsubseteq w\}$
admits a four-nonterminal unambiguous CNF grammar (start $S$,
``after-1'' state $T$, helpers $A \to 0$, $B \to 1$, with productions
$S \to AS \mid BT \mid 0 \mid 1$, $T \to AS \mid 0$). By a standard
Zeckendorf argument, $|F \cap \{0,1\}^N| = \mathrm{Fib}(N+2)$, with
asymptotic rate $\log_2((1+\sqrt 5)/2) \approx 0.6942$ bits/symbol.
At $N = 20$, $|F \cap \{0,1\}^{20}| = \mathrm{Fib}(22) = 17{,}711$;
the algorithm transmits $\lceil \log_2 17711 \rceil = 15$ bits per
length-20 message for rate $0.75$ bits/symbol, within $1/N = 0.05$
of the finite-$N$ optimum $\log_2(17711)/20 \approx 0.7056$. The
script \texttt{scripts/verify/op3\_cfl\_enumerative\_encoder.py}
implements all three operations $(\mathrm{count}, \mathrm{rank},
\mathrm{unrank})$ and verifies (i) count consistency with the
closed-form Fibonacci value for $N \in \{1, \ldots, 30\}$, (ii) full
lex-enumeration round-trip for $N \leq 15$, (iii) 100 random
round-trips at $N = 20$, (iv) rank uniqueness and range, and (v) the
empirical $0.75$-bits/symbol rate against the optimum.

\emph{What this closes and what remains.} Lemma 13.4.3a constructively
closes OP3$''$ for the sub-case of \emph{unambiguous CFG-definable
data classes}. The outer envelope of OP3$''$ --- natural data classes
such as natural images, video, audio, where no explicit generative
grammar is known --- remains open; constructing unambiguous CFG (or
PCFG) approximations to these classes is the remaining engineering
challenge that the Lemma reduces the open problem to. The complexity
constant $O(N^4 \sigma k |P_2|)$ in the rank step is the
straightforward bound; standard incremental-CYK and Earley-set
optimisations can reduce this to $O(N^3 \sigma k |P_2|)$ at the cost
of more intricate bookkeeping (sharing pattern-CYK cells across the
$N$ prefix-positions, since consecutive patterns differ in exactly
one symbol). The optimised bound is not required for the
``polynomial-time'' guarantee that establishes the closure.

\textbf{13.4.4 OP4 (Approximability of optimal support selection for
general dependent joint laws of $M_v$).}

\emph{What has been tried (failed paths):}
\begin{itemize}
\item NP-hardness via CLIQUE reduction --- works (Lemma 5.7c), but
the positive-affine reduction $\mathbb{E}_\mu[A_s] = A_s(1) -
(A_s(1) - A_s(0)) \cdot n_0(S)/m$ does not preserve multiplicative
approximation ratios.
\item Dense-DkS partial closure (Lemma 5.7d) --- gives polynomial gap
in the dense regime; Manurangsi 2017 sparse-DkS ETH-hardness does
not transfer (gap dilutes to $1 + o(1)$ under natural log-binomial
$A_s$).
\item \emph{$E_{12}$ sub-locus (Hamming-restricted labeling; see §13.2
Conj.\ 4.3e and Remark 4.3d for the full treatment).} Gap-amplification
from NCP / CVP (Arora-Babai-Stern-Sweedyk 1997; Dinur-Kindler-Raz-Safra
2003), $0$-extension / metric-labeling (Karzanov 1998), and a
row-stochastic-source / multiway-cut reduction were all considered ---
none lands a gap-preserving map into the injective-into-$Q_m$ geometry.
\emph{Corrected obstacle framing (per Remark 4.3d, validated; supersedes
the earlier ``row-stochastic $F$ is the obstacle'' reading):} the
blocker is \emph{not} row-stochasticity (the ratio $\rho(F)=\mathrm{OPT}/2w$
is rescaling-invariant, Remark 4.3d(ii)) and \emph{not} affine dilution
(the cost is a quadratic assignment functional, gap survives
normalization), but the absence of a \emph{gap-amplified hypercube-embedding
source} exhibiting a YES dilation-$1$ instance together with a
constant-fraction robust-NO instance (Remark 4.3d(iii)--(iv)). The natural
candidates fail for structurally distinct reasons: partial cubes are
poly-time recognizable (so embeddability alone is not NP-hard;
Djokovi\'c 1973 / Winkler 1984 / Eppstein 2008), and dense
constant-stretch graphs (e.g.\ $K_{2,3}$) are never hypercube subgraphs
(so they have no YES baseline). Do not re-attempt these routes; the
precise remaining step is constructing the robust-NO-promise embedding
source.
\end{itemize}

\emph{Attack vectors not yet tried:}
\begin{enumerate}
\item \emph{Use a different $A_s$ structure where the gap transfers.}
Search for monotone $A_s$ families where DkS gap multiplies (not
dilutes) into the support selection problem.
\item \emph{Reduction from Densest-$k$-Subhypergraph} (Chuzhoy-Manurangsi-
Schroeppel 2020), which has different gap structure than DkS and
may transfer better.
\item \emph{Average-case hardness via planted-clique.} If support
selection on random graphs with a planted dense subgraph is
average-case hard (under planted-clique conjecture, hardness for
$k \leq \sqrt{N}$ subgraphs), this gives conditional hardness.
\end{enumerate}

\emph{Status update 2026-05-29 (attack vector 2 closed negatively;
the DkSH gap dilutes \textbf{worse} than DkS, and a citation correction).}
We pursued attack vector~2 (Densest-$k$-Subhypergraph, DkSH) and
verified the literature, with one correction. The attribution
``Chuzhoy--Manurangsi--Schroeppel 2020'' above is a mis-citation: the
DkSH inapproximability is due to \textbf{Applebaum} (\emph{Pseudorandom
generators with long stretch and low locality from random local one-way
functions}, SIAM J.\ Comput.\ 42(5):2008--2037, 2013; ECCC TR11-007;
STOC 2012), and the nontrivial algorithmic upper bounds are due to
Chlamt\'ac--Dinitz--Konrad--Kortsarz--Rabanca (\emph{The Densest
$k$-Subhypergraph Problem}, APPROX-RANDOM 2016 / SIAM J.\ Discrete
Math.\ 32(2), 2018, arXiv:1605.04284). We adopt the corrected citations.

\emph{(i) Verified DkSH hardness.} Applebaum's $p$-Densest-Subhypergraph
($p$-DSH) is the promise problem on $d$-uniform hypergraphs ($d$ constant,
$n<m\le\mathrm{poly}(n)$) distinguishing: \textbf{YES} --- some node-set
$S$ of density $p$ (size $pn$) contains a fraction
$\ge p^{\,d-1}(1-o(1))$ of the hyperedges; \textbf{NO} --- every density-$p$
set contains a fraction $\le p^{\,d}(1+o(1))$. By Applebaum's
Thm.\ 7.1/1.2, for $m\ge n^{c}$ ($c>3$ constant) and every
$n^{-(c-3)/2d}\le p\le \tfrac12$, $p$-DSH is intractable assuming
$F_{m,Q}$ is $o(1/(\sqrt n\log n))$-pseudorandom (a cryptographic
assumption, following from one-wayness of $F_{m_0,Q}$ for a large
polynomial $m_0$). The multiplicative inapproximability factor is
$g=p^{\,d-1}/p^{\,d}=1/p$, maximised at $p=n^{-(c-3)/2d}$ to
$g=n^{(c-3)/2d}=n^{\varepsilon}$. So the hardness is conditional
(cryptographic, not $\mathrm P\neq\mathrm{NP}$ nor ETH) and the strong
$n^{\varepsilon}$ gap lives \emph{only} in the small-$p$ regime
$p=n^{-\Theta(1)}\to 0$; at $p=\tfrac12$ one recovers only the
constant-factor Feige--Khot ``quasi-random PCP'' hardness.

\emph{(ii) Reduction sketch and its dichotomy.} A DkSH instance maps to
Type-II by taking each hyperedge $T$ ($|T|=d$) as a $d$-ary atom and
identifying $n_0(S):=|\{T:T\subseteq S\}|$ (the contained-hyperedge
count, exactly the $p$-DSH objective) with the Type-II support
substructure, with $s:=pn$. For $\mathbb E_\mu[A_s(|T\setminus S|)]$ to
be a function of $n_0(S)$ \emph{alone} (the affine identity of
Lemma 5.7c, the only structure that lets a cost-minimiser decide the
combinatorial problem), one needs $A_s(1)=A_s(2)=\dots=A_s(d)$, i.e.\
$A_s$ \emph{flat above $0$}; then
$\mathbb E_\mu[A_s]=A_s(1)-(A_s(1)-A_s(0))\,n_0(S)/m$ as before. A
\emph{genuinely affine} $A_s(j)=a+bj$ ($b>0$, the natural
monotone-increasing code-length direction, since more out-of-support
coordinates cost more) instead yields
$\mathbb E_\mu[A_s]=a+b\,(dm-\sum_{v\in S}\deg v)/m$, whose minimiser is
the $s$ highest-degree vertices --- a trivially polynomial
max-degree-sum problem (greedy top-$s$) carrying no hardness (verified
by brute force, \texttt{scripts/verify}). Thus only the flat-above-$0$
$A_s$ encodes DkSH, and it lands in exactly the Lemma 5.7c/5.7d affine
regime.

\emph{(iii) Quantitative gap (does not survive; dilutes worse than DkS).}
Plugging $D_{\mathrm{YES}}/m=p^{\,d-1}$, $D_{\mathrm{NO}}/m=p^{\,d}$ into
the Lemma 5.7d dilution formula with normalised $A_s$
($A_s(0)=0,A_s(1)=1$, the most gap-favourable choice) gives
\[
\rho_{\mathrm{Type\text{-}II}}
=\frac{1-D_{\mathrm{NO}}/m}{1-D_{\mathrm{YES}}/m}
=\frac{1-p^{\,d}}{1-p^{\,d-1}}
=1+\frac{p^{\,d-1}(1-p)}{1-p^{\,d-1}}
=1+\Theta\!\bigl(p^{\,d-1}\bigr)
=1+\Theta\!\bigl(g^{-(d-1)}\bigr).
\]
The $n^{\varepsilon}$ DkSH gap is \emph{inverted and collapsed}: as the
DkSH gap $g=1/p\to\infty$ (the regime where it exceeds DkS), the Type-II
gap $\to 1$ \emph{polynomially}, and the larger the source gap the
smaller the transferred gap. This is the same affine-dilution obstacle
as Lemma 5.7d(b1), \emph{strengthened} by the exponent $d-1\ge 2$: the
$d=2$ (DkS) case gives $1+\Theta(g^{-1})$, whereas $d$-uniform DkSH gives
$1+\Theta(g^{-(d-1)})$, so moving from edges to hyperedges makes the
dilution \emph{worse}, not better. (Verified numerically and
symbolically across $d\in\{2,3,4,6\}$ and $p\in\{10^{-1},\dots,10^{-9}\}$,
\texttt{scripts/verify}.) The intuition is structural: Applebaum's gap is
already a \emph{fractional} (edge-density $p^{d-1}$ vs $p^{d}$) gap with
both branches $o(1)$ in the hard regime, exactly the sparse regime the
affine cost flattens; DkSH does not escape because, unlike DkS where a
large count-ratio can coexist with $D_{\mathrm{NO}}=1$, here both
densities are tied to the \emph{same} $p\to 0$.

\emph{(iv) Verdict: negative, with a unification.} DkSH fails to close
OP4 for the same reason DkS did, and quantitatively worse. The hyperedge
structure does \emph{not} produce a different functional form for the
Type-II objective: the contained-hyperedge count is just $n_0(S)$ with
$d$-ary atoms, so the affine identity and its dilution apply verbatim,
with the dilution exponent rising from $1$ to $d-1$. Attack vector~2 is
therefore closed negatively. The unified obstacle now reads: \emph{any
count-of-contained-substructures objective (DkS edges, DkSH hyperedges)
maps to Type-II through the affine $n_0(S)$ identity, so the transferred
multiplicative gap is always $1+\Theta(D_{\mathrm{YES}}/m)$ --- governed
by the YES-instance \emph{density} $D_{\mathrm{YES}}/m$, not by the source
gap factor. Whenever the hard instances are sparse ($D_{\mathrm{YES}}/m
\to 0$), the transferred gap collapses to $1$.} For Applebaum's
fractional-gap $p$-DSH family this density is exactly
$D_{\mathrm{YES}}/m=p^{\,r-1}=g^{-(r-1)}$ for $r$-uniform atoms (gap
$g=1/p$), so the larger the source gap the smaller the transferred gap,
and the $r\ge 3$ hypergraph case dilutes worse than the $r=2$ DkS case
by the extra exponent. Escaping OP4 hence requires a Type-II hardness
source that is \emph{not} a contained-substructure count --- e.g.\ an
entropy- or partition-functional objective whose minimiser does not
reduce to $\max n_0(S)$ --- as already flagged in Lemma 5.7d, route~(c).
Attack vectors~1 (gap-amplifying $A_s$) and 3 (planted-clique average
case) remain open; the present finding makes vector~1 look unpromising,
since flatness-above-$0$ is forced by the affine identity and a non-flat
$A_s$ trivialises the objective.

\emph{Status update 2026-05-29 (non-count objective; route~(c) closed
negatively \emph{within} Type-II, with a structural invariance
characterisation).} We pursued the route flagged immediately above and in
Lemma 5.7d route~(c): replace the contained-substructure \emph{count} by a
\emph{non-count} (entropy / partition-functional) objective, in the hope
that a concave/log functional escapes the affine $n_0(S)$ dilution. We
verified the relevant entropy-hardness literature and reached a negative
verdict for Type-II, together with a clean structural reason that
\emph{extends the unified obstacle from counts to all partition
functionals}.

\emph{(i) The candidate non-count hard problems are genuinely non-affine.}
\textsc{Max-Entropy Sampling} (MESP; Ko--Lee--Queyranne, \emph{Oper.\ Res.}\
43(4):684--691, 1995) maximises $\tfrac12\log\det\Sigma_S$ over size-$s$
principal submatrices of a covariance $\Sigma$ and is NP-hard; the closely
related \textsc{Log-Determinant Maximisation} / DPP-MAP is NP-hard to
approximate within $5/4$ (unconstrained) and within $1+10^{-10^{13}}$
(size-constrained monotone) under $\mathrm P\neq\mathrm{NP}$, and the
\emph{unconstrained} DPP-MAP is inapproximable within $2^{\beta n}$,
$\beta=10^{-10^{13}}$ (Ohsaka, \emph{J.\ Artif.\ Intell.\ Res.}, 2022;
PMLR 130, 2021, improving the $9/8-\varepsilon$ of Kulesza--Taskar 2012).
Crucially these gaps live on a \emph{spectral} functional $\log\det K_S$
that does \emph{not} reduce to any count: subsets $S,S'$ with the same
induced-edge (correlated-pair) count have different $\log\det K_S$
(witnessed in \texttt{scripts/verify}, check~B), so the affine identity of
Lemma 5.7c simply does not apply and the multiplicative gap \emph{would}
survive. So a log-det objective is exactly the kind of non-count functional
sought.

\emph{(ii) But the Type-II cost is provably a count-functional, so no such
objective is the Type-II cost.} The actual Type-II support-selection cost
is the Theorem 5.3 enumerative residual
$\mathbb E_\mu[A_s(|T\setminus S|)]$ with
$A_s(\alpha)=\log_2\binom{s}{k_v-\alpha}+\log_2\binom{N-s}{\alpha}$. Each
atom $T$ contributes $A_s(|T\setminus S|)$, a function of the integer
\emph{overlap} $|T\cap S|$ \emph{alone}; hence \emph{any} functional of the
per-atom Type-II contributions --- the expected code length, the Shannon
entropy of the induced within-support overlap law, or any other
``entropy of the partition restricted to $S$'' --- is a function of the
(weighted) overlap histogram $(W_0,\dots,W_d)$, $W_\alpha:=\sum_{T:\,
|T\setminus S|=\alpha}\mu(T)$, and \emph{therefore} of $n_0(S)$ for
$2$-ary atoms at the Lemma 5.7c collapse point. Equivalently, the Type-II
cost is \emph{invariant} under every relabelling of positions/atoms that
preserves the overlap histogram, whereas a DPP/MESP $\log\det K_S$ is
\emph{not} invariant under such relabellings (it depends on the off-diagonal
kernel entries). The two functionals have \emph{different invariance
groups}, so no kernel $K$ satisfies $\text{Type-II-cost}(S)=\log\det K_S$
identically (\texttt{scripts/verify}, checks~A and~C). The discrete
entropy-of-partition objective the route proposed thus collapses to the
same affine $n_0(S)$ dilution as the count: its transferred gap is again
$1+\Theta(D_{\mathrm{YES}}/m)$, vanishing on sparse instances.

\emph{(iii) Where the non-count gap really lives: Type-III-C, not Type-II.}
The log-det functional \emph{does} arise in the RNR framework --- it is
exactly the Gaussian total correlation $\mathrm{TC}(X)=\tfrac12[\sum_i
\log_2(1+a_i)-\sum_j\log_2(1+b_j)]$ of Lemma 5.6j, a \emph{differential}
entropy of a continuous latent (Cover--Thomas 2006, Thm 8.4.1). By the
type taxonomy (Definitions 3.2 / 3.5; the discussion at the end of \S 1.2:
``Type-II as defined in \S 3.2 is partition-based,'' while a continuous
latent ``fits into Type-III-C''), a log-det cost is a Type-III-C
(continuous-dictionary / variational-posterior) cost, \emph{not} a Type-II
partition cost. Importing MESP/DPP hardness as a ``Type-II objective''
would therefore require silently replacing the discrete partition cost by an
arbitrary spectral functional --- precisely the substitution the anti-
tautology requirement forbids. The non-count gap is real, but it belongs to
the continuous (Type-III-C) optimal-support question, which is a
\emph{different} open problem from OP4.

\emph{(iv) Verdict: negative for Type-II, extending the unified obstacle to
all partition functionals.} Route~(c) does not close OP4. The unified
obstacle now reads in its strongest form: \emph{the Type-II cost is a
function of the overlap histogram, hence invariant under
histogram-preserving relabellings; every objective with this invariance
(every count, and every entropy/partition-functional built from the
discrete partition cost) maps through the affine $n_0(S)$ identity and
inherits the density-governed dilution $1+\Theta(D_{\mathrm{YES}}/m)$,
collapsing to $1$ on sparse hard instances. The only functionals that
escape (MESP/DPP $\log\det K_S$) lack this invariance and are therefore not
the Type-II cost --- they are continuous Type-III-C costs.} This does
\emph{not} weaken the prospect of OP4 closure by other means (a non-affine
gap-amplifying $A_s$ that nonetheless encodes a hard count, vector~1; or an
average-case route, vector~3, both still open); it only rules out the
``swap the count for an entropy'' route by showing the swap leaves the
Type-II invariance class. Verified in
\texttt{scripts/verify/op4\_noncount\_entropy\_objective.py} (checks A, A$'$,
B, B$'$, C: discrete entropy/enumerative objective is a count-functional and
affine in $n_0(S)$; MESP $\log\det$ is a genuine non-count functional whose
gap survives but is a continuous cost; the Type-II cost is a sum of
log-binomials of integer counts with the wrong invariance group to be any
DPP $\log\det$).

\emph{Status update 2026-05-29 (the dense-vs-sparse residual of
Lemma 5.7d is resolved \textbf{negatively for the positive branch}: no
natural Type-II family yields a polynomial inapproximability).} Lemma 5.7d
left precisely one residual open: ``open whether any natural Type-II family
lies in [the dense] regime'' where the multiplicative gap survives. We close
it. The resolution has two prongs and a structural witness.

\emph{(i) Two independent reasons the positive (dense) branch carries no
polynomial gap for natural $A_s$.} \emph{Reason~A (unconditional, the primary
resolution).} Reading the exact dilution formula of Lemma 5.7d(b),
\[
\rho_{\mathrm{Type\text{-}II}}
= 1 + \frac{(D_{\mathrm{YES}} - D_{\mathrm{NO}})/m}
            {\frac{A_s(1)}{A_s(1) - A_s(0)} - D_{\mathrm{YES}}/m},
\]
a \emph{polynomial-in-$N$} gap can arise only if the denominator $\to 0$,
i.e.\ $D_{\mathrm{YES}}/m \to A_s(1)/(A_s(1) - A_s(0))$. For the
\emph{idealised} normalised cost $A_s(0) = 0$ this threshold is $1$ (a
diverging gap as $D_{\mathrm{YES}}/m \to 1$), but for the \emph{natural}
log-binomial $A_s(0) = \log_2\binom{s}{2} > 0$ the denominator is bounded
below by $A_s(0)/(A_s(1) - A_s(0)) = \Theta(\log s) > 0$ \emph{for all}
$D_{\mathrm{YES}}/m \le 1$, so $\rho_{\mathrm{Type\text{-}II}}$ saturates at
the \emph{constant} $A_s(1)/A_s(0) = 1 + O(1/\log s)$ \emph{regardless of the
density}. Under the genuine code cost the gap is therefore $O(1)$ in
\emph{every} regime, dense or sparse --- the dense regime cannot produce a
polynomial gap at all. \emph{Reason~B (the idealised-$A_s$ escape hatch is
also closed).} Even granting the non-physical $A_s(0) = 0$, the regime where
its formula diverges is $D_{\mathrm{YES}}/m \to 1$ --- which is
\emph{disjoint} from Manurangsi's hard regime $D_{\mathrm{YES}}/m \to 0$, so
the diverging idealised ratio is never transferred from a Manurangsi-hard
source. Concretely, $D_{\mathrm{YES}}/m \to 1$ splits into two sub-cases,
both gap-free. \emph{(B1)} If additionally the graph is \emph{everywhere
dense} ($m = \Theta(N^2)$, $k = s = \Theta(N)$, average degree $\Omega(N)$),
then DkS \emph{itself} admits a polynomial-time $(1+\varepsilon)$-
approximation: Arora--Karger--Karpinski (STOC 1995; JCSS 58(1):193--210,
1999, \emph{Polynomial time approximation schemes for dense instances of
NP-hard problems}) give an additive-$\varepsilon N^2$ PTAS for dense graph
maximisation (average degree $\Omega(N)$, by exhaustive sampling), and since
the dense optimum is $\max n_0(S) = \Theta(N^2)$ the additive guarantee
converts to a multiplicative $(1+\varepsilon)$ factor (the DkS survey
arXiv:2303.14467 and Manurangsi--Raghavendra, STACS 2016, arXiv:1507.04391,
state DkS with $m = \Omega(N^2)$, $k = \Omega(N)$ has a PTAS). \emph{(B2)} If
instead $m = o(N^2)$ (a \emph{sparse} edge set) yet $D_{\mathrm{YES}}/m \to
1$, then a $(1-o(1))$-fraction of all edges concentrates inside the chosen
$S$ --- a configuration to which \emph{no known super-constant DkS
inapproximability applies}: Manurangsi's $N^{1/(\log\log N)^c}$ gap
\emph{provably} requires the opposite regime $D_{\mathrm{YES}}/m \to 0$ (i.e.\
$D_{\mathrm{YES}} = O(k^2)$, $m = \Theta(N^2)$, where the affine cost flattens
it), and the Bhaskara et al.\ $O(N^{1/4+\varepsilon})$ algorithm is itself a
$(1+o(1))$-approximation once $D_{\mathrm{YES}} = (1-o(1))m$ (the trivial
upper bound $D_{\mathrm{NO}} \le m$ then forces ratio $D_{\mathrm{YES}}/
D_{\mathrm{NO}} \le m/D_{\mathrm{NO}}$, and greedy recovery of the
edge-concentrated set is observed to attain gap $\approx 1$, see
\texttt{scripts/verify} PART~I Reason~B2). So in both sub-cases the idealised
gap is transferred from a source carrying no \emph{established} super-constant
gap. Either way: \emph{the b2 dilution formula is correct but transfers a gap
that is non-existent (under natural $A_s$, Reason~A) or unestablished
anywhere it could be non-trivial (under idealised $A_s$, Reason~B)}.

\emph{(ii) Where natural instances actually sit: both regimes occur, and
both are gap-free.} The OP1.5 optimisation object is the support of the
\emph{joint law} of $M_v$ (the predictor's positive-probability position
pairs), \emph{not} a single realised block. We measured its edge-density
$m/\binom{N}{2}$ and scaling $m \sim N^\beta$ on real-ish $n$-gram streams.
For any source with \emph{full marginal support} (the generic noisy case ---
i.i.d.\ Zipf, order-1 Markov: $\beta \approx 2$, density $\to 1$) the support
graph fills toward the complete graph as the corpus grows, so the natural
instance \emph{is dense} ($b2$) --- resolved by prong~(i)'s PTAS. For a
genuinely low-diversity finite corpus (the measured text-byte stream: $\beta
\approx 0.18$, density $\to 0.02$ at $N = 256$) the support is \emph{sub-dense}
($b1$) --- resolved by the Lemma 5.7d(b1) dilution to $1 + o(1)$. The
$b1$-sparse \emph{hard} regime of Manurangsi requires the joint law to vanish
on almost all pairs, a measure-zero structural condition not produced by a
natural source. Either way the natural multiplicative gap is capped at the
constant $1 + O(1/\log s)$.

\emph{(iii) Structural witness (conditional).} Within a single realised
block the positions are \emph{partitioned} by symbol class ($M_v = \{i : x_i
= v\}$ pairwise disjoint, §5.3), so the realised ``co-occupied-position''
graph is a \emph{disjoint union of cliques} (cluster graph), on which DkS is
poly-time \emph{exact} (maximise $\sum_i \binom{x_i}{2}$ s.t.\ $\sum_i x_i =
s$, $0 \le x_i \le c_i$; convexity of $\binom{\cdot}{2}$ puts the optimum at
the largest-cliques-first extreme point). This is the \emph{conditional}
object, distinct from the distributional optimisation, but it confirms that
the partition structure of Type-II is intrinsically benign; the NP-hard
instances of Lemma 5.7c require atoms that do \emph{not} form an equivalence
relation (a class spanning a position-\emph{pair}, e.g.\ the $C_5$ test),
which a single partition-based Type-II layer on natural data does not
generate.

\emph{Verdict.} The dense-vs-sparse dichotomy is resolved, and for the
natural cost it is in fact \emph{moot}: \emph{no natural Type-II
support-selection family yields a polynomial inapproximability}. The decisive
fact is Reason~A --- under the genuine log-binomial $A_s$ the transferred
ratio saturates at the constant $A_s(1)/A_s(0) = 1 + O(1/\log s)$ for
\emph{every} density, so the question of which regime natural instances
occupy does not even need to be answered to rule out a polynomial gap. The
regime characterisation then confirms the picture concretely: the dense
regime (generic full-support sources) is independently
$(1+\varepsilon)$-approximable by Arora--Karger--Karpinski; the sub-dense
regime (low-diversity sources) dilutes to $1 + o(1)$ by Lemma 5.7d(b1); the
per-realisation structure is a cluster graph (poly-time exact). The Lemma
5.7c \emph{decision} NP-hardness is unaffected (it lives on adversarial
intermediate-density non-cluster graphs); the multiplicative
\emph{inapproximability} does not bite for natural data. This narrows
OP4/OP1.5: for the natural log-binomial $A_s$ the worst-case multiplicative
gap is provably constant on natural inputs (Reason~A holds unconditionally),
so the only remaining route to a genuine natural-instance hardness is
average-case (attack vector~3, planted-clique), not worst-case multiplicative
inapproximability. (Attack vector~1, a contrived gap-amplifying non-natural
$A_s$, is not closed here, but is irrelevant to the \emph{natural} family by
Reason~A.) Verified in
\texttt{scripts/verify/op1\_5\_natural\_typeii\_density\_regime.py} (PART~I:
dense-regime greedy/brute gap stays $\approx 1$, PTAS-consistent, the
natural-$A_s$ saturation is a constant $1 + O(1/\log s)$, and Reason~B2
sparse-concentrated instances are greedily recoverable with gap $\approx 1$;
PART~IIa: per-realisation cluster-graph + exact convex-knapsack DkS $=$ brute
over all small $(N,s)$; PART~IIb: distributional support density on
i.i.d./Markov/text streams spans $b1$ and $b2$; PART~III: adversarial $C_5$
non-cluster boundary).

\emph{What would constitute closure:} either a polynomial-time
$O(\log N)$-approximation for general support selection (matching
DkS upper bound), or APX-hardness with explicit constant under
P $\neq$ NP or a standard cryptographic assumption.

\emph{Status update 2026-05-29 (the non-count gap's genuine home:
\textbf{Type-III-C continuous} support selection \emph{is} APX-hard via
MESP in the fixed-field regime --- a positive transfer in the regime that
mirrors Theorems 6.7/6.8, with free-field absorption re-blocking the rest;
\textsc{Partial}).} The previous update closed the non-count route
\emph{within Type-II} and observed that the genuine non-count gap
``belongs to the continuous (Type-III-C) optimal-support question, which is
a \emph{different} open problem from OP4.'' We now attack that question
directly. \emph{This is not a relabelling of the Type-II finding:} there
the log-det was a foreign functional that does not match the partition cost
(the forbidden substitution); here the log-det \emph{is} the actual
Type-III-C coding cost, so the import is legitimate.

\textbf{Lemma 6.9 (Type-III-C continuous support selection is APX-hard via
MESP in the fixed-field regime).} \emph{Let $X\sim\mathcal N(0,\Sigma)$ be
the Gaussian entropy field of a Type-III-C source on $\mathbb R^N$
(Definition 3.5; the continuous-latent / bits-back limit of \S 1.2 and
Theorem 6.6), with $\Sigma$ supplied by the model description (the
continuous analogue of the fixed predictor/dictionary of Theorems
6.7/6.8). The \emph{fixed-field support-selection problem} is: choose a
coded set $S$, $|S|=s$, coding $X_S$ and reconstructing $X_{\bar S}$ from
the field model, so as to minimise the per-block residual rate. Then:}
\begin{enumerate}
\item \emph{(The cost is a log-det.) By Theorem 6.6 the bits-back
Type-III-C rate is $H_D(X\mid Y)+\mathrm{KL}$, and for the Gaussian field
the residual of coding $X_S$ given $X_{\bar S}$ is the conditional
differential entropy}
\[
h(X_S\mid X_{\bar S}) \;=\; \tfrac12\log_2\det\!\bigl(2\pi e\,
\Sigma_{S\mid\bar S}\bigr),\qquad
\Sigma_{S\mid\bar S}=\Sigma_{SS}-\Sigma_{S\bar S}\Sigma_{\bar S\bar
S}^{-1}\Sigma_{\bar S S}
\]
\emph{(Schur complement; Cover--Thomas 2006, Thm 8.4.1). The marginal
selection variant $\min_{|S|=s}\tfrac12\log_2\det\Sigma_S$ is the Gaussian
differential entropy, equal to the total-correlation log-det of Lemma 5.6j
($\mathrm{TC}=\tfrac12[\sum_i\log_2\Sigma_{ii}-\log_2\det\Sigma]$).}
\item \emph{(The selection is MESP / $0$-$1$ D-optimal.) Minimising this
conditional entropy over size-$s$ supports is exactly the maximum-entropy
sampling / $D$-optimal design objective on principal submatrices of
$\Sigma$ (Ko--Lee--Queyranne, \emph{Oper.\ Res.}\ 43(4):684--691, 1995),
which is NP-hard, and which the unconstrained log-determinant maximisation
form makes NP-hard to approximate within $5/4$ and the size-constrained
monotone form within $1+10^{-10^{13}}$ (Ohsaka, \emph{JAIR} 2022; AISTATS
2021). Hence the fixed-field Type-III-C support-selection problem inherits
the MESP/Ohsaka multiplicative inapproximability: it is APX-hard with
constant $5/4$ (unconstrained log-det form) under $\mathrm P\neq\mathrm{NP}$.}
\item \emph{(The transfer is legitimate, unlike Type-II.) The gap is
genuinely non-count: two supports with identical induced correlated-pair
counts have different $\log\det\Sigma_S$, so the Lemma 5.7c affine $n_0(S)$
identity does not apply and the gap does not dilute. Crucially the log-det
here \emph{is} the Type-III-C cost (item 1), not an arbitrary functional
substituted for the partition cost; the anti-tautology objection that
sinks the Type-II import (no kernel $K$ makes the Type-II cost a $\log\det
K_S$) does not arise, because no substitution is performed.}
\item \emph{($\Sigma$-restriction does not kill the hardness.) For the
Lemma 5.6j manifold form $\Sigma=U\Lambda U^T+\sigma^2 I_N$, the MESP gap
survives in the informative regime $\bar\lambda\gg N\sigma^2$ (high SNR):
the per-coordinate spread of $\log\det\Sigma_S$ is $\Theta(1)$ and collapses
only as $\sigma^2\to\infty$ (when $\Sigma\to\sigma^2 I$ and all subsets tie).
This is consistent with MESP being NP-hard even for rank-deficient
covariance, so the low-rank-plus-noise restriction does not reduce the
complexity at high SNR.}
\end{enumerate}

\emph{Why this is \textsc{Partial} and not a full closure of the free-field
problem (the predictor-absorption boundary).} Lemma 6.9 establishes
hardness for the \emph{fixed-field} problem --- $\Sigma$ given, the encoder
selecting only $S$. The full RNR Type-III-C cost is $L(M)+L(D)+L(R\mid Y)$
with $M$ \emph{free}, and there the predictor-absorption obstacle of Remark
6.8b re-enters in a precise, quantified form. A free $M$ may condition on a
learned side channel $Y$ ($X\mid Y$ Gaussian with conditional covariance
$\Sigma_{\mid Y}$), shrinking the per-block residual $\tfrac12\log_2\det
\Sigma_{S\mid Y}$; but by the mutual-information identity $h(X_S)=h(X_S\mid
Y)+I(X_S;Y)$ \emph{every bit removed from the residual is paid, one for one,
in the information $I(X;Y)$ that $M$ must describe} (no free lunch /
data-processing). The consequence is a clean dichotomy:
\begin{itemize}
\item \emph{Single-block / non-amortised regime} (\S 3.4, single-block
case): $L(M)$ is paid in full, so the total never drops below the marginal
entropy floor, and the only lever is the \emph{choice} of $S$ --- which is
MESP. The $5/4$ gap survives. Lemma 6.9 holds verbatim.
\item \emph{Amortised regime} (\S 3.4, the regime RNR assumes for the main
results): $L(M)/m\to0$, so a free $M$ can pay the one-time $I(X;Y)$ to drive
the per-block residual towards zero \emph{symmetrically across all field
realisations}, diluting the multiplicative gap --- exactly the Remark 6.8b /
Lemma 13.4.2a absorption, now realised continuously as ``$M$ learns the
field covariance once.'' The fixed-field reduction does not transfer to the
free-field amortised cost.
\end{itemize}
This is the \emph{same} restricted-versus-unrestricted boundary already
present for the discrete dictionary: Theorem 6.8 is APX-hard with the fixed
predictor, while the free-predictor OP2(b) is open precisely because of
absorption (Remark 6.8b). Lemma 6.9 is the continuous (Type-III-C, log-det)
counterpart of Theorem 6.8's discrete (SGP, symbol-count) restricted
APX-hardness: \emph{the genuine non-count gap that L showed cannot live in
Type-II does live here, and it is captured exactly when --- and only when ---
the field model is fixed.} Whether the free-field amortised Type-III-C
log-det problem is APX-hard is therefore \emph{the OP2(b) question for the
continuous objective}, and is blocked by the same absorption obstacle; it is
not closed by the present lemma.

\emph{Verdict: \textsc{Partial}.} A genuine Type-III-C APX-hardness (MESP,
constant $5/4$ under $\mathrm P\neq\mathrm{NP}$; Ohsaka 2022) holds for the
fixed-field support-selection problem and survives the Lemma 5.6j
$\Sigma$-restriction at high SNR, legitimately (the log-det is the actual
cost, not a substituted functional). It does \emph{not} extend to the
free-field amortised regime, where continuous predictor-absorption (Remark
6.8b, Lemma 13.4.2a) dilutes the gap. The result locates the non-count gap's
true home (Type-III-C, as L predicted), realises it as a hardness theorem in
the fixed-field regime, and pins the residual obstacle to free-field
absorption --- the continuous mirror of OP2(b). Verified in
\texttt{scripts/verify/lemma\_6\_9\_typeiiic\_mesp\_support\_selection.py}
(C1: the cost is a Schur-complement log-det and the Lemma 5.6j TC identity
holds; C2: the selection has a genuine multiplicative MESP spread; C3: the
gap is non-count, defeating the affine identity; C4: the no-free-lunch
conservation $h_{\mathrm{marg}}=h_{\mathrm{cond}}+I(X;Y)$, giving the
single-block-survives / amortised-re-blocks dichotomy; C5: the Lemma 5.6j
restriction keeps the gap at high SNR).

\textbf{Remark 6.9a (Continuous OP2(b) characterised: the free-field
amortised Type-III-C rate problem is EASY (in P); Lemma 6.9's fixed-field /
single-block restriction is the sharp worst-case hardness boundary).}
\emph{Lemma 6.9 pinned the residual obstacle to ``the continuous mirror of
OP2(b)'': whether the free-field amortised Type-III-C log-det problem is
APX-hard. We resolve that status. With the predictor $M$ free and amortised
($L(M)/m\to0$, the regime of \S 3.4 and the main results), the problem is in
$\mathrm P$, and the optimum is achieved by an explicit poly-time
decoder-reproducible encoder.}

\emph{Formalisation.} The free-field continuous OP2(b) problem is: over all
computable predictors $M$, dictionaries $D$, and repairs $\{R_i\}$ for a
Gaussian field source $X\sim\mathcal N(0,\Sigma)$ on $\mathbb R^N$
(Definition 3.5), minimise the amortised Type-III-C cost
$L_{\mathrm{amort}}(M,D,\{R_i\}) := L(M)+L(D)+\frac1m\sum_{i=1}^m L(R_i\mid
Y,M,D)$ of Lemma 13.4.2a, where the $m$ blocks are i.i.d.\ draws from the
field. (``Free $M$'' is the most permissive class --- arbitrary computable
predictors --- so an \textsc{Easy} verdict here is the strongest form: no
restriction of the predictor class can make the problem harder.)

\emph{Verdict: \textsc{Easy} (in $\mathrm P$), unconditionally for the
amortised rate objective.} The optimal encoder is the \emph{full
sample-covariance maximum-likelihood plug-in}: on a shared prefix of $m_0$
blocks visible to both encoder and decoder, both run the identical
deterministic estimator $\hat\Sigma=\frac1{m_0}\sum_i X_iX_i^\top$ (the
\emph{unstructured} Gaussian MLE: the unconstrained maximiser of the Gaussian
log-likelihood is the sample covariance in closed form --- Anderson,
\emph{An Introduction to Multivariate Statistical Analysis}, 3rd ed.\ 2003,
Thm 3.2.1 --- computable in $O(N^2m_0)$ time; for the structured linear-model
case the log-likelihood is moreover essentially convex once $m_0\gtrsim 14N$,
Zwiernik--Uhler--Richards, \emph{J.\ R.\ Stat.\ Soc.\ B}\ 79(4):1269--1292,
2017), then code each subsequent block under $M^\star=\mathcal N(0,\hat\Sigma)$. Three facts make this optimal and the
problem easy, each verified numerically:
\begin{enumerate}
\item \emph{(The MESP combinatorial object dissolves under a free field.)}
Lemma 6.9's $5/4$ hardness \emph{is} the MESP support selection
$\min_{|S|=s}\tfrac12\log_2\det\Sigma_{S\mid\bar S}$. That object exists only
because the fixed-field problem imposes a hard budget $|S|=s$ and \emph{forbids
transmitting} $X_{\bar S}$ (the complement must be reconstructed from the
frozen model). A free amortised $M$ has no such budget: it codes the whole
block at the support-\emph{invariant} rate $h(X)=\tfrac12\log_2\det(2\pi e\,
\Sigma)$, a single scalar with no subset to choose. By the chain rule
$h(X_S)+h(X_{\bar S}\mid X_S)=h(X)$ for \emph{every} $S$, so the
combinatorial selection that carried the gap is gone (D1).
\item \emph{(Learning $\Sigma$ is poly-time and decoder-reproducible.)}
The plug-in rate is the cross-entropy $\mathbb E[-\log_2 p_{\hat\Sigma}(X)]
=h(X)+\mathrm{KL}(\mathcal N(0,\Sigma)\,\|\,\mathcal N(0,\hat\Sigma))\to
h(X)$ as $m_0\to\infty$, attaining the Slepian--Wolf / differential-entropy
floor. Both parties run the \emph{same} deterministic MLE on the \emph{same}
prefix, so $\hat\Sigma$ is bit-identical at encoder and decoder (a
universal-Slepian--Wolf shared model; the decoder-reproducibility constraint
adds no complexity --- Tan--Kosut 2014), and $L(M^\star)=O(N^2\log(1/
\varepsilon))$ bits amortises away ($L(M^\star)/m\to0$) (D2).
\item \emph{(The only candidate residual hardness --- sparse PCA --- does not
obstruct the rate.)} Hardness in covariance learning lives only in the
\emph{sparse} spiked regime of the Lemma 5.6j manifold $\Sigma=U\Lambda
U^\top+\sigma^2I$ with $U$ $k$-sparse and the signal in the band
$\sqrt{k\log N/m}<\theta<k/\sqrt m$ ($\log N<k<\sqrt N$): the average-case
sparse-PCA / spiked-covariance statistical--computational gap, conjecturally
hard via planted clique (Berthet--Rigollet, \emph{COLT} 2013, PMLR 30:1046--1066;
Brennan--Bresler--Huleihel, \emph{COLT} 2018, PMLR 75:48--166). This is \emph{average-case detection}
hardness, not worst-case multiplicative APX-hardness, and it vanishes above
the Baik--Ben~Arous--P\'ech\'e threshold $\theta>\sqrt{N/m}$ and for dense $U$
(where plain PCA recovers the spike in poly time --- consistent with Lemma
6.9's own item (4), the gap dying as $\Sigma\to\sigma^2 I$) (D3).
Decisively, the sparse-PCA gap concerns efficiently \emph{recovering the
sparse support of $U$}, which the coding objective never requires: the
\emph{full} MLE drives the rate to $h(X)$ with no sparse recovery even inside
the hard band (D4). The gap is therefore irrelevant to the amortised rate.
\end{enumerate}

\emph{Why this is not a relabelling of Lemma 6.9, and where the two meet.}
Lemma 6.9 is the \emph{fixed-field / single-block} statement and stands
verbatim: with $\Sigma$ supplied and the $|S|=s$ budget enforced (or
equivalently $L(M)$ paid in full on a single block), the MESP $5/4$ gap is
real. Remark 6.9a is the complementary \emph{free-field / amortised}
statement: lifting the budget collapses the same problem to the entropy
floor. The boundary is exactly Lemma 6.9's stated dichotomy (single-block
survives / amortised re-blocks), now resolved on the easy side --- the
``re-block'' is not a hardness but the trivial fact that a free $M$ pays the
one-time $I(X;Y)$ and thereafter codes at $h(X)$.

\emph{Why the continuous problem resolves while discrete OP2(b) stays open.}
The contrast is structural, not merely a statement that one is currently
proved and the other is not. For the discrete objective the only known route
to \emph{any} hardness past the Lemma 13.4.2a absorption is a
Kolmogorov-incompressible instance family, available solely through a
cryptographic PRG and yielding average-case \emph{decision} hardness
(\S 13.4.2, 2026-05-29 update), never worst-case APX --- so even the lone
surviving lever does not deliver APX-hardness, leaving the status open. The
continuous problem has \emph{no} such lever at all: the ``family'' is i.i.d.\
samples from one Gaussian law, and a PSD covariance cannot be made
computationally incompressible from its own samples (the sample covariance is
a consistent poly-time estimator for every $\Sigma$). The single complexity
mechanism on which discrete OP2(b) hardness would have to rest is therefore
structurally unavailable in the continuous setting, and continuous OP2(b)
collapses unconditionally to $\mathrm P$. (This does not claim a complexity
separation between the two problems --- discrete OP2(b) may itself be in
$\mathrm P$; it claims that the continuous problem is \emph{resolved} in
$\mathrm P$ whereas the discrete one is open precisely because its only
candidate hardness source has no continuous counterpart.)

\emph{Verdict: \textsc{Easy}.} Free-field amortised continuous OP2(b) is in
$\mathrm P$; the optimal encoder is the full sample-covariance MLE plug-in,
poly-time and decoder-reproducible, achieving rate $h(X)$ for every PSD
$\Sigma$ (including the sparse-spiked manifold). Lemma 6.9's fixed-field /
single-block restriction is therefore the \emph{sharp} worst-case hardness
boundary for Type-III-C support selection: hardness exists exactly on its
restricted side and nowhere in the free-field amortised regime. Verified in
\texttt{scripts/verify/remark\_6\_9a\_continuous\_op2b\_easy.py} (D1: the MESP
object dissolves under a free field via the chain rule; D2: the poly-time MLE
plug-in achieves rate $\to h(X)$ and is encoder/decoder bit-identical; D3: the
only candidate hard core is the average-case sparse-PCA gap, with dense /
above-BBP easy; D4: even that gap is irrelevant to the coding rate ---
full-MLE attains $h(X)$ without support recovery).

\textbf{Corollary 6.9b (operational Gaussian-on-manifold RNR: an
\emph{additive-regret} end-to-end encoder, fixed-field).} \emph{Assemble Lemma
5.6j (rate), Lemma 6.9 (fixed-field hardness), and the matching poly-time
support-selection approximation (the monotone-submodular greedy MESP selector,
whose calibrated $(1-1/e)$ guarantee is certified in Remark 6.9d). Let $X\sim\mathcal N(0,\Sigma)$ be a Gaussian-on-manifold
source with $\Sigma=U\Lambda U^T+\sigma^2 I_N$ satisfying the Lemma 5.6j
delocalisation / reach hypotheses (Haar-spread tangent $U$,
$\rho=\bar\lambda/(N\sigma^2)>0$), with $\Sigma$ supplied by the model (the
fixed-field regime of Lemma 6.9; the regime that is non-trivial precisely
because Remark 6.9a places the \emph{free}-field problem in $\mathrm P$), and
assume the monotonicity noise floor $\sigma^2\ge 1/(2\pi e)\approx 0.0585$ (the
isotropic floor makes the differential-entropy set function $M(T):=h(X_T)=
\tfrac12\log_2\det(2\pi e\,\Sigma_T)$ non-negative and monotone submodular ---
$\log\det$ is submodular by Kelmans--Kimelfeld / Krause--Guestrin 2008). Fix a
coding budget $s$, write $t:=N-s$, and recall the complementation identity: the
per-block residual $C(S)=h(X_S\mid X_{\bar S})=h(X)-M(\bar S)$, so
$\min_{|S|=s}C(S)$ is exactly $\max_{|T|=t}M(T)$ over the reconstruction set
$T=\bar S$ (an exact bijection). Let $S^\star$ be the optimal coded support and
$\widehat S$ the support produced by greedy submodular maximisation of $M$ on
$\widehat T=\overline{\widehat S}$. Then the end-to-end encoder} ((i)
estimate/read $\Sigma$; (ii) select $\widehat S$ by greedy MESP; (iii) code
$X_{\widehat S}$ at its conditional entropy and reconstruct the complement by
linear MMSE) \emph{runs in $\mathrm{poly}(N)$ time, and its residual obeys the
\emph{additive-regret} bound}
\[
C(\widehat S)\;-\;C(S^\star)\;=\;M(T^\star)-M(\widehat T)\;\le\;
\tfrac1e\,M(T^\star)\;\le\;\tfrac1{e-1}\,M(\widehat T),
\tag{6.9b.1}
\]
\emph{i.e.\ the encoder's excess cost over the information-theoretic optimum is
at most the Nemhauser--Wolsey--Fisher 1978 greedy slack: the multiplicative
guarantee $M(\widehat T)\ge(1-\tfrac1e)M(T^\star)$ gives $M(T^\star)-M(\widehat
T)\le\tfrac1e M(T^\star)$ (equivalently $\le\tfrac1{e-1}M(\widehat T)$ in terms
of the achieved value). The total per-block rate is
$C(S)+M(\bar S)=h(X)$ for every $S$ by the chain rule; the joint code realising
$h(X)$ carries the Lemma 5.6j compression advantage $\mathrm{TC}(X)=\tfrac12
N\,c_d(\rho)+o(N)$ over independent-coordinate coding.}

\emph{Why the composition is additive-regret, not the na\"ive $(c\cdot 5/4)$.}
The two factors do \emph{not} multiply. (a) The Lemma 6.9 hardness ($5/4$,
Ohsaka 2022) and the greedy ratio ($e/(e-1)$, Nemhauser--Wolsey--Fisher 1978;
submodular monotone $M$) are an inapproximability lower bound and an algorithmic
upper bound \emph{on the maximisation objective $M$}, not stages to be composed.
(They are not even on the identical objective form: the $5/4$ is Lemma 6.9's
\emph{unconstrained} log-det hardness, whereas the $e/(e-1)$ greedy is for the
\emph{size-constrained monotone} maximisation Lemma 6.9 item~2 shows is hard only
within $1+10^{-10^{13}}$. The two bracket the difficulty of selection but do not
multiply into one ratio.) (b) The multiplicative $e/(e-1)$ ratio holds on $M$
but \emph{not} on the residual $C$: since $C=h(X)-M$ is a \emph{difference}, a
multiplicative bound on $M$ transfers to $C$ only \emph{additively} (subtraction
does not preserve ratios; Ohsaka 2022 notes $\log$-det and $\det$ hardness are
incomparable), which is precisely (6.9b.1). A clean multiplicative rate factor
$(c\cdot 5/4)$ on the residual would require a multiplicative $\log$-det
approximation of $C$ directly, for which no constant-factor algorithm is known
in general (\c{C}ivril--Magdon-Ismail 2013, \emph{Algorithmica} 65:159--176; the
greedy guarantee is purchased by the RNR-specific monotone-submodular structure
of $M$, not by a general $C$-approximation). The honest end-to-end statement is
therefore the additive-regret (6.9b.1).

\emph{Hypothesis compatibility (the substantive check).} The pieces are
\emph{simultaneously} satisfiable. Lemma 5.6j wants delocalised $U$ and $\rho>0$
(and high SNR, $\bar\lambda\gg N\sigma^2$, for the Lemma 6.9 selection gap to be
$\Theta(1)$ per coordinate); the greedy approximation wants $\sigma^2\ge
1/(2\pi e)$ for monotonicity. These coexist: take $\sigma^2=0.25\ (\ge 0.0585)$
and $\bar\lambda=N\rho\sigma^2$ with $\rho=4$, so SNR is high \emph{and} the
noise floor holds. (Below $\sigma^2=1/(2\pi e)$ the entropy objective can be
non-monotone; a modular shift restores monotonicity without changing greedy
choices, so the corollary extends there at the cost of the shift bookkeeping.)
The concrete witness $N=14$, $d=3$, $s=7$, $\rho=4$, $\sigma^2=0.25$ satisfies
all hypotheses; see check~A of the verification script.

\emph{Fixed-field scope; sharpened by Remark 6.9a.} Corollary 6.9b is a
\emph{fixed-field} statement: $\Sigma$ is given (step~(i) reads or estimates a
\emph{fixed} field, not a learned per-instance side channel). This is exactly ---
and, by Remark 6.9a, \emph{only} --- where the synthesis has content: in the
free-field amortised regime Remark 6.9a shows the support-selection object
dissolves (the full sample-covariance MLE plug-in codes at rate $h(X)$ for every
$\Sigma$, in $\mathrm P$), so the regret (6.9b.1) and the selection lever vanish
identically. Thus Corollary 6.9b is the operational content of the \emph{hard}
side of the Lemma 6.9 / Remark 6.9a dichotomy, and is silent on (because subsumed
by the $h(X)$-floor on) the easy free-field side.

\emph{Novelty: synthesis; the value-add is the composition and the
compatibility check.} Corollary 6.9b introduces no new information-theoretic or
complexity-theoretic content beyond Lemma 5.6j, Lemma 6.9, the greedy-submodular
approximation, and Remark 6.9a. Its contribution is (i)~the \emph{compatibility
verification} --- a single concrete source provably meeting all hypotheses,
including the previously-unstated interaction between Lemma 5.6j's high-SNR
requirement and the greedy selector's noise-floor $\sigma^2\ge1/(2\pi e)$ ---
and (ii)~the \emph{correct composition}: the end-to-end guarantee is the
\emph{additive regret} (6.9b.1), \emph{not} the multiplicatively-composed
$(c\cdot 5/4)$ the surface forms suggest, because the residual is a difference of
the entropy objective. The practical payload (\S 10/\S 11) is a concrete
$\mathrm{poly}(N)$ encoder for fixed-field manifold-structured data with a stated
additive regret over the optimum; the natural-image classes of \S 13.4.5 fall
under it \emph{only} insofar as their tangent statistics are (a)~well-approximated
by a fixed Gaussian field and (b)~delocalised --- the two empirical conditions
\S 13.4.5 leaves open --- so the corollary does \emph{not} close \S 13.4.5.
Verified in \texttt{scripts/verify/cor\_6\_9a\_operational\_manifold\_rnr.py}
(A: one source satisfies all hypotheses, including $\sigma^2\ge 1/(2\pi e)$ and
high SNR; B: the residual transfer is additive ---
$C_{\mathrm{greedy}}-C_{\mathrm{opt}}=M_{\mathrm{opt}}-M_{\mathrm{greedy}}$ ---
and the $(c\cdot5/4)$-multiplicative prediction on $C$ is shown false on a
greedy-suboptimal instance; C: the per-coordinate regret is governed by the
greedy slack, vanishing in the near-modular (curvature $c\to0$) regime --- the
dependence on SNR is \emph{non-monotone}, so the honest control is curvature,
not SNR; D: the end-to-end encoder attains total per-block rate $h(X)$ with
residual within the additive regret of optimal, in $\mathrm{poly}(N)$, greedy
near-optimal on the benign low-rank-plus-isotropic instance).

\textbf{Remark 6.9c (the Type-III-C support-selection approximability
\emph{threshold}: two-sided, by budget regime; tight only by over-sampling
past the rank, a constant gap otherwise).} Lemma 6.9 pins one side --- the
fixed-field Type-III-C selection is the log-determinant / maximum-entropy
sampling (MESP) problem, hence APX-hard. It is worth pinning the
\emph{matching algorithm} so the approximability is bracketed on both sides,
because the threshold is genuinely \emph{regime-dependent} and the regime RNR
occupies is not the tight one. Throughout, ``ratio'' is the multiplicative
approximation factor on the \emph{log-determinant value} $f(S)=\log\det
\Sigma_S$ --- the same scale on which Ohsaka (\emph{JAIR} 2022) states his
hardness (his proof of Thm.\ 3.2 fixes $f(S):=\log_2\det A_S$, monotone when
$\lambda_{\min}(A)\ge 1$ by Sharma et al.\ 2015, and proves $k$-\textsc{LogDetMax}
hard to approximate ``within a factor $>1+10^{-10^{13}}$'' \emph{on $f$}),
\emph{not} the multiplicative-on-$\det$ scale of his Thm.\ 3.1 ($5/4$) nor the
$(\det)^{1/s}$ rooted-volume scale of \c{C}ivril--Magdon-Ismail
(\emph{Algorithmica} 65, 2013), whose exponential $2^{-ck}$ inapproximability
is for a \emph{different} (rooted, not additive) objective and must not be
imported here --- that would be the like-for-unlike move the Lemma~6.9
anti-tautology guards against. The complete two-sided picture from the
2015--2026 literature, like-for-like:
\begingroup\footnotesize\sloppy
\setlength{\tabcolsep}{5pt}
\noindent\begin{tabular}{@{}p{3.7cm}p{3.6cm}p{4.7cm}@{}}
\hline
\textbf{regime / objective} & \textbf{hardness (log-det value)} &
\textbf{best algorithm $\Rightarrow$ status}\\
\hline
Unconstrained log-det
& $5/4$ (Ohsaka 2022, Thm 3.1)
& $2$ (Buchbinder--Feldman 2018) $\Rightarrow$ gap $[5/4,2]$\\
Size-constrained \emph{monotone}, $|S|{=}s$, $\Sigma\succeq I$
& $1+10^{-10^{13}}$ (Ohsaka 2022, Thm 3.2)
& greedy $\tfrac{e}{e-1}\!\approx\!1.58$ (Han--Gillenwater 2020;
Sharma 2015) $\Rightarrow$ gap $[1{+}\epsilon_0,1.58]$\\
Large budget, $s\ge d+d/\epsilon$ (over-sampling)
& \emph{(no hardness $>1$)}
& $1+\epsilon$ (Allen-Zhu 2017; Madan 2019, Thm 1; NST 2019:
const.\ $e$ at $s{=}d$) $\Rightarrow$ \textsc{tight}\\
\hline
\end{tabular}
\par\endgroup
Two facts make this a genuine characterisation and not a relabelling of
Lemma~6.9.

\emph{(a) MESP and D-optimal design are the \underline{same} log-det selection
problem, separated only by the budget-vs-rank crossover.} For a rank-$d$
covariance ($\Sigma=V^TV$, $V\in\mathbb R^{d\times n}$), $\det\Sigma_S=
\prod(\text{$s$ largest eigenvalues of }\sum_{i\in S}v_iv_i^T)$ (Li--Xie,
\emph{Oper.\ Res.}\ 2023, Obs.\ 1): the MESP principal-submatrix objective
\emph{is} the D-optimal information-matrix objective on the same selection. The
non-trivial MESP regime is $s\le d$ (for $s>d$ the value is $-\infty$); ``when
$d\le s\le n$, MESP becomes the well-known D-optimal design problem''
(\emph{ibid.}, after eq.\ (2)). So the \emph{tight} $(1+\epsilon)$ guarantee
lives strictly in the \emph{over-sampled} regime $s\ge d+d/\epsilon>d$ ---
budget \emph{exceeding} the rank, reached only by repetitions (multiset
designs). The MESP / principal-submatrix selection (no repetition, $s\le d$)
\emph{cannot enter it}. (The Allen-Zhu/Madan $(1+\epsilon)$ is stated on the
rooted $\det(M)^{1/d}$ objective; on the log-det \emph{value} scale used here it
reads $1+\tfrac{\log(1+\epsilon)}{\log(s/d)}\to1$ as the over-sampling ratio
$s/d\to\infty$, so the tightness holds on the same scale as the gap below ---
but is driven precisely by $s/d>1$, vacuous at the MESP cap $s\le d$.)

\emph{(b) In the budget-capped regime $s\le d$ the gap is real and
rank-localised, not closed by any known algorithm.} On the monotone
($\Sigma\succeq I$) log-det value scale the best provable factor is the greedy
$e/(e-1)\approx1.58$, against Ohsaka's $1+10^{-10^{13}}$ lower bound: a
\emph{constant gap} $[1+\epsilon_0,\,1.58]$ that no result in the surveyed
literature closes for size-constrained log-det. The verification
(\texttt{scripts/verify/lemma\_6\_9c\_typeiiic\_threshold.py}) measures the
achieved greedy factor on synthetic Lemma~5.6j fields: it stays $\le e/(e-1)$
per instance (the monotone-submodular guarantee), is maximised at budgets
$s\approx d_{\mathrm{eff}}$ (the informative rank) and decays towards $1$ as
$s\to K$, confirming the gap is \emph{rank-localised}; a $1$-swap local-search
post-pass drives the achieved factor to $\approx1$ on tractable $K$, an honest
caveat that the worst-case hardness is asymptotic, not generic at small sizes.

\emph{RNR mapping (the decisive item).} In the Type-III-C encoder the
selection acts inside a sub-block of size $K=\Theta(\sqrt N)$ (Theorem 7.5),
on a covariance $\Sigma=U\Lambda U^T+\sigma^2 I$ (Lemma 5.6j) whose informative
spectrum has effective rank $d_{\mathrm{eff}}\ll K$ in the high-SNR regime
where the Lemma~6.9 gap lives. The principal-submatrix selection admits no
repetition and codes $s\le d_{\mathrm{eff}}$ informative coordinates. Hence
\emph{RNR sits in the budget-capped (gap) regime, not the tight one}: the
$(1+\epsilon)$ large-budget guarantee is structurally unreachable because it
demands a budget past the rank, which $s\le d_{\mathrm{eff}}$ forbids. The
sharpened verdict is therefore: \emph{Type-III-C fixed-field support selection
is APX-hard (Lemma 6.9) and its approximability is bracketed in
$[\,1+10^{-10^{13}},\ e/(e-1)\,]$ on the monotone log-det value scale (and in
$[5/4,2]$ unconstrained); it becomes tight $(1+\epsilon)$ only under
over-sampling past the rank ($s\ge d+d/\epsilon$), a regime the RNR
principal-submatrix encoder does not occupy.} The residual open question is
correspondingly sharp and \emph{purely worst-case-algorithmic}: closing the
constant gap $[1+\epsilon_0,\,e/(e-1)]$ between Ohsaka's hardness and the greedy
bound for size-constrained monotone log-det. This is orthogonal to the
free-field question: Remark 6.9a places the \emph{free-field amortised}
Type-III-C rate problem in $\mathrm P$ (predictor absorption makes it easy),
so the only Type-III-C hardness that survives is the fixed-field /
single-block one of Lemma 6.9, and the only slack in \emph{its}
characterisation is the present hardness-vs-algorithm gap $[1+\epsilon_0,
\,e/(e-1)]$ --- not predictor absorption, which is now resolved. Verified
in \texttt{scripts/verify/lemma\_6\_9c\_typeiiic\_threshold.py} (T1:
MESP$=$D-optimal, same log-det problem; T2: greedy gap real and $\le e/(e-1)$;
T3/T4: gap rank-localised at $s\approx d_{\mathrm{eff}}$, tight regime
$s\ge d+d/\epsilon$ out of range when $s\le d_{\mathrm{eff}}$; T5:
like-for-like guard, log-det value scale vs $\det$ vs rooted volume; T6: the
rooted $(1+\epsilon)$ bridges to a $(1+o(1))$ value-scale factor only as the
over-sampling ratio $s/d\to\infty$).

\textbf{Remark 6.9d (The $\Sigma\succeq I$ / noise-floor calibration of the
greedy guarantee is \emph{necessary}, not cosmetic: a concrete failure of the
raw $e/(e-1)$ bound below the entropy threshold).} The greedy
$\tfrac{e}{e-1}$ guarantee of Remark 6.9c and the additive-regret transfer of
Corollary 6.9b are both stated on a \emph{monotone, non-negative} log-det set
function --- in Remark 6.9c via the calibration $\Sigma\succeq I$ (Ohsaka's
Thm.\ 3.2 frame, monotone by Sharma et al.\ 2015 Prop.\ 2), in Corollary 6.9b
via the noise floor $\sigma^2\ge\tfrac1{2\pi e}$ that makes the
differential-entropy objective $M(T)=\tfrac12\log_2\det(2\pi e\,\Sigma_T)$
non-negative. It is worth recording, against the temptation to drop the
calibration as a harmless normalisation, that it is \emph{load-bearing}: the
Nemhauser--Wolsey--Fisher ratio is multiplicative and a multiplicative bound is
not invariant under the additive shift that calibration performs, so on the
\emph{raw} (uncalibrated) entropy the $(1-\tfrac1e)$ guarantee can fail
outright.

\emph{Counterexample (single block, $N=3$, $s=1$, $t=2$).} Take
$A:=2\pi e\,\Sigma=\bigl(\begin{smallmatrix}1.1&0&\frac1{\sqrt2}\\
0&1.1&\frac1{\sqrt2}\\ \frac1{\sqrt2}&\frac1{\sqrt2}&1.2\end{smallmatrix}\bigr)$,
with spectrum $\{0.149,\,1.1,\,2.151\}$, so $\Sigma\succ0$ but
$\lambda_{\min}(\Sigma)=0.149/(2\pi e)\approx0.0087<\tfrac1{2\pi e}$ --- the
sub-threshold regime. Maximising the raw $M(T)=\tfrac12\log_2\det A_T$ over
$|T|=2$: greedy first takes index $3$ ($M(\{3\})=\tfrac12\log_2 1.2$, the
largest singleton), then is forced into $\det A_{\{1,3\}}=\det A_{\{2,3\}}=
1.1\cdot1.2-\tfrac12=0.82$, giving $M(T_{\mathrm{greedy}})=\tfrac12\log_2
0.82\approx-0.143$; but the optimum is $T^\star=\{1,2\}$ with $\det A_{\{1,2\}}=
1.21$, $M(T^\star)=\tfrac12\log_2 1.21\approx+0.137$. The raw bound
$M(T_{\mathrm{greedy}})\ge(1-\tfrac1e)M(T^\star)$ \emph{fails}: the left side is
$-0.143$, the right $+0.087$. Equivalently the na\"ive additive transfer with
the \emph{raw} optimum, $C_{\mathrm{greedy}}-C^\star\le\tfrac1e M(T^\star)$,
also fails ($0.281\not\le0.050$). After the calibration shift $g(T)=M(T)+\tfrac{|T|}2
\log_2\alpha$ with $\alpha=1/\lambda_{\min}(A)$ (so $\lambda_{\min}(\alpha A)=1$,
$g\ge0$ monotone) the guarantee is restored --- $g(T_{\mathrm{greedy}})\approx
2.606\ge(1-\tfrac1e)g(T^\star)\approx1.825$, and the additive deficit obeys
$C_{\mathrm{greedy}}-C^\star=g(T^\star)-g(T_{\mathrm{greedy}})\le\tfrac1e
g(T^\star)$ with the \emph{calibrated} $g(T^\star)$ --- confirming that the
deficit is shift-invariant but the \emph{bound} must use the non-negative
$g(T^\star)$, never the raw $M(T^\star)$.

\emph{Two consequences, both honest caveats on the surrounding results.}
\emph{(i)} The monotonicity condition is \emph{sufficient, not necessary}: a
strongly-correlated $\Sigma$ with $\lambda_{\min}<\tfrac1{2\pi e}$ can still
have every conditional variance $\mathrm{Var}(X_i\mid X_T)\ge\tfrac1{2\pi e}$
and hence a monotone $M$ (e.g.\ $\Sigma=\bigl(\begin{smallmatrix}100&99.95\\
99.95&100\end{smallmatrix}\bigr)$ has $\lambda_{\min}=0.05<\tfrac1{2\pi e}$ yet
$\mathrm{Var}(X_2\mid X_1)=0.0999>\tfrac1{2\pi e}$, so $M$ is monotone). The
exact condition is ``all conditional variances $\ge\tfrac1{2\pi e}$'';
$\Sigma\succeq\tfrac1{2\pi e}I$ is the clean sufficient proxy used in
Corollary 6.9b / Remark 6.9c. \emph{(ii)} The calibration is always available
(scale up; it preserves greedy choices and the size-$t$ optimum because it adds
a constant to every size-$t$ set and to every per-step marginal gain), so the
guarantees of Corollary 6.9b and Remark 6.9c hold as stated --- the present
remark only certifies that their stated hypothesis cannot be silently removed.
Verified in
\texttt{scripts/verify/remark\_6\_9d\_calibration\_necessity.py} (R1: the
$N=3$ raw bound fails and the calibrated bound holds; R2: across random
sub-threshold $\Sigma$ the raw $(1-\tfrac1e)$ bound fails on a positive
fraction while the calibrated one always holds; R3: the conditional-variance
condition is the exact monotonicity boundary, with the $\lambda_{\min}\ge
\tfrac1{2\pi e}$ proxy sufficient but not necessary; R4: the additive deficit
is shift-invariant and bounded by $\tfrac1e g(T^\star)$ on the calibrated
scale).

\emph{Discrete-deployment refinement and a new open problem (Remarks
7.15a--7.15b).} The §6.9 picture above is for the Gaussian (bits-back-limit)
field. Remarks 7.15a--7.15b port it to the \emph{discrete} autoregressive
neural-LM field of Theorem 7.15 of [RNR-II] and record two effects with no Gaussian
analogue. First, the submodularity is \emph{unconditional}: discrete entropy
$H(X_T)$ is a polymatroid rank function for every law (Fujishige 1978), so the
Remark 6.9d calibration is not needed --- the discrete result is strictly
cleaner. Second, a \emph{value-oracle} obstacle appears that the Gaussian case
(where the oracle is the $O(s^3)$ log-det) entirely lacks: the support-selection
greedy needs $H(X_T)$, which for a discrete field is marginal inference, and is
poly for order-$1$ arbitrary subsets and bounded-order contiguous blocks but
\#P-hard (exact) for general high-order fields (Roth 1996). So the discrete
deployment carries \emph{two} obstacles --- NP-hard selection (Max-Coverage,
Feige 1998) \emph{and} a \#P-hard exact oracle --- versus the Gaussian one. This
exposes a \textbf{new open problem}: the \emph{noncontiguous bounded-order}
oracle is exactly the \emph{exact finite-block Shannon entropy of a hidden
Markov / function-of-Markov-chain process}, for which neither a polynomial
algorithm nor a \#P-/NP-hardness theorem is known (the adjacent results ---
entropy \emph{rate} with no closed form, Blackwell 1957; most-likely string
NP-hardness, Lyng\o{}--Pedersen 2002; automata relative-entropy
PSPACE-completeness, Cortes--Mohri--Rastogi--Riley 2008 --- do not settle exact
finite-block Shannon entropy, which is moreover tractable for unambiguous PFA).
Remark 7.15b then shows this hardness is \emph{orthogonal} to the deployed
encoder: causal coding and byte-granular random access condition only on the
realised prefix (one forward pass), and the \#P-hard marginal arises solely in
the optional non-contiguous selection lever. Remark 7.15c makes \emph{partial
progress} on the open exact problem: under a filter-stability hypothesis (FS:
geometric ergodicity of the order-$w$ hidden chain with positive Doeblin
emissions, the discrete analogue of Lemma 5.6e's GE\,+\,EP), the induced HMM's
prediction filter forgets geometrically, so the order-$d$ truncation
approximates the block entropy with error $O(r\rho^d)$ and the oracle is
\emph{$\varepsilon$-approximable} in slice-wise-polynomial time
($d=O(\log(r/\varepsilon))$); robust submodular greedy then selects within
$(1-\tfrac1e)-O(\varepsilon)$. So the \#P-hardness is an artifact of
\emph{exact} computation --- it is suspected to persist even under mixing, but
the operative approximate oracle is tractable on filter-stable sources. Remark
7.15d then maps the complexity \emph{landscape} via the R\'enyi spectrum: every
\emph{integer}-order R\'enyi entropy $H_\alpha$ ($\alpha\ge2$) is \emph{exactly}
poly --- unconditionally --- because $\sum_y P(y)^\alpha$ is the agreement
probability of $\alpha$ independent copies, a product-chain forward computation;
the min-entropy ($\alpha=\infty$) is NP-hard (HMM consensus, Lyng\o{}--Pedersen
2002); and the operative Shannon point ($\alpha=1$) is not reached by either,
though $H_2$ furnishes a free exact poly \emph{lower bound}. The exact Shannon
complexity remains the open problem (now bracketed: poly $H_2$ below, FS
approximation, NP-hard min-entropy above). Remark 7.15e then \emph{maps} this wall:
the exact oracle lies in $\mathrm{FP}^{\#P}$ (a weighted exponential sum of
poly-computable terms), and both natural attacks are provably obstructed --- the
predicate-encoding \#P-hardness route (Remark 7.15a/V5) is self-defeating for
bounded-order chains (they encode only read-once predicates, whose model-counting is
in $\mathrm P$, and the gap-marginalisation adds only $\#\mathrm L\subseteq\mathrm{FP}$
walk-counting), and the obvious filtering algorithm is blocked by the non-collapsing
Blackwell measure. So a resolution needs a \emph{non-predicate-encoding} hardness
construction or a \emph{non-filtering} algorithm --- the wall is pinned, not broken.

\textbf{13.4.5 Tight TC scaling for natural-data classes (refined
after Lemma 5.6i).}

\emph{What has been established:}
\begin{itemize}
\item Lemma 5.6a-d: $\Theta(N)$ TC for four explicit dependent
families (uniqueness-constrained, hidden-class iid, stationary
processes, HMMs).
\item Lemma 5.6e-h: extensions to continuous-state HMMs, AMS-mixing,
tree-MRFs, chordal junction-tree decomposable.
\item Lemma 5.6i: \emph{sub-linear} TC for $K$-sparse combinatorial
uniform; $\Theta(N)$ TC requires pairwise dependence (permutations),
not just low-cardinality support.
\end{itemize}

\emph{Remaining open: rigorous $\Theta(N)$ TC for less combinatorially
structured natural-data classes.}

\emph{Attack vectors:}
\begin{enumerate}
\item \emph{Empirical TC measurement on real data.} For specific
data classes (natural images at fixed resolution, audio waveforms,
genome sequences), compute empirical TC via plug-in estimators
(Kraskov-Stögbauer-Grassberger 2004 nearest-neighbor MI estimator)
and check $\Theta(N)$ scaling. This is empirical (§11) but provides
data for theory.
\item \emph{Manifold-plus-regularity assumptions.} Show that for
data classes satisfying explicit manifold-and-density assumptions
(bounded reach, Lipschitz density on the manifold, bounded
parameter complexity of the manifold itself), TC is $\Theta(N)$.
The proof would require quantitative geometric measure theory
beyond the informal manifold intuition.
\item \emph{Bounded-relative-entropy classes.} For data classes
$D_N$ satisfying $D(P_{X_i \mid X_{<i}} \,\|\, P_{X_i}) \geq \delta$
for all $i$ (each coordinate provides $\delta$ bits of conditional
information beyond marginal), TC $\geq N \delta = \Theta(N)$. Show
that natural classes satisfy this with measurable $\delta$.
\end{enumerate}

\emph{Remark (attack vector~3 reduces to the chain rule; the open
part is empirical $\delta$-estimation).} The inequality in
attack vector~3 is \emph{not} an independent mathematical lemma: it
is the multi-information chain rule
$\mathrm{TC}(D_N) = \sum_{k=2}^{N} I(X_k; X_{1:k-1})$ (Cover \&
Thomas 2006, already used in \S5.6 --- e.g.\ the tree-MRF bound
``$I(X_u; X_v) \geq \varepsilon$ for every edge $\Rightarrow
\mathrm{TC} \geq (N-1)\varepsilon$'' of Lemma 5.6g is the identical
termwise-lower-bound move applied to sequential conditionals)
combined with the trivial observation that a sum of $N-1$ terms each
$\geq \delta$ is $\geq (N-1)\delta$. Two points of rigour, neither of
which yields new theory. \emph{(i) Random vs.\ expected.} The
quantity $D(P_{X_i \mid X_{<i}=x} \,\|\, P_{X_i})$ is a \emph{random}
pointwise relative entropy (it depends on the realisation $x$ of
$X_{<i}$, and can lie above or below its mean); the quantity that
enters the chain rule is its expectation,
$\mathbb{E}_{X_{<i}}\!\bigl[D(P_{X_i \mid X_{<i}=x} \,\|\, P_{X_i})\bigr]
= I(X_i; X_{<i})$. The correct hypothesis is therefore
$I(X_i; X_{<i}) \geq \delta$ (expected/mutual-information form);
``$D \geq \delta$ for every realisation $x$'' is a strictly stronger
\emph{sufficient} condition and is not what natural-data
$\delta$-estimates measure. \emph{(ii) No gain over Lemma 5.6c.} For
\emph{stationary} sources this lower bound is already subsumed by
Lemma 5.6c ($\mathrm{TC} = N E - \delta_N$ with $E = H(X_1) - h(X)$),
and the genuinely \emph{non-stationary} (e.g.\ AMS) cases are covered
by Lemma 5.6f; in every case the inequality is the chain rule, so the
bounded-$\delta$ framing contributes no information-theoretic content
beyond \S5.6. The substantive open part of attack vector~3 is purely
\emph{empirical}: exhibiting a non-combinatorial natural-data class
with a \emph{measured} $\delta = \inf_i I(X_i; X_{<i}) > 0$ bounded
away from $0$ (the Shannon-1951 English-text estimates
$E_{\mathrm{bigram}} \approx 0.71$, $E_{\mathrm{en}} \approx 3.08$
bits/letter of Lemma 5.6c are illustrative but not a proof for a
fixed class). This is the §11 estimation question, not a missing
lemma.

\emph{What would constitute closure:} an explicit construction of a
natural-data class (not combinatorially defined) with rigorously
proven $\mathrm{TC}(D_N) = \Theta(N)$, plus an explicit polynomial-time
rank coder achieving rate $\log_2 |D_N| / N + o(1)$.

\textbf{13.4.6 Methodological observation.}

The pattern across these open problems is consistent: \emph{the
standard reduction templates (Berman-Karpinski-Charikar for SGP,
CLIQUE for support selection, multi-way-cut for joint problems)
have all been tried at the natural target regimes and shown to
fail by quantitative gap-dilution arguments}. Closing any of these
problems likely requires either (a) a structurally novel reduction
template suited to the target regime, (b) cryptographic conditional
hardness via Liu-Pass-style indirect connections, or (c) a positive
algorithmic result (PTAS or polynomial-time exact algorithm)
contradicting the current intuition. We list these as the natural
next research directions, and explicitly retract any premature
claims of closure that did not survive adversarial validation (notably: K-restricted ``closures'' that excluded the
obstacle by fiat, and constant-substitution ``closures'' that
conflated different complexity-theoretic quantities).

\textbf{References}

\textbf{{[}1{]}} Banner, R., Hubara, I., Hoffer, E., and Soudry, D.
(2018). Scalable Methods for 8-bit Training of Neural Networks. NeurIPS.
\textbf{{[}2{]}} Bellard, F. (2019). Lossless Data Compression with
Neural Networks. Technical report and software,
https://bellard.org/nncp/.
\textbf{{[}3{]}} Bernstein, J., Wang, Y.-X., Azizzadenesheli, K., and
Anandkumar, A. (2018). signSGD: Compressed Optimisation for Non-Convex
Problems. ICML.
\textbf{{[}4{]}} Burrows, M., and Wheeler, D. J. (1994). A Block-Sorting
Lossless Data Compression Algorithm. SRC Research Report 124, Digital
Equipment Corporation.
\textbf{{[}5{]}} Carlsson, G. (2009). Topology and Data. Bulletin of the
American Mathematical Society, 46(2), 255-308.
\textbf{{[}6{]}} Carlsson, G., Ishkhanov, T., de Silva, V., and
Zomorodian, A. (2008). On the Local Behavior of Spaces of Natural
Images. International Journal of Computer Vision, 76(1), 1-12.
\textbf{{[}7{]}} Chen, X., Mishra, N., Rohaninejad, M., and Abbeel, P.
(2018). PixelSNAIL: An Improved Autoregressive Generative Model.
International Conference on Machine Learning.
\textbf{{[}8{]}} Charikar, M., Lehman, E., Liu, D., Panigrahy, R.,
Prabhakaran, M., Sahai, A., and Shelat, A. (2005). The Smallest Grammar
Problem. IEEE Transactions on Information Theory, 51(7), 2554-2576.

\textbf{{[}8a{]}} Casel, K., Fernau, H., Gaspers, S., Gras, B., and
Schmid, M. L. (2021). On the Complexity of the Smallest Grammar Problem
over Fixed Alphabets. Theory of Computing Systems, 65(1), 158-203.

\textbf{{[}8b{]}} Ga\'nczorz, M. (2018). Entropy bounds for grammar
compression. arXiv:1804.08547 (also presented at DCC 2018/2020).

\textbf{{[}8c{]}} Kozma, L., and Voderholzer, J. (2024). The
Computational Complexity of Tokenisation. arXiv:2411.08671.

\textbf{{[}8d{]}} Je\.z, A. (2015). A really simple approximation of
smallest grammar. Theoretical Computer Science, 616, 141-150.

\textbf{{[}8e{]}} Karzanov, A. V. (1998). Minimum 0-Extensions of Graph
Metrics. European Journal of Combinatorics, 19(1), 71-101.

\textbf{{[}8f{]}} Schwartz, R., and Tur, S. (2023). Improved hardness
of approximation for graph problems via Boolean PCPs.
arXiv:2104.11670 (also conference version, ICALP 2023).
\textbf{{[}9{]}} Cover, T. M. (1973). Enumerative Source Encoding. IEEE
Transactions on Information Theory, 19(1), 73-77.
\textbf{{[}10{]}} Cover, T. M., and Thomas, J. A. (2006). Elements of
Information Theory, 2nd edition. Wiley.
\textbf{{[}11{]}} Donoho, D. L. (2000). High-Dimensional Data Analysis:
The Curses and Blessings of Dimensionality. AMS Math Challenges Lecture.
\textbf{{[}12{]}} Duda, J. (2009). Asymmetric Numeral Systems.
arXiv:0902.0271.
\textbf{{[}13{]}} Frey, B. J. (1997). Bayesian Networks for Pattern
Classification, Data Compression, and Channel Coding. Ph.D. thesis,
University of Toronto.
\textbf{{[}14{]}} Frey, B. J. (1998). Graphical Models for Machine
Learning and Digital Communication. MIT Press.
\textbf{{[}15{]}} Hinton, G. E., and van Camp, D. (1993). Keeping Neural
Networks Simple by Minimizing the Description Length of the Weights.
Conference on Learning Theory.
\textbf{{[}16{]}} Hoogeboom, E., Peters, J. W. T., van den Berg, R., and
Welling, M. (2019). Integer Discrete Flows and Lossless Compression.
NeurIPS.
\textbf{{[}17{]}} Huffman, D. A. (1952). A Method for the Construction
of Minimum-Redundancy Codes. Proceedings of the IRE, 40(9), 1098-1101.
\textbf{{[}18{]}} Kingma, F. H., Abbeel, P., and Ho, J. (2019).
Bit-Swap: Recursive Bits-Back Coding for Lossless Compression with
Hierarchical Latent Variables. ICML.
\textbf{{[}19{]}} Knoll, B. (ongoing). Cmix: a lossless data compression
program. https://www.byronknoll.com/cmix.html.
\textbf{{[}20{]}} Levina, E., and Bickel, P. J. (2004). Maximum
Likelihood Estimation of Intrinsic Dimension. NIPS.
\textbf{{[}21{]}} Liveris, A. D., Xiong, Z., and Georghiades, C. N.
(2002). Compression of Binary Sources with Side Information at the
Decoder Using LDPC Codes. IEEE Communications Letters, 6(10), 440-442.
\textbf{{[}22{]}} Mahoney, M. (ongoing). Data Compression Programs and
the PAQ family. http://mattmahoney.net/dc/.
\textbf{{[}23{]}} Mentzer, F., Agustsson, E., Tschannen, M., Timofte,
R., and Van Gool, L. (2019). Practical Full Resolution Learned Lossless
Image Compression. CVPR.
\textbf{{[}24{]}} Pradhan, S. S., and Ramchandran, K. (2003).
Distributed Source Coding Using Syndromes (DISCUS): Design and
Construction. IEEE Transactions on Information Theory, 49(3), 626-643.
\textbf{{[}25{]}} Reed, I. S., and Solomon, G. (1960). Polynomial Codes
over Certain Finite Fields. Journal of SIAM, 8(2), 300-304.
\textbf{{[}26{]}} Rissanen, J. (1978). Modeling by Shortest Data
Description. Automatica, 14(5), 465-471.
\textbf{{[}27{]}} Rissanen, J., and Langdon, G. G. (1979). Arithmetic
Coding. IBM Journal of Research and Development, 23(2), 149-162.
\textbf{{[}28{]}} Salimans, T., Karpathy, A., Chen, X., and Kingma, D.
P. (2017). PixelCNN++: Improving the PixelCNN with Discretized Logistic
Mixture Likelihood and Other Modifications. ICLR.

\textbf{{[}29a{]}} Crutchfield, J. P., and Feldman, D. P. (2003).
Regularities Unseen, Randomness Observed: Levels of Entropy
Convergence. Chaos, 13(1), 25-54. DOI: 10.1063/1.1530990.
\textbf{{[}29{]}} Sahni, S., and Gonzalez, T. (1976). P-Complete
Approximation Problems. Journal of the ACM, 23(3), 555-565.
\textbf{{[}30{]}} Shannon, C. E. (1948). A Mathematical Theory of
Communication. Bell System Technical Journal, 27, 379-423 and 623-656.

\textbf{{[}31a{]}} Shannon, C. E. (1951). Prediction and Entropy of
Printed English. Bell System Technical Journal, 30(1), 50-64.
\textbf{{[}31{]}} Slepian, D., and Wolf, J. K. (1973). Noiseless Coding
of Correlated Information Sources. IEEE Transactions on Information
Theory, 19(4), 471-480.
\textbf{{[}32{]}} Tenenbaum, J. B., de Silva, V., and Langford, J. C.
(2000). A Global Geometric Framework for Nonlinear Dimensionality
Reduction. Science, 290(5500), 2319-2323.
\textbf{{[}33{]}} Townsend, J., Bird, T., and Barber, D. (2019).
Practical Lossless Compression with Latent Variables Using Bits Back
Coding. ICLR.
\textbf{{[}34{]}} Townsend, J., Bird, T., Kunze, J., and Barber, D.
(2020). HiLLoC: Lossless Image Compression with Hierarchical Latent
Variable Models. ICLR.
\textbf{{[}35{]}} van den Berg, R., Gritsenko, A. A., Dehghani, M.,
Sonderby, C. K., and Salimans, T. (2021). IDF++: Analyzing and Improving
Integer Discrete Flows for Lossless Compression. ICLR.
\textbf{{[}36{]}} van den Oord, A., Kalchbrenner, N., and Kavukcuoglu,
K. (2016). Pixel Recurrent Neural Networks. ICML.
\textbf{{[}37{]}} van den Oord, A., Kalchbrenner, N., Vinyals, O.,
Espeholt, L., Graves, A., and Kavukcuoglu, K. (2016). Conditional Image
Generation with PixelCNN Decoders. NeurIPS.
\textbf{{[}38{]}} Rhee, H., Jang, Y. I., Kim, S., and Cho, N. I. (2022).
LC-FDNet: Learned Lossless Image Compression with Frequency
Decomposition Network. CVPR, 6023-6032.
\textbf{{[}39{]}} Witten, I. H., Neal, R. M., and Cleary, J. G. (1987).
Arithmetic Coding for Data Compression. Communications of the ACM,
30(6), 520-540.
\textbf{{[}40{]}} Wu, S., Yan, J., Wang, J., Maxson, A., Xu, K., and He,
K. (2018). Training and Inference with Integers in Deep Neural Networks.
ICLR.
\textbf{{[}41{]}} Wyner, A. D., and Ziv, J. (1976). The Rate-Distortion
Function for Source Coding with Side Information at the Decoder. IEEE
Transactions on Information Theory, 22(1), 1-10.
\textbf{{[}42{]}} Ziv, J., and Lempel, A. (1977). A Universal Algorithm
for Sequential Data Compression. IEEE Transactions on Information
Theory, 23(3), 337-343.
\textbf{{[}43{]}} Ziv, J., and Lempel, A. (1978). Compression of
Individual Sequences via Variable-Rate Coding. IEEE Transactions on
Information Theory, 24(5), 530-536.
\textbf{{[}44{]}} Watanabe, S. (1960). Information Theoretical
Analysis of Multivariate Correlation. IBM Journal of Research and
Development, 4(1), 66-82.
\textbf{{[}45{]}} Studen\'y, M., and Vejnarov\'a, J. (1998). The
Multiinformation Function as a Tool for Measuring Stochastic
Dependence. In: \emph{Learning in Graphical Models}, Jordan, M. I.
(ed.), MIT Press, pp. 261-297.
\textbf{{[}46{]}} Lund, C., and Yannakakis, M. (1994). On the
Hardness of Approximating Minimization Problems. Journal of the ACM,
41(5), 960-981.
\textbf{{[}47{]}} Lehman, E., and Shelat, A. (2002). Approximation
Algorithms for Grammar-Based Compression. Proceedings of the
Thirteenth Annual ACM-SIAM Symposium on Discrete Algorithms (SODA),
205-212.
\textbf{{[}48{]}} Berman, P., and Karpinski, M. (1999). On Some
Tighter Inapproximability Results. Proceedings of the 26th
International Colloquium on Automata, Languages and Programming
(ICALP), Lecture Notes in Computer Science 1644, Springer, 200-209.
\textbf{{[}49{]}} Storer, J. A., and Szymanski, T. G. (1982). Data
Compression via Textual Substitution. Journal of the ACM, 29(4),
928-951.
\textbf{{[}50{]}} Ar{\i}kan, E. (2009). Channel Polarization: A Method
for Constructing Capacity-Achieving Codes for Symmetric Binary-Input
Memoryless Channels. IEEE Transactions on Information Theory, 55(7),
3051-3073.
\textbf{{[}51{]}} Azuma, K. (1967). Weighted Sums of Certain Dependent
Random Variables. T\^ohoku Mathematical Journal, 19(3), 357-367.
\textbf{{[}52{]}} Bellare, M. (2006). New Proofs for NMAC and HMAC:
Security Without Collision-Resistance. In: Advances in Cryptology ---
CRYPTO 2006, Lecture Notes in Computer Science 4117, Springer, 602-619.
\textbf{{[}53{]}} Bellare, M., Canetti, R., and Krawczyk, H. (1996).
Keying Hash Functions for Message Authentication. In: Advances in
Cryptology --- CRYPTO '96, Lecture Notes in Computer Science 1109,
Springer, 1-15.
\textbf{{[}54{]}} Bellare, M., and Namprempre, C. (2000). Authenticated
Encryption: Relations Among Notions and Analysis of the Generic
Composition Paradigm. In: Advances in Cryptology --- ASIACRYPT 2000,
Lecture Notes in Computer Science 1976, Springer, 531-545.
\textbf{{[}55{]}} Berger, T. (1971). Rate Distortion Theory: A
Mathematical Basis for Data Compression. Prentice-Hall, Englewood
Cliffs, NJ.
\textbf{{[}56{]}} Breslau, L., Cao, P., Fan, L., Phillips, G., and
Shenker, S. (1999). Web Caching and Zipf-like Distributions: Evidence
and Implications. Proceedings of IEEE INFOCOM, 126-134.
\textbf{{[}57{]}} Cover, T. M. (1975). A Proof of the Data Compression
Theorem of Slepian and Wolf for Ergodic Sources. IEEE Transactions on
Information Theory, 21(2), 226-228.
\textbf{{[}58{]}} Cover, T. M., and King, R. C. (1978). A Convergent
Gambling Estimate of the Entropy of English. IEEE Transactions on
Information Theory, 24(4), 413-421.
\textbf{{[}59{]}} Csisz\'ar, I., and K\"orner, J. (2011). Information
Theory: Coding Theorems for Discrete Memoryless Systems, 2nd edition.
Cambridge University Press. (First edition: Akad\'emiai Kiad\'o, 1981.)
\textbf{{[}60{]}} Cybenko, G. (1989). Approximation by Superpositions
of a Sigmoidal Function. Mathematics of Control, Signals, and Systems,
2(4), 303-314.
\textbf{{[}61{]}} Del\'etang, G., Ruoss, A., Duquenne, P.-A., Catt, E.,
Genewein, T., Mattern, C., Grau-Moya, J., Wenliang, L. K.,
Aitchison, M., Orseau, L., Legg, S., and Veness, J. (2024). Language
Modeling Is Compression. Proceedings of ICLR. arXiv:2309.10668.
\textbf{{[}62{]}} Doshi, V., Shah, D., and M\'edard, M. (2007). Source
Coding with Distortion through Graph Coloring. Proceedings of the IEEE
International Symposium on Information Theory (ISIT), 1501-1505.
\textbf{{[}63{]}} Doshi, V., Shah, D., M\'edard, M., and Effros, M.
(2010). Functional Compression through Graph Coloring. IEEE
Transactions on Information Theory, 56(8), 3901-3917.
\textbf{{[}64{]}} Equitz, W. H. R., and Cover, T. M. (1991). Successive
Refinement of Information. IEEE Transactions on Information Theory,
37(2), 269-275.
\textbf{{[}65{]}} Fagin, R. (1977). Generalized First-Order Spectra and
Polynomial-Time Recognizable Sets. In: Complexity of Computation,
SIAM-AMS Proceedings 7, 43-73.
\textbf{{[}66{]}} Ferragut, A., Rodr{\'\i}guez, I., and Paganini, F.
(2018). Optimizing TTL Caches under Heavy-Tailed Demand. Queueing
Systems, 88(3-4), 207-241.
\textbf{{[}67{]}} Grossi, R., and Vitter, J. S. (2005). Compressed
Suffix Arrays and Suffix Trees with Applications to Text Indexing and
String Matching. SIAM Journal on Computing, 35(2), 378-407.
\textbf{{[}68{]}} Hinton, G., Vinyals, O., and Dean, J. (2015).
Distilling the Knowledge in a Neural Network. arXiv:1503.02531.
\textbf{{[}69{]}} Hoeffding, W. (1963). Probability Inequalities for
Sums of Bounded Random Variables. Journal of the American Statistical
Association, 58(301), 13-30.
\textbf{{[}70{]}} Hoffmann, J., Borgeaud, S., Mensch, A., Buchatskaya,
E., Cai, T., Rutherford, E., et al. (2022). Training Compute-Optimal
Large Language Models. arXiv:2203.15556.
\textbf{{[}71{]}} Hornik, K., Stinchcombe, M., and White, H. (1989).
Multilayer Feedforward Networks Are Universal Approximators. Neural
Networks, 2(5), 359-366.
\textbf{{[}72{]}} Kaplan, J., McCandlish, S., Henighan, T., Brown,
T. B., Chess, B., Child, R., Gray, S., Radford, A., Wu, J., and
Amodei, D. (2020). Scaling Laws for Neural Language Models.
arXiv:2001.08361.
\textbf{{[}73{]}} Korada, S. B., and Urbanke, R. L. (2010). Polar
Codes Are Optimal for Lossy Source Coding. IEEE Transactions on
Information Theory, 56(4), 1751-1768.
\textbf{{[}74{]}} McDiarmid, C. (1989). On the Method of Bounded
Differences. In: Surveys in Combinatorics 1989, London Mathematical
Society Lecture Notes 141, Cambridge University Press, 148-188.
\textbf{{[}75{]}} McLeish, D. L. (1974). Dependent Central Limit
Theorems and Invariance Principles. Annals of Probability, 2(4),
620-628.
\textbf{{[}76{]}} Muthitacharoen, A., Chen, B., and Mazi\`eres, D.
(2001). A Low-Bandwidth Network File System. Proceedings of the 18th
ACM Symposium on Operating Systems Principles (SOSP), 174-187.
\textbf{{[}77{]}} Orlitsky, A., and Roche, J. R. (2001). Coding for
Computing. IEEE Transactions on Information Theory, 47(3), 903-917.
\textbf{{[}78{]}} Polyanskiy, Y., Poor, H. V., and Verd\'u, S. (2010).
Channel Coding Rate in the Finite Blocklength Regime. IEEE Transactions
on Information Theory, 56(5), 2307-2359.
\textbf{{[}79{]}} Rabin, M. O. (1981). Fingerprinting by Random
Polynomials. Technical Report TR-15-81, Center for Research in Computing
Technology, Harvard University.
\textbf{{[}80{]}} Richardson, T., and Urbanke, R. (2008). Modern Coding
Theory. Cambridge University Press.
\textbf{{[}81{]}} Rimoldi, B. (1994). Successive Refinement of
Information: Characterization of the Achievable Rates. IEEE Transactions
on Information Theory, 40(1), 253-259.
\textbf{{[}82{]}} Rio, E. (2000). Th\'eorie Asymptotique des Processus
Al\'eatoires Faiblement D\'ependants. Math\'ematiques \& Applications 31,
Springer.
\textbf{{[}83{]}} Rissanen, J. (1996). Fisher Information and
Stochastic Complexity. IEEE Transactions on Information Theory, 42(1),
40-47.
\textbf{{[}84{]}} Sadakane, K. (2003). New Text Indexing Functionalities
of the Compressed Suffix Arrays. Journal of Algorithms, 48(2), 294-313.
\textbf{{[}85{]}} Said, A. (2003). Arithmetic Coding for Data
Compression. Technical Report, Hewlett-Packard Laboratories.
\textbf{{[}86{]}} Samson, P.-M. (2000). Concentration of Measure
Inequalities for Markov Chains and $\Phi$-Mixing Processes. Annals of
Probability, 28(1), 416-461.
\textbf{{[}87{]}} Shannon, C. E. (1959). Coding Theorems for a Discrete
Source with a Fidelity Criterion. IRE National Convention Record, Part 4,
142-163.
\textbf{{[}88{]}} Sleator, D. D., and Tarjan, R. E. (1985). Amortized
Efficiency of List Update and Paging Rules. Communications of the ACM,
28(2), 202-208.
\textbf{{[}89{]}} Steinberg, Y., and Merhav, N. (2004). On Successive
Refinement for the Wyner-Ziv Problem. IEEE Transactions on Information
Theory, 50(8), 1636-1654.
\textbf{{[}90{]}} Witsenhausen, H. S. (1976). The Zero-Error Side
Information Problem and Chromatic Numbers. IEEE Transactions on
Information Theory, 22(5), 592-593.
\textbf{{[}91{]}} Zhang, Y., Huang, Z., Chen, K., Tao, D., et al.
(2024). L3TC: Leveraging RWKV for Learned Lossless Low-Complexity Text
Compression. arXiv:2412.16642.
\textbf{{[}92{]}} Krichevsky, R. E., and Trofimov, V. K. (1981). The
Performance of Universal Encoding. IEEE Transactions on Information
Theory, 27(2), 199-207.
\textbf{{[}93{]}} Baum, L. E., and Petrie, T. (1966). Statistical
Inference for Probabilistic Functions of Finite State Markov Chains.
Annals of Mathematical Statistics, 37(6), 1554-1563.
\textbf{{[}94{]}} Lauritzen, S. L. (1996). Graphical Models. Oxford
University Press.
\textbf{{[}95{]}} Lauritzen, S. L., and Spiegelhalter, D. J. (1988).
Local Computations with Probabilities on Graphical Structures and
Their Application to Expert Systems. Journal of the Royal Statistical
Society, Series B, 50(2), 157-194.
\textbf{{[}96{]}} Dinur, I. (2007). The PCP Theorem by Gap
Amplification. Journal of the ACM, 54(3), Article 12.
\textbf{{[}97{]}} Dinur, I., Kindler, G., Raz, R., and Safra, S.
(2003). Approximating CVP to Within Almost-Polynomial Factors is
NP-Hard. Combinatorica, 23(2), 205-243.
\textbf{{[}98{]}} Barron, A. R. (1985). The Strong Ergodic Theorem
for Densities: Generalized Shannon-McMillan-Breiman Theorem. Annals
of Probability, 13(4), 1292-1303.
\textbf{{[}99{]}} Gray, R. M. (1990). Entropy and Information
Theory. Springer.
\textbf{{[}100{]}} Pradhan, S. S., and Ramchandran, K. (1999).
Distributed Source Coding Using Syndromes (DISCUS): Design and
Construction. Proceedings of the Data Compression Conference (DCC),
158-167.
\textbf{{[}101{]}} Stratonovich, R. L. (1960). Conditional Markov
Processes. Theory of Probability and Its Applications, 5(2), 156-178.
\textbf{{[}102{]}} Wainwright, M. J., and Jordan, M. I. (2008).
Graphical Models, Exponential Families, and Variational Inference.
Foundations and Trends in Machine Learning, 1(1-2), 1-305.
\textbf{{[}103{]}} Arora, S., Babai, L., Stern, J., and Sweedyk, Z.
(1997). The Hardness of Approximate Optima in Lattices, Codes, and
Systems of Linear Equations. Journal of Computer and System Sciences,
54(2), 317-331.
\textbf{{[}104{]}} Birch, B. J. (1962). Approximations for the Entropy
for Functions of Markov Chains. Annals of Mathematical Statistics,
33(3), 930-938.
\textbf{{[}105{]}} Manurangsi, P. (2017). Almost-Polynomial Ratio
ETH-Hardness of Approximating Densest $k$-Subgraph. Proceedings of the
49th Annual ACM Symposium on Theory of Computing (STOC), 19-23.
\textbf{{[}106{]}} Pakzad, P., and Anantharam, V. (2002). Estimation
and Marginalization Using Kikuchi Approximation Methods. Neural
Computation, 17(8), 1836-1873.
\textbf{{[}107{]}} Kieffer, J. C., and Yang, E.-H. (2000). Grammar-Based
Codes: A New Class of Universal Lossless Source Codes. IEEE Transactions
on Information Theory, 46(3), 737-754.
\textbf{{[}108{]}} Merlev\`ede, F., Peligrad, M., and Rio, E. (2009).
Bernstein Inequality and Moderate Deviations under Strong Mixing
Conditions. In: High Dimensional Probability V: The Luminy Volume,
IMS Collections 5, 273-292.
\textbf{{[}109{]}} Bradley, R. C. (2005). Basic Properties of Strong
Mixing Conditions: A Survey and Some Open Questions. Probability
Surveys, 2, 107-144.
\textbf{{[}110{]}} Rio, E. (2017). Asymptotic Theory of Weakly
Dependent Random Processes. Probability Theory and Stochastic
Modelling 80, Springer.
\textbf{{[}111{]}} Kontorovich, L., and Ramanan, K. (2008).
Concentration Inequalities for Dependent Random Variables via the
Martingale Method. Annals of Probability, 36(6), 2126-2158.
\textbf{{[}112{]}} Davydov, Yu.\ A. (1968). Convergence of
Distributions Generated by Stationary Stochastic Processes. Theory
of Probability and Its Applications, 13(4), 691-696.
\textbf{{[}113{]}} Bradley, R. C. (2007). Introduction to Strong
Mixing Conditions, Volumes 1--3. Kendrick Press.
\textbf{{[}114{]}} Doukhan, P. (1994). Mixing: Properties and
Examples. Lecture Notes in Statistics 85, Springer.
\textbf{{[}115{]}} Brown, P. F., Della Pietra, V. J., Mercer,
R. L., Della Pietra, S. A., and Lai, J. C. (1992). An Estimate of
an Upper Bound for the Entropy of English. Computational Linguistics,
18(1), 31-40.
\textbf{{[}116{]}} Lezaud, P. (1998). Chernoff-Type Bound for Finite
Markov Chains. Annals of Applied Probability, 8(3), 849-867.
\textbf{{[}117{]}} Che, H., Wang, Z., and Tung, Y. (2002). Analysis
and Design of Hierarchical Web Caching Systems. Proceedings of IEEE
INFOCOM, 1416-1424.
\textbf{{[}118{]}} Aldous, D., and Diaconis, P. (1987). Shuffling
Cards and Stopping Times. American Mathematical Monthly, 93(5),
333-348.
\textbf{{[}119{]}} Knudsen, J. B., and Tsitsiklis, J. N. (1990).
On the Time Required to Drain a Network. Mathematics of Operations
Research, 15(4), 685-708.
\textbf{{[}120{]}} Levin, D. A., Peres, Y., and Wilmer, E. L. (2017).
Markov Chains and Mixing Times, 2nd edition. American Mathematical
Society.
\textbf{{[}121{]}} Gillman, D. (1998). A Chernoff Bound for Random
Walks on Expander Graphs. SIAM Journal on Computing, 27(4),
1203-1220.
\textbf{{[}122{]}} Paulin, D. (2015). Concentration Inequalities
for Markov Chains by Marton Couplings and Spectral Methods.
Electronic Journal of Probability, 20(79), 1-32.
\textbf{{[}123{]}} Aldous, D., and Fill, J. A. (2002).
Reversible Markov Chains and Random Walks on Graphs.
Unfinished monograph, available at
https://www.stat.berkeley.edu/\textasciitilde aldous/RWG/book.html.
\textbf{{[}124{]}} Diaconis, P., and Stroock, D. (1991). Geometric
Bounds for Eigenvalues of Markov Chains. Annals of Applied
Probability, 1(1), 36-61.
\textbf{{[}125{]}} King, W. F. (1971). Analysis of Demand Paging
Algorithms. Proceedings of the IFIP Congress, 485-490.
\textbf{{[}126{]}} Kesidis, G., Shan, Y., Fumagalli, A., Gulati,
A., and Singh, A. (2017). An Analysis of LRU Replacement Caches
under Independent Reference Models. arXiv:1704.04849.
\textbf{{[}127{]}} Niyogi, P., Smale, S., and Weinberger, S. (2008).
Finding the Homology of Submanifolds with High Confidence from
Random Samples. Discrete and Computational Geometry, 39(1-3),
419-441. (Introduces the ``reach'' parameter of a sub-manifold and
gives quantitative manifold-learning sampling bounds.)
\textbf{{[}128{]}} Muirhead, R. J. (1982). Aspects of Multivariate
Statistical Theory. Wiley.
\textbf{{[}129{]}} Diaconis, P., and Freedman, D. (1987).
A dozen de~Finetti-style results in search of a theory. Annales de
l'Institut Henri Poincar\'e (B) Probabilities et Statistiques, 23(S2),
397-423.
\textbf{{[}130{]}} Gersho, A., and Gray, R. M. (1992). Vector
Quantization and Signal Compression. Kluwer Academic Publishers.
(Bennett high-resolution quantization theory, Chapter 5.)
\textbf{{[}131{]}} Talagrand, M. (1995). Concentration of Measure
and Isoperimetric Inequalities in Product Spaces. Publications
Math\'ematiques de l'IH\'ES, 81, 73-205.
\textbf{{[}132{]}} Bennett, W. R. (1948). Spectra of Quantized
Signals. Bell System Technical Journal, 27(3), 446-472.
\textbf{{[}133{]}} Ryabko, B. Ya. (1984). Twice-Universal Coding.
Problems of Information Transmission, 20(3), 173-177.
\textbf{{[}134{]}} Willems, F. M. J., Shtarkov, Y. M., and Tjalkens,
T. J. (1995). The Context-Tree Weighting Method: Basic Properties.
IEEE Transactions on Information Theory, 41(3), 653-664.


\textbf{{[}RNR-II{]}} Ratushnyi, P. (2026). Ratushnyi Neural Repair
Coding, Part II: Random Access and Deviation Theory. Companion paper (RNR
series, Part II).

\textbf{{[}RNR-III{]}} Ratushnyi, P. (2026). Ratushnyi Neural Repair
Coding, Part III: Operational Rate--Distortion Dispersion. Companion paper
(RNR series, Part III).

\end{document}
